\documentclass[10pt,a4paper,twoside,twocolumn,openany,bg=print,nodeprecatedcode]{dndbook}
\usepackage[utf8]{inputenc}
\usepackage{tabulary}
\usepackage[normalem]{ulem}
\usepackage{color,soul}
\usepackage{float}
\usepackage{caption}
\usepackage{wrapfig}
\usepackage{tocloft}
\usepackage[hidelinks]{hyperref}
\raggedbottom
\extrafloats{2000}
\cftsetindents{section}{1em}{0em}
\cftsetindents{subsection}{2em}{0em}
\cftsetindents{subsubsection}{3em}{0em}
\setcounter{tocdepth}{1}
\begin{document}
\addcontentsline{toc}{chapter}{Contents}
\tableofcontents
\newpage
\chapter{Welcome to Lairs of Etharis}
\begin{figure}[h]
\includegraphics[width=0.5\textwidth]{images/adventure/GHLoE/000-00-001.welcome-to-etharis.png}
\end{figure}
\section{Introduction}
From its deep and forbidding forests, where horrific aberrations lurk, to the hulking and impenetrable mountain ranges where lost legacies lie, the landscape of Etharis is formidable. However, the true monsters of this land are found in the degenerating cities of old, where humanity clings to survival. While nefarious figures entice those who seek power, the devotees of dead gods punish the impious. There are few heroes in a world as morally corrupted as this. The only hope in Etharis is that which you bring, yourself.
Etharis is the official world for the Grim Hollow setting. The adventures found herein are an introduction to this gritty and dangerous realm. The rules added in the Grim Hollow tomes are designed for weaving dark fantasy tales. Blood magic, undeath, lycanthropy, fey bargains, fell sorcery, and other temptations: would-be heroes are always one wrong step away from becoming the villains they claim to fight.
\section{Monster Hunting}
Grim Hollow: Lairs of Etharis offers a wealth of new creature stat blocks, allowing GMs to challenge, frighten, and delight their players. From new characters cutting their teeth on relatively weak threats, to the most advanced monster slayers dealing with world-shattering events, the monsters in this book offer something for everyone.
The monsters are drawn from every corner of horror and dark fantasy: gothic intrigue, existential dread, gruesome fairy tales, slasher horror, undead hoards, body horror, killer dolls, and so much more.
\section{Horror Triggers Horror}
With so many horrible monsters and threats at your fingertips, there's a good chance that something presented here will bring into your players' minds something that they're uncomfortable with. Before you use these monsters, or run any sort of horror-based campaign, we urge you to discuss with your players their preferences, desires, and limits. Gore, child endangerment, loss of free will, nightmares: all of these and more infest these pages.
The previous Grim Hollow books talk about safety tools and methods you can use to ensure that everyone involved in your game is having a good time and avoiding problematic material. We suggest you review those guidelines—or use other safety tool procedures and guidelines—to learn more about running games, especially games steeped in horror and dark fantasy, for any player.
\section{Capturing the Grim}
Etharis provides a great setting for a grim, gritty campaign. One tool provided in Lairs of Etharis is salvage rules. These rules encourage a type of play that invites characters to make the most of their knowledge and proficiencies, as well as their downtime, to have the best chance of surviving the threats they face. The blood, bones, brains, and bile of the creatures they kill might provide characters advantages later. Weapons, potions, components, and other magic items might not be readily available, but the characters can make their own from salvaged parts of monsters. And in some cases, monsters might be nigh invulnerable unless the characters use their wits.
\section{A Little Knowledge Goes a Long Way}
Like the salvage rules, lore plays an important part in a dark fantasy setting. Characters may be in trouble if they don't take the time to learn more about the monsters they face. Local legends about a beast may be myth or folly, but even within the folk tales, there could be an iota of truth—and that truth could be the difference between surviving an encounter and suffering an ultimate defeat.
The ability checks associated with each monster are meant to provide details on what a character might know about a monster—or what a character might find out. Use these checks to dole out important information. Some checks might require characters to do some research, or even face the creature to observe its traits and behavior. Failed checks could result in characters gaining incorrect or partial bits of information, which could lead to dire consequences.
\section{Ultimately, the Goal is Fun}
Regardless of how you run your games, or what your players want from a game, it's important to remember that the goal of playing this game is to have fun and tell great stories. The monsters in this book are there to challenge characters, but they're also there to enhance stories. No bit of lore or game mechanical number is more important than the fun and the story. Change whatever you must to ensure everyone at the table enjoys themselves, contributes to the story, and leaves the table satisfied.
With that said, prepare yourself to encounter some of the most terrifying, challenging, entertaining, and bizarre creatures to ever stalk the lands of your imagination!
\section{You are a Game Master}
As Game Master (GM) of a Lair, you have the power of a storyteller, spinning tales of heroism and adventure to a rapt audience. Imagine yourself dressed in the clothes of a traveling bard, your face cast in glowing firelight as you weave a tapestry of words. Your words draw those who sit around the fire with you into another world—but they aren't just the audience of these tales, but actors and storytellers themselves.
The GM's role is to establish the setting of each scene in the Lair, introduce and portray its nonplayer characters (NPCs), and play monsters and villains who long to bring a swift end to the heroes' adventures. The players at your table are storytellers in their own right. Though they embody a single character each while you embody many, the players have the power to make decisions that will change the world forever. The best GMs are willing to improvise, even ignoring the text when necessary, to ensure the characters' actions have consequences—for both good and ill.
If you don't plan to run a Lair, stop reading now. What follows is for the Game Master's eyes only.
\section{Monsters, Magic, and More}
When you see text in bold, that's our way of telling you a scene or location contains a monster or NPC that's ready to fight; that monster's game statistics are found in the fifth edition core rules. Text in bold* followed by an asterisk indicates a Grim Hollow monster; that monster's game statistics are found at the end of that chapter and in Grim Hollow: The Monster Grimoire.
When you see text in italics, that's to tell you that you're looking at the name of a spell (like Tasha's Hideous Laughter) or a magic item (like winged boots). Most spells and magic items are found in the fifth edition core rules. Text in Italics* followed by an asterisk are Grim Hollow spells or magic items and are described in are described in Appendix A of Lairs of Etharis and in Grim Hollow: The Players Guide or Grim Hollow: The Monster Grimoire.
\begin{DndReadAloud}{}
When you see boxed text like this, the characters have likely entered a new area or triggered an event, such as an NPC entering the scene. This text is meant for the GM to read or paraphrase aloud to set the scene.
\end{DndReadAloud}

\DndQuote%
{}
{''lore plays an important part in a dark fantasy setting...''}
{}
\chapter{1: Bone House}
By Tom "Evhelm" Donovan
Bone House is a bone trader lair for four or five 1st-level characters.
\section{Background}
The bone trader Fibekruk has built its "estate" into the cool caverns of this forest-adjacent hillside. The caverns include its bone manse (a haunting negotiations chamber), a nest for its tamed knifewings, a storeroom, and private chambers. Outside the caverns is the forest, which hides the huts of Fibekruk's faevlin minions, their refuse pit, and a ruined cottage used as a storage space.
The bone trader is busy polishing and cleaning its bone collection; its recent efforts have brought it a measure of success, allowing it to move from animal bones to much more prestigious (and tasty) human bones. The creatures it has pressed into service rest comfortably in their nest, content that the faevlin bring them enough snacks—and that the faevlin themselves are always available as snacks if necessary. A small tribe of nine faevlin are spread out across the Bone House. One is assisting the bone trader with its work. Two are moving supplies into the ruined cabin. Two more are searching the refuse pit for trinkets and gold. Two are out hunting (and can return if the characters take too long exploring or dealing with the lair) while two are sleeping in the huts at the camp. If the faevlin have recently abducted anyone, the abductees are imprisoned in the animals' nesting chamber.
\DndQuote%
{}
{''The bone trader is busy polishing and cleaning its bone collection...''}
{}
\subsection{Set the Hook}
Fibekruk has lately ordered the faevlin to collect bones from humans living in the nearby villages. The characters have been hired by a local village to deal with the faevlin spotted in the area. Perhaps someone close to the characters has also been abducted, and the trail led here. If the characters are returning to the area after a long absence (perhaps after receiving training in a distant land), the ruined cottage could hold significance to one or more of them. Of course, it's also possible the characters simply run across the Bone House in their explorations.
\section{Lair Overview}
The lair is described as it is during nighttime hours, when the faevlin do their best work. During the day, two sleepy faevlin stand watch with the knifewings while the rest of the faevlin sleep in their huts. The bone trader also rests during the day. The only light sources at night are the cooking fire in the faevlin camp and the braziers (kept burning constantly) in the bone manse. Otherwise, the forest and cavern are dark. Characters observing the scene see the faevlin eventually swap places or jobs, but other than the day-night crew change, the faevlin always divide into work crews.
Faevlin who encounter characters try to incapacitate characters rather than kill, hoping to offer them to Fibekruk. Even if all the characters are incapacitated, it may not be the end of the adventure: set up a meeting between the characters and the bone trader. They can make a deal for their lives—though the cost should be suitably dark.
Dealing with the bone trader is making a proverbial deal with the devil: you likely get what you want in the here and now, with regrettable consequences down the line. For example, bone traders are willing to give up immediate gratification (snacking on a character or one of the character's NPC friends) in exchange for long-term satisfaction (such as a blood oath that the character will bring two humanoids for it to feast on before the next full moon). Captured characters unwilling or unable to make a deal may find themselves fed on, one at a time, over the next few days. Escape is always possible, of course!
\subsection{1. Ruined Cabin}
\begin{DndReadAloud}{}
Rotting wooden bookcases, chairs, and a table are covered in damp moss. The southeast corner of the roof has collapsed, taking the nearby walls down with it. Three stout barrels rest in the northwest corner.
\end{DndReadAloud}

Two faevlin* are moving some new barrels into the cottage, one at a time, from outside. The bookshelves are empty, its items were moved to Fibekruk's bedroom (area 6). The barrels contain fresh water and pickles.
\subsection{2. Faevlin Camp}
\begin{DndReadAloud}{}
Six tiny huts, poorly constructed, are loosely spaced around a cooking fire among the trees.
\end{DndReadAloud}

Two faevlin* are sleeping in separate huts. The cooking fire is burning low; it has not been attended for some time.
\paragraph{Gold Pouch.}
The faevlin have meaningless trinkets and talismans, but one of them had the guts to secret away 3 gp in a bladder hidden in the hay coating the floor of one of the huts.
\subsection{3. Refuse Pit}
\begin{DndReadAloud}{}
A 15-foot-deep hole has been roughly dug out of the forest floor; the dirt removed from it sits in haphazard mounds nearby. The stench emanating from it identifies it as an outhouse and refuse pit.
\end{DndReadAloud}

Two faevlin* were rooting around at the base of the pit for valuables when the stake they used to hold their rope snapped and left them stuck at the bottom. They have noticed this, but they're too busy hunting for gold to panic yet.
\paragraph{Treasure.}
One of their many now-boneless victims at the bottom of the pit carried a small leather purse with 9 gp inside. If the characters watch the faevlin search for more than 10 minutes, the faevlin find it. The characters can find it themselves with a successful DC 10 Intelligence (Investigation) check.
\begin{figure}[h]
\includegraphics[width=0.5\textwidth]{images/adventure/GHLoE/003-01-001.faevlin.png}
\end{figure}
\subsection{4. Knifewing Den}
\begin{DndReadAloud}{}
This large cavern smells of blood and refuse. Two piles of straw are placed opposite each other, and four rough metal cages are arrayed about the chamber. One end of the cavern opens into the forest, the other deeper into the cave.
\end{DndReadAloud}

\begin{figure}[h]
\includegraphics[width=0.5\textwidth]{images/adventure/GHLoE/004-01-002.knifewing.png}
\end{figure}
Three Knifewing* are sleeping lightly in their beds. They are friendly to the faevlin (if they've been fed) but are hostile to other creatures not escorted by either the bone trader or faevlin.
\paragraph{Cages.}
Two of the cages are empty and two contain sheep stolen from nearby farms. As GM, you may decide if any prisoners are trapped in the cages, such as villagers from the nearby town or an NPC the characters have been sent to find. If you need a default humanoid prisoner, use Dirk Wheatfield, a human farmer who made the mistake of investigating sounds in his barn the previous night before being beaten senseless and dragged here. He has 1 hp and is suffering a level of exhaustion. His left leg is missing its bones below the knee, hobbling him. Magical healing restores his ability to walk.
His cage can be broken open with a DC 10 Strength (Athletics) check or unlocked with a DC 10 Dexterity (Thieves' Tools) check. If saved, he knows nothing about the interior of the Bone House but conveys that there are at least half a dozen faevlin in addition to the three knifewings and the bone trader. He refuses to go further inside and wishes to be taken home.
\subsection{5. Bone Manse}
\begin{DndReadAloud}{}
This large square pagoda has walls comprised of thousands of interlocking bones. Through the open doorway you can see a macabre set of chairs and a desk, also made of interwoven bones, crested with skulls of woodland creatures. A humanoid skull features prominently on the desk's center. Metal braziers built into the corners cast the space in flickering shadows.
\end{DndReadAloud}

By default, this space is empty, but if the characters have arranged a meeting with the bone trader, this is where he meets them—from behind his desk in his bone manse.
\paragraph{Unsettling Hollow.}
Between the bones and the "artwork," this space is designed to unsettle and disturb the bone trader's guests (or enemies). Charisma ability checks or saving throws made by humanoids other than Fibekruk and the faevlins suffer disadvantage while in this room.
\paragraph{Skull of Cohesion.}
Bones are notoriously difficult to bind together without copious binding agent. Most of the time, bone-builds are merely bone-accented! Skilled bone traders are capable of crafting a common magic item known as bones of cohesion, providing an aura of stability to structures in a 30 foot radius. This skull of cohesion lends its strength to the pagoda. If removed from the space (30 feet away from the structure), the pagoda collapses like a house of cards. Characters may use this item themselves to quickly set up or take down an encampment, but more usefully it makes donning (and subsequently doffing) armor a bit faster. Characters bearing the skull of cohesion may halve the time to don heavy armor and (if they set the skull no more than 30 feet away) halve the time to doff that same armor.
\subsection{6. Bedroom and Showcase}
\begin{figure}[h]
\includegraphics[width=0.5\textwidth]{images/adventure/GHLoE/005-01-003.bone-trader.png}
\end{figure}
\begin{DndReadAloud}{}
Bones arranged in a herringbone pattern line the floor of this smooth-walled chamber. A canopied bed, framed paintings, and sturdy wardrobes furnish this otherworldly room. Three bony pedestals are each topped with a different skull—one large and tusked, one so tiny as to be awkward for the size of the pedestal, and one that is just the lower jawbone of a medium-sized creature.
\end{DndReadAloud}

One faevlin* is busy polishing the tusked skull, a trade made by the bone trader some time ago. Even he does not recognize the creature from which it came. The bone trader* fusses over the jawbone; it was once a fully intact skull, but Fibekruk has been stress-eating it, and the jawbone is all that remains. The tiny skull belonged to a fairy.
\paragraph{Wardrobes and Linens.}
The bone trader acquired its linens and the costumes in its wardrobes over several years. Most are beautiful, but fragile. Collected, the linens and clothes weigh around 30 pounds, are extremely bulky, and would be worth 20 gp to a collector. Traveling with them is likely to make them disintegrate: for each day a character moves the linens, roll a d10. On a roll of a 1, the linens are damaged beyond repair and become worthless. (Characters who have proficiency with a tailor's kit cope with the fragility and are not subject to the roll.)
\paragraph{Coin Purse.}
A search of the area with a successful DC 12 Intelligence (Investigation) check reveals a section of the herringbone floor that doesn't match the pattern. Underneath is a small hollow where the bone trader hides the coins and other bargaining chips stolen or traded for from others. Its coffers are low, having recently acquired the tusked skull at a great price. Now, 5 gp and 49 sp are in the purse.
\subsection{7. Storeroom}
\begin{DndReadAloud}{}
Stacks of bones of varying sizes have been meticulously organized and stacked into 6-and-7-foot square piles. Barrels, crates, jars of oils, foodstuffs, and sundries are marked with a variety of symbols belonging to farms and villages throughout the region.
\end{DndReadAloud}

In addition to the bones, the other supplies are mostly kept by Fibekruk to feed itself and pacify the faevlin and knifewings, though it does occasionally offer to trade these supplies for more bones.
\paragraph{Sundry Supplies.}
If hauled back to civilization, most of these supplies can be returned to their original owners. If characters are willing to return the supplies, the villagers are extremely grateful and owe the characters a future favor. Alternately, the characters may sell the sundries, netting 15 gp profit. Moving this many barrels and crates requires a wagon or other means of overland transport.
\section{Conclusion}
Lucky or particularly skilled characters who manage to kill, defeat, or otherwise drive off Fibekruk and its minions become folk heroes in the nearest village, gaining offers of free meals and lodging for the next two weeks in the village.
The characters may tussle with and eventually make a deal with the bone trader. Villagers don't usually have a lot of coin to use as payoff, so rewards they offer may be in information (perhaps leading to a new adventure) or admiration. If necessary, the villagers can give the characters 10 gp each to convince the bone trader to stop preying on them. This is reduced to 5 gp each if they are only able to convince the bone trader to agree not to hunt the people, leaving the animals as fair game. Obviously, if the characters make a deal that saves their own bones but doesn't account for nearby villages' needs, the villages won't pay them anything!
Dealing with the bone trader could be the first step toward helping a local ruler regain control of their fiefdom, leading the characters back to them for another quest against sinister forces in these lands. The bone trader may mention that its own actions are limited by an even greater bone trader in another vale—one who has several bone trader vassals of its own! Or, perhaps, the bone trader is annoyed because the characters have been sent in breach of an existing contract, and the next adventure for the characters investigates who sent them into a deadly situation to get rid of Fibekruk!
\DndQuote%
{}
{''Stacks of bones of varying sizes have been meticulously organized and stacked into six-and-seven-foot square piles.''}
{}
\section{Creatures}

\begin{itemize}
\begin{multicols}{3}
\item Bone Trader
\item Faevlin
\item Knifewing
\end{multicols}

\end{itemize}

\chapter{2: Faren's Rest}
By Tom "Evhelm" Donovan
Faren's Rest is a vitebriate lair for four or five 2nd-level characters.
\section{Background}
Faren's Rest is a roadside village that grew up along a major trading route. Once, it was just an inn and tavern. Over time, it became an important waystation: a stable and a small hamlet grew up around the inn. A few generations ago, a local baron (Faren) got the bright idea to put a tax collector at the waystation. He quickly abused the tax, building a house across from the inn to keep an eye on things—along with a statue of himself as a heroic warrior. Eventually one of the baron's many enemies did him in, and Faren's Rest became independent once again.
Recently, the newest mayor of Faren's Rest, Rolant Joss, began acting strangely. In recent weeks, he's taken in an orphan child, hired bodyguards, and started talking about reviving Faren's dreaded tax.
The vitebriate is playing the role of the poor orphan child calling herself "Rose," taken in by the mayor. Rose has grown greedy in her years of avoiding discovery and sees in Rolant an opportunity to play the long game and accumulate wealth from the old tax. Eventually, she figures, the town will get mad at Rolant (or she will use up her host), and either way she will be seen as the "innocent victim" rather than an unworthy beneficiary of the wealth, and she can disappear again. Unfortunately, her domination of Rolant is making him behave very strangely, and some in town fear something is wrong.
\subsection{Set the Hook}
There are a few ways the characters might be involved. The easiest is that the inn and tavern (The Outrageous Pig) is a recommended waystation for the party. Perhaps they're on the way to another adventure, or perhaps they are returning from one. When they arrive for their well-earned rest (from questing or traveling), the owner Betha Mathildis asks the party for help. She needs outsiders to talk to the mayor and try to convince him not to revive Faren's road toll and tax.
If you're worried about the players not taking the hook, have the mayor be a friend of someone in the party or even send the party to Faren's Rest to meet with the mayor. "Remind" the players that the mayor is known for being social, for his generosity, and for his verbosity.
\section{Lair Overview}
The first time the characters see Rolant and Rose is in the mansion (area 3), but after that, they can be in Rolant's friend's home (area 8). If the characters spend more than 24 hours before solving the mystery and dealing with Rose, two bodyguards join Rolant and Rose.
\subsection{1. The Outrageous Pig (Inn and Tavern)}
\begin{DndReadAloud}{}
The first floor of The Outrageous Pig is a subdued but cozy roadside establishment. The fireplace roars and crackles, and the barkeep clinks glasses as he buses the bar. A matronly woman leans over the bar and beckons you toward her.
\end{DndReadAloud}

Betha Mathildis (NG, female human commoner) is the matronly woman. Five regulars sit at tables waiting for food, two more stand at the bar, the cook and waiter work in the kitchen, and upstairs three traveling merchants rest in two of the five rooms (all Commoner except the noble Ardel).
Betha invites the characters to dine for free if they'll handle a small errand: she has been hoping for some competent-looking outsiders to give Rolant Joss (LG male human commoner, dominated by the vitebriate*, Rose) a reminder that Faren's tax is a gross injustice by delivering him only a portion of the meals he ordered. Betha is more than willing to share her positive opinion about Rolant (and how he has changed recently), and if the characters ask what might have changed him, the only thing she can think of is adopting Rose. She insists the characters investigate. If the characters want to haggle or insist on further payment for the errand, she informs them that The Outrageous Pig is the only inn for miles, and she won't give them a room if they won't help.
She's serious about the food and the room; once the characters agree to help, she sells them two or three discounted rooms (two beds per room) at a rate of 3 sp per room along with the promise that if they figure out what's wrong with Rolant, she'll comp the rooms entirely.
\paragraph{Ardel Bigring.}
One of the people getting a drink at the bar is Ardel Bigring. He and his family live next to the inn. If the characters talk to him, he invites them back to his home to meet his wife and talk there—he thinks Betha's mistrustful talk might be bad for the public to hear, but he's happy to talk in private.
\DndQuote%
{}
{''The fireplace roars and crackles, and the barkeep clinks glasses...''}
{}
\subsection{2. Stables}
\begin{DndReadAloud}{}
The strong, sharp scent of horse dung and hay permeate these stables. A few post-horses and some work animals are quartered here.
\end{DndReadAloud}

The stable boy, Merk (N, male human commoner child) tends the horses and has a blooming black eye. He is reluctant to say where he got it, but a successful DC 10 Charisma (Intimidation or Persuasion) check reveals that Rolant hit him this morning for being "too slow" to prepare a riding horse for his little girl Rose. He's never been hit by Rolant before, and he thinks it must be the fluke of a bad mood. Merk is authorized to rent workhorses at 5 sp per day.
\paragraph{Helpful Merk.}
Merk is a nice kid, but extremely gullible—and he idolizes Rolant. If the characters tell him that Rolant said to do something, he believes them without question and does his best to carry it out.
\subsection{3. Faren's Mansion}
\begin{DndReadAloud}{}
The outside of this three-story brick house is overwrought with fine details. The four corners each have decorative edges and statues of imposing angels, the windows are tall and uniform, and the bricks themselves are of a deeper shade of bruised purple than the bright red of the inn. The interior is compact but aggressively opulent, as if someone decided to furnish a home three times this size and put all the furniture in the same room.
\end{DndReadAloud}

Rolant Joss (N, male human commoner) and the vitebriate* Rose are both on the first floor. While meeting the characters, Rose tries to achieve her goals without completely eroding the people's trust in Rolant as their mayor.
Characters who arrive at the "mansion" with dinner from Betha are admitted immediately. It's easy to see that Rolant doesn't match Betha's description of him. The first time they meet him and Rose, have the characters make a Wisdom (Insight) check:

\begin{description}[nolistsep]
\item[0+:.]
Rolant is not as he was described.
\item[5+:.]
He is gaunt and taciturn.
\item[10+:.]
He seems lethargic, even sickly.
\item[15+:.]
His eyes are glazed and dead. Everything he does seems to be stiff and deliberate, like a puppet on a string.
\item[20+:.]
Something is clearly controlling his words and actions.
\end{description}

No amount of talk persuades him to change his mind about the tax. Though he's a bit sluggish in his responses, his words are concise and unyielding.
\paragraph{Rose, the Vitebriate.}
Rose has managed to secret a gem worth 67 gp in a purse hidden on her person. Wherever she is killed, this purse remains even after she crumbles into a swarm of harmless roaches.
A successful DC 15 Wisdom (Perception) check reveals two details about Rose that are striking. Although she wears fine silk clothes, flecks of moss cling to the cuffs of her blouse. Additionally, although she is very clean and tidy, flecks of rust sit beneath her nails. The moss is from hiding a ring beneath the statue (area 4) and the rust is from the bucket in the well (area 5).
If the characters want to focus on Rose, have them make a Wisdom (Insight) check:

\begin{description}[nolistsep]
\item[0–13:.]
She is a frightened little girl.
\item[14–18:.]
She may be hiding behind Rolant, but she is fierce and focused.
\item[19–23:.]
She neither holds nor carries herself like a child. She seems annoyed that you're here.
\item[24+:.]
Something is clearly off about her—she's not human, and she's armed with a hidden dagger.
\end{description}

\paragraph{Orb of the World.}
An expensive-looking orb, representing the world as an Arch Seraph wants it to be, is on the dining table as a centerpiece. The orb is made of solid steel and can be used as a club. If threatened here, Rose orders Rolant to use it to defend them. (It uses the same stats as a dagger but change the damage from piercing to bludgeoning.) If taken as treasure, the orb can be sold for 10 gp.
\begin{figure}[h]
\includegraphics[width=0.5\textwidth]{images/adventure/GHLoE/012-02-001.vitebriate.png}
\end{figure}
\subsection{4. Defaced Statue of Faren}
\begin{DndReadAloud}{}
This covered gazebo houses a disproportionately muscular stone statue of a human man. The base of the statue has been inscribed with: "Faren, Conquering Wart—" The words clearly once read "warrior," but have been chiseled out. The face of the statue has been similarly broken.
\end{DndReadAloud}

\paragraph{The Statue.}
Characters who investigate the statue might recall two things. With a successful DC 15 Intelligence (History) check, characters learn the information from the first paragraph of the "Background" section. With a successful DC 10 Intelligence (Investigation) check, characters learn that the statue has been defaced for decades—even the exposed defacements show signs of erosion and a lack of dust or debris.
\paragraph{Hidden Cache.}
A successful DC 20 Intelligence (Investigation) check reveals a hole at the base of the statue covered with moss. Rose stole a silver toe ring (worth 20 gp) from Zila Bigring and hid it here to use in case she needed to flee the area quickly.
\subsection{5. Well}
\begin{DndReadAloud}{}
This deep well serves as a source of water for the whole hamlet. A simple wooden winch system attached to a metal bucket draws out the water. It is dozens of feet deep.
\end{DndReadAloud}

A successful DC 15 Intelligence (Investigation) check reveals that the underside of the metal bucket is a bit rusty. Characters checking the underside of the bucket find a silver dagger secured there, hidden by Rose after she took it from Merk. The characters can keep it, but if they return it to Merk, he is grateful. The stablemaster, if the characters do this and save the town from Rose as well, offers the characters a single riding horse as payment.
\DndQuote%
{}
{''Characters checking the underside of the bucket find a silver dagger secured there, hidden by Rose after she took it from Merk.''}
{}
\subsection{6. Outhouses}
\begin{DndReadAloud}{}
This tiny wooden structure with a half-moon carved out of its door reeks of excrement.
\end{DndReadAloud}

Two outhouses are marked on the map. In one of the outhouses, Rose has secured a Alchemist's Fire (flask) to use in an emergency.
\subsection{7. Wealthy Merchant's Home}
\begin{DndReadAloud}{}
This comfortable living space has three plush chairs seated around a roaring hearth fire. A small dining area waits behind. Doors leading to a pair of bedrooms exit to the north and south, and a door at the end of a hallway leads outside.
\end{DndReadAloud}

This home belongs to Ardel Bigring (LN male noble) and his family. He is a staunch supporter of Rolant. At his friend's (dominated) urging, he sent for two Thug to guard Rolant.
\paragraph{Ardel and Zila Bigring.}
Ardel and his wife Zila are easy to talk to and willing to discuss Rolant. Zila believes something is wrong; Ardel admits his friend is behaving strangely but believes Zila and Betha are overreacting. If the characters try to convince Ardel that something is wrong (and/or that Rose is a threat), he can be persuaded with a successful DC 15 Charisma (Persuasion) check. If Zila is present when the check is made, the characters have advantage on the roll. If Ardel is persuaded, he calls off the thug support for Rolant.
Zila also mentions that small things have gone missing from her home. Just in the last two days a small silver toe ring has gone missing. She doesn't remember when she saw it last, so it might have fallen off. But it's just as likely that someone stole it.
\subsection{8. Ambitious Merchant's Home}
\begin{DndReadAloud}{}
The common room of this small home has a large dining table and a corner workspace covered in ledgers and bolts of cloth. Beyond a large hearth and around the corner is the home's only bedroom.
\end{DndReadAloud}

Morwa Curwen (N female commoner), a down-on-her-luck textile merchant desperate to get ahead, has agreed to help Rolant restart Faren's Tax for a small cut of the toll. Once Rolant and Rose have talked to the characters at Rolant's mansion, this is where they come to hide out and gather strength. Morwa has invited some of her family and employees to come (armed with clubs) to protect their mayor in case Betha tries to threaten him. If the characters come here to talk to Morwa or Rolant (or Rose), they must deal with a small mob of five angry Commoner outside first!
However they deal with the commoners, Morwa, Rolant, and Rose remain inside; Morwa follows Rolant's lead, and Rose has no problem using Morwa to guilt the characters into staying away from her and Rolant, if possible.
\paragraph{Bolts of Cloth.}
The fabric on the corner table can be used as makeshift bindings to tie someone up in a pinch.
\subsection{9. Assorted Homes}
\begin{DndReadAloud}{}
This small home has a bed, desk, chair, wardrobe, and large fireplace.
\end{DndReadAloud}

Each of the buildings labeled area 9 contain similar furnishings and belong to the villagers of Faren's Rest.
\subsection{10. Open Road}
\begin{DndReadAloud}{}
The road through Faren's Rest is well-traveled and packed solid despite being a glorified, widened dirt path. Various buildings line the road, their chimneys lazily dusting the sky with smoke.
\end{DndReadAloud}

Characters who stay on the road for a prolonged period see wagons, travelers, and even the occasional caravan clop through, roughly one each hour. Every third group passing through stops at the tavern for a meal and/or at the inn and stables for the night. Wagons are parked behind the stables when necessary.
\section{Conclusion}
Chasing the vitebriate "Rose" out of Faren's Rest only delays the inevitable: only killing her releases Rolant from his domination and saves the town from a slow death at the hands of the greedy creature. Assuming the characters eventually kill the vitebriate, the question is not whether they have any success, but rather what that success looks like.
If the characters fail to intervene soon enough, it's possible that Betha will be killed by Rolant's borrowed thugs. Similarly, if they kill lots of innocents who are protecting Rolant, is that really a victory? The characters may claim the title of the town's heroes, or leave a subdued, somber town behind them as it deals with the chaos wrought by their own trusting natures. If the characters kill anyone without sufficient evidence of wrongdoing, they might also find themselves wanted by the law, a bounty put on their heads.
\DndQuote%
{}
{''Chasing the vitebriate "Rose" out of Faren's Rest only delays the inevitable:''}
{}
\section{Creatures}

\begin{itemize}
\begin{multicols}{3}
\item Vitebriate
\end{multicols}

\end{itemize}

\chapter{3: The Leatherhollow}
By Tom "Evhelm" Donovan
The Leatherhollow is a skinweaver lair for four or five 3rd-level characters.
\section{Background}
Once upon a time, a cave belonging to morbus kobolds was attacked by a skinweaver. The creature slew their leader, installing itself as "chief" over the kobolds. It expects them to deliver fresh supplies of food and (just as importantly) hides for it to tan and work. When they fail to deliver sufficient hides, the skinweaver uses the kobolds for that purpose.
This cave system has three major areas: the deepest belongs to the skinweaver itself. The midrange depths belong to a spythronar swarm. The caves closest to the outside belong to the skinweaver's kobold minions.
\DndQuote%
{}
{''The characters are hired to investigate the mysterious disappearances from farms on the outskirts of the village.''}
{}
\subsection{Set the Hook}
The characters are hired to investigate the mysterious disappearances from farms on the outskirts of the village. First it was just livestock, but now people are vanishing too. In the latest abduction, a child saw his parents carried off by "little scaly men" (kobolds). To increase the suspense, have one of the missing abductees be a family member or friend of the characters.
Alternately, the characters could be investigating why kobolds, who used to fear coming too close to human settlements, have suddenly become more brazen, stealing livestock and abducting people.
\section{Lair Overview}
No natural sources light the caves, so only the initial cave opening has daylight. The tunnels between the three cave sections are cramped enough that only Small creatures (like kobolds) can walk them upright, but larger creatures need to squeeze. The caves are damp and cool with smooth walls and floors. In the spythronar and skinweaver sections, spiderwebbing is prevalent on walls and ceilings and make the floors tacky.
At least once a day, a pair of kobolds (whoever draws the short straws) must bring their latest catch, whether livestock or humanoid, to the skinweaver. Doing so is dangerous both for the capricious nature of the skinweaver and because of the spythronar in the caves between—both the skinweaver and the spythronar sometimes take a kobold as a meal. The kobolds have not dared to bring a separate offering to feed the spythronar, being more afraid of the skinweaver's knives.
\subsection{1. First Cave (Entrance)}
\begin{DndReadAloud}{}
Beyond the partially barricaded mouth, this cave opens to vaulted ceilings. A large cooking fire burns brightly in the center of the space. Two openings lead deeper into the cave system to the north and east.
\end{DndReadAloud}

\begin{figure}[h]
\includegraphics[width=0.5\textwidth]{images/adventure/GHLoE/016-03-001.morbus-kobold-r.png}
\end{figure}
Six Morbus Kobold* are ringing the fire, telling stories instead of paying attention to the cave entrance. Their passive Perception is 6. A fight here draws the attention of the kobolds sleeping in area 3, who arrive as soon as they can. They travel through area 2 to reach area 1 to avoid the pit trap.
\paragraph{Pit Trap.}
The northern exit to this chamber is trapped. A concealed 10-foot-square spiked pit trap, is on area 1's side of the opening. Passive Perceptions of 15 or higher detect it just before stepping on it; a successful DC 14 Intelligence (Investigation) check while looking for traps find it as well. When one character steps onto the trap, that character must succeed at a DC 12 Dexterity saving throw or take 2d10 piercing damage from the spikes in the 10-foot-deep pit.
\subsection{2. Lavatory Pit \& Cages}
\begin{DndReadAloud}{}
This chamber reeks of rotting waste. A large depression in the center buzzes and writhes with the sounds and movements of thousands of insects feasting on the assembled dung and trash. Four crude wooden cages are arrayed on the far side. Three tunnels and an opening lead beyond.
\end{DndReadAloud}

The lavatory pit is smelly and gross but not dangerous. The cages contain at least two sheep and a single humanoid prisoner. If the characters have been sent to find and rescue anyone, all but one of the NPCs they've been sent to rescue are in the cages. The last prisoner is in area 7. Any prisoners here are suffering from two levels of exhaustion from malnutrition and have 1 hp.
\paragraph{Cages.}
The cells are rudimentary and can be broken open with a successful DC 13 Strength (Athletics) check or the locks picked with a successful DC 13 Dexterity (Thieves' Tools) check. There are two keys. The skinweaver carries one and the other is on one of the kobold corpses in the spythronar webs in area 5.
\subsection{3. Barracks}
\begin{DndReadAloud}{}
The stench of malodorous sweat fills this cave. Clusters of rough sleeping mats are radially arranged around an enormous stalagmite pillar in the center. The pillar has been carved to resemble a strange, malformed creature. A few battered crates serve as footlockers for the creatures who dwell here. Two tunnels lead east, larger tunnel grades downward to the north, and a larger opening leads south.
\end{DndReadAloud}

Eight Morbus Kobold* are sleeping in this chamber if they weren't already drawn to fighting in area 1. A fight here brings the three Morbus Kobold* who have snuck off to the root cellar in area 4 to gamble.
\paragraph{Stalagmite.}
A successful DC 15 Intelligence (History) check reveals that this stalagmite has been carved into a crude representation of the Filth Grazer, which the morbus kobolds worship. They've been forced to denounce their allegiance to that entity while serving the skinweaver.
\subsection{4. Root Cellar}
\begin{DndReadAloud}{}
A short down-grade levels out in this musty, earthen space. Poorly piled stacks of potatoes, beets, and onions are heaped around. Most of these dirt-covered foods have sprouted green roots and tendrils. Only one tunnel leads up and out of this space.
\end{DndReadAloud}

The three Morbus Kobold* here, if they haven't been drawn to fighting elsewhere in the complex, are gambling with dice over a beautiful pair of elven boots.
\paragraph{Magical Boots.}
The elven boots are quickheels, magical boots (requiring attunement) that allow the wearer to use a bonus action to double their speed for a single round. This effect requires a long rest before the item can be used again. They have an Elven inscription written across the tongues that, when taken together reads, "You only need to cross the finish line once to win the race."
\subsection{5. Spider Lair}
\begin{DndReadAloud}{}
The expansive cave is obscured by omnipresent webbing: walls, ceiling, and floor are all coated in the tendrils of pale web. Large cocoons, some intact and some ripped open, hint at the hungry nature of the cavern's inhabitants.
\end{DndReadAloud}

The webs are tacky, making the floor (and walls and ceiling) here difficult terrain. One spythronar swarm* is concealed in normal webs on the ceiling at the northwest end of the chamber. A successful DC 17 Wisdom (Perception) check is required to see them before they surprise any creature passing through their territory. At the southeast end of the chamber is a spythronar web*.
\paragraph{Cocoons.}
Among the dozens of entombed animals are several kobolds who had the misfortune to cross the chamber while the spythronar were hunting. One of these cocooned (and deceased) kobolds has the key to the cages in area 2 as well as a small pouch containing a potion of healing. The characters automatically find this kobold and his loot after thoroughly searching the area for ten minutes. Characters who succeed on an DC 13 Intelligence (Investigation) check can speed up this process and do it in five minutes.
\subsection{6. Skinweaver's Blinds}
\begin{DndReadAloud}{}
Emerging from the tunnel, you find yourself in a forest of leather strips and webbing. Spider webs crisscross with patches and strands of leather, making it difficult to see more than a few feet. The strands are attached to the ceiling and form an opaque mass before you, like an overabundance of fabric fringe.
\end{DndReadAloud}

\begin{figure}[h]
\includegraphics[width=0.5\textwidth]{images/adventure/GHLoE/017-03-002.skinweaver.png}
\end{figure}
The skinweaver set up this blind as an early warning system for itself. The leather and webbing move only when disturbed (there is no breeze this far into the caves), disturbing the skinweaver's webs, alerting it to a visitor or intruder. If the characters move directly forward through it or follow the wall to their right, they end up in the skinweaver's chamber. If they follow the wall to the left instead, they find one of the narrow tunnels leading to the tanning chamber in area 8.
\paragraph{Leather and Web Blind.}
Characters particularly worried about being ambushed in the blind may want to try hacking their way through it or burning it. Hacking is reasonably effective (though slow). Each strand of webbing acts as webbing: AC 10; hp 5; vulnerability to fire damage; immunity to bludgeoning, poison, and psychic damage. The leather is similar: AC 12; hp 7; vulnerability to acid damage; immunity to bludgeoning, poison, and psychic damage. Destroying a web and a leather section exposes and clears a 5-foot-square section of the blind. If the characters take more than three rounds to clear the blind or move through it, the skinweaver attacks them through the blind. Because of its web-connections, treat the skinweaver as having blindsight against any characters in contact with the blind.
\subsection{7. Skinweaver Chamber}
\begin{DndReadAloud}{}
Aside from the grotesque splotches of spider web, this cave looks like the most opulent of bedrooms. Leather tapestries hang on the walls and from the high ceiling; a large crude statue of malformed monstrosity peers out from the north wall; a circular mattress rests on a frame of webbing and leather suspended from the ceiling; and bookcases, wardrobes, and mirrors all give the impression of culture and refinement. Atop a three-tiered dais rests a large throne made of different-colored leather strips.
\end{DndReadAloud}

This is the skinweaver's inner sanctum. If not for how all the clothing and furniture is sized for a Large creature with several more limbs than the average humanoid, it could pass for a royal bedroom suite. All the shine is from bronze and brass, however, not gold or silver. Unless it has been fetched by the kobolds to deal with threats elsewhere, the characters probably run into the skinweaver* and three morbus kobold* attendants here. The kobolds recently brought a victim for the skinweaver to flay and eat, and they receive new orders as the characters arrive.
The skinweaver fights to the death not because it's desperate or ideologically driven, but because it's so full of itself that it cannot conceive of its own defeat.
\paragraph{Filth Grazer Statue.}
This crude statue is worthless but belonged to the kobold's original chieftain. She stashed a potion of growth down its throat. The vial can be discovered with a successful DC 13 Intelligence (Investigation) check.
\paragraph{Key.}
The skinweaver carries one of the keys to the cages in area 2.
\paragraph{Tapestries, Wardrobes, Bookshelves, and Desk.}
The majority of the tapestries are sewn to anchor points in the walls and are immovable without destroying them. The clothes in the wardrobe are tailored to the skinweaver and are similarly worthless except for their leather value. The books are on a variety of topics including languages, geography, and, of course, tailoring. Characters who spend at least an hour perusing the books gain advantage on their next Intelligence ability check. The small collection can be sold for 35 gp. The desk is filled with strange notations; even comprehend languages does not reveal the meaning of the scribbles. A character proficient with Leatherworker's Tools or weaver's tools who spends four hours researching them realizes they're patterns for exquisite leatherworking. The patterns can be sold to a collector (once translated out of the skinweaver's shorthand) for 100 gp.
\DndQuote%
{}
{''Hissing, popping, and sizzling sounds...''}
{}
\subsection{8. Tanning Area}
\begin{DndReadAloud}{}
This chamber is stiflingly hot and filled with pungent, eye-watering, lung-burning air. Hissing, popping, and sizzling sounds come from half-a-dozen large metal vats, each containing gallons of liquid. A metal rack, large enough to restrain a human-sized creature, is attached to the east wall. Metal tables covered in a wide range of tools and hides in various stages of tanning interweave between the vats.
\end{DndReadAloud}

If the characters were sent to rescue someone, they find their quarry here, unconscious but stable (0 hp) and strapped to the metal rack. Some of the NPC's skin has been removed with surgical precision in disturbing patches of specific shapes; from a distance, these patches appear as monochromatic red tattoos. If the characters have not been sent to rescue anyone—or if they took days to explore the complex—the NPC on the rack is dead when they arrive.
\paragraph{Vats.}
The vats contain the various chemical agents and acids the skinweaver employs to create its insane masterpieces. A creature who touches a substance takes 1d4 acid damage; a creature submerged in a vat takes 2d6 acid damage per round it is submerged and an additional 2d6 acid damage on the round after it's no longer submerged.
\paragraph{Rack.}
The metal rack is effective at restraining a creature but can easily be opened by any creature not restrained on it.
\section{Conclusion}
Defeating the skinweaver, rescuing the prisoners, and determining the fate of the person in the tanning room ends this adventure and disperses any remaining kobolds. The characters return to the local farms as heroes. Future adventures might seek to answer the question of where the skinweaver came from. Perhaps it only targeted people from a particular extended family or even a particular village: why only them? Who was really pulling the strings?
\DndQuote%
{}
{''Aside from the grotesque splotches of spider web, this cave looks like the most opulent of bedrooms.''}
{}
\section{Creatures}

\begin{itemize}
\begin{multicols}{3}
\item Morbus Kobold
\item Skinweaver
\item Spythronar Sac
\item Spythronar Web
\item Spythronar Swarm
\end{multicols}

\end{itemize}

\chapter{4: Apostate Temple}
By Tom "Evhelm" Donovan
Apostate Temple is a bloodbonded lair for four or five 4th-level characters.
\begin{figure}[h]
\includegraphics[width=0.5\textwidth]{images/adventure/GHLoE/021-03-004.end-of-chapter-splash.png}
\end{figure}
\section{Background}
After turning on both Arch Seraph Solyma and Arch Daemon Venin, a bloodbonded (known to its followers as "the Holy One") has used its twisted lies to gather desperate villagers to the ruins of this forgotten temple. With its cult growing, the bloodbonded orders raids on shrines of the "false" gods, preaching the end of the dominion of Seraphs and Daemons alike. The temple's remote location proves perfect for the Holy One's efforts. Its followers have constructed makeshift storehouses, barracks, and even special chambers for the Holy One itself.
In the temple's sanctuary, the bloodbonded and its followers have knocked over the statue of the discarded god once revered here and marked the sanctuary with both sacred and profane symbols. In its madness, the bloodbonded hopes to summon and bind minor celestials and infernals to its cause. While its conjuring is so far ineffective, given sufficient time and effort it will succeed. Prisoners and willing volunteers wait in cages to be sacrificed.
\subsection{Set the Hook}
A local village, temple, or ruler may hire the characters to seek out the bloodbonded's temple and deal with its inhabitants because of the danger they pose. The characters may have received a vision from a seraphic (or daemonic) entity asking them to deal with the threat. A priest or other holy person might mandate the characters to silence the apostate and scatter its followers. The characters may be exploring a dense and unpopulated forest or a lofty and barren mountain, stumbling across the ruins. If a more personal reason is needed, have the characters track the bloodbonded's followers after an abduction; the bloodbonded is stealing people to sacrifice for its dark purposes.
The temple is designed to be viable in any location: mountain, tundra, forest, mire, etc.—wherever it would make sense that temple could be tucked away undisturbed and forgotten! If it's near a civilized area, the cages in the sanctuary are likely more populated. If not, fewer cages have prisoners.
\section{Lair Overview}
Assume that most of the roof that once covered the stylobate has fallen in, leaving all but the raised sanctuary platform partly exposed to the elements and sunlight during the day. The colonnades are almost entirely intact. Only a pair of large supporting columns in the northwest of the temple have collapsed outward, causing part of the stylobate and roof to collapse along with it.
The site is divided into two major areas: the main temple stylobate (where all the bloodbonded's followers live and where they've built their ramshackle shelters), and the temple's sanctuary. The stylobate level is 15 feet from ground level. The sanctuary is an additional 25 feet above (for a total of 40 feet from ground level). The stylobate level is open on all sides. The sanctuary level is enclosed in a thin stone wall except for the stairs that lead in from the west and a set of eleven 5-foot-diameter circular openings that rest 20 feet above the sanctuary floor. The steeply pitched roof of the sanctuary (intact after the collapse of the main roof over the whole temple) is another 20 feet above the top of the openings.
\subsection{1. Stylobate Base}
\begin{DndReadAloud}{}
At the top of the stairs, the stylobate of this ancient temple stretches before you. A few watchfires and four makeshift structures occupy the space. Up another stair flight several dozen feet away, an intact inner sanctuary rises from the base.
\end{DndReadAloud}

Three Cultist are on guard here near the watchfire. They are supposed to be actively watching the stairs leading up the stylobate, but they're sitting around the fire instead. At the first sign of trouble, they run to the barracks and toward the sanctuary to fetch help.
\paragraph{Watchfires.}
These low fires shed 30 feet of bright light and an additional 30 feet of dim light. They are the only consistent light source on the stylobate level.
\subsection{2. Barracks}
\begin{DndReadAloud}{}
Thirty-six bunked beds are crammed into this tight space, triple-stacked. The stench of sweat and unwashed bedding spill out of these barracks.
\end{DndReadAloud}

Twelve Cultist are resting and relaxing at any given time in the barracks. Their armor is in the armory (area 4), meaning if fighting here they only have an AC of 10, but they do keep their scimitars with them.
\paragraph{Footlockers.}
Small lidless fruit crates have been nailed to the bunks to form open-shelved footlockers for the cramped cultists. Most have a change of clothes, eating utensils, or hardtack, but one has secreted away a small gold bracelet worth 25 gp.
\subsection{3. Storeroom}
\begin{DndReadAloud}{}
This confined shed is stuffed to the brim with barrels, jars, and crates. It smells of flour and brine.
\end{DndReadAloud}

The camp cook (N female commoner named Uliri) is taking inventory in this space. She is unarmed and unarmored but does carry a manifest of the space's contents. She wears a blood-and-vegetable-stained apron. She carries one of the two keys to the bloodbonded's chambers. (The other is on the bloodbonded itself.)
\paragraph{Food Supplies.}
There are more than enough nonperishable supplies here to replenish up to twenty days' worth of rations.
\subsection{4. Armory}
\begin{DndReadAloud}{}
A motley collection of leather armor is stacked in piles on the floor. Piles of pebbles and strips of fabric sit alongside a few crossbows, blades, and shields along the walls.
\end{DndReadAloud}

An eldritch priest* named Blarrie who's given his allegiance to the bloodbonded is stationed here to parcel out armor to cultists. If fighting lasts for more than a round in area 1, he leaves the armory to investigate.
\paragraph{Weapons and Armor.}
The arms and armor here are of poor quality. Twenty sets of leather armor, a dozen Sling and a hundred pellets, four Longsword, five Scimitar, and two Light Crossbow (no bolts) are here for the taking.
\begin{figure}[h]
\includegraphics[width=0.5\textwidth]{images/adventure/GHLoE/024-04-001.eldritch-priest-r.png}
\end{figure}
\subsection{5. Bloodbonded Quarters}
\begin{DndReadAloud}{}
Unlike the rest of the camp, this shed's floor is covered in a patterned rug, its walls lined with tapestries, and a central candelabra quietly illuminates a desk and chair, bookcase, and soft-looking bed.
\end{DndReadAloud}

The door to these quarters is locked. The bloodbonded (if not here, see area 7) carries one of the keys and the camp cook (see Area 3) carries the other. Breaking down the door requires a successful DC 17 Strength (Athletics) check. Picking the lock requires a successful DC 14 Dexterity (Thieves' Tools) check.
The bloodbonded* sleeps here for four hours each day, preferring to do most of its work during nighttime hours. Almost all its other waking time is spent in the sanctuary.
\paragraph{Books and The Mad Treatise.}
Most of the books the bloodbonded has acquired are written in Celestial and center on superstitions and folktales that are 90 percent fiction and 10 percent fact. They are dry grimoires with no literary value or illustrations. Only one book on the shelves stands out from the rest: The Mad Treatise. When the bloodbonded hits a dead end in its summoning attempts, it spends some time adding to its experimental log. The Mad Treatise is a lightly enchanted magic item (conjuration) and can be used as an arcane focus. Any spellcaster who uses the tome as its arcane focus and casts a conjuration spell that conjures at least one creature can choose one of the summoned creatures to be summoned with 5 temporary hp.
\paragraph{Chest.}
The collected valuables of the camp are stored in an unlocked chest: 2,800 cp, 1,300 sp, and 60 gp.
\paragraph{Furnishings.}
The wall tapestries are faded and have been splattered with blood, defacing the sacred and profane themes that once played across them. The rug is similarly dirty and ripped, though still soft. The desk and chair are rudimentary, having neither arms nor drawers.
\subsection{6. Sanctuary Stairs}
\begin{DndReadAloud}{}
These taller-than-average stairs were made for creatures twice or three times the size of a human. They ascend to the sanctuary above.
\end{DndReadAloud}

A pair of Eldritch Priest* guard these stairs from inside the sanctuary. If they see fighting in area 1, they warn the bloodbonded and then go assist their allies two rounds later.
\paragraph{Large Stairs.}
These extra-tall, extra-wide stairs are almost too large to ascend without climbing. They are treated as difficult terrain for any creature smaller than Huge.
\subsection{7. The Sanctuary}
\begin{DndReadAloud}{}
This enormous space is a perfectly crafted, artificial primordial forest. Lit only by torch set in sconces at intervals along the walls, the lights flicker and dance across the smooth walls and ceiling. An array of columns, symmetrical but staggered, forego the entasis of the exterior columns in favor of attenuation that leaves large bases and shrinks as they ascend. As they approach the ceiling, the columns branch into funnel shapes like great abstract trees. The ruined rubble of a toppled colossal statue that once stood in the chamber's center fades away into the darkness in the northeast.
\end{DndReadAloud}

\begin{figure}[h]
\includegraphics[width=0.5\textwidth]{images/adventure/GHLoE/025-04-002.bloodbonded.png}
\end{figure}
The bloodbonded* and its six closest followers (Commoner), as well as two other sycophants (Cultist), are gathered in this space, creating a summoning circle for the next conjuring attempt.
\paragraph{Profane Summoning Circle.}
A complicated set of artistic shapes are inscribed in blood on the floor in a 15-foot-diameter area. Characters who succeed on a DC 14 Intelligence (Arcana) check recognize the general arrangement and layout of the shapes as a circle of daemon summoning. However, this one has druidic symbols merged with it. Stepping voluntarily onto this circle without having dealt damage to another creature in the last round causes a character to take 7 (2d6) necrotic damage unless it succeeds on a DC 13 Charisma saving throw. This effect can trigger once per creature per round.
\paragraph{North-Side Altar.}
This altar has been horribly defaced and smashed with a large heavy object. Only the top of the altar remains intact. A copper chalice with gold filigree rests atop it, filled with dried blood. The chalice (cleaned of blood) is worth 25 gp.
\paragraph{Statue Ruins.}
It is difficult to determine what this statue looked like before its fall. A successful DC 15 Intelligence (History) check determines that this statue depicted a four-winged seraphic figure wearing a crown and holding a scythe. Characters proficient with mason's tools have advantage on this check.
\paragraph{East Main Altar.}
This altar has been horribly defaced with a caustic agent. It was once lovingly painted but is now pocked and discolored. A silver ewer filled with acid rests on the altar. The acid from the ewer acts as using a Acid (vial) in combat. If the acid is cleaned out, the ewer is worth 30 gp.
\paragraph{Sacred Summoning Circle.}
A set of overlapping geometric shapes are roughly chiseled into the floor. The shapes are obscured by four dismembered body parts and the blood that accompanies them: two hands and two feet are spaced as if the body that connected them has been removed. The shapes cover a 15-foot-diameter area. Characters who succeed at a DC 14 Intelligence (Religion) check recognize the general shape and layout of the shapes as a circle of seraphic summoning. However, the body parts and blood would not normally be part of such a ritual. Stepping voluntarily onto this circle without having dealt damage to another creature in the last round causes a character to take 7 (2d6) radiant damage unless it succeeds on a DC 13 Charisma saving throw. This effect can trigger once per creature per round.
\paragraph{South-Side Altar.}
This altar has been left entirely intact but is frustratingly bereft of art or pigmentation that would help divine its purpose or dedication. A pair of cloth-of-gold vestments have been folded meticulously on the altar, each worth 20 gp.
\paragraph{Cages.}
Depending on your comfort level and the comfort level of your players, these cages can be stocked with animals the cultists have caught, by dead prisoners they've recently killed, by living hostages, by living voluntary sacrifices, or by living involuntary sacrifices. There are four cages, each one large enough to comfortably contain a Medium-sized creature.
\section{Conclusion}
Killing the bloodbonded is the only way to guarantee that its cult does not survive to trouble civilization again. Chasing it off is a half-measure. If the characters have been contracted to rescue captives (or if they merely witnessed a kidnapping and set out to track the abductors), freeing the prisoners in the cages and getting them to safety qualifies as success as well. If the prisoners are killed (perhaps in a failed hostage negotiation), that reflects poorly on the characters. Similarly, the bloodbonded's closest followers (the commoners) are brainwashed. To earn the most respect from their employers, the characters should allow as few commoners as possible to perish defending their monstrous leader.
Depending on which path led the characters to the ruins, this can start several other adventures: perhaps the temple was built over a dungeon even more ancient than the ruins, and some of the cultists flee into its depths with a prisoner. Perhaps this bloodbonded, with its dying breath, reveals that its plans are already in motion and its death won't stop them—what could those plans be? Other, more powerful bloodbonded may still be out there, and tracking them down to save a church or temple the characters value could also be a follow-up!
\DndQuote%
{}
{''Killing the bloodbonded is the only way to guarantee that its cult does not survive to trouble civilization again. Chasing it off is a half-measure.''}
{}
\section{Creatures}

\begin{itemize}
\begin{multicols}{3}
\item Bloodbonded
\item Eldritch Priest
\item Eldritch Herald
\end{multicols}

\end{itemize}

\chapter{5: Gloomrock Caverns}
By Tom "Evhelm" Donovan
Gloomrock Caverns is a dawndrinker lair for four or five 5th-level characters.
\section{Background}
This cave system houses a natural underground ecosystem. If not for the dawndrinker's presence, it's unlikely the characters would have found their way into it: there is no long-forgotten dungeon, no veins of precious metal, no dark and evil civilization below it—unless you want there to be for a future adventure!
The dawndrinker lairs in the deepest part of the cavern, but several monstrous creatures have long resided in the caves, hunting the rats, fish, worms, and sightless creatures that feed on the fungi around the waterways. The only ones who even notice the difference caused by the dawndrinker's presence are the grimlocks. Before their night foraging was restricted, now they have discovered freedom from the sun's heat and easier hunting in its absence—their blindness is no longer a curse but a blessing.
\subsection{Set the Hook}
This dawndrinker has had months to (fairly literally) soak up the sun. The dawndrinker's aura of advanced darkness is in full effect. For a mile in any given direction, the sun is as pale as a cloud-shrouded moon, and, beneath even the slightest shade, total darkness reigns. Trees and plants have begun to wither, and docile animals have fled the area in droves, driven away by the persistent darkness. The characters have either been hired to determine the source of this darkness and restore the light, or have undertaken this quest on their own.
\section{Lair Overview}
Carefully read over the dawndrinker's abilities, as it establishes the light-source rules for the lair. Notably, darkvision does not function in the Gloomrock Caverns! Because the area is in perpetual night, other than the foraging habits of the grimlocks (which are cyclic and focused on the underwater lake in areas 2 and area 3), the creatures' lives are almost static.
\subsection{1. Caverns Entrance}
\begin{DndReadAloud}{}
After hundreds of feet of trudging downward in the dark and uniform tunnel, the cavern opens a bit. The height of the ceiling triples and the stale smell of raw fish and manure interrupts your journey.
\end{DndReadAloud}

Unless the characters are extremely careful and stealthy the entire length of the corridor, the two Grimlock, who have climbed up from their island to eat some captured fish, are aware that there are intruders. They attempt to hide and, if spotted, dive down into area 2 to fetch the other grimlocks as reinforcements. Smart characters would normally find this as easy as shooting fish in a barrel, but in the advanced darkness it may not be so easy!
\paragraph{Rough Ascent.}
Hidden 25 feet up on the northeast wall is a large opening to another tunnel (area 4). From ground level, it is difficult to see. Characters who succeed on a DC 15 Wisdom (Perception) check (and have light sufficient to see it) notice it. The stones on the wall are jagged and easy to climb unless the characters are in combat. In combat, characters must succeed on a DC 10 Strength (Athletics) check to advance up the wall.
\paragraph{Smooth Descent.}
Characters without sufficient light sources who are moving too quickly may pitch themselves over the 25-foot drop and into the water below. The wall here is wet, requiring a DC 15 Strength (Athletics) check to successfully scale without the help of a rope. Characters who climb up from the water in area 2 and reach the top automatically notice the handholds the grimlocks made leading to area 4.
\DndQuote%
{}
{''The height of the ceiling triples and the stale smell of raw fish and manure interrupts your journey.''}
{}
\subsection{2. Grimlock Island}
\begin{DndReadAloud}{}
A small island rises out of the cool, languid water. To the west, the cave moves sharply up. To the east, it narrows and follows the water down into another cavern.
\end{DndReadAloud}

In addition to the two in area 1, six more Grimlock fish and swim on and around the island here. If they fight and begin to lose, they make a break to the east, hoping to lure their foes into the path of the sea drakes that live in the pool in area 3.
\paragraph{Garbage Island.}
The small island is covered in the detritus of sightless living. Beneath the filth and bones are the remains of a spelunker the grimlocks killed, who carried a climber's kit and Silk Rope (50 feet).
\subsection{3. Drake Lake}
\begin{DndReadAloud}{}
The water is eerily still and frigid here. The only way in or out of this domed chamber is through the waterway back to the small island.
\end{DndReadAloud}

Four Sea Drake* make their home in this tiny freshwater lake. They spar with the grimlocks, occasionally catching them and just as often getting caught themselves. They attack any humanoids (including the characters and grimlocks) who enter their lair. The sea drakes pursue enemies back to the small island but ignore targets that escape onto land or who climb the wall back to area 1.
\paragraph{Buried Treasure.}
One previous, tentative expedition to deal with the dawndrinker was fated to end here. Two skeletons sit at the bottom of the lake, twenty feet below the surface. While most of their gear is waterlogged, rusted, and worthless, a pristine +1 flail sits with the bones.
\subsection{4. Slow Rise}
\begin{DndReadAloud}{}
There is a 20-foot climb down behind you. Ahead of you, a passage opens and meanders away into the darkness.
\end{DndReadAloud}

The walls here are high compared to the rest of the cavern—likely beyond the range of any light sources the characters are carrying (though they may get creative). The slope rises five feet by the time it reaches the base of area 5, then another five feet as it encounters the first of the stalagmites and stalactites in area 6.
Unless the characters are entirely silent, they draw the attention of the horror flits in area 5 when they pass by. Those with a passive Perception of 15 or higher (or who succeed at an active DC 15 Wisdom (Perception) check) note that bat guano, a common occurrence in caves, is thicker and more common here than it has been in the rest of the cavern.
\paragraph{Scale the Wall.}
Characters who want to scale the 40-foot-tall wall to get to area 5 can do so with successful DC 13 Strength (Athletics) checks.
\subsection{5. Horror Flit Nest}
\begin{DndReadAloud}{}
The bones of countless fish and small game are nested among clear signs of larger kills.
\end{DndReadAloud}

Three Horror Flit Hunter* here have gotten by on the natives of the caverns, including poaching the occasional explorer and grimlock, after they feasted on bats that once made their home in the caverns. The dead bats are the source of the guano in area 4.
\paragraph{Treasure.}
Characters who seek out the horror flits' nest discover several spiked bone clubs (from grimlock meals) along with a belt pouch containing a handful of pearls (worth 250 gp total), assorted loose coins totaling 82 sp and 57 cp, and a spell scroll of burning hands.
\begin{figure}[h]
\includegraphics[width=0.5\textwidth]{images/adventure/GHLoE/030-05-001.horror-flits-r.png}
\end{figure}
\subsection{6. Unsteady Stalagmites}
\begin{DndReadAloud}{}
The passage opens wider here, but the ceiling drops to a comparatively claustrophobic ten or twelve feet overhead. Stalagmites and stalactites make this stone forest resemble an enormous gaping, crooked-toothed maw. The steady plip of water dripping from the ceiling and splashing on the floor is nearly constant in this damp cavern, muffled only by the plentiful fungi that cover much of the floor and stalagmites.
\end{DndReadAloud}

There are no monstrous dangers here, but a significant section of the floor that crosses over the waterway in areas 2 and 3 has grown unstable from all the water draining through the floor. Creatures with blindsight automatically detect the weakened floor and can go around it. Hugging the west wall of the tunnel (roughly a 10-foot walkway in the section over the water below) is safe and stable. The rest of the passage in the section over the water below is fragile. Characters with a passive Perception of 16 or higher notice this just before stepping on it; otherwise, succeeding on a DC 15 Intelligence (Nature) or DC 20 Wisdom (Perception) check is needed to notice the weak floor.
Beyond the faulty floor, the passage narrows in width, but the height of the passage rises to more than thirty feet. It is here that the dawndrinker's abilities begin to extinguish magical light sources. If the characters still have a way to perceive it, they notice the passage branching off and leading to area 8. (See area 8 for access details.)
\paragraph{Fungi.}
The characters may be wary of the fungi, but it's harmless and edible. When the dawndrinker is not around, they are luminescent, but their light is one of the first things the dawndrinker consumed! A successful DC 10 Intelligence (Nature) check tells a character that this fungi should be glowing, but isn't, for some reason that must be magical.
\paragraph{Floor Collapse.}
If two or more characters walk over the faulty section, it collapses. Characters who succeed at a DC 18 Dexterity saving throw grab the stable sides as the floor gives way and stay in area 6. Those who fail take 7 (2d6) bludgeoning damage from the floor debris as it crashes into the water around them. (These characters fall 30 feet into the water below, putting them between areas 2 and 3.)
\subsection{7. Dawndrinker Glimmercave}
\begin{DndReadAloud}{}
The sounds of water trickling into the cavern echo faintly and persistently, giving you the impression of a grand space.
\end{DndReadAloud}

The steep cliff from area 6 requires successful DC 15 Strength (Athletics) checks to navigate. A small pool of water has collected on the west wall. The passage to area 8 is gradual.
\begin{figure}[h]
\includegraphics[width=0.5\textwidth]{images/adventure/GHLoE/031-05-002.dawndrinker.png}
\end{figure}
The dawndrinker lairs here. Its blindsight means that it's aware of characters whether they fly, climb, or fall into the cavern. It uses hit-and-run tactics to claw each one to death as quickly as possible, prioritizing those with light sources first.
\paragraph{Pool.}
The pool of water is not deep—only a few feet—but it's large enough that smart characters may retreat into it to both slow the dawndrinker's movement (wading counts as difficult terrain) and to hear the dawndrinker's approach before it strikes.
\paragraph{Retreat Up.}
Characters may think to climb away from the dawndrinker. In such a case, it uses the passage to area 8 to get the jump on the characters from behind.
\paragraph{Retreat East.}
Characters may discover that the passage to area 8 narrows considerably—and that retreating into it might force the dawndrinker to attack head on. Depending on the characters' abilities, this may be a useful way to contain the dawndrinker's strengths and force it to face only the character with the best AC at any given time!
\subsection{8. Sloped Tunnel}
\begin{DndReadAloud}{}
This narrow tunnel is sloped steeply.
\end{DndReadAloud}

At the crook of the passage is the remains of the only other adventurer to get this far in search of the dawndrinker: a dragonborn paladin of Miklas. Her sword broke in the fight, and she tried to retreat here, but the dawndrinker followed and slew her.
Climbing up to this passage from area 6 (or back down to area 6 from this passage) requires succeeding at a DC 12 Strength (Athletics) check.
\paragraph{Miklas's Gear.}
The dragonborn corpse wears its suit of Plate Armor armor and shield. Her rations are spoiled, but her backpack still holds 18 gp, 3 moonstones (each worth 25 gp) and a spell scroll of daylight. There are also three empty vials, as well as an intact one holding a potion of greater healing.
\DndQuote%
{}
{''The dragonborn corpse wears its suit of plate mail armor and shield.''}
{}
\section{Conclusion}
Killing the dawndrinker allows all light sources to return to their original, normal strength. Failing to penetrate deeply enough into the lair to face the dawndrinker or allowing it to escape causes the darkness to persist and more of the region to fall into perpetual twilight.
Who sent this abomination to the region? Is an evil wizard loosing them on the world? Did something otherworldly or extraplanar crawl up from the bowels of the earth? By the time the characters return to civilization, a villain may have already taken credit for darkening the skies themselves as a show of its power and reach. If your characters are adventuring in Castinellan territory, perhaps an inquisitor will seek them out to see what they know of this arcane abomination. If they're adventuring in the Ostoyan Empire, perhaps the villagers whisper that this must be a creature that escaped The City Below.
\section{Creatures}

\begin{itemize}
\begin{multicols}{3}
\item Dawndrinker
\item Horror Flit Hunter
\item Sea Drake
\end{multicols}

\end{itemize}

\chapter{6: Boulderwood Path}
By Tom "Evhelm" Donovan
Boulderwood Path is a lupilisk lair for four or five 6th-level characters.
\section{Background}
The Boulderwood was named for its large stone outcroppings, but now that name takes on a more sinister tone for the stony fate that befalls travelers on its dangerous paths. A pack of lupilisks recently moved into this dense forest, stopping trade and cutting off communication between realms on either side.
Recently, a very important set of three signed treaties were being sent by disguised courier from one regent to another. These signed documents refreshed a century-old peace treaty, and if they are not delivered soon, these regents may find themselves at war.
The courier was sent through the Boulderwood and has not returned; it is feared she was attacked. Now someone must go and retrieve the documents (and deliver them to the eastern regent) before war breaks out.
\DndQuote%
{}
{''Recently, a very important set of three signed treaties were being sent by disguised courier from one regent to another.''}
{}
\subsection{Set the Hook}
A desperate western regent (or her agent, perhaps a captain of the guard or trusted advisor) tells the characters a rough outline of the information laid out in the Background. The characters' primary mission is to recover the signed and sealed documents (if possible) and the secondary mission is to determine what happened to the courier. Is the courier a traitor? Is there foul play involved? Is this just an accident?
Alternately, the characters may be traveling through the Boulderwood on another quest altogether; perhaps they are hunting lupilisks to make Periapt of Proof against Poison, or maybe this is just the fastest way to get where they're going, and they stumble onto the lupilisks and the remains of the courier's disastrous mission.
\section{Lair Overview}
The characters arrive on the road from the west at area 1. This deep in the Boulderwood, the tree canopy is very thick. The rocky outcroppings here have led to fewer taller, stronger trees, so the canopy is still all encompassing but light does pass through.
The lupilisks stumbled across a great bounty with the attack on the courier. Between the horses, the courier, and the courier's guards, the entire pack has been eating well for days, supplemented with the occasional stray forest creature. As such, most of the lupilisks are lazing in their respective nests and hangouts.
The eye crows in area 6 have peers all around the lair, watching to ensure the treaty pieces are not recovered. Characters with a passive Perception of 17 or higher notice the unusual number of birds watching them as they travel the lair, as do any characters who actively search the trees and succeed on a DC 17 Wisdom (Perception) check.
When the characters have investigated any two of the areas and are on their way to a third, a pack of gnolls (four Beast Gnoll*, two Venomous Gnoll*, and two Gnoll Brute*) arrives at area 3, traveling the road toward area 1. They make a ruckus, destroying the ruined cart and devouring the rations and supplies the forest animals haven't gotten into unless stopped by the characters. The gnolls carry their plunder from other travelers and adventurers they've killed. The treasure they carry includes 147 gp of random items and coins, as well as a bronze coronet with the word "blessed" inscribed in elegant Celestial worth 250 gp.
\subsection{1. Forest Path}
\begin{DndReadAloud}{}
After hours of trudging down this dark forest path, you find yourself staring down at the statue of a horse that's been tipped over onto its side in the middle of the road. Strangely, it's missing a leg and tail, but the missing pieces aren't nearby.
\end{DndReadAloud}

One of the two cart horses from the courier's wagon was turned to stone when the lupilisks attacked. They ate its foreleg and tail, leaving the rest for later. The rest of the wagon, along with the other horse, courier, and two guards made it to area 3 before succumbing.
\paragraph{Stone Horse.}
Characters who investigate the horse with a successful DC 12 Intelligence (Investigation) check discover that the statue is of impeccable workmanship, displaying the animal rearing back. It even has bridles and a broken yoke carved with symbols of the western regent. Those who suspect the horse was turned to stone (rather than being an artificial statue) can attempt a DC 15 Intelligence (Arcana). A success recalls stories of fantastic fiends called Lupilisk whose venom—which grows stronger as the creature ages—petrifies its prey. Characters who succeed by 5 or more also recall that lupilisks tend to hunt in packs and are often led by the oldest and most powerful among them.
\paragraph{Path.}
Characters who succeed at a DC 15 Wisdom (Survival) check can pick out a set of wagon wheel tracks that indicate the wagon picked up speed here. Those who succeed by 5 or more know that the speed increase was uncontrolled and unsafe.
\subsection{2. Watering Hole}
\begin{DndReadAloud}{}
Past the rocks, the air here is slightly warmer thanks to the break in the trees afforded by the pool in this swale.
\end{DndReadAloud}

Six Lupilisk Whelp*, a lupilisk*, and a lupilisk elder* are here: three whelps are curled up napping while the others are drinking at the pool. The lupilisk elder was injured in the fight against the courier's group and is at half its hit points. Even with the elder's weakened condition, this can be a very difficult encounter!
\paragraph{Rocky Area.}
The raised area that separates the pool's surroundings from the forest path is difficult terrain. Characters can use this to their advantage to slow the lupilisks!
\paragraph{Pool.}
The water is cool and clean to drink. It is 10 feet deep at its deepest point.
\paragraph{Treasure.}
The napping whelps had just finished a stony meal. A stone hand and a metal scroll case (one of the three pieces of the treaty) are all that remain of this courier guard.
\begin{figure}[h]
\includegraphics[width=0.5\textwidth]{images/adventure/GHLoE/037-06-006.lupilisk-elder.png}
\end{figure}
\subsection{3. Ruined Wagon}
\begin{DndReadAloud}{}
A capsized wagon with a broken wheel blocks the path.
\end{DndReadAloud}

A cursory search confirms that this is the missing courier's wagon, bearing the emblems of a trading house instead of the regent's seal to avoid suspicion. Flight of any kind is a great way to stay out of reach of the lupilisks; however, staying airborne or sticking to the treetops for too long draws the attention of two Sky Drake*.
\paragraph{Wagon.}
To establish the courier's cover, she was given several crates filled with jars of flour. The wreck of the wagon has broken most of the jars, and the crates are leaking flour. A successful DC 14 Wisdom (Survival or Perception) check reveals that the wagon capsized after being chased by at least ten quadrupedal creatures. A successful DC 14 Intelligence (Investigation) check reveals that some of the rocks on the road are not rocks at all, but pieces of another horse statue. Succeeding at either of these checks, or spending 10 minutes searching, reveals no lockbox or other suitable safe place for the treaty; the treaty pieces are not here and were likely carried by the courier or her guards.
\subsection{4. Elder Mount}
\begin{DndReadAloud}{}
From this high vantage point, you can see down the path to where you first encountered the stone horse and began your investigation. A rocky projection juts up from this hillock, blocking your view of the rest of the path.
\end{DndReadAloud}

A lupilisk elder* and two Lupilisk Whelp* are resting here and attack the characters on sight. Loud fighting here draws the creatures from area 5 into the fight after 1d3 + 1 rounds.
\paragraph{Climbing.}
Climbing or descending the gentle slope on the eastern side of the hillock does not require a check but is considered difficult terrain. To ascend or descend the steeper three sides of the hillock requires a successful DC 15 Strength (Athletics) check. The hillock is 30 feet above the forest floor.
\paragraph{Treasure.}
Near a petrified boot lies a necklace (an amulet of proof against detection and location) and the second piece of the treaty in a scroll case. This is the final resting place of the courier.
\subsection{5. Hillside Nest}
\begin{DndReadAloud}{}
The detritus of a large animal haunch is nestled under a large tree in the lee of this hillock.
\end{DndReadAloud}

This is the original and main nest of the lupilisks in the area: two Lupilisk* and two Lupilisk Whelp* are here. All four defend the nest to the death. Loud fighting here draws the creatures from area 4 into the fight after 1d3 + 1 rounds.
\paragraph{Treasure.}
Searching the nest for treasure yields the third metal scroll case containing the third part of the treaty. A successful DC 16 Intelligence (Investigation) check also yields a ring of jumping that was hidden in the filth.
\subsection{6. Ancient Altar}
\begin{DndReadAloud}{}
Atop this knoll rests a large worked-stone altar covered in ivy and moss. The unnerving sound of whispering seems to be all around you.
\end{DndReadAloud}

\begin{figure}[h]
\includegraphics[width=0.5\textwidth]{images/adventure/GHLoE/038-06-001.eye-crow.png}
\end{figure}
The altar's original purpose is long lost, and the gods who were worshipped here long gone. However, it is a beacon to a murder of Eye Crow that watches the characters. If the characters come here before finding any of the treaty scroll cases, the dozens of eye crows are in the trees all around the altar when they arrive, whispering. A DC 10 Wisdom (Insight) check determines that the birds are the source of the whispering sound.
\paragraph{Climbing.}
Climbing or descending the gentle slope on the western side of the knoll does not require a check but is considered difficult terrain. To ascend or descend the steeper southeast side of the hillock requires a successful DC 15 Strength (Athletics) check. The knoll's crest is 30 feet above the forest floor.
\paragraph{Eyes on the Altar.}
If the characters arrive after collecting any or all the treaty scroll cases, a group of 12 Eye Crow* fly onto the altar and stare at the characters. Characters who approach the 12 swear they hear the word "meddler" being whispered by the birds—but each character hears it in their native tongue. Attacking or moving threateningly towards the eye crows causes them to scatter to the trees (or further if pursued). However, if only a single character is present at the altar, the 12 attack, trying to pluck out the characters' eyes.
Characters who can talk to animals find that these crows have little to say to them; if compelled, the crows can reveal that they're watching to see if anyone takes the "metal tubes" (scroll cases) and who it is that takes them. They can describe a woman falconer who feeds them and trains them and to whom they deliver this information.
\section{Conclusion}
Characters who unseal and open the scroll cases find that the treaty was written, signed, and then cut vertically, meaning that only someone who has all three pieces can truly understand all the treaty's articulations and stipulations. Characters who put the treaty together and read it find that it is exactly what the western regent said it was: a renewal of a treaty between the western and eastern regions. It is up to the characters what to do with it. Taking it back to the western regent is a halfway measure but taking it to the eastern regent may be a gamble (after all, what if the eastern regent is responsible for the attack on the courier?). In any case, opening the seals should diminish (but not negate) the reward.
This quest ends with more questions than answers: were the gnolls directed to the courier's path, or were they merely a wandering pack? Were they after the characters? The treaty? Who do the eye crows ultimately report to, and why were they whispering "meddlers"? How did so many lupilisks take up residence so quickly in an area that had been cleared of dangers sufficiently for the courier to think this route was safe? Who doesn't want this peace treaty to be re-ratified?
\DndQuote%
{}
{''Who do the eye crows ultimately report to, and why were they whispering "meddlers"?''}
{}
\section{Creatures}

\begin{itemize}
\begin{multicols}{3}
\item Eye Crow
\item Beast Gnoll
\item Gnoll Brute
\item Venomous Gnoll
\item Lupilisk
\item Lupilisk Whelp
\item Lupilisk Elder
\item Sky Drake
\end{multicols}

\end{itemize}

\chapter{7: Flamegrit Iron Mine}
By Tom "Evhelm" Donovan
Flamegrit Iron Mine is a mjork lair for four or five 7th-level characters.
\section{Background}
In Stehlenwald province, the small city of Flamegrit is an iron mine that has a problem with elementals. Mjorks have been moving up through the lower levels of the mine, and now they've taken over. The dwarves have sealed off the mine, but they know it won't hold the mjorks forever: the mjorks seek to escape the mine's confines and spread their environment of ash and fire to the whole of Etharis.
\subsection{Set the Hook}
The easiest way for the characters to be involved is to be in Flamegrit when the mine is sealed; the dwarven council recruits them to clear the top level of the mine and rescue miners trapped there. Perhaps the dwarves have a magic item or resource the characters need, or perhaps the dwarven council's permission is needed to travel through their lands or to explore. Whatever the characters need, the dwarves have it and will give it to the characters if they defeat the mjorks who have shut down their mine.
If you remove the miner-survivors (area 4), it's also possible for the characters to stumble into the mine unwittingly while exploring a dungeon or other subterranean complex (though they have less of a reason to fight through the mjorks in this case). Just age the equipment and give the mine the feel of being abandoned rather than recently used—perhaps block off the area beyond the double doors in area 1 with a cave-in and have the characters enter from somewhere else.
\section{Lair Overview}
The mine is hot. While rocky surfaces won't cause damage, it is exceptionally warm, like the inside of a sauna. Prolonged exposure without a buffer or cooler causes characters to accrue levels of exhaustion at a rate of one level per eight hours. The air is stale and sooty, even in areas lacking mjorks. The northern area of the cavern is largely static, but if the characters delay too long in area 2, two Mjork Sootling Swarm* climb from the pit and harass them. The mjork charger* in area 7 is on patrol across the southern area. Exceptionally loud noise (as from a thunderwave or knock spell) draws the charger to almost any location in the mine in 2d4 rounds.
\subsection{1. Entrance}
\begin{DndReadAloud}{}
Beyond the doors, the air is thick, a fog cloud of ash and soot making it difficult to see more than a few feet. The doors close behind you.
\end{DndReadAloud}

A mjork asher* is here, razing the supplies in the east corner after being unable to get through the doors out of the mine. It attacks immediately. The supplies are ruined slag. The doors have been barred from the other side; characters must communicate with the guards to get them to open the doors.
\paragraph{Asher Fog.}
The mjork asher's smoke shroud is in effect, emanating from its starting position near the supplies.
\subsection{2. North Side}
\begin{DndReadAloud}{}
The mine's central excavation stretches before you. Directly to the south, a cavernous pit with a stone bridge leading up to the south side. To the west, a cave. To the east, supplies, stacks of raw ore, and a large mining elevator.
\end{DndReadAloud}

The ore is iron, brought up on the elevator before the mjorks destroyed it.
\paragraph{Pit.}
The central scar is between fifty and sixty feet deep, leading to the next level of the mine. Due to the mjorks' interference, the floor of the lower level is extremely hot; characters who contact it directly take 2 (1d4) fire damage per round. The branching pathways have been caved in, allowing only Tiny creatures to squeeze through. The ruined mining scaffold on the western side of the pit and the collapsed elevator on the eastern side can be climbed with a successful DC 10 Strength (Athletics) check, or the sides of the pit walls can be scaled with a successful DC 15 Strength (Athletics) check. Characters who spend more than a minute on the bottom level are attacked by three Mjork Sootling Swarm* that emerge from the side tunnels.
\subsection{3. Smoking Crates}
\begin{DndReadAloud}{}
Dozens of crates marked with symbols for climbing gear and mining equipment are emitting tendrils of smoke. Many appear scorched.
\end{DndReadAloud}

Three Mjork Sootling* hide among the crates, destroying their contents. Patient observation of the crates (five or more minutes) reveals them, as does a successful DC 13 Wisdom (Perception) or Intelligence (Investigation) check. If they survive two rounds of combat with the characters, they realize they won't win and flee toward area 7 for help.
\paragraph{Crates.}
The supplies are mostly destroyed. Characters who spend at least ten minutes salvaging from the crates find all the items in a dungeoneer's pack minus the rations and waterskin.
\begin{figure}[h]
\includegraphics[width=0.5\textwidth]{images/adventure/GHLoE/045-07-006.branches.png}
\end{figure}
\subsection{4. Trapped Miners}
\begin{DndReadAloud}{}
This space has been recently but precisely excavated, forming a narrow tunnel. The muffled sound of hammering is offset by the roar of a fire.
\end{DndReadAloud}

\begin{figure}[h]
\includegraphics[width=0.5\textwidth]{images/adventure/GHLoE/046-07-008.the-mjork-burner.png}
\end{figure}
Four miners (dwarven Commoner) were cut off from the entrance by the mjorks and retreated into this unfinished excavation tunnel. They managed to throw up a make-shift barricade, but a mjork burner* has discovered them and is slowly smashing its way through the barricade. Fighting here draws weak cries and coughs for help from the trapped miners.
\paragraph{Barricade.}
If the characters don't interfere, the mjork burner breaks through the barricade 6 rounds after the characters see it. It proceeds to slaughter the miners. If the mjork is defeated, the miners happily help the characters disassemble the barricade.
\paragraph{Miners.}
The miners are siblings from Clan Izim (Drik, Klup, Gred, and Trilk) and are extremely thankful to be saved. If the characters successfully save all four from the mjorks and escort them safely back to area 1, they promise a boon of 50 gp from their clan treasury. Drik also gives the characters a pretty emerald she recovered while mining; unbeknownst to Drik, it is an Elemental Gem, Emerald. The elemental gem and 12 sp can be recovered from the miners' corpses if the characters allow the burner to finish its grim task. The miners willingly share what they know of the mine: they can identify the three paths from the north to south areas, the bunkroom, the general store, and can confirm that there are no survivors below this level of the mine.
\DndQuote%
{}
{''The muffled sound of hammering is offset by the roar of a fire.''}
{}
\subsection{5. West Cave}
\begin{figure}[h]
\includegraphics[width=0.5\textwidth]{images/adventure/GHLoE/047-07-001.torcheater.png}
\end{figure}
\begin{DndReadAloud}{}
This cool cave ascends sharply from north to south.
\end{DndReadAloud}

This cave slopes sharply up 25 feet from the north entrance to the south entrance. The slope is considered difficult terrain.
\paragraph{Torcheaters.}
A flight of 6 Torcheater* dwell in this cave, perched among the stalactites 15 feet overhead. They remain motionless and practically invisible unless the characters are carrying a fire-based light source. In that case, they swoop down to consume the light. They have no interest in fighting the characters unless attacked and flee if more than three of their number are slain.
\subsection{6. General Store}
\begin{DndReadAloud}{}
A cramped vestibule is lined with stone chairs. Beyond a stone countertop are rows of shelves, knocked over and singed.
\end{DndReadAloud}

Two Mjork Sootling Swarm* and 3 Mjork* are busy turning the general store into general slag. If interrupted in their destruction, they attack. Combat here that spills out of the store's entrance draws the attention of the mjork charger* in area 7 in 1d3 rounds.
\paragraph{Till.}
In a locked stone box behind the counter are 20 gp, 110 sp, and 255 cp. The box's lock can be opened with a successful DC 14 Dexterity (Thieves' Tools) check.
\paragraph{Supplies.}
As elsewhere, valuable supplies have been ruined. A successful DC 15 Intelligence (Investigation) check reveals a matching set of three silk handkerchiefs that escaped the destruction, each worth 15 gp.
\subsection{7. South Side}
\begin{DndReadAloud}{}
From this higher vantage point, the pit seems much deeper. A ruined mining scaffold creaks ominously to the west. A stone bridge ramps down to the north side of the cavern, and a multitude of cave entrances pocket the south wall.
\end{DndReadAloud}

A mjork charger* patrols the south side of the cavern. It takes it roughly five minutes to traverse the length of the cavern, stopping to spend about a minute at the entrance to each of four places: areas 6, 8, 10, and 11. It attacks any non-mjorks on sight and attempts to push enemies near the edge of the pit off into the depths. If it spots non-mjorks in any of its patrol areas, it moves to that new location to attack. Exceptionally loud noise (as from a thunderwave or knock spell) draw the charger to almost any location in the mine in 2d4 rounds.
\paragraph{Pit and Scaffold.}
See area 2 for details on the pit. Falling from the north side to the bottom is an 80-foot drop.
\subsection{8. Bunkroom}
\begin{DndReadAloud}{}
This worked stone chamber is filled with uncomfortable-looking stone bunkbeds and reeks of sweat and soot.
\end{DndReadAloud}

For miners who didn't want to waste time between long shifts, a series of stone slabs with minimal bedding was carved out. Four Mjork* have destroyed all the linens and are on their way to finishing the job with the footlockers. Fighting that spills out of this area draws the attention of the mjork charger* in area 7.
\paragraph{Footlockers.}
More apt for resting mining picks than safe keeping for belongings between shifts, there are nevertheless a few valuables in these containers. Searching all the footlockers yields a pair of engraved bone dice (25 gp), three emeralds (each worth 50 gp), and 47 sp.
\subsection{9. Abandoned Excavation}
\begin{DndReadAloud}{}
This narrow tunnel is partly excavated and partly natural. Rubble has sealed off what used to be a second entrance.
\end{DndReadAloud}

The dwarves gave up on this tunnel after their initial efforts collapsed. None of the mjorks are interested in it, making it a superb place to catch a short rest or to hide from pursuing mjorks.
\DndQuote%
{}
{''Four mjorks have destroyed all the linens...''}
{}
\subsection{10. East Cave}
\begin{DndReadAloud}{}
The steady drip of water echoes through this sloped cave, almost drowned out by the sizzle of liquid hitting a hot surface and the sounds of a roaring, hammering fire.
\end{DndReadAloud}

A mjork burner* is flinging itself at the water-coated walls here, trying to dry them out. Each body-slam stings the burner a bit and leads it to be enraged even further. The floor slopes up from the north to the south side at a steep angle, making the floor difficult terrain. Combat that spills out the south entrance to the cave automatically draws the attention of the mjork charger* in area 7.
\subsection{11. Stone Bridge}
\begin{DndReadAloud}{}
A sharply ascending stone bridge crosses the chasm.
\end{DndReadAloud}

Two Mjork Sootling Swarm* are underneath the bridge. Characters looking for danger beneath the bridge have advantage on seeing them with a successful DC 15 Wisdom (Perception) check. The second time a character crosses the bridge, the swarms climb out from below and attack. Combat here draws the attention of the mjork charger* in area 7.
\paragraph{Collapse.}
The swarms have been eating away the wooden support beams holding the bridge in place. While it is still strong enough for a single character to cross, it collapses if two or more characters cross the bridge at the same time. Characters can spot the damage with a successful DC 12 Intelligence (Investigation) check to determine the integrity of the bridge before crossing. It's a 60-foot drop off the bridge. If the bridge collapse is the reason for the fall, add 3 (1d6) additional bludgeoning damage from the debris to anyone who hits the pit floor.
\section{Conclusion}
Clearing the mjorks from this level of the mine satisfies the dwarven council and completes the characters' quest. The characters can be rewarded and move on, but questions remain unanswered: What lurks below the mine? What caused the mjorks to suddenly emerge from the depths and wreak havoc? Perhaps another nearby mine has also gone dark, and this is the beginning of something far more sinister than a random attack by elementals? While the characters were clearing the mine, perhaps the dwarves above found themselves embattled by another foe—were the mjorks just a convenient distraction? The dwarven council may have more missions for the characters, even if they're not exploring more subterranean spaces.
\section{Creatures}

\begin{itemize}
\begin{multicols}{3}
\item Mjork
\item Mjork Charger
\item Mjork Burner
\item Mjork Asher
\item Mjork Sootling
\item Mjork Sootling Swarm
\item Torcheater
\end{multicols}

\end{itemize}

\chapter{8: Shadowsteel Citadel}
By Tom "Evhelm" Donovan
Shadowsteel Citadel is a blightscale dragon lair for four or five 8th-level characters.
\section{Background}
Before the dragon's arrival, this mountain stronghold was an open secret. Built near a large settlement, readily supplied with sundries and patrons, it's far enough from civilization that the powers-that-be didn't need to be threatened by its presence. The citadel was devoted with religious fervor to one pursuit: curses. The priests and mages who built the citadel did so through the favors of the common people: curses in exchange for a year of service building, guarding, maintaining, or cooking for the stronghold. Of course, payment in Shadowsteel or sufficient gold was also acceptable, allowing the rich to buy their curses as the poor pledged their service.
Then Chuldroth, a young blightscale dragon, driven mad by his ruined body and knowing distantly of the citadel's work with curses (and blaming it for his condition), attacked it. After causing a rockslide at the entrance and trapping the inhabitants, he tore a hole in the mountainside and made his new lair in the ruins of their cathedral, killing the guardians sent to chase him away. Now only the undead dwell in the citadel, waiting for the dragon to move on to a new lair.
\subsection{Set the Hook}
There are three easy reasons the characters may be investigating the Shadowsteel Citadel. First, a powerful figure in the nearby city may have sought a curse, but found the entrance blocked—such a figure may hire the characters to investigate. Second, Chuldroth is wreaking havoc on the countryside; the characters may be tracking Chuldroth and followed him back to the citadel. Finally, the characters may stumble on the path leading to the citadel and investigate the caved-in entrance.
\section{Lair Overview}
The citadel was prepared for a siege (after all, in the curse business you never know who might come knocking) but was not ready for a dragon to drop into their cathedral. Rooms meant for the living are generally tidy though dusty, with torches ready to be lit and braziers filled with oil. Rooms built for the dead are tomblike; no sources of light are prepared. When the dragon is absent, the Shadowsteel ghasts occasionally leave their tomb to get materials from the rest of the complex for their fell experiments.
\subsection{1. Path and Cliffs}
\begin{DndReadAloud}{}
The path leads to a forbidding cliff face, once smooth but now choked with rubble from a collapse high above. Even at a distance, an opening between the rubble and the cliff is visible a few dozen feet up.
\end{DndReadAloud}

Climbing the rubble is easy but time consuming. The hole leading to the interior and area 2 is large enough for a Medium creature; Large creatures can fit by squeezing.
\subsection{2. Collapsed Entrance}
\begin{DndReadAloud}{}
This hemispheric room has three doors, each with inscriptions cut into their frames. Small piles of rubble from the collapse have been hastily stacked in two piles between the doors.
\end{DndReadAloud}

The soldiers trapped by the collapse made a short-lived effort to move the debris before realizing there was too much of it. Only the passage of time has caused the debris to settle enough for an opening to appear.
\paragraph{Three Doors.}
Each inscription is in highly stylized Castinellan. The south door's inscription reads: "Death and Curses." The north door reads: "Long Death." The middle door reads: "Curses Worse than Death." Opening the north or south doors activates the enchantments animating the undead in both spaces labeled area 3. The triggering runes on both doors can be discovered by a successful DC 15 Intelligence (Investigation) check. Dispel magic suppresses the rune for an hour, knock triggers the rune, and it can be deactivated permanently through a successful DC 16 Intelligence (Arcana or Thieves' Tools) check. The middle door is locked from this side but not trapped; it can be picked with a successful DC 14 Dexterity (Thieves' Tools) check.
\subsection{3. Undead Guards}
\begin{DndReadAloud}{}
This roughly square-shaped room with two doors is unfurnished.
\end{DndReadAloud}

In each of these rooms, five Skeleton Rifler* are lined up along the wall opposite the door to area 2, five more are lined up along the east wall, and a skeleton commander* stands at the ready in the corner between them. The riflers have their firearms trained on the door to area 2 but are immobile unless the rune on the door activates them (see area 2) or someone enters from areas 4 or 5 without speaking the password ("cursed steel" in Castellan). Once activated, the skeleton troopers defend themselves and the rest of the complex until destroyed, pursuing fleeing enemies and fighting to the death.
\paragraph{Locked Door.}
The door to area 5 is locked from both sides. The key is in the dragon's nest in area 13. The lock can be picked with a successful DC 16 Dexterity (Thieves' Tools) check.
\begin{figure}[h]
\includegraphics[width=0.5\textwidth]{images/adventure/GHLoE/054-08-006.skeleton-rifler-r.png}
\end{figure}
\subsection{4. Empty Room}
\begin{DndReadAloud}{}
This room is entirely empty except for the thick layer of dust on the floor. Two doors lead out.
\end{DndReadAloud}

\paragraph{Clawed Footprints.}
Characters with a passive Perception of 14 or higher note a single trail through the dust from door to door. A successful DC 10 Wisdom (Survival) check determines they were made by a single Medium-sized biped, while DC 15 reveals that the creature had clawed feet that walked from the west door to the east door and returned the same way.
\subsection{5. Bunkroom}
\begin{DndReadAloud}{}
This crowded bunkroom has eight sets of bunkbeds and footlockers; doors lead out to the east and west.
\end{DndReadAloud}

Half the beds are perfectly made, the other half are in disarray. A few open footlockers have had their contents shredded and strewn about the space. A sitri cat* naps in a bunk. It has been trapped here since the collapse and is starving after having eaten everything it could from the soldiers' footlockers.
\paragraph{Locked Door.}
The door to area 3 is locked. See area 3 for details.
\paragraph{Treasure.}
The indentured soldiers' belongings are in the footlockers. A search of the space finds 27 gp and a small mirror set into a painted wooden frame (worth 25 gp).
\DndQuote%
{}
{''This hemispheric room has three doors, each with inscriptions cut into their frames.''}
{}
\subsection{6. Guarded Negotiations Chamber}
\begin{DndReadAloud}{}
Two half-walls divide this room. In the west, comfortable chairs surround a pair of circular tables. Two of the chairs are knocked over. In the eastern part of the room, three doors lead out.
\end{DndReadAloud}

\paragraph{Secret Door.}
Characters may notice that the way the chairs fell looks like whoever knocked them over was running toward the west wall. Any investigation into why the chairs are tipped over grants advantage on subsequent investigations to find the western secret door leading to area 9. A successful DC 15 Wisdom (Perception) or Intelligence (Investigation) check discovers the door and the stone button to activate it.
\paragraph{Locked Door.}
The southern door is locked and leads to the armory. The key is in area 13. A successful DC 18 Dexterity (Thieves' Tools) check opens it.
\DndQuote%
{}
{''The few insects here are harmless, but their yellow, lumpy hives are massive.''}
{}
\subsection{7. Armory}
\begin{DndReadAloud}{}
This narrow armory is filled with assorted weapons and armor, gleaming on wooden racks.
\end{DndReadAloud}

\paragraph{Arms and Armor.}
Six Shortbow, eight Longsword, four Pike, and three sets of leather armor hang on racks. Two Quiver with 20 Arrows (20) each hang from hooks on the wall; five of the arrows in the first quiver are +1 Arrow.
\subsection{8. Dead in Waiting}
\begin{DndReadAloud}{}
Six stone sarcophagi, each depicting a featureless robed humanoid form, are arranged in two columns of three.
\end{DndReadAloud}

The builders of the citadel lie on these stone coffins, each one twisted by their excessive use of Shadowsteel. Four are currently here, reanimated as Shadowsteel Ghoul*, resting in their coffins until more shadowsteel is available. The northeastern and southwestern sarcophagi are empty. Any creatures passing through this room without speaking the password ("cursed steel" in Castellan) draw them out. They fight to the death and pursue fleeing enemies throughout the complex.
\paragraph{Treasure.}
One of the ghouls is wearing a cape of the mountebank.
\subsection{9. Dining Room}
\begin{DndReadAloud}{}
This area was once a dining hall with two large refectory tables and accompanying benches, now tipped on their sides. The stench of decay is overpowering, emanating from several yellow insect hives scattered around the room creating the soft buzz of insect wings. The northern door is open. The double doors to the west hang slightly ajar.
\end{DndReadAloud}

The few insects here are harmless, but their yellow, lumpy hives are massive.
\paragraph{Secret Door.}
A successful DC 15 Wisdom (Perception) or Intelligence (Investigation) check discovers the secret eastern door leading to area 6 and the stone button to activate it.
\subsection{10. Kitchen}
\begin{DndReadAloud}{}
The rough outline of workstations and preparatory tables with hanging pots, pans, and vegetable baskets is visible under an enormous yellow insect hive that covers most surfaces in the room. Two doors exit the room; the northern door is completely covered by the hive.
\end{DndReadAloud}

\begin{figure}[h]
\includegraphics[width=0.5\textwidth]{images/adventure/GHLoE/055-08-001.mold-spider.png}
\end{figure}
The insects here are harmless, but a colony of Mold Spider* has claimed the space as well. All the cooking aids have been corroded and rotted, making them useless.
\paragraph{Moldy Hitchhikers.}
If the characters disturb the hive at all (as by clearing the north door, for example), the insects leave their hives and envelop the characters. Disturbing as this is, they retreat after two rounds, leaving behind 3 Mold Spider* per character.
A meticulous character who picks off every insect eventually sees them, aggravating the spiders, causing all of them to attack the characters they've hitched a ride on. If no one notices immediately, they lie dormant until the characters enter combat or take a rest, chomping down on their fresh prey then.
\subsection{11. Storage}
\begin{DndReadAloud}{}
This narrow area is filled with crates, caskets, barrels, and chests marked with the names of foodstuffs and supplies.
\end{DndReadAloud}

The door to this chamber kept out the insects, preserving food and supplies.
\paragraph{Supplies.}
Characters willing to peer into every container find 40 days of iron rations, 100 feet of hempen rope, barrels of fresh water, and crates of dry baking ingredients. A successful DC 13 Intelligence (Investigation) check discovers one of the chests has a false bottom with 100 gp inside.
\subsection{12. Sacristy}
\begin{figure}[h]
\includegraphics[width=0.5\textwidth]{images/adventure/GHLoE/056-08-003.eye-crawler.png}
\end{figure}
\begin{DndReadAloud}{}
Wall-to-wall shelves and a central credence table reveal that this vesting room is untouched by the rot just outside its door.
\end{DndReadAloud}

An eye crawler* lurks among the books and jars on the table, observing.
\paragraph{Treasure.}
Three sets of black vestments inlaid with gold can be sold for 30 gp each. Most of the books are worth little, but a few tomes on the art of cursing and the magic of Shadowsteel would sell for a total of 50 gp to a collector.
\subsection{13. Ruined Cathedral}
\begin{DndReadAloud}{}
This massive space is lit by a hole smashed in the ceiling. The light is shifting and green, tinted by branches and leaves that make it difficult to see the sky. The north end of the room is covered in debris crushing a raised platform. Wooden pews are cracked and smashed, shoved to the sides. Pockets of smelly puss-filled insect hives litter the ground, and a nest of broken pews is in the southwest corner.
\end{DndReadAloud}

Chuldroth, a young blightscale dragon*, has made this his temporary lair. He is cunning, but just as desperate as any other blighted creature—far more likely to fly into a murderous rage than to reason calmly.
\paragraph{Pews and Rubble.}
The debris around the room makes the area difficult terrain.
\paragraph{Treasure.}
Scattered around the area are the remains of the people who made their last stand against the dragon, as well as the precious articles dedicated to the space. This includes the keys to doors in areas 3-5 and 6-8 along with two silver ewers (25 gp each), a potion of lightning resistance, a +1 shortbow, 638 sp, and 220 gp.
\DndQuote%
{}
{''An eye crawler lurks among the books and jars on the table, observing.''}
{}
\subsection{14. Undead Alchemy Lab}
\begin{DndReadAloud}{}
This area has another six sarcophagi that match those you saw earlier. On the eastern end of the area is a set of worktables with glass beakers, various jars, and shelves full of dried herbs.
\end{DndReadAloud}

Two Shadowsteel Ghast* work tirelessly to perfect their curses and magical defenses. The dragon's intrusion destroyed their support network, so they are slowly planning their escape. The sarcophagi here are empty. They fight to the death if forced to, but if given the chance to flee, they take as much gold as they can carry and use potions of levitate to escape through the hole in the roof of area 13.
\paragraph{Alchemy Lab.}
The materials here function as a fully-stocked alchemical lab. By spending an hour breaking down the supplies and taking only essentials, a character proficient in their use can create a complete set of Alchemist's Supplies. Two completed potions of levitate are already prepared.
\paragraph{Treasure Hoard.}
The ghasts have been slowly collecting everything of value from the complex when the dragon goes hunting and storing it in the eastern alcove. Along with a chest containing 316 gp, characters may find three silver chalices (30 gp each), four gold bracelets (45 gp each), and a platinum-studded velvet mask (100 gp), all thematically linked to curses and Shadowsteel.
\section{Conclusion}
Defeating the dragon ensures its attacks on the region stop and provide information as to why the citadel is no longer functioning. If the Shadowsteel ghasts escape, they may set up a new "cursed citadel" the characters must find. Similarly, someone was gathering information through the eye crawler in area 12: who wanted to keep an eye on things here? Also, perhaps someone who doesn't know the whole story blames the characters for the citadel's fall—they may be seeking to curse the characters now!
\DndQuote%
{}
{''Defeating the dragon ensures its attacks on the region stop and provide information as to why the citadel is no longer functioning.''}
{}
\section{Creatures}

\begin{itemize}
\begin{multicols}{3}
\item Young Blightscale Dragon
\item Eye Crawler
\item Mold Spider
\item Shadowsteel Ghoul
\item Shadowsteel Ghast
\item Sitri Cat
\item Skeleton Rifler
\item Skeleton Cannoneer
\item Skeleton Commander
\end{multicols}

\end{itemize}

\chapter{9: Storm Sanctuary}
By Tom "Evhelm" Donovan
Storm Sanctuary is an ithjar lair for four or five 9th-level characters.
\section{Background}
Long ago, worshippers built a temple and shrine to the gods, who they imagined in the form of dragons. The temple and shrines are long abandoned, but a small cult of panjaian ilharans found the ruins and turned them to their own purposes: luring and protecting stormborn ithjars. The ilharans have built a small worship space and created an enchanted playground for their darling ithjars to frolic in. They are about to achieve their grandest goal: the stormborn ithjar who roosts here has lain eggs. The ilharans have taken great pains to make sure that the roost is safe—even from curious members of their own sect—and are awaiting the hatchlings. Wild animals have taken over the nearby ruins, but through long months of confrontation, the ilharans stay out of the ruins and the beasts don't venture into the ilharans' sanctuary.
\subsection{Set the Hook}
The ilharans' plans are at odds with civilized society: hungry ithjar (both regular and stormborn) have been raiding roads and farms for decades, but lately they've grown bolder, wiping out villages and even attacking fortified settlements and castles. The characters have been hired by the governor or regent of the region to find the ithjar lair and wipe out the threat. Diviners have tracked the ithjar to the Storm Sanctuary, and now the characters must take them on. Alternatively, the characters may need an ithjar egg for a ritual, craft, or contract. Their efforts to locate an egg ultimately lead them to the Storm Sanctuary.
\section{Lair Overview}
The Storm Sanctuary is high in the forested mountains and surrounded by natural formations and artificial constructions; many partitions close off parts of the Storm Sanctuary from other areas. Any flight—whether magical or natural—draws the immediate attention of two ithjars, and combat in the sky risks bringing all the ithjars in area 6 and the stormborn ithjar from area 8. Panjaian ilharans* and beasts present in areas where aerial combat is occurring get involved as well. Similarly, climbing the cliffs and crags on the eastern half of the map can be done with a successful DC 12 Strength (Athletics) check. When a character reaches the zenith of a crag or cliff, the ithjars are alerted as if the character were flying.
The descriptions of each area assumes a daylight approach. At night, halve the number of panjaian ilharans and double the number of beasts in each area; the number of ithjar stays the same.
\subsection{1. Open Pagoda}
\begin{DndReadAloud}{}
A simple wooden pagoda, elegantly constructed, dominates the landscape here. Two rows of wooden benches, little more than felled trees, face the pagoda. Tall vertical wooden walls capped with sharpened tops flank the area to the north and east.
\end{DndReadAloud}

Six Panjaian Ilharan* await the stormborn ithjar's hatchlings here. One stands in the pagoda, giving what amounts to a sermon to her peers about the glories of the storm and the power of the plane of air. Her name is Arial (CN, desires all sentient life to acknowledge the superiority of Ilhara over the other Primordials). The other five are sitting and listening or inattentively daydreaming. Characters have advantage on checks to surprise these creatures.
\begin{figure}[h]
\includegraphics[width=0.5\textwidth]{images/adventure/GHLoE/065-09-005.panjaian-ilharan.png}
\end{figure}
\paragraph{Negotiating with the Ilharans.}
The panjaian stationed here are guardians of the ithjar. Respectful and knowledgeable characters may be able to convince the panjaian they're pilgrims come to adore the ithjar. The panjaians start as hostile but not so hostile they'll attack on sight. Arial is gullible and willing to believe that her own reputation has brought potential pilgrims; however, she has a hatred of elves. This can be overlooked unless an elf is the characters' spokesperson!
Through interaction and diplomacy, the characters may shift Arial and her peers' attitude to indifferent. Doing so avoids combat with the panjaians (for now) and even means they'll share an important secret: no one can see the ithjars now because the key to the gate that leads to their nests was lost in the ruins to the north (and flying or climbing enrages the ithjar).
Characters who manage to do particularly well in negotiations may even be able to shift Arial (if not her peers') to a friendly attitude. If friendly, Arial confirms that a stormborn ithjar has lain eggs. If the negotiations break down or the characters decide to attack first, the panjaians fight to the death.
\subsection{2. Air Elemental Statue}
\begin{DndReadAloud}{}
At the convergence of the two wooden walls, the ground is replaced by plate-sized natural stones polished and placed to appear as pavers in a plaza. At the center of this plaza is a moving statue, an aery elemental that is anchored to the ground but which wafts and flutters in the wind. To the north, ruined walls are all that remain of a once impressive structure. To the east, cliffs rise fifty feet or more above your head. The cliffs disappear beyond your view to the north and south.
\end{DndReadAloud}

The statue is a complicated illusion created by the panjaians. Over time, some of them have come to believe it is a real air elemental that has been trapped, but it's still just an illusion. The stone it is anchored to radiates faint illusion magic when subjected to detect magic.
\paragraph{Beneath the Statue.}
Casting dispel magic on the statue suppresses it for one hour. When the statue is suppressed, it becomes clear the stone it's anchored to is really the lid to a chest! A search of the base with a successful DC 18 Intelligence (Investigation) check also reveals this to be true. In either case, once the truth about the "stone" is realized, it's a simple matter to pry it open. Inside are two moonstones (worth 50 gp each) and four Potion of Greater Healing.
\subsection{3. Ruined Hall}
\begin{DndReadAloud}{}
Up some crumbling stairs is a jagged hall. Its roof has long since rotted away and only a few stone columns remain to provide a sense that the walls were once two or even three times higher than their ruined remains would suggest. A deep growl is quickly joined by several others from the west!
\end{DndReadAloud}

\begin{figure}[h]
\includegraphics[width=0.5\textwidth]{images/adventure/GHLoE/066-09-001.sloth-galloper.png}
\end{figure}
The dragon-god worshippers who built the original temple were replaced by druids. These potent spellcasters put curses on the dragon-god temple to make it an attractive dwelling place for beasts. Eight Corpsejaw*, eight Shieldhead*, and four Sloth Galloper* arrive over three rounds. In the first round, half of the animals arrive: corpsejaws arrive from the east and south, shieldheads smash through the ruined walls to the north and west, and the sloth gallopers climb over the northwest and northeast corners. In the second round, a quarter of the animals arrive in the same manner. In the third round, the last quarter of the animals arrive.
\paragraph{Beast Behavior.}
Whenever a humanoid creature enters, if it is not a druid, the beasts are aggressively drawn from the nearby woods to attack it after one round has passed. Any characters with at least one level in druid are completely ignored by the attacking animals unless they attack the animals or aid the animals' enemies. A group entirely composed of druids still draws the animals in, but they merely circle the characters, sniffing curiously before leaving the way they came.
These animals are bloodthirsty and enraged. Animal Handling checks made to deal with them have disadvantage, but calm emotions or similar magic breaks the spell; "calmed" creatures wander off confused if they are not attacked. The beasts fight to the death and pursue fleeing characters.
The magical lure only functions once every 24 hours. Once the characters have dealt with the beasts, they are safe until the next day.
\subsection{4. Ruined Study}
\begin{DndReadAloud}{}
The purpose of this room is unclear: mold and choked trees grow amid the debris of the ruin, but no furniture remains. In the southeast corner, concealed under a fallen chunk of stone, a half-eaten corpse is rotting.
\end{DndReadAloud}

The corpse belongs to an adventurous panjaian who mistakenly thought that the gate key would protect him from the beast lure. If the characters climbed into this space before exploring area 3, the same animal attack that would have happened in area 3 happens here instead.
\paragraph{Panjaian Corpse.}
Characters who drag the corpse from under the rock or who lift the rock off of him discover an ornate dragon-headed key needed to open the gate in area 5, a potion of hill giant strength, and a coin purse with a handful of gemstones, each worth 50 gp: a citrine, a carnelian, three onyx, and two quartz.
\subsection{5. Magic Gate}
\begin{DndReadAloud}{}
This disproportionally tall gate spans the ten-foot gap between two rocky outcroppings. It appears to be composed entirely of a smooth white stone, is spotless, and seamless except for a small keyhole in its center.
\end{DndReadAloud}

The gate itself is not strong, but it is loud. The stone is akin to porcelain in its smoothness and sheen but is less brittle.
\paragraph{Gaining Entry.}
The easiest way in is with the key (see area 4). Smashing the gate open requires dealing 20 damage to it (AC 17, damage threshold 5, 40 hp), but each time damage is done to it, it reverberates, creating a shockwave that deals 7 (2d6) thunder damage to creatures within a 20-foot radius of the door. This alerts the ithjar (and the panjaians) that someone is breaking through the gate. The gate is nearly impossible to scale, requiring a successful DC 20 Strength (Athletics) check. The cliffs surrounding it can be climbed more easily. Picking the lock requires a successful DC 15 Dexterity (Thieves' Tools) check.
\subsection{6. The Sacred Pool}
\begin{DndReadAloud}{}
A small pool of water collects at the base of a ten-foot waterfall before trickling away into the rocks behind. A series of small cairns have been constructed on the path beside the water. The path ascends to the north and descends to the south.
\end{DndReadAloud}

The cairns were built by the panjaian pilgrims before the gate key was lost. Three Ithjar* sunbathe on the rocks above the waterfall, keeping an eye on the pool below. They are used to panjaian pilgrims, and if the characters create their own cairns, they're ignored and can investigate the water or continue up the path. If, however, the characters touch the water or continue along the path north without first making cairns, the ithjars attack.
\paragraph{Wishing Pool.}
Hopeful pilgrims have quietly offered riches to the pool. The water is frigid but clear, allowing characters to see down to the bottom of the 15-foot-deep pool. Characters who search its depths with a successful DC 12 Intelligence (Investigation) check over several dives may recover 74 sp, a zircon (50 gp), and 23 gp.
\subsection{7. Nesting Cave}
\begin{DndReadAloud}{}
This narrow cave halfway up the cliffside shelters a grass nest from wind, sun, and weather.
\end{DndReadAloud}

The cave is 15 feet up the cliffside from area 8. The stormborn ithjar's three eggs rest in this nest. If tended by the mother, they hatch in 1d6 + 3 days. If taken from the nest, they cannot hatch without supernatural or magical intervention. The stormborn ithjar* itself is here guarding its brood, unless it was dealt with elsewhere. The stormborn ithjar defends the nest to the death and (if they haven't already been dealt with), its cries summon the ithjars from area 6.
\paragraph{Timing.}
The stormborn ithjar rests during the day but hunts briefly during the nighttime hours. Characters who wait for nightfall have a chance to steal the eggs without fighting the stormborn ithjar (if that's their objective). Observing the nest unseen requires waiting in area 8 and succeeding at a group DC 16 Dexterity (Stealth) check as they stay hidden for hours.
\paragraph{The Nest.}
In addition to the eggs, a suit of Mithral Half Plate Armor has been broken up and its constituent pieces used to hold precious stones, each worth 50 gp: three sardonyx, two star rose quartz, and one cracked garnet.
\subsection{8. Dragon Sanctuary}
\begin{DndReadAloud}{}
This mountain haven is preternaturally peaceful. Butterflies flutter silently in the cool breeze, and the two ancient carvings of dragon heads set among the larger monoliths seem inviting rather than foreboding.
\end{DndReadAloud}

Characters who wish to sneak up on the stormborn ithjar in area 7 must be hidden or to have disguised themselves in panjaian pilgrim clothing while entering this area. If not, it waits for an opportunity to strike.
\paragraph{Dragon Heads.}
Once, this was a holy sanctuary for followers of the dragon gods. Its sacred energy is still present: divine spells that strengthen, bolster, or guard allies (such as bless and spirit guardians) are cast as if one spell level higher while in this sanctuary. Cure wounds and other healing spells are unaffected. Spellcasters who can cast an affected spell immediately realize this property and intuitively understand how it works.
\section{Conclusion}
Characters who capture the stormborn ithjar or recover the stormborn ithjar's eggs and fulfill the requirements of their quest can receive their reward. If they were hired to recover the creature or its eggs, they may find themselves facing an ithjar-riding army in the future! If the egg was a reagent in a ritual or spell, what else do they need to complete it? Does the ritual or spell do what they set out to do? Perhaps the Primordial Ilhara or her followers are upset about their desecrated sanctuary—they may hunt down the characters in the future.
\DndQuote%
{}
{''The gate itself is not strong, but it is loud.''}
{}
\section{Creatures}

\begin{itemize}
\begin{multicols}{3}
\item Corpsejaw
\item Ithjar
\item Stormborn Ithjar
\item Panjaian Ilharan
\item Shieldhead
\item Sloth Galloper
\end{multicols}

\end{itemize}

\chapter{10: Sardonyx Necropolis}
By Tom "Evhelm" Donovan
Sardonyx Necropolis is a corpse walker lair for four or five 10th-level characters.
\section{Background}
At the height of this city's glory and power, a necropolis was built below the streets to commemorate the city's honored war-dead and the general who valorously led their forces before falling in battle. Over time, as the city's power began to wane, the necropolis was sealed and forgotten.
Now, as the city's future is uncertain, a young and ambitious necromancer has rediscovered the Sardonyx Necropolis, so named for the rare and precious red and white marble and gemstones decorating the central mausoleum. In his quest to enhance his city's power, the necromancer began animating the dead, whose corpses have lain undisturbed for so long. Unfortunately for him, his power was insufficient to control them, and now the city above is attacked each night by bands of roving undead soldiers wearing, horrifically enough, the tattered remnants of the city's own uniforms from its glory days.
\subsection{Set the Hook}
As the city panics, the characters are brought in to deal with the threat. The city's leadership has been reminded of the Sardonyx Necropolis's existence, but the city's watch has its hands full guarding the sewers, caves, and basements where the undead have broken through to the surface. Adventurers are needed to descend to the necropolis, destroy the undead, and stop whoever is causing the problem.
Alternatively, the necropolis may still be sealed, but important people (perhaps including the characters) are receiving horrifying visions of the dead rising from their graves and destroying the city. The visions all end with an image of a red-and-white stone structure embossed with sardonyx stones. Research in the city library reveals the existence of the necropolis, and the characters must go there to confront the source of the visions before the worst comes to pass. If your adventures have taken your characters far from any city, feel free instead to put this necropolis above the first level of a dungeon or just below a ruined wilderness temple.
\section{Lair Overview}
The necropolis is underground, where natural light never penetrates. While there are torch sconces and braziers for oil spread throughout the area, none are stocked. Characters who put torches in the sconces or oil in the braziers and light them draw the undead like moths to a flame. Clever characters may use this strategy to draw undead where they want them, but since the dead are drawn to the light, they may get more than they bargain for! Descriptions assume the characters have some way to see but not necessarily in color.
Without a necromancer controlling them, most of the dead are listless; they "patrol" a small area around their graves but lack a sense of past or future that would cause them to change that routine without outside stimuli. If using the city-under-siege model, the characters should not have time to descend at night but should be strongly encouraged to enter during the day. If they spend too long, the corpse walkers gather the undead hordes and lead them to the surface to pillage and destroy—possibly creating even more undead—before returning to the necropolis.
\subsection{1. Cobblestone Walkway}
\begin{DndReadAloud}{}
The natural caverns you've traversed give way to precisely worked cobblestones, a twenty-foot-wide road leading to an enormous stone structure, studded with gemstones, that is covered in angular depictions of soldiers drilling, angels floating, and families crying. Gravestones are spaced to the north and south.
\end{DndReadAloud}

Characters enter the necropolis from this location. Characters who choose to approach quietly and without a light source can enter and move about unmolested if they succeed at a group DC 15 Dexterity (Stealth) check. Those using a light source or failing to move quietly draw undead from area 7. Those who use light are able to see the stone structure to the east is red and white, confirming its identity.
\paragraph{Sardonyx Gemstones.}
Though the structure is magnificent, some expense was spared in its construction: only one in every 10 gemstones that sparkle on the structure are real sardonyx, worth 50 gp each. The rest are colored glass baubles. Characters proficient with jeweler's tools or glassblower's tools notice the difference immediately if attempting to remove one from the wall; otherwise, there is only a 10 percent chance that a given gemstone is real. There are 100 stones across each of the three entrances to the necropolis's main structure, 10 of which are real in each case. Prying a stone from the wall takes 5 minutes and an appropriate tool (such as a crowbar or piton and hammer). Glass fakes shatter after a minute of attempted extraction.
\paragraph{Graves.}
Gravestones and crypt niches throughout have a first initial of a first name, first three letters of a last name, and a rank abbreviation inscribed on them. They are made of white stone and sanded or filed smooth but not shiny.
\subsection{2. Angel of Health}
\begin{DndReadAloud}{}
An angelic statue holding a staff in one hand and a chalice upraised in the other watches over the graves in this area. A spongy fungus grows over the base of the statue and over many of the graves in this area.
\end{DndReadAloud}

No creatures threaten the characters here, but the "fungus" hides a danger. The soldiers buried here died from disease they brought home from battle, not from battle wounds.
\paragraph{Angel of Health.}
The statue depicts Julara, an angel of health and former servant of Aurelia (now an ally of Miklas). Characters who clear off the fungus from the base of the statue discover a plate inscribed with the following prayer: "Julara, angel of health, protect us from disease and death." The blessing of Julara is potent: cleaning the statue grants characters who participate in the cleaning advantage on saves against disease while in the necropolis.
\paragraph{Spongy Fungus.}
Closer inspection (and knowledge) can reveal that what appears to be fungus is really a yellow, disease-bearing moss. Druids and those proficient with Medicine automatically recognize it. Others must succeed at a DC 15 Intelligence (Nature) check to discern that the spongy stuff is really moss infected with the Weeping Pox. Physical contact with the moss is sufficient to spread it (though if this contact is in service of cleaning the statue, the character gains advantage on their checks against the disease!).
\subsection{3. Unmarked Mausoleum}
\begin{DndReadAloud}{}
This gated and walled mausoleum bears no identifying markers.
\end{DndReadAloud}

One of the general's humbler lieutenants requested an unmarked grave, but the people couldn't bear the thought of not honoring him; the compromise was this unmarked but magnificent mausoleum. Characters can climb the walls, unlock the gate with a successful DC 15 Dexterity (Thieves' Tools) check, or use another means to get to the door of the mausoleum proper.
\paragraph{Riddle of the Mausoleum.}
The door has no handle or keyhole but bears a riddle in archaic Celestial: "If thou art better, thou dost not possess me; yet they who do art the best of men. I am never shouted, often overlooked, and cannot be bragged of. Speak my name and enter." The best answer to the riddle is "humility," but "humble," "meek," "modest," or some other similar trait works as well. Players who struggle with riddles can still gain access if a character makes a successful DC 15 Wisdom (Insight) check.
\paragraph{Inside.}
This tomb's interior contains a simple stone platform where a corpse in tarnished armor reposes. Gleaming in his hands is a +1 greataxe that, when attuned, gives the wielder a +2 bonus to initiative checks.
\subsection{4. Obelisk}
\begin{DndReadAloud}{}
A towering obelisk shaped like a saber looms over the graves here.
\end{DndReadAloud}

One corpse walker* and eight Skeleton Rifler* lurk in the dark. They attack living creatures on sight.
\paragraph{Obelisk.}
The obelisk bears a simple inscription in Higher Bürach: "By our lives have we secured yours."
\subsection{5. Crypt Niches}
\begin{DndReadAloud}{}
These dense crypt niches are stacked from the floor to the ceiling in two rows beside the central tomb and a smaller mausoleum.
\end{DndReadAloud}

\begin{figure}[h]
\includegraphics[width=0.5\textwidth]{images/adventure/GHLoE/075-10-001.corpse-walker.png}
\end{figure}
The floor here is made of smoother, red pavers compared to the white cobblestones of the road.
\paragraph{Crypt Niches.}
The niches have not yet been defiled or destroyed. The remains inside repose peacefully. Even a systematic search of the niches (which requires breaking their seals or smashing them open) reveals nothing of value interred with the corpses.
\subsection{6. Statue of Lady Gilrain}
\begin{DndReadAloud}{}
A life-size statue of an armored woman wielding a blade guards these graves.
\end{DndReadAloud}

Lady Gilrain died leading her troops in battle, one of the only nobles to lead from the front. Her remains were never recovered, but her sword found its way back to friendly hands and is held by her statue.
\paragraph{Gilrain's Blade.}
Characters with a passive Perception of at least 13 or who succeed at a DC 13 Wisdom (Perception) check, recognize that the sword itself is metal (not stone) and removable. It is a +1 longsword.
\subsection{7. Assorted Graves}
\begin{DndReadAloud}{}
Scattered gravestones, precisely planned and placed, mark the resting places of the fallen.
\end{DndReadAloud}

A pair of corpse walker* trundle through these graves looking for other undead agents. They attack on sight.
\DndQuote%
{}
{''These dense crypt niches are stacked from the floor to the ceiling in two rows beside the central tomb and a smaller mausoleum.''}
{}
\subsection{8. General's Tomb}
\begin{DndReadAloud}{}
A simple wooden box, mostly rotted, sits on a stone pedestal in the center of this mausoleum. The lid is askew. A man wearing a once-fine but now-bloodied cloak stands over the coffin, unmoving.
\end{DndReadAloud}

The general's spirit is long gone, but when the necromancer who raised the first corpse walkers and skeleton soldiers came to raise the general's body, the magics protecting the tomb severed his control over his undead minions. They slaughtered him. He is now a single mindless zombie, doomed to stand watch over the tomb he tried to defile. Characters who call out to him get no response, and any who approach him realize he's a zombie. If they move to interact with the coffin, he attacks; he similarly defends himself if attacked from range.
\paragraph{Noise.}
Loud combat or opening the general's coffin draws the creatures from area 9.
\paragraph{General's Coffin.}
The coffin is behind bars that stretch to the ceiling, but the necromancer's zombified corpse is not. Getting inside the bars and opening the general's coffin unleashes a wave of potent disruptive magic on anyone or anything standing inside the interior walls containing area 8: spellcasters concentrating on spells automatically lose their concentration. Magical effects granted by potions or spells immediately end. Summoned creatures are banished. If the lid is replaced, the effect can be repeated after 24 hours. The general's corpse is skeletal and his clothes colorless and crumbling; nothing of value remains.
\paragraph{The Zombie.}
Unlike the other undead and corpses in the necropolis, the zombie's clothes are modern, and though torn and ripped are reasonably new. A pin bearing a seal of office marks him as a low-level bureaucrat in the city above.
\subsection{9. Destroyed Crypts}
\begin{DndReadAloud}{}
This half of the mausoleum has been utterly destroyed. What were once row upon row of crypt niches have been pulverized. The corpses within are missing, but the musty dust of the crypt has settled over the copious debris.
\end{DndReadAloud}

An ancient corpse walker* and a skeleton commander* are here and attack on sight.
\paragraph{Untouched Mausoleum.}
At the southwest end of area 9 is a mausoleum with double doors. This is where the ancient corpse walker retreats to if defeated (and where it can ultimately be destroyed). The doors are generally closed but unlocked. Inside the characters find cast off armor including gauntlets of ogre power.
\paragraph{Rough Terrain.}
The floor in this area is difficult terrain.
\subsection{10. Niches \& Mosaic}
\begin{DndReadAloud}{}
Densely stacked crypt niches reach to the ceiling along the western and northern walls. The eastern side of this area has been repeatedly smashed with something large and heavy, but it's clear that once these were more crypt niches too. Beyond the niches to the east, a twenty-foot-diameter floor mosaic of a radiant, angelic figure is visible in the gloom.
\end{DndReadAloud}

One of the dead bodies from the niches has found itself flung onto the mosaic, but otherwise most of the debris is far from it. The mosaic was under the influence of a hallow spell that kept out the corpse walkers and other undead, but when the necromancer opened the general's coffin in area 8, it destroyed the effect. The undead still shy from this area, however.
\paragraph{Niche Search.}
The destroyed niches have been looted, but the sealed ones have not. A thorough, systematic search of these niches (which requires breaking their seals and/or smashing them open) reveals one of the soldiers was interred with an outer garment that hasn't succumbed to time: a cloak of elvenkind.
\section{Conclusion}
Destroying the corpse walkers and their skeletal legion marks the conclusion of the characters' journey to the Sardonyx Necropolis. Who was the lone zombie before he was turned into an undead puppet? Have some undead already escaped to wreak havoc and raise more undead legions elsewhere?
Now that the immediate danger posed by the corpse walkers is vanquished, some powerful figures in the city above might be upset that the characters slew them rather than tamed them or turned them against the enemies of the city. Others will cry "desecration!" after learning that these creatures were once revered soldiers and leaders from the city's glorious past. This adventure can easily lead to intrigue in the city above as the characters defend their actions and the enemies of the city (both within and without) make their move!
\DndQuote%
{}
{''Unlike the other undead and corpses in the necropolis, the zombie's clothes are modern, and though torn and ripped are reasonably new.''}
{}
\section{Creatures}

\begin{itemize}
\begin{multicols}{3}
\item Corpse Walker
\item Ancient Corpse Walker
\item Skeleton Rifler
\end{multicols}

\end{itemize}

\chapter{11: Foulest Truths}
By Greg Marks
Foulest Truths is a harvester of lies lair for four or five 11th-level characters.
\section{Background}
Mauro Cavaliere began life as an entertainer, plying his trade as an actor, singer, and grifter. Eventually he accumulated enough money to branch into usury, and later, legitimate banking. Throughout his career, fast talking and silver lies lined his purse. Unfortunately, a few weeks ago, he met a violent end when former acquaintances murdered him for his fat purse. After pushing his body into the canal, it got caught in an effluent pipe near his home, floating there unnoticed for a day before rising again the next night as a harvester of lies.
Greed still claws at Mauro, but it's a hunger to collect the tongues of liars and sew himself the finest coat of tongues Morencia has ever seen. The recently born harvester of lies haunts low areas of the city: half-flooded basements, sewage pipes that drain to the canals, and walkways under the docks—dragging liars to their deaths. In a city of merchants' and bankers' guilds, liars are easy to come by.
\subsection{Set the Hook}
In one week, it will be midsummer—and the Black Carnival comes to the city. Morencians dress in costumes, lie about who they are, and indulge in all sorts of base activities. If allowed free reign during the Black Carnival, Mauro may feast and his pack of undead grow out of control.
\begin{figure}[h]
\includegraphics[width=0.5\textwidth]{images/adventure/GHLoE/078-11-001.harvester-of-lies.png}
\end{figure}
But a member of the Bankers' Guild, Vissia Giliverti, noticed that many in her neighborhood have gone missing. She fears for her person and fortune. Money is replaceable but her life is not, so she hired a group of troubleshooters to get to the bottom of what the city's watch could not. Many of the missing are con artists, grifters, and even a few politicians. All the disappearances have occurred in the same neighborhood, so there is bound to be some lair within. The characters have seven nights before the Black Carnival and Doge Giliverti's ball. If her ball suffers, the characters will suffer, but if the murderer is brought to heel before, the reward will be great.
\paragraph{Alternate Hook.}
Magistrate Lucente was as crooked as a dog's leg, taking bribes to issue official writs, change a judicial ruling, or arrange for the jailing of political enemies. An enemy of the characters pays Magistrate Lucente to issue a warrant for the arrest of the characters. Before it can be served, however, his lifeless body is found with its tongue cut out. Using the corrupted warrant as leverage, Magistrate Savina orders the characters to solve Lucente's murder or be imprisoned.
\section{Lair Overview}
Much of the city on the canals rests on pilings that keep the buildings from sinking into the water. However, many older buildings have sunk a story or more, creating a maze of partially submerged basements, docks, and forests of large wooden poles.
Less a lair and more a hunting ground, Mauro Cavaliere prowls the wet, flooded undercity that surrounds his now-empty home. The more his coat of tongues has grown, the more victims have risen from their own deaths to follow him. He has put them to work clearing out some of the basements, knocking down walls, and attaching planks to pylons under the buildings. This allows him to crawl above the water but under the streets. Everything is wet, dark, and foul smelling, but the undead don't mind. It gives them a quiet place to feast without being noticed by the city watch.
\subsection{1. Underdock}
\begin{DndReadAloud}{}
Algae-covered stairs lead below street level to a half-submerged gate beneath a darkened house. A private boathouse is under the mansion, complete with dock and a poleboat.
\end{DndReadAloud}

Mauro's home has a dock underneath, which he uses to access the canal. With his death, the key has gone missing and his boat remains untended. The rest of his home has been emptied by creditors and heirs.
\paragraph{Gate.}
The stairs lead down to a ledge intended for dropping off passengers. There is no way to open the gate without a boat or getting into the water. The rusted iron gate's lock is beneath the dirty water. Though not locked, a successful DC 12 Strength check is needed to pull it open.
\paragraph{Boat and Dock.}
A poleboat large enough for four passengers and a boatman is tied to a dock. While inspecting the boat, a successful DC 15 Wisdom (Survival) check notes from the growth of algae between the boat and the dock that no one has used or cared for the boat in several weeks.
\paragraph{Hole in the Wall.}
Hard to see from the gate, a hole has been knocked into one of the foundation walls. This leads to planks nailed to pilings that support a neighboring building (area 2).
\paragraph{Locked Door.}
A locked ironbound timber door leads up into the house. The door can be opened with a successful DC 20 Dexterity (Thieves' Tools) check or forced open with a DC 20 Strength check. Stairs lead up to a three-story, narrow home that once belonged to Mauro Cavaliere. It is empty.
\subsection{2. Path of Planks}
\begin{DndReadAloud}{}
The hole in the wall reveals boards nailed to the neighboring building's pilings: a secret path running under the buildings and streets!
\end{DndReadAloud}

The boards are soaked through and bend precariously when any weight is placed on them, dipping to touch the water. The path passes through holes punched through walls, under streets and buildings, avoiding the canals themselves. It passes several areas where a nimble climber could pass up to an alley or waterfront without being seen.
Mauro left a surprise for anyone who might stumble across his secret path. He captured two Shambling Mound and four Swarm of Quippers, leaving them trapped in a flooded basement, occasionally throwing in the bodies of his victims. The shambling mounds look like floating masses of kelp and trash, but as soon as a leg dangles, they try to pull the victim underwater. The quippers are likewise attracted to motion, swarming anyone in the water. All the creatures are starving and desperate for a fresh meal.
\paragraph{Planks.}
The slippery, warped boards can easily be crawled across, but any creature that tries to walk on them must succeed on a DC 15 Dexterity (Acrobatics) check or slip into the water. If more than 250 pounds is placed on any plank at a time, the boards crack, dumping everyone into the water.
\subsection{3. Pumping Room}
\begin{DndReadAloud}{}
The sound of splashing water and a rhythmic thumping can be heard before anything can be seen. Six waterlogged corpses move with a sick squishing and tinkling of glass as they endlessly work two pumps, fighting the encroachment of the sea into this low basement-like room. Several chests are piled atop an overturned wardrobe, out of the water, near the exit out of this room.
\end{DndReadAloud}

\begin{figure}[h]
\includegraphics[width=0.5\textwidth]{images/adventure/GHLoE/081-11-002.scream-thief.png}
\end{figure}
This room appears to be from an older building that has sunk further into the mud. The hole from neighboring buildings break through the near the ceiling of this room. There are no planks crossing this basement. Six Shatter Corpse* are manning pumps, forcing water out of the room. They take no action unless the characters enter the room or attack them, in which case they stop pumping, silently advance. They open their mouths as if to scream but no noise comes out. All of the zombies have had their tongues ripped out. A successful DC 20 Wisdom (Medicine) check can tell that it appears their tongues were removed while they were still alive, before being thrown from windows. Some of the zombies are recognizable as some of the missing grifters and con artists.
The ceiling of this room is eight feet high and is only one foot above the water in the neighboring areas.
\paragraph{Chests.}
The three locked (DC 20) chests rest on an overturned wardrobe. The first chest contains Mauro Cavaliere's business records. If studied over a short rest, a character can determine that many of Mauro's business deals were based on lies, and he was stealing from many of his clients. The records also detail a money laundering scheme for the Red Mask gang. The documents also include details of his personal bank accounts. Most are empty now, but one in a local account box should have 2,700 gp in it, and the key is in the book! If the chest is underwater, the characters must carry the entire chest carefully and dry it out for at least a day in order to rescue the documents. The second chest includes 750 gp worth of jewelry and twenty-two copies of the same fake diamond necklace. A successful DC 15 Intelligence (Investigation) check is needed to determine the forgery. Otherwise, the necklaces seem to be worth 50 gp each. The third chest is the smallest and has spell scroll of lesser restoration in an old scroll tube, a potion of greater healing, and the mummified hand of a child. Each round this chest is underwater, roll 1d20. If the roll is 8 or less, the scroll is destroyed.
\paragraph{Rising Water.}
Every round, two feet of water rushes into the room and each pump maintains the balance by each pushing out one foot of water. If the shatter corpses stop pumping, the balance is disrupted, and the room quickly starts to fill. A character who succeeds on a DC 15 Strength (Athletics) check can use an action to operate a pump. If the water rises to four feet, the chests go under water. If the water gets to seven feet, the Shambling Mound and four Swarm of Quippers from area 2 can get inside, if they are still alive. The room becomes difficult terrain if the water is two feet deep or more.
\subsection{4. The Liar's Test}
\begin{DndReadAloud}{}
The planks end at a flooded, round room. From several different walls, pipes drip vile waste into a cesspit that stands between you and a set of stairs that rise to the opposite wall. The iron door has been cast with a bas relief of two women, one smiling and coy and the other stoic and chaste. Both are puckering their lips as to offer a kiss. Markings around the door frame could be writing.
\end{DndReadAloud}

The last obstacle before locating Mauro's lair is the Liar's Test. Several broken pipes that should empty into the main canal now deposit their contents into a sunken room below the water level, trapping the waste. If someone is willing to brave the filth, they still must bypass a magical door.
\paragraph{Cesspool.}
The pool of filth is 10 feet deep, and can be swum across normally. Any creature doing so must succeed on a DC 11 Constitution saving throw or become infected with Sewer Plague. Any open flame in this room causes the room to explode, dealing 21 (6d6) fire damage. A successful DC 20 Dexterity saving throw halves the damage.
\paragraph{Liar's Door.}
An iron door with no obvious hinges, handle, or lock is at the top of stairs that rise out of the waste pool. The iron is cast to depict two women with lips pursed as if to offer a kiss with small holes in the door between their lips. There appears to be a small switch inside the mouths of both women, just in reach of a creature's tongue. One woman is poised in a flirtatious manner offering a smiling kiss while the other's kiss appears more honest or virtuous. The door radiates abjuration and evocation magic if checked for.
Writing around the frame reads: "She'll make you a crier, But I've never been slyer. You should have cursed me prior, For now, you kiss me sire." Placing anything other than a tongue inside one of the holes causes the switch to magically disappear and the creature is shocked with 22 (4d10) necrotic damage. A successful DC 16 Constitution saving throw halves the damage. The solution is to kiss the "honest" woman, using one's tongue to flip the switch which causes the door to disappear for 1 minute. If a creature kisses the "lying" woman, a blade slashes down on the liar's tongue doing 33 (8d6 + 5) slashing damage, and the target must succeed on a DC 16 Constitution saving throw or have their tongue ripped out. A creature without a tongue cannot speak or cast spells with verbal components. A creature is immune to this effect if it is immune to slashing damage, doesn't have or need a tongue, or has legendary actions. A creature losing its tongue can have it reattached with a successful DC 20 Wisdom (Medicine) check and a week of rest. Higher level healing magic can also do the trick, at the GM's discretion.
\subsection{5. Mauro's Lair}
\begin{DndReadAloud}{}
The iron door disappears, revealing the inside of a rickety tall house that has partially sunk into the water. Its windows are boarded, and it appears as if most of the floors and interior walls have collapsed, leaving rubble and a pile of corpses rising out of the water. Pieces of masonry collapse from above, splashing into the water. Scaffolding has been affixed to the unsteady walls, climbing several stories.
\end{DndReadAloud}

The building has partially sunk into the water, and only four stories remain above the water with parts of crumbling floor at the ends every 10 feet. Sunlight (or moonlight) can be seen coming through a hole in the ceiling. Mauro Cavaliere (harvester of lies*), 2 Psychic Fragment Swarm*, and 4 Scream Thief* lair here. The scream thieves have no tongues,
If the creatures have forewarning of the characters' approach, the scream thieves join the corpse pile to be surprisingly close to invaders when they enter the room. The psychic fragment swarms hover atop scaffolding on opposite walls, with thier maddening voices kept a low murmur until the characters enter. Mauro carries a Alchemist's Fire (flask), which he lobs into area 4, trying to catch foes there in a secondary explosion. He then casts invisibility and uses its spider climb trait to hang 15 feet above the door.
\paragraph{Boarded-Up Windows.}
Every story has a window on each side of the building. They are heavily boarded to prevent the curious or foolish from entering the condemned building. A window can be broken through with a successful DC 25 Strength check or by doing 25 points of damage to the timbers.
\paragraph{Collapsing Masonry.}
Every round that violent actions occur in the building causes a chunk of masonry to fall from overhead. A randomly determined character must succeed on a DC 15 Dexterity saving throw or get hit by the masonry, taking 14 (4d6) bludgeoning damage. Any creature hit by the masonry must succeed on a DC 15 Strength saving throw or be knocked into the water below. If the saving throw is failed by 5 or more, the chunk of wall pins the character underwater. The character is restrained until the victim succeeds on a DC 15 escape check. A character adjacent to the victim can move the masonry with a DC 15 Strength (Athletics) check.
Large explosions such as spells that do thunder damage, or a fire detonating the gas in area 4, cause a rain of masonry targeting all creatures in the room. Those with scaffolding over their heads gain advantage to the saving throw. If this happens three times, the entire building starts to collapse and all creatures still in the building after three rounds suffer 55 (10d10) bludgeoning damage and are pinned by the rubble. Such a creature is restrained until they succeed on freeing themselves with a successful DC 20 Strength or Dexterity check to escape.
\paragraph{Corpse Pile.}
On the first floor, the only solid ground above water is a pile of loose rubble and corpses. Treat this ground as difficult terrain.
\paragraph{Scaffolding.}
The scaffolding has plank floors attached to the walls every 10 feet vertically, roughly where each story previously was. There are ladders at alternating ends. Moving at more than half speed on the scaffolding requires a successful DC 15 Dexterity (Acrobatics) check or the creature falls off into the water below.
\section{Conclusion}
The defeat of Mauro Cavaliere fulfills any obligation the characters may have to Doge Vissia Giliverti or Magistrate Lucente, including a reward of 500 gp.
If the characters recovered the records from the pumping room (area 3), they can claim the money in Mauro's secret account, but the party also becomes a potential target of the Red Masks, the gang that killed Mauro when he stole their money in a fake deal to launder their ill-gotten gains. The gang wants their money, and they want to silence anyone who has seen Mauro's books.
\section{Creatures}

\begin{itemize}
\begin{multicols}{3}
\item Harvester of Lies
\item Scream Thief
\item Psychic Fragment
\item Psychic Fragment Swarm
\item Shatter Corpse
\end{multicols}

\end{itemize}

\chapter{12: The Tower of Flicker and Shadow}
By Greg Marks
The Tower of Flicker and Shadow is a candlelight daemon lair for four or five 12th-level characters.
\section{Background}
Seeking power, the wizard Henri L'ombre de la Nuit knowingly walked into the Dark Mists of Charneault and came out changed. Any trace of mercy or kindness was scrubbed away, and only a thirst for power remained. In addition, the wizard carried out The Dark-Hearted Oath, a spellbook containing vile rituals, including the steps needed to summon a candlelight daemon.
Henri L'ombre de la Nuit is now a necromancer of peerless skill with dreams of greatness. His study of The Dark-Hearted Oath only made him stronger. While gathering power and developing his skills, he offended several other wizards and important nobles. Most suspect his ambition, but none realized how dark his heart truly became. He had secretly plagued the city for years, but only recently has his thievery of the dead and use dark magic been uncovered. Several of his rivals and detractors met bloody ends, despite being locked inside their homes. Henri knew his time in the shadows would have to end and his rivals would send assassins after him, just as he did for them. He prepared a death trap in his own tower, hoping to eliminate any would-be heroes standing between him and his final goals.
\subsection{Set the Hook}
Following a chain of clues that include missing cats, the missing body of a murderer, and a mysterious fire that killed a family of seven, the characters have arrived at a tower in the Charneault Kingdom near Castle Lamesdhonneur. The tower is bathed in shadow, and occasional flickering arcane lights are seen through its narrow windows. The locals shun the Tower of Flicker and Shadow.
\paragraph{Alternate Hook.}
If the party has a patron in your campaign, they were the rival targeted by Henri L'ombre de la Nuit. Their patron has asked the characters to bring the wizard to justice by going to the Tower of Flicker and Shadow to find proof of Henri's misdeeds and arrest the wizard.
\begin{figure}[h]
\includegraphics[width=0.5\textwidth]{images/adventure/GHLoE/085-12-001.candlelight-daemon.png}
\end{figure}
\section{Lair Overview}
The tower cannot be entered via teleportation and the walls are enchanted to prevent effects that transmute them, such as passwall or stone shape. There is one door that provides entrance to the first floor (area 1), and there are two small slit windows on the top floor. The glass windows are similarly enchanted against transmutation or breakage. Each story is 20-feet tall with 5-foot thick stone floors separating them.
\paragraph{The Death Trap.}
Henri has left via the portal mirror in area 5 and rigged the tower death trap to activate when the characters enter the tower. When all the characters enter the tower, the following changes occur:

\begin{itemize}
\item The windows and exterior door disappear, replaced with sections of wall.
\item The candle in area 5 is lit, summoning a candlelight daemon* with orders to kill the intruders.
\item Teleportation and planar travel are impossible through the walls of the tower, except by using the mirror portal in area 5, though one can still teleport within the tower.
\item A guards and wards spell triggers. In addition to the corridor fog, magically locked doors, and Web in the stairs, Stinking Cloud are conjured in areas 2 and 3 and a suggestion is placed in area 1. See the listed areas for specific details.
\end{itemize}

\paragraph{Shadow Pillar.}
A 5-foot-wide pillar of impenetrable flickering darkness rises from the bottom of the tower through holes in each floor all the way to the very top of the tower. Any creature that is not an undead or fiend that enters or starts its turn in the pillar takes 22 (4d10) necrotic damage and has its speed is reduced by 10 feet until the end of its next turn. A successful DC 15 Constitution saving halves the damage and negates the decrease in speed. The hole in the floors that the pillar passes through are large enough that a creature could pass through them, provided they are willing to enter the darkness. The candlelight daemon uses the shadow pillar to move from floor to floor without being seen. It utilizes hit and run tactics when the party triggers other combats or traps, taking advantage of the teleportation effect of its horns attack to flee back into the pillar.
\paragraph{Stairways.}
Each floor has an enclosed stair that ascends to the next floor. At the top and bottom of each is a locked, reinforced wooden door. Each door has an AC 15 and 25 hp. The doors are immune to poison and psychic damage but vulnerable to fire. The locks can be opened with a successful DC 15 Dexterity (Thieves' Tools) check or a DC 20 Strength check to force it.
\subsection{1. First Floor–Grand Hall}
\begin{DndReadAloud}{}
A pillar of opaque darkness rises from floor to ceiling in the center of the room. Banners displaying a silver skull over crossed candlesticks hang on the walls, flanking a large chair. A door appears to lead to stairs going up.
\end{DndReadAloud}

The grand hall is where Henri would receive guests, seated in his throne-like chair. After the exterior door disappears, four Oblivion Brute* and two Oblivion Whistler* emerge from the banners. While the creatures attack, they attempt to push living creatures into the shadow pillar.
\paragraph{Suggestion.}
A suggestion from the guards and wards affects the area in front of the stairs. A creature passing through that area must succeed on a DC 17 Wisdom saving throw or stand quietly next to the shadow pillar for 8 hours.
\subsection{2. Second Floor - Library}
\begin{DndReadAloud}{}
A yellow cloud and numerous bookshelves make it difficult to see across the room. A large, closed tome sits on a reading table in front of a single chair. Another tome is in a glass case over the mantle. A low fire in the hearth provides dim illumination. As below, the pillar of darkness rises through the ceiling.
\end{DndReadAloud}

The library is filled with a maze of bookshelves. The collection contains local history, anatomy texts, political treatises, natural history, recipes for candle making, arcana from the real to weird, and a wide selection of fiction dominated by mysteries. Many of the books are rare and the entire collection could be sold for 750 gp to the right collector.
\paragraph{Stinking Cloud.}
This entire area is filled with a stinking cloud (spell save DC 17) from the guards and wards spell.
\paragraph{Dark-Hearted Oath.}
Henri's copy of The Dark Hearted Oath rests on the reading table. Inside the tome are copies of all the spells that Henri has memorized, as well as a complete description of a ritual to summon a candlelight daemon and step by step directions for the construction and enchantment of their candle. A successful DC 17 Intelligence (Investigation) notices a nearly invisible rune on the cover. Touching or moving the tome causes the rune to summon six Shadow and three Wraith that attack the creature that triggered the rune. The rune can be dispelled as a 6th-level spell. Any creature of non-evil alignment who completes a long rest within 10 feet of the book must succeed on a DC 20 Wisdom saving throw or suffer 10 (3d6) psychic damage from nightmares filled with images of them reading The Dark Hearted Oath. If a character fails the saving throw and you are using the Influence dice optional rule from the Grim Hollow's Player's Guide, add one die to the Beast Pool.
The Dark Hearted Oath cannot be damaged or destroyed by the characters at this time.
\paragraph{Trapped Tome.}
The tome in the glass case is a trap. A glass box encloses a tome whose cover is embossed with a skull. If the case is opened and exposed to air, an alchemical coating on the tome bursts into flame and fills the library with burnt othur fumes (DC 13 Constitution saving throw). A successful DC 20 Intelligence (Investigation) check made while the glass case is sealed notices a glistening coating on the book. The trap cannot be disarmed while the case is closed and opening it triggers the trap. The book itself is blank.
\DndQuote%
{}
{''A yellow cloud and numerous bookshelves make it difficult to see across the room.''}
{}
\begin{figure}[h]
\includegraphics[width=0.5\textwidth]{images/adventure/GHLoE/088-12-002.oblivion-stalkers.png}
\end{figure}
\subsection{3. Third Floor–Living Quarters}
\begin{DndReadAloud}{}
A sleeping area, cooking hearth, and other household essentials take up this floor. A yellow cloud drifts throughout the room. The pillar of darkness goes through the ceiling and floor.
\end{DndReadAloud}

The detritus of Henri's life and biological needs take up this floor. It's clear he doesn't spend much time on this floor beyond what is needed.
\paragraph{Stinking Cloud.}
This entire area is filled with a stinking cloud (spell save DC 17) from the guards and wards spell.
\paragraph{Treasure.}
A successful DC 15 Wisdom (Perception) check locates a potion of supreme healing hidden under the bed's mattress.
\subsection{4. Fourth Floor–Laboratory}
\begin{DndReadAloud}{}
Framed stain glass windows, with a niche in the wall behind each, decorate an operating theater. A candle burns in each niche, illuminating the stained glass. Bodies under sheets rest upon stone slabs, with trays of surgical tools next to some of them. A counter against the wall is filled with bottles and jars. The ever-present pillar of darkness is here as well.
\end{DndReadAloud}

This room is where Henri plies his necromantic trade.
\paragraph{Bodies.}
Seventeen humanoid bodies lie on slabs, each covered by a plain white sheet. There are symbols painted on each, recognizable with a successful DC 15 Intelligence (Arcana) check as preparation for animating the undead. None are animated when the characters arrive in the room.
\paragraph{Bottles.}
The table of solutions contains spell components, organs in embalming fluid, blood, bile, and a host of other disgusting substances. Each is labeled with cramped handwriting.
\paragraph{Stained Glass Windows.}
Four large stained-glass windows decorate the cardinal points in the room. A shallow niche with a magically burning candle is enclosed behind each. The only way to reach the candle is to break the window. The windows depict the following:

\begin{itemize}
\item A dark mist populated with ghostly faces rising over a large tome
\item A cowled wizard with a ritual dagger raised over a cat, the blood dripping off the blade to a candle
\item A horned daemonic creature sliding from under a bed while an elderly nobleman sleeps unaware
\item A group of adventures in a room that looks like this one being overwhelmed by undead and daemons.
\end{itemize}

Breaking any of the stained-glass windows causes all of them to explode, doing 66 (12d10) piercing to all creatures in the room. Note that area of effect spells might cause the windows to explode as well, such as when targeting the candlelight daemon. This extinguishes the candles and plunges the room into darkness. A successful DC 18 Dexterity saving throw halves the damage.
The round after the explosion 12 Shatter Corpse* rise from the slabs. A successful DC 20 Wisdom (Perception) notices arcane runes etching into the lead between the panes of glass. Each window can be disabled with a successful DC 20 Intelligence (Arcana) check or a Dexterity check with proficiency in thieves' tools. If a skill check is failed by 5 or more, or a dispel magic fails, the windows explode. A window can also be dispelled as a 6th-level spell. For each disabled window, decrease the damage of the explosion by 4d10 and remove 4 Shatter Corpse*.
\subsection{5. Fifth Floor–Summoning Chamber}
\begin{DndReadAloud}{}
A black candle covered in foul runes burns atop a pedestal next to the column of darkness. A mirror is mounted opposite it on the wall.
\end{DndReadAloud}

This is the daemon's candle, and if the party makes it this far without having killed the candlelight daemon*, it does everything within its power to keep them away from the candle, foregoing its hit-and-run tactics. In addition, Henri (archmage) is scrying on the summoning chamber from afar. If the characters have made this far, he uses the mirror portal to join the fight at an advantageous point, bringing his staff of frost with him.
\paragraph{Extinguishing the Candle.}
The candlelight daemon's candle can be extinguished with an action while adjacent to the candle. Strong winds, splashed fluids, or other attempts to smother it from a distance require the character to succeed on a DC 15 Dexterity check. The candle can also be destroyed by doing 15 points of damage to it. It is immune to necrotic, poison, and psychic damage and resistant to fire damage. If the candle is extinguished, the candlelight daemon is instantly banished.
\paragraph{Mirror Portal.}
Henri has set up a portal to a cave outside the city, where there is a matching mirror. With a command word, the mirrors fill with darkness in which a single candle flame flickers. By touching the flame in the mirror, and speaking the command word a second time, a creature can step between mirrors as part of their movement.
\section{Conclusion}
With the candlelight daemon banished, the characters can kill or capture Henri L'ombre de la Nuit. There's ample evidence of his wrongdoing in the tower, and he is quickly tried and put to death. For their aid, the characters might be rewarded with the Tower of Flicker and Shadow, on the condition that they purify it. Who knows what other evils or traps await them in the tower? Is there a strong source of evil beneath it? Finding a way to banish the shadow pillar and determining a way to destroy The Dark Hearted Oath are both entire quests on their own.
\section{Creatures}

\begin{itemize}
\begin{multicols}{3}
\item Candlelight Daemon
\item Oblivion Brute
\item Oblivion Juggernaut
\item Oblivion Leaper
\item Oblivion Whistler
\item Shatter Corpse
\end{multicols}

\end{itemize}

\chapter{13: Swamp of Fate}
By Greg Marks
Swamp of Fate is an ordeal tree lair for four or five 13th-level characters.
\section{Background}
In times now forgotten, a small temple to Typharia, the goddess of truth, was raised on a series of small hills. Called Astranag, the holy shrine was a place of learning, focused on uncovering natural wonders and divining fate from the stars. After the God's End, the clear and fertile land about Astranag sank into more, and what was bright became foul. The burgeoning swamp swallowed most of the temple, and thick mists closed off the stars from sight. The place of clarity became obscured and forgotten by most until it was put to darker purpose.
What remains became the lair of an ordeal tree, a mystical evil plant that tempts the foolish with corrupted wishes. For decades, the ordeal tree has rested atop one of the few hills that rise out of the swamp, feasting on the foolish and granting power to those of the darkest hearts. The latest, a dark werebear known as Folmar Duskscape, has camped in the ruins. Folmar seeks to gain more and more power from the ordeal tree, but the tree taunts him. It granted him power once, but despite surviving the ordeal three times, his further pleas have gone unanswered. Folmar is systematically seeking loopholes in the mystical bargain.
His latest attempt is to awaken animals and have them undergo the ordeal while under the effects of dominate beast, such that they might wish for more power for him. None of the animals have survived or granted him power despite the promises of the ordeal tree. Hungry for power at any cost, Folmar is not ready to give up.
\subsection{Set the Hook}
The characters hear tales of a mystical tree in the swamp that grants wishes to those judged worthy. Hoping to receive their hearts' desire, they set out on a quest to find the tree. After much research, the characters locate a map drawn by a madman guilty of a series of increasingly outlandish murders. The map depicts the location of the tree among the ruins of some place called Astranag. After traveling through a dangerous swamp, they come across the swamp.
\paragraph{Alternate Hook.}
Galia Feykin, a druid, has received a vision of a dryad in trouble. In her vision the dryad stands among the ruins of a temple sinking into a swamp, tears streaming down her face as a great shadow passes over her. The last thing she sees is the symbol of Typharia upon the temple's altar. Galia begs the characters to find this ancient temple and aid the dryad.
\begin{figure}[h]
\includegraphics[width=0.5\textwidth]{images/adventure/GHLoE/091-13-001.oozing-vulture.png}
\end{figure}
\section{Lair Overview}
Several awakened animals patrol the ruins of Astranag. They are listed where they're most likely to lair, but they move about. A delicate treaty exists between three groups of animals, and the characters might be able to take advantage of this animosity when dealing with the inhabitants of the ruins.
In addition, the entire area is filled with a thick fog that heavily obscures all vision beyond 30 feet and lightly obscures everything within that radius.
\subsection{1. The Hanging Tree}
\begin{DndReadAloud}{}
Your boat bumps into a decrepit dock jutting out of the dense fog. Beyond, you can make out a small, tree-covered patch of land. A huge tree is just far enough away that you can't make out any details.
\end{DndReadAloud}

Folmar has taken the bodies of several petitioners to the ordeal tree and hung them upon a large tree, hoping to distract or delay future visitors. The rotting meat frequently attracts many of the awakened animals who come to feast.
Yacka (awakened giant ape) and Bludfeather (awakened oozing vulture*) are most frequently found here. Neither has yet undergone the ordeal, and both are charmed by Folmar. The two have a wary truce and are considering joining together, particularly since they fear and dislike both Snapjaw (area 6) and the Hollow Sisters (area 3). Neither is especially hungry at the moment. If approached by adventurers, they try to convince them that this is the real ordeal tree and that to be judged worthy, they must bind themselves to it.
\paragraph{Hanging Tree.}
A massive, entirely natural tree grows on the island. Six humanoid bodies of various races hang from its branches, fraying weathered rope tied around their necks or waists. A successful DC 13 Intelligence (Investigation) check notes that the bodies appear to have been tied to the tree after they died from various forms of trauma and are unlikely to have bound themselves. All the bodies have been nibbled on since dying. A successful DC 15 Intelligence (Arcana) check notes that several of the bodies show damage from necrotic energy.
\subsection{2. Standing Stones}
\begin{DndReadAloud}{}
Several standing stones crown a hill that rises steeply more than twenty feet out of the water.
\end{DndReadAloud}

A natural observatory for stargazing, the island radiates a feeling of peace and calm which the awakened animals find repellant. The animals are never found here without being drawn here, making it the safest place to rest in Astranag.
\paragraph{Standing Stones.}
Each of the stones is covered in complicated scrollwork; many of the spirals have small holes carved into them. A successful DC 15 Intelligence (Nature) check recognizes the patterns as the paths of planets and stars through the heavens, though in positions from long ago.
The stones radiate evocation magic. Anyone that completes a short or long rest within the circle heals an additional hit point per hit die spent, and any effect causing the creature to be charmed or frightened is suppressed while they remain within the circle. Once a creature benefits from the circle of standing stones, they cannot do so again until a new moon has passed.
\subsection{3. Ruins of Astranag}
\begin{DndReadAloud}{}
Ruined stone walls, blocks of marble, and broken pillars dot this patch of dry land.
\end{DndReadAloud}

The ruins of the temple proper stand on this rise. Weeds grow between the stones of a walkway that passes between broken pillars to a mostly intact temple. The Hollow Sister, a pack of seven awakened night wolves (treat as Winter Wolf, but they are Medium-sized, and replace their cold immunity and cold breath with necrotic) lair in the other ruined buildings. Folmar tried binding an entire pack of black-furred wolves to the ordeal tree. All the males died, but the females grew larger and more intelligent. Now calling themselves the Hollow Sisters, the pack has claimed the ruins of Astranag as their own. They are the most ardent supporters of Folmar and use their numbers to bully Yacka and Bludfeather while avoiding Snapjaw.
\paragraph{Ruins.}
The Hollow Sisters sleep in the ruins, now desecrated by the remains of other animals and humanoids they have partially eaten. The Hollow Sisters are cruel and torture their prey before consuming it. There are ample places one could hide and take cover.
\paragraph{Temple.}
At night, Folmar sleeps in the small temple to Typharia that is still mostly intact beyond a large hole in the roof. A plain room with a simple stone slab altar is all that remain of the former faith. The wall opposite the doors is badly chipped and damaged where Folmar has tried to desecrate Typharia's holy symbol. A bedroll and backpack are behind the altar. In the backpack, beyond food and personal effects, is a pouch with 270 gp and an elixir of health.
\subsection{4. Wooded Hill}
\begin{DndReadAloud}{}
A dense group of trees top a small mound. The trees appear healthier than the other foliage you've seen.
\end{DndReadAloud}

The dryad Dyssodia hides within her tree in the center of the small mound. Once a follower of Typharia and a scholar concerned with the movements of the night sky, the dryad has drowned in melancholy as the world slipped into darkness. When the ordeal tree first came to Astranag, she tried to warn visitors away, but inevitably they were drawn to the tree. Most died. Those that didn't possessed the darkest hearts, and Dyssodia knew they would further spread corruption. Tied to her tree, there is nowhere to go, so she simply watches corruption spread and despairs. She avoids visitors now, but if the characters can find her, or show true hearts when confronting Folmar and the ordeal tree, they might be able to coax her out to offer aid.
Neither Folmar nor the awakened animals know of her. The ordeal tree recognizes the dryad's tree, but with no mortal soul to corrupt and with her wallowing in sadness, it simply doesn't care.
\subsection{5. The Ordeal Tree}
\begin{DndReadAloud}{}
A massive hoary tree rises from the island. Its bark is scarred rope, spikes, and deep gashes.
\end{DndReadAloud}

\begin{figure}[h]
\includegraphics[width=0.5\textwidth]{images/adventure/GHLoE/094-13-002.ordeal-tree.png}
\end{figure}
The ordeal tree* is rarely alone. Folmar Duskscape (werebear ascetic*; the ordeal tree has given Folmar resistance to all damage) spends most of his time here in supplication or studying the tree, hoping to divine some way to gain more power. Between fits of coughing (from his curse of foul blight), he orders visitors to flee before he attacks them with his staff of thunder and lightning. If they refuse, he summons all the awakened animals that remain and launches his attack. The ordeal tree itself uses its heart sight on the characters. If they're more powerful and more flawed than Folmar, it waits for the druid's defeat so that it might curse the characters. If the adventurers are pure of heart, it joins Folmar when the party least expects it and tries to kill them. This emboldens Folmar, and he stops at nothing to murder the characters.
\paragraph{Sunken Boat.}
A foot of the rickety dock is a sunken boat that succumbed to Snapjaw's attack. While the passengers have been eaten, a chest has gone unnoticed in the remains. Inside the locked (DC 20) chest is the treasure a group of mercenaries had intended to offer the ordeal tree in exchange for power: 1,000 gp, a goat's horn with the blasphemous names of daemons carved into it, and a ring of free action that is cursed to bring misfortune to the wearer.
\subsection{6. Shallows}
\begin{DndReadAloud}{}
The water is shallow enough between the hills that you can see the bottom in some places and possibly walk across.
\end{DndReadAloud}

The shallows are the domain of Snapjaw, an awakened giant crocodile. Snapjaw buries itself in the mud and watches, occasionally feasting on something that draws too close. Only Folmar need not fear the crocodile, though as Folmar's first experiment, the druid's charm has worn off and the crocodile merely tolerates the werebear now. Most of the time Snapjaw sleeps or waits, but if hungry, it rouses to feast. When buried in the mud of the shallows, a successful DC 20 Wisdom (Perception) check or comparable passive Perception is needed to distinguish Snapjaw from the surrounding foliage and debris. The ordeal tree's blessing has made Snapjaw as large as its hunger. Snapjaw is a gargantuan giant crocodile with double hit points, +3 AC, +3 to hit, +10 damage on all attacks, and gains an additional bite attack when using Multiattack).
The shallows range from 3 to 15 feet deep and are difficult terrain for anyone walk across.
\begin{figure}[h]
\includegraphics[width=0.5\textwidth]{images/adventure/GHLoE/095-13-003.werebear-ascetic.png}
\end{figure}
\section{Conclusion}
If Folmar and his animals are defeated, the characters can choose to undergo the ordeal, but it most likely leads to death at the hands of the ordeal tree, one way or another. Any boons granted by the ordeal tree fade once the tree is destroyed, and only by destroying the tree can the world be made safer.
If the ordeal tree is defeated, Dyssodia comes forth, grateful for the weight lifted from her heart. She sets about restoring the ruins, though Typharia no longer speaks to her. The dryad may become a great font of knowledge for the party, particularly about history from before the God's End or about the stars and the natural world. If the party aids her in restoring her home, she considers them friends and allow them to visit when they have need of her.
Followers of the Arch Seraph Zabriel and the Arch Daemon Venin would be greatly interested in the ruins, and in particular, the standing stones that retain a fragment of divine power. Protecting, or selling, the location may lead to future adventures.
\section{Creatures}

\begin{itemize}
\begin{multicols}{3}
\item Werebear Ascetic
\item Oozing Vulture
\item Ordeal Tree
\end{multicols}

\end{itemize}

\chapter{14: Street Prophet}
By Greg Marks
Street Prophet is a doomcaller lair for four or five 14th-level characters.
\section{Background}
Kazilor the doomcaller came to a poor neighborhood to preach that the end-times are at hand and that the city's leadership has done nothing to prepare the populace. The neighborhood increasingly fell under the doomcaller's sway, and throngs gather each day to hear the prophet speak in the square in front of his storefront chapel. Daemon attacks, civil unrest, and general lawlessness are on the rise as the apparent end of the world approaches. The fall of the city seems more likely.
Kazilor claims to have been gifted a local furrier's shop and has moved his "church" inside, though he preaches from the fountain in the square outside on most days. No one knows where the furrier is or moved on to, if he even still lives. He welcomes all to share in the lamentations he leads: The world is ending. There is no future or afterlife. Enjoy what little time you have how you will and make ready for the end!
The Hearthkeepers are not happy about the formation of this cult. They have ordered Wachtmeister Nikos Hiedler to expel him and his followers. The guard sergeant has only three patrolmen under his command and doesn't think they'll be able to handle the crowd, much less the street prophet and his acolytes. Worse, he has seen the prophet's proclamations of curses and daemonic invasion come true. It's beyond the guard, but maybe, just maybe, it's something heroes could put an end to before the Hearthkeepers blame him.
\subsection{Set the Hook}
The street prophet has gained a popular following. Dozens of commoners come out to see him preach while his two attractive acolytes move through the crowd handing out bread and meat. The local clergy have asked the guard to remove the street preacher, but the guards are worried about the reaction of the crowds and turn to instead to adventurers to handle the situation.
\paragraph{Alternate Hook.}
If the characters are friendly with a reoccurring low-station NPC, they have gone missing while in the vicinity of the doomcaller and are currently being held in area 4. If the party doesn't rescue them soon, they become the next free meal prepared for the crowd!
\section{Lair Overview}
During the day, Kazilor, wrapped in bandages and robes, climbs the fountain in the street and begins haranguing passersby. By midday, a crowd forms. Kazilor's assistants Klanssa and Alezc hand out bread and meat while accepting donations. When night falls, a few are invited to join them inside for a worship service, though no one can say what they worship, and those who go rarely tell the tale (since few get out alive).
The descriptions below assume the characters arrive while it's light, and the fiends are not aware the characters are coming. The GM should adjust if that is not the case.
\subsection{1. Fountain Square}
\begin{DndReadAloud}{}
A man wrapped in rags and a dirty robe stands atop the square's fountain, preaching to throngs of commonfolk. Decrepit buildings, some boarded up but a few still sporting businesses struggling to hold on, ring the square overflowing with people.
\end{DndReadAloud}

During the day, Kazilor the doomcaller* stands atop the fountain wall, walking around as he preaches to the crowd. His acolytes Klanssa (succubus) and Alezc (incubus) polymorphed into human forms move through the crowd, feeding them and offering words of kindness. Both Klanssa and Alezc carry Potion of Supreme Healing.
If anyone challenges Kazilor or threatens to disperse his crowd, he has a shaking fit and warns that evil approaches before pointing to the roof of one of the neighboring buildings, and summoning 1d4 Lenchtahg* (area 6). Kazilor and his acolytes use the confusion the attack creates to flee into their storefront chapel. Once out of sight, Kazilor's first action is to use gate to summon a fiend listed in his book (area 3) into the square.
\begin{DndReadAloud}{}
The Words of the Prophet Kazilor
"Listen well, children, for time grows short. The world's end draws near, and the veil between worlds thins. Already fiends of the underworld stalk us day and night, hungering for your souls. But in forsaking hope there is a chance to live in joy for but a little longer. There is no afterlife. There is nothing else but here and now. Should you escape the fiends' hunger that the nobles do not protect you from, at world's end your soul shall fray and wither away. Live as you would. Those in power would tell you your place is in the dirt, is to toil for their gain. It is not so! Take what you desire. Live as you wish. For what is chaos but the freedom to live in elation. Embrace pleasure for pleasure's sake, and care not for the rigid society those in power would impose. Nothing matters but what you desire!"
\end{DndReadAloud}

\begin{figure}[h]
\includegraphics[width=0.5\textwidth]{images/adventure/GHLoE/100-14-002.doom-caller.png}
\end{figure}
\paragraph{Crowd.}
On an average day, 5d20 Commoner come to listen to Kazilor preach. They crowd into the small square and fill the surrounding streets. Some climb onto surrounding buildings or hang from windows. Consider the entire area to be difficult terrain while the crowd is present. Three Downcast Mercenary* in the crowd are charmed by Klanssa, Alezc, and Ilsoa to act as bouncers, keeping the crowd from being too disruptive during the sermon.
If combat breaks out, the crowd confusedly tries to flee in every direction for five rounds. In addition to the difficult terrain, any creature in the square or surrounding streets must succeed on a DC 15 Strength saving throw if they move on their turn or be knocked prone. Creatures that are prone at the start of their turn suffer 7 (2d6) bludgeoning damage. Several commoners are trampled to death.
\paragraph{Fountain.}
The fountain is surrounded by a 4-foot-high wall made of marble and can be crouched behind to take cover. Water cascades out of the top and over three successively larger tiered bowls in a classic stacked configuration. The water is pure and drinkable.
\paragraph{Neighboring Shops.}
Several shops front onto the square and surrounding streets. Half are boarded up and abandoned. Other nearby shops include a chandler, a baker, and canvas maker. The surrounding blocks are a significant fire hazard.
\subsection{2. Storefront}
\begin{DndReadAloud}{}
The front area of the former shop is now largely abandoned. Empty racks and a lone counter are all that remain of the former furrier's establishment. A door leads to a back room.
\end{DndReadAloud}

When this was an active shop, this room displayed wares and had a corner where the furrier would take orders or accept pieces to be repaired. The cracked window still reads "Schilnd's Furrier, Sales, Commissions, and Repairs" in flaking gold paint.
\paragraph{Counter.}
Under the counter is a small, sealed wooden box that has been nailed to the underside of the counter and filled with an alchemical device. A successful DC 15 Wisdom (Perception) check notes that a wire connects the box to the door leading to area 3. A character that beats the DC by 5 or more also locates a hidden switch also under the counter. If the door to area 3 is opened without flipping the hidden switch, the wire opens the box and exposes it to air causing it to explode and do 22 (4d10) fire damage to all creatures in area 2. The trap can be disabled with a successful DC 20 Dexterity (Thieves' Tools) check.
\paragraph{Front Door.}
The front door is locked (DC 20 Dexterity (Thieves' Tools) check to unlock). Kazilor carries the only key. The door can be forced open with a successful DC 15 Strength check. The window can be easily smashed, though it makes a lot of noise and is easily heard for blocks.
\subsection{3. Chapel}
\begin{DndReadAloud}{}
A symbol is crudely painted on the southern wall above a low dais and podium. Offerings are piled around the podium. There are two doors exiting the small chapel.
\end{DndReadAloud}

This former workroom has been cleaned out to make room for a small chapel. After sundown Kazilor leads worship services, where he preaches persuasively to throw off the shackles of morality and order, encouraging receptive audiences to do anything they want since the world is about to end. Anyone not receptive is asked to stay after the indoctrination session and invited for additional discussions in area 4. When not preaching or otherwise causing havoc, the fiends rest in this room.
\paragraph{Offerings.}
The offerings of the desperate are packed around the podium. Sacks of coins and jewelry (worth 273 gp), food, crafts, clothing, and the detritus of lives abandoned in anticipation of the end of the world.
\paragraph{Podium.}
On a shelf under the podium is a small book written in Abyssal containing lists of names: congregants that have fully embraced the doomcaller's message, missing people that have been killed and fed to the crowds, and a short list of fiendish true names including Malfikan, Allonassa, Rawnivour, and Ixixxitollan. These can be any kind of fiend the GM wishes to use in the campaign.
A successful DC 20 Wisdom (Perception) check while searching the podium locates a hidden switch that deactivate the trap in area 2 before opening the door.
\paragraph{Symbol.}
The runes painted on the wall are Abyssal runes associated with chaos. Looking closely and succeeding on a DC 18 Intelligence (Investigation) check notes that a Symbol (spell save DC 18) worked into it. The symbol triggers when a humanoid is in the room for more than three rounds without a fiend present.
\subsection{4. Abattoir}
\begin{DndReadAloud}{}
Hacked and bloody bodies are piled in front of cooking hearth, and butchering tools hang on the wall over a low cabinet. A reinforced and barred door exits the building to the west.
\end{DndReadAloud}

Ilsoa the infernal tormentor* plies her trade as a chef here, making the bread and meat that Klanssa and Alezc hand out. All the food is made from, or otherwise tainted by, the bodies of those they abduct. If Ilsoa knows that non-fiends are in the building, she casts alter self to appear as a tied-up victim so she can either surprise would-be rescuers or escape if things turn against Kazilor. Ilsoa has a rod of security that is sometimes brought out for special services in the chapel.
\paragraph{Alley Door.}
The door to the alley is locked. It can be unlocked with a successful DC 20 Dexterity (Thieves' Tools) check or broken down with a successful DC 20 Strength check. Both Kazilor and Ilsoa have keys to this door.
\paragraph{Bodies.}
Recently collected street people, congregants that would not be missed, or those that asked too many questions, have been killed and stacked here. Several have been partially butchered and are currently cooking in the hearth.
\subsection{5. Abandoned Shop}
\begin{DndReadAloud}{}
The original purpose of this abandoned shop is no longer determinable. Three bedrolls are spread out on the floor.
\end{DndReadAloud}

At night, the three charmed Downcast Mercenary* sleep in the abandoned shop. A night one of them stands watch in case Kazilor needs them.
\subsection{6. Summoning Shop}
\begin{DndReadAloud}{}
Inside the boarded-up shop, light pours in through a hole in the ceiling illuminating a circle of magical runes and half burned candles.
\end{DndReadAloud}

Kazilor uses this shop for summoning fiends, and potentially binding them within the prepared circle. If there is trouble, he summons Lenchtahg* to this location.
\section{Conclusion}
\begin{figure}[h]
\includegraphics[width=0.5\textwidth]{images/adventure/GHLoE/101-14-006.lenchtahg.png}
\end{figure}
Defeating Kazilor and his fiends brings a measure of stability to the neighborhood, but unrest remains and eventually leads to riots by the disaffected. Reading through Kazilor's book identifies those who are fully converted to his apocalyptic teachings, including at least one highly placed noble. The characters find the name of Ulsam March among the list of the fiend's converts. Ulsam is the son of Countess March who is well placed in court. Attempting to bring him to justice is a difficult task indeed if the characters don't wish to run afoul of the powerful noble.
Among the bodies in area 4, one carries a map to the fabled ordeal tree, potentially connecting this adventure to the Swamp of Fate.
\section{Creatures}

\begin{itemize}
\begin{multicols}{3}
\item Doomcaller
\item Downcast Mercenary
\item Downcast Apostate
\item Infernal Tormentor
\item Lenchtahg
\end{multicols}

\end{itemize}

\chapter{15: The Thundering Hills}
By Greg Marks
The Thundering Hills is a lindwyrm lair for four or five 15th-level characters.
\section{Background}
Generations-old legends speak of rumbling beneath the Thundering Hills. The stories change over time, but one thing never changes: there is always a great beast called a lindwyrm, a serpent with one pair of claws and an endless appetite. The stories say, when the lindwyrm wakes, devastation always follows. The stories are wrong. There are two lindwyrms under the Thundering Hills.
Digging through the earth, the Lickscale kobold tribe came across the vast beasts and took to worshipping the century-long sleeping monstrosities. The kobolds piled small offerings around the mated pair in hopes that when they awoke, the draconic creatures grant blessings and favor to the tribe. A few weeks ago, two-thirds of the tribe was eaten the first day the pair woke, but that didn't stop the kobolds. The tribe brought more offerings, capturing wild game, farmers, and livestock to appease the lindwyrm pair which, with their hunger sated, began to tolerate the kobolds. Other than a few raids of their own against nearby villages or farms, the lindwyrms have turned their attention toward their hatching eggs. Shortly after they hatch, the lindwyrms will outpace the ability of the kobolds to feed them, and the rest of the surrounding villages will be endangered. Should an entire brood of lindwyrms spawn, the entire north will be imperiled.
\subsection{Set the Hook}
The Lickscales have raided a farm, abducted the farmers, and carried off livestock. Nearby fishermen heard the commotion but, by the time they arrived, the Tolferson family had been carried away. It's clear from tracks in the snow that some humanoids attacked the farm and that at least some of the family was alive when they were taken. The Tolfersons were popular and the eldest son, Ulfrick Snow-Swimmer, was a local hero for saving two children from drowning when they feel through the ice. The adventurers are begged to follow the tracks, rescue the family, and wreck a terrible vengeance upon those who would raid the clan. Many whisper of creatures beneath the Thundering Hills, but the folk tales are rarely consistent beyond claiming a huge feral beast sleeps beneath there, its snores or restless sleep making the hills shake.
\paragraph{Alternate Hook.}
Young Ingrind Settur, barely nine winters, was the only survivor of an attack on her village by one of the lindwyrms. No one believes her tale of a giant serpent burrowing out of the ground, though something clearly devastated the village. She seeks the aid of heroes. She claims that before being swallowed, her father buried the family's ancestral blade into the beast's scales, and Ingrind would have it returned that she might wield it when she grows to be a warrior like her father.
\section{Lair Overview}
The trail leads deep into the snow-covered hills to a cave that shows significant recent traffic, including something quite large coming in and out of the cave. There is no light in the caves, but a foul smell emanates from within. While rough, the ceilings are generally 20 feet high in most places.
\subsection{1. Cave Entrance}
\begin{DndReadAloud}{}
A dark, natural tunnel descends below the hills. The ground is disturbed by the passage of many feet.
\end{DndReadAloud}

While the lindwyrms can and do burrow, they frequent the natural entrance used by the kobolds. The smell of rot and offal is strong. The path proceeds down at a steep angle.
\paragraph{Tracks.}
A successful DC 15 Wisdom (Survival) check identifies humanoid and kobold tracks, as well as some great serpentine beast with huge claws. The traffic is so great that it is impossible to determine exact numbers or if they last went in or out of the cave.
\subsection{2. Windy Ravine}
\begin{DndReadAloud}{}
The tunnel is abruptly bisected by a crevasse whose bottom is out of sight. A strong wind whips from somewhere above, making crossing the distance even more difficult. Strange graffiti is scrawled upon the walls on the other side of the crevasse.
\end{DndReadAloud}

This natural fissure blocks passage. The lindwyrms are long enough to reach themselves across it, while the kobolds keep a log around the corner for crossing.
\paragraph{Crevasse.}
The crevasse is 15 feet across and continues out of sight in all directions. Should a creature fall, the bottom is 200 feet down. A creature with a 15 Strength can make the jump without a check, provided the violent winds don't stop them. A creature with a lower score must succeed on a Strength (Athletics) check equal to 10 + (15–the character's Strength score). Should a character miss the jump, they may attempt a DC 20 Dexterity saving throw to catch the opposite ledge before falling.
The Lickscales keep a long tree trunk and ropes out of sight around the corner for crossing the crevasse.
\paragraph{Graffiti.}
The graffiti scrawled on the walls shows crude drawings of two draconic serpents with one pair of claws surround by several small worshipping draconic creatures. Writing in poor Draconic spells out "Gods of the Lickscales" amid various rude statements.
\paragraph{Violent Winds.}
Random violent winds blow in all directions, making jumping or flying across the crevasse difficult. When a creature attempts to cross the crevasse, roll a 1d10. If the result is odd, the wind pushes down, back, or sideways, impeding progress. Subtract the result of the roll to the distance jumped. If the creature is flying, they must succeed on a Strength (Athletics) check with a DC equal to 10 plus the number rolled or take 10 (3d6) bludgeoning damage as they bounce along crevasse's walls before righting themselves. If the result is even, it blows the creature up or forward, aiding progress. Add the result of the roll to the distance jumped or flown.
\subsection{3. Burrow of the Lickscale}
\begin{DndReadAloud}{}
The flickering light of campfires gives away inhabitation even before the smells. Small burrows and niches, some covered by dirty blankets or scraps of leather, house many small creatures.
\end{DndReadAloud}

\begin{figure}[h]
\includegraphics[width=0.5\textwidth]{images/adventure/GHLoE/108-15-001.morbus-kobold.png}
\end{figure}
The Lickscales live here. When the characters arrive, 20 Morbus Kobold* are cooking, caring for young, sharpening weapons, and sorting through stolen goods. They try to hide if the adventurers make themselves known. Only Quitak Darktongue (mage) is left to speak for the tribe. Quitak is the apprentice to Shuuna Crowbiter (area 4) and as such carries some measure of respect among the kobolds. Quitak realizes that without Shuuna and her bodyguards, the tribe could easily be exterminated by most adventuring bands, so he stalls for time, welcoming them to the burrow of the Lickscales and share a campfire. He claims to have no knowledge of any monsters, thievery, or kidnappings but is happy to sit and discuss it, asking the party to tell him everything that has happened.
While he talks, one of the kobolds sneaks through the small tunnel to area 4 to alert Shuuna, who in turn goes to warn the lindwyrms. If the Tolfersons are dropped into the Nursery Pit, their screams are easily heard in this room.
\subsection{4. Nursery Pit}
\begin{DndReadAloud}{}
As you approach a foul sulfurous smell adds to the fetid rotten meat smell that permeates the tunnels. The sounds of cracking and hissing can be heard echoing from the chamber.
\end{DndReadAloud}

The lindwyrms have lain seventeen eggs and so far, two newborn lindwyrms (treat as Giant Crocodile) have hatched. Shuuna Crowbiter (kobold archmage) and her bodyguards Schminge and Schux (kobold Assassin) are in the process feeding the newborns. They have just fed each a sheep and are about to toss the captured farmers to the newborns. If the kobolds see or hear the characters coming, or are warned, they start cutting the ropes and drop the captives as quickly as possible. If the newborn lindwyrms are harmed, their cries attract their parents (area 5).
\paragraph{Captives.}
There are five Tolfersons, all in bad shape from their recent violent capture: Rune, Vigdis, Ulfric (called Snow-Swimmer), Asgiern, and Natala. They are individually bound and hanging over the ledge by ropes, five feet from the ledge and ten feet above the pit floor. If they are dropped, a newborn lindwyrm can kill one per bite, but then spends a round to swallow its prey.
\subsection{5. Lair of the Lindwyrms}
\begin{DndReadAloud}{}
The floor of this rot-stinking room is covered in bones, many of them with pieces of spoiling meat upon them. A mound of treasure and goods are heaped in the center around a makeshift altar.
\end{DndReadAloud}

\begin{figure}[h]
\includegraphics[width=0.5\textwidth]{images/adventure/GHLoE/109-15-002.lyndworm.png}
\end{figure}
Two Lindwyrm*, called Bloodthirst and Eater of Men by the kobolds, lounge here feasting on a pile of livestock next to the altar. Eater of Men has a longsword of Valikan make stuck in his scales, behind his left claw. This is the sword Ingrind Settur seeks. The lindwyrms are happy to feast on the adventurers, yet another meal delivered directly to them.
\paragraph{Altar.}
The kobolds have piled many years' worth of stolen armor, weapons, trade goods, coins, and other knickknacks on a large flat stone. Several lit candles decorate the altar cause light to scatter off the treasure for 30-feet. Most of the treasure is not especially valuable, but 9,000 gp worth of coins and goods can be recovered. In addition, there is a potion of flying, a staff of healing that is cursed to give the owner an unhealthy obsession with blood, and a dented brass clock that only counts backwards and tolls ominously at noon.
\paragraph{Bones.}
Bones of all types litter the floor, many of them centuries old, to a depth of three feet, making it difficult to move through the room. Treat the entire room as difficult terrain for any creature smaller than Large size.
\section{Conclusion}
If the lindwyrms and the Lickscales are defeated, a great reoccurring evil is prevented—though the legends of the beasts beneath the Thundering Hills are likely to remain for many generations. Oddly though, the rumbling doesn't stop. What horrors might bubble up from the crevasse that the presence of the lindwyrms had kept in check?
If the Tolfersons are rescued, they have little to reward the party with because of the recent damage to their farm. But after a season, they get back on their feet and start sending goat cheese, dried fruits, and mead. A break in their regular deliveries might mean the family is in trouble once again.
If Ingrind Settur's father's sword is recovered, she swears a blood oath to the characters to serve the heroes that avenged her father, though she cannot yet lift his sword. There is a dark anger in her eye, that if left unchecked, might blossom into sinister corruption. If the characters leave her on her own, the oath weighs heavily upon the girl, and she becomes a murderer killing those she believes to be the characters' enemies or detractors. As bodies accumulate everywhere the characters go, rumors begin to circulate that the party is cursed, or possibly evil.
\section{Creatures}

\begin{itemize}
\begin{multicols}{3}
\item Lindwyrm
\item Morbus Kobold
\end{multicols}

\end{itemize}

\chapter{16: Safe in the Stillborn}
By Shawn Merwin
Safe in the Stillborn is a hraptnon lair for four or five 16th-level characters.
\section{Background}
A downcast apostate called Master Krumm wanted to become immortal again, and he planned to perform dark rituals to make that happen. After being driven away from two different lairs by interfering adventurers, he sought a place where no one would bother him: the Stillborn Forest. Krumm found abandoned lodgings upon which to build his new home, including a small manor and a house for servants. Within a month, Krumm and his followers had created a well for water, a garden for food, abundant game for meat, and Krumm knew privacy like they never had before.
Once established in their new dwellings, Krumm's entourage scoured their area of the forest, finding a few woodcutters, hunters, and scouts. Krumm and his lackeys captured these folks in preparation to attempt yet another ritualistic sacrifice to reestablish Krumm's immortality.
All seemed well for Krumm's plan until some roving xakalonuses noticed the new settlement. Krumm repelled their first attack, although they destroyed the home of his servants and many of the trees in the area. The xakalonuses can smell the delicious, magic-infused brain of the downcast apostate, and they know it will make a wonderful meal for their progenitor, the hraptnon.
When the characters arrive at the lair, they find themselves caught in the middle of an ongoing struggle between Master Krumm and several xakalonuses, with the hraptnon about to arrive on the scene to take what it feels belongs to it: the life of any living creature in the Stillborn Forest.
\subsection{Set the Hook}
A small but thriving hamlet called Oakendale rests at the edge of the Stillborn Forest. It acts as a gathering point for those brave souls who ply their trade in the forest as woodcutters, loggers, hunters, and treasure seekers. The leader of the town, a woman named Yvell Stowin, oversees the mercantile businesses that operate in and through Oakendale, acting as a de facto guildmistress. When many more foresters than normal fail to return from their work in a timely manner, Yvell reaches out to adventurers to enter the dangerous forest and see what the problem might be.
Yvell can provide a map of the main logging roads at the edges of the forest, showing the area that supplies Oakendale with its resources. (Krumm's new estate rests one of those logging roads, and Yvell is unaware of its presence.) She can also regale the adventurers with the legend of the hraptnon, which contains an interesting mixture of truth and wild speculation. Discerning which of these tales are true and which are false is up to the characters.
Yvell can truthfully avow to an increase in sightings of creatures called xakalonuses in the area in the last several days. These creatures are very dangerous, but the local militia is experienced at driving away threats. The only time a xakalonus didn't flee after taking some heavy longbow fire was when there was a mage in town doing research on the hraptnon. In that case, the xakalonuses did not retreat until one of them had attacked the mage, ripping off the top of his head and eating his brain. Yvell offers the characters a total of 500 gp to put an end to the attacks and find the cause of the missing foresters.
\paragraph{Alternate Hook.}
Xakalonus attacks have increased dramatically all over Etharis, with most of those attacks happening against people who use magic. After attacking magic-users and eating their brains, the xakalonuses immediately retreat. More than one report has pointed to the Stillborn Forest as the ultimate destination of these brain-eating creatures, so the characters are prompted to investigate.
\section{Lair Overview}
The Stillborn Forest is thick with trees, underbrush, wild game, and even some lesser fey creatures. Moving off established paths means entering difficult terrain, and the countless trees offer much concealment and cover from ranged attacks.
\DndQuote%
{}
{''Also, the simple trodden path is replaced with a road paved with stones.''}
{}
\subsection{1. Logging Road}
\begin{DndReadAloud}{}
Thick, luxuriously leafed oak and pine trees teem in this area of the Stillborn Forest. As you travel, however, you notice some changes: the trees on one side of the road thin out a bit, where they appear sickly and stripped of leaves. Fewer forest-dwelling creatures chitter at you from their dens and perches. Felled logs form a rough barrier on each side of the path, slowing plant growth over the path.
In the distance, on either side of the path, are a pair of wooden structures. Although they are difficult to see through the thick foliage, you can discern a few details. The structure to the west of the path, on the same side as the sickly trees, has obviously sustained structural damage. The structure on the eastern side of the path is in better shape.
\end{DndReadAloud}

The characters get a better view only if they approach or use magic. The current lull in the battle between Krumm's contingent and the xakalonuses means there's nothing to hear from the road.
\paragraph{Tracks.}
A successful DC 15 Wisdom (Survival) check identifies tracks of Medium-sized and booted humanoid creatures, along with some domesticated animals like horses. More recent tracks to and from the area were made by xakalonuses, which look like large wolf tracks but with much sharper claws.
\paragraph{Corrupted Foliage.}
If the characters take a moment to examine the corrupted trees along the west of the path, they can attempt a DC 20 Intelligence (Arcana or Nature) check. On a success, they learn that these trees were at one point magical, infused with a type of fey magic. Something drained the magic from them, leaving them alive but severely damaged.
The magic of the trees was drained by the presence of the xakalonuses as they moved through the area recently. Unless the characters have explicit knowledge of the xakalonus, however, they have no way of knowing their presence was the cause of this corruption of the trees. However, when the characters see magic fail to affect those creatures, they might figure it out.
\paragraph{Lurking Monsters.}
Two sets of five Xakalonus* each wait to the north and south of the Krumm's Manor (Area 3). They watch the situation, waiting for the hraptnon to arrive and take part in a final assault against the occupants of the manor. The characters' presence throws their plan into question. If the characters notice either of these sets of monsters and attack, the group of xakalonuses they attack defends themselves.
The other group of creatures lingers, waiting to see what happens next. They may be tempted to join any fracas between the characters and Krumm's entourage (see Area 3).
\subsection{2. Destroyed Dwelling}
\begin{DndReadAloud}{}
The outer walls and foundation of this dwelling are cracked, the glass windows shattered and the doors torn from their hinges. Three human corpses litter the lawn outside the house.
\end{DndReadAloud}

This building housed Krumm's mundane servants: cooks, farmers, animal keepers, etc. They locked themselves in the building when the xakalonuses attacked, but the monsters entered the building and killed everyone inside, chasing down those who attempted to flee.
\paragraph{Building.}
A character succeeding on a DC 15 Intelligence (Arcana or Nature) check determines that while it looks like the damage to the building was caused by an earthquake, the damage is too localized and specific. Highly concentrated thunder damage was the cause of the damage to the building.
\paragraph{Corpses.}
The three corpses in the front lawn belonged to members of Master Krumm's staff: a cook, a farmer, and a woodcutter. A successful DC 15 Wisdom (Medicine) check determines that the bodies where slain by claws and bites from a rather larger creature. On the body of the cook is a key, which opens the back door of Krumm's manor.
\paragraph{Survivor.}
Inside the house are more corpses like the ones in the front yard. They bear the same wounds as well, although some of the bodies were killed by thunder damage. Unconscious but stable, hidden by an overturned mattress, is a stable hand, a young human named Feder Estes. If roused from his current state, Feder speaks only Ostoyan and has little information to impart, save for the following:

\begin{itemize}
\item In the last two months, Krumm and his staff of about 20 people created this small settlement. Most of the staff is now dead in the servants' dwelling. Krumm paid well but insisted on secrecy, insisting that Feder tell no one where he was going or who he was serving.
\item They were attacked today by large wolf-like creatures lacking hair. Their howls shattered windows and buckled doors.
\item Krumm is probably in his manor house, assuming he survived the attack. The Master is a powerful arcanist who sought this remote place to perform his research.
\item [A successful DC 20 Charisma (Intimidation or Persuasion) check is needed to prompt the following information.] The staff catches glimpses of strange creatures (the Chapped Brute) moving through the woods, and in and out of the manor house late at night. Occasionally a scream is heard from the manor as well, but no one is brave enough to ask questions or investigate.
\end{itemize}

\subsection{3. Krumm's Estate}
\begin{DndReadAloud}{}
Unlike the other building, this one is still intact. The bodies of several creatures, looking like some sort of wolf-lizard hybrid, lie dead around the house's perimeter. There are no signs of life within or outside the house.
\end{DndReadAloud}

\begin{figure}[h]
\includegraphics[width=0.5\textwidth]{images/adventure/GHLoE/114-16-001.awakened-chapped-brute.png}
\end{figure}
Master Krumm (downcast apostate*) is locked within his home, waiting for the next barrage from the xakalonuses. He is unaware of the presence of the hraptnon nearby, although he knows of its mythical existence and fears that if it attacks, he and his bodyguards won't be able to withstand the assault. Those bodyguards are five Awakened Chapped Brute*, who currently hide in the woods behind the manor, awaiting the call of their master.
\paragraph{Krumm's Motivation.}
Master Krumm's rituals to regain his immortality failed for one simple reason: none of his sacrifices were powerful enough in spirit or magic to ignite the divine spark in him. With the characters, he sees several potential sparks. This leaves him in a quandary: he may need the characters to help hold off further attacks, but if he can kill one to complete the ritual, he can regain immortality, so defending against the attack won't matter.
\paragraph{The Manor.}
The manor is full of Krumm's personal items, as well as tomes on rituals that could return immortality to once-immortal creatures. Components for those rituals are stored throughout the manor and could be sold for 1,500 gp to the right buyer. In the shed near the garden, three bound woodcutters await their fate as sacrifices to Krumm's ritual.
\subsection{Hrapton's Advance}
\begin{figure}[h]
\includegraphics[width=0.5\textwidth]{images/adventure/GHLoE/115-16-002.hrapton.png}
\end{figure}
Whenever you feel is the best time to throw the entire situation into even more chaos, unleash the hraptnon*!
\begin{DndReadAloud}{}
A horrifying creature the size of a house smashes through the forest, uprooting trees and roaring out a challenge.
\end{DndReadAloud}

When the hraptnon arrives, its instinct is to target the creature wielding the strongest magical ability. That is likely Master Krumm, who's an 18th-level spellcaster. However, the creature may go after a more accessible target who is almost as powerful.
If the hraptnon is being easily handled by the characters, a squad of five Xakalonus enters the fray to support their progenitor. They can act as a buffer, or the hraptnon can devour them to regain hit points.
\section{Conclusion}
If what could turn into a three-way battle results in the characters' victory, they learn that the foresters' disappearance was caused by both the hraptnon and xakalonuses, as well as the machinations of Master Krumm. Bringing Krumm back to Oakendale for justice settles the affair in the eyes of Yvell, who gratefully provides the promised reward.
Unfortunately for the characters, unless they find the hraptnon's magical Stillborn Heart at the grove in the center of the Stillborn Forest, a new hraptnon appears within 24 hours of the original's demise. It is likely that the new hraptnon gathers a large force of xakalonuses, and maybe some other deadly forest creatures, and assaults Oakendale directly, hoping to gorge on the brains of any arcane characters.
\section{Creatures}

\begin{itemize}
\begin{multicols}{3}
\item Chapped Brute
\item Chapped Brute Abomination
\item Awakened Chapped Brute
\item Downcast Mercenary
\item Downcast Apostate
\item Hraptnon
\item Xakalonus
\end{multicols}

\end{itemize}

\chapter{17: Cave of Rising Sands}
By Shawn Merwin
Cave of Rising Sands is an hourglass widow lair for four or five 17th-level characters.
\section{Background}
Jacoty Endak traveled the length and breadth of Etharis. He sang songs, told tales, collected stories, and survived countless adventures. Jacoty thought he'd seen and experienced everything, until he fell madly in love with a chef named Hosken Kwilleau. The two moved to a remote region of the Castellan Province to live out their lives together, sharing their love with each other and those around them.
Hosken had to return to his hometown in the Charneault Kingdom to attend to family matters, but he disappeared after leaving. Jacoty did everything in his immense power to find his lost love, but it was all for naught. Jacoty spent the rest of his life pining for Hosken, wracked with guilt and anger. When Jacoty finally died of wasting grief, even that passing failed to reunite him with his love. Jacoty was reborn as an hourglass widow.
As an undead creature that despises hope, Jacoty moves throughout the countryside, snuffing out hope and joy wherever he finds it. He's recently found an opportunity to corrupt the hopes and dreams of a werebear ascetic named Zeragonia. Zeragonia found this region suffused with lycanthropy after a visit from the Great Beast, and she established a cave complex to harbor lycanthropes trying to learn to control their affliction. An acolyte of Miklas, Zeragonia calls upon that Arch Seraph's power to aide her in healing her charges, or at least helping them control themselves.
When Jacoty learned of Zeragonia's ministrations, he found a new focus for his evil intentions. He wormed his way into Zeragonia's cavern-home, and he is using his powers to subvert Zeragonia and corrupt her followers. He's already taken control of the lair and of those whom the werebear attempted to help, and he's almost completed the process of turning Zeragonia herself to the side of darkness.
\begin{figure}[h]
\includegraphics[width=0.5\textwidth]{images/adventure/GHLoE/119-17-001.hourglass-widow.png}
\end{figure}
\subsection{Set the Hook}
A sighting of the Great Beast in the area has led to a severe and dangerous increase in lycanthropy. While the Beast has vanished from the area, the citizens are left to deal with the aftermath. Lycanthropes ravage the area. To make matters worse, people in the region are disappearing. Others talk about experiencing lost time.
Scouts tracked some lycanthropes to a cave opening on the edges of civilization. The scouts who entered the cave never emerged, and one of the scouts left behind to watch the cave opening went mad, claiming that she'd been gone for a year when she was actually gone for no more than a couple of days. Local authorities beg the characters to investigate.
\section{Lair Overview}
Jacoty appears in different areas of the cavern complex, depending on how he has interacted with the characters in other areas. The boxed text describes the area and the occupants, minus Jacoty's possible presence. Adjust the boxed text as needed to describe the situation you're presenting.
\paragraph{Sealed Inside.}
Once the characters enter the mouth of the cave, Jacoty's magic takes effect, and the mouth of the cave disappears, trapping the characters inside. The following is true until the characters manage to destroy the shrine and statue of Miklas in Area 7:

\begin{itemize}
\item Teleportation and planar magic cannot move characters or other creatures into or out of the cave complex, although they can magically travel within the cavern complex.
\item As a bonus action, Jacoty can pop into or out of any location in his lair in the blink of an eye.
\item Mentioning the name "Hosken" to Jacoty has an obvious effect on the hourglass widow*. When that name is mentioned to him, Jacoty is stunned until the end of this next turn with the pain of his loss. In the next turn after the stun wears off, he must physically attack the person who said the name, eschewing any tactical or magical strategy.
\item As he can affect the passage of time within his lair, Jacoty can take one long rest at any point while the characters are trapped within the caverns in less than a minute's time. He may only do this once, however. Alternately, if the characters attempt to take a long rest, Jacoty teleports in to make sure they cannot get a full rest.
\end{itemize}

\DndQuote%
{}
{''The stench of this room assaults the senses.''}
{}
\subsection{1. Midden}
\begin{DndReadAloud}{}
The stench of this room assaults the senses. Three monstrous beings hurl flayed corpses down a deep hole, while others sit at a table feasting on raw, bloody flesh.
\end{DndReadAloud}

Seven fzeglaich* occupy this room, throwing the remains of their last meals down the refuse hole. They attack immediately upon seeing intruders.
\paragraph{Refuse Hole.}
The stench coming out of the hole is so terrible that it has taken on magical properties. Living creatures (aside from the fzeglaichs) who start their turn within 20 feet of the refuse hole must succeed on a DC 20 Constitution saving throw or be stunned until the start of their next turn.
\paragraph{Jacoty's Presence.}
Early in this combat, Jacoty (hourglass widow*) teleports into the area and spends a round or two attacking the characters, telling them that they are never going to leave this lair alive. After testing their mettle and taking in their tactics and capabilities, he pops away to deal with them later elsewhere in the cavern lair. He repeats this process in other areas, hounding the characters and then disappearing before he is seriously injured.
\subsection{2. Tainted Pool}
\begin{DndReadAloud}{}
This chamber contains a large pool filled with clear water. The basin is about ten feet deep, fed by tiny cracks in the bottom that allow the water to bubble up.
\end{DndReadAloud}

No creatures currently occupy the chamber, although if the characters wait here long enough, a few of the residents of the lair arrive to take a drink.
\paragraph{Pool.}
The pool contains fresh water that the residents of the lair use for survival. Jacoty's presence has caused the pool to take on corrupted characteristics in certain circumstances.
If the characters drink the water without using magic to test its potability first, it's very refreshing: provide them with inspiration. A successful DC 20 Intelligence (Nature) or Wisdom (Survival) check determines that the water is clean and pure.
If the characters use magic on the water, even if just to check for magic on it, the magic taints the water. Any spells determine that the water is fine, but the magic taints the water going forward. Characters that drink the water once it's tainted must succeed on a DC 20 Charisma saving throw or be unable to regain hit points through magical means for 24 hours.
\subsection{3. Ill-Gotten Gains}
\begin{DndReadAloud}{}
Chests, barrels, and sacks line the walls of this chamber. A writing desk and a table also occupy the room, covered with papers and scrolls.
\end{DndReadAloud}

Treasure taken from the bodies of victims rest here. If the characters take at least an hour to sort through the collected goods, they find the following:

\begin{itemize}
\item 750 gp worth of coins, gems, jewelry, and other goods
\item 3 Potion of Supreme Healing
\item A spell scroll of heartseeker.
\end{itemize}

\paragraph{Writing Desk.}
The writing desk contains poems, notes, half-completed stories, and sketches. All of these were written by Jacoty, although he signed none of the works. The documents all convey a sense of loss, and the subject of the writings is a person named Hosken, who is apparently dead, although none of the work describes exactly who he was or how he died. Whoever wrote and drew these works is obviously heartbroken by the loss of Hosken.
\DndQuote%
{}
{''Treasures taken from the bodies of victims rest here.''}
{}
\subsection{4. Larder}
\begin{DndReadAloud}{}
A floor-to-ceiling stone barricade, formed from interlocking bricks, seals off a section of this cavern. Arcane runes decorate the bricks. Faint voices echo from the other side of the brick wall.
\end{DndReadAloud}

\begin{figure}[h]
\includegraphics[width=0.5\textwidth]{images/adventure/GHLoE/122-17-002.werewolf-ravager.png}
\end{figure}
Jacoty installed this magical brick wall to hold prisoners, used as food for the creatures who he's corrupted into becoming evil monsters. He'll occasionally come in and tell the prisoners that he's going to let them go, and then uses his magic to make them think they've been locked away for years as a way to taste their deliciously crushed hope.
\paragraph{Wall.}
The stone wall can magically expand or retract to block or allow access. A successful DC 15 Intelligence (Arcana) check reveals that the runes on the wall function as a lock. When a spell of level 1 or higher is cast on the wall, it opens for 10 seconds to allow passage. Then it closes, needing magic to re-open again. Magic cast on the wall is instantly absorbed, so spells that would affect the door do not. Each five-foot section of the wall has an AC of 25 and 200 hit points.
\paragraph{Prisoners.}
Currently six prisoners occupy the larder. Three are Scout who tracked the lycanthropes back to this lair. The other three are lycanthropes (Werewolf Ravager*) who refused to bow to the corruption of Jacoty. All of them are disoriented by lack of sustenance, as well as by Jacoty's time-distorting magic. Each can, however, describe Jacoty and call him "the devil with the hourglass body."
\subsection{5. Barracks}
\begin{DndReadAloud}{}
This three-tiered room contains living quarters.
Six vaguely lupine figures rest here.
\end{DndReadAloud}

These six Werewolf Ravager*, firmly under Jacoty's control, rest while waiting for further instructions. They attack immediately unless given other instructions by their new master.
\subsection{6. Study}
\begin{DndReadAloud}{}
This chamber holds furnishings that reveal it to be a living area for a single creature.
\end{DndReadAloud}

Before she fell under the sway of Jacoty, Zeragonia came here to rest and reflect. The area contains utilitarian furnishings, personal effects, and a journal.
\paragraph{Zeragonia's Journal.}
Written in code, the journal spans the last two years and talks about Zeragonia's personal struggle with her lycanthropy, how Miklas helped her overcome the affliction, and how she created this place to safely assist others who suffer from lycanthropy to control their more violent instincts. The last entry talks about the arrival of a stranger named Jacoty who seemed troubled. Zeragonia was interested in what those troubles might be, sensing a deeper darkness dwelling within him.
In order to decipher the journal, characters must either use magic or succeed on a DC 20 Wisdom (Insight) or Intelligence (Investigation) check.
\subsection{7. Shrine to Miklas}
\begin{DndReadAloud}{}
A statue and shrine to Miklas dominates this chamber. Three stone columns surround a dais facing the statue. An elderly woman with white streaks in her otherwise dark hair kneels before the altar, flanked by a pair of humanoid creatures with lupine forms.
\end{DndReadAloud}

The woman is Zeragonia, a werebear ascetic*, in her human form. The two flanking her are Werewolf Ravager* who've already succumbed to Jacoty's control. The ravagers attack as soon as the characters approach the shrine, but Zeragonia spends the first round of combat focused on the shrine. She attacks during the next round unless the characters find a way to help her.
\paragraph{Zeragonia's Plight.}
If any of the characters focus on seeing what's happening with Zeragonia, they can attempt a DC 15 Wisdom (Insight) check (no action required). On a success, they see that the woman is highly distressed and confused as she stares at the statue of Miklas.
\paragraph{Shrine.}
Jacoty corrupted the shrine to slowly turn the creatures in the cavern from those who wanted to be peaceful into killers. Zeragonia is the last holdout, but she is about to be turned to the dark side. The shrine radiates strong magic, and a successful DC 20 Intelligence (Arcana or Religion) reveals that the shrine is corrupting both the attitudes of the creatures in the area, as well as the time inconsistencies that have plagued the region.
\paragraph{Deactivating the Shrine.}
To deactivate the shrine, the characters must perform these actions, in order. Unless otherwise noted, each of these steps requires anyone attempting them to use an action to perform the step:

\begin{itemize}
\item Figure out what magic is affecting the shrine (no action, noted above). This reveals the rest of the steps as well.
\item Succeed on a DC 20 Intelligence (Arcana) check to begin to sever the connection to the dark forces.
\item Succeed on a DC 20 Dexterity (Sleight of Hand) check to undo the subtle physical changes that Jacoty made to the statue.
\item Succeed on a DC 20 Intelligence (Religion) check to reconsecrate the statue to Miklas.
\end{itemize}

Failing the checks in steps 2 - 4 by 5 or more results in the statue lashing out and doing 8d8 force damage to the person who failed the check.
Once the final step is completed, all of the magic of the lair stops working, making Jacoty less powerful and more susceptible to defeat by the characters. In addition, Zeragonia regains her senses and can be convinced into helping the characters fight the rest of the lycanthropes and Jacoty.
\section{Conclusion}
If the characters can rescue and redeem Zeragonia rather than killing her, Miklas smiles upon them for showing mercy. The statue speaks, and promises to answer three questions for the characters, or provide them with one wish (as per the wish spell).
If the characters defeat Jacoty, the hourglass falls out of his chest, spilling the sand that begins to flow upward. As the spirit of the poor man departs, it whispers a thanks for the release, and then asks the characters to find the fate of his love Hosken. But that's another story for another time.
\DndQuote%
{}
{''An elderly woman with white streaks in her otherwise dark hair kneels before the altar, flanked by a pair of humanoid creatures with lupine forms.''}
{}
\section{Creatures}

\begin{itemize}
\begin{multicols}{3}
\item Fzeg
\item Fzeglaich
\item Hourglass Widow
\item Werewolf Ravager
\item Werebear Ascetic
\end{multicols}

\end{itemize}

\chapter{18: Mayhem at the Crossroads}
By Shawn Merwin
Mayhem at the Crossroads is a caprathorn lair for four or five 18th-level characters.
\section{Background}
Floss's Cross was a prosperous hamlet set at the crossroads of two important trade routes. The people of Floss's Cross worked hard to provide the best food, drink, and services to the many travelers who passed through their village, and they grew accustomed to the prosperity that came with a respected reputation. The place earned a high standing among even the most powerful people in Etharis, making it a place where political, religious, social, and economic leaders would meet in secret to discuss, negotiate, barter, entreat, and strategize.
Unfortunately, the hamlet's success and reputation also drew some unwanted attention. An avarice seraph called Rusk and a gluttony seraph called Yort realized that by infiltrating the town, they might be able to corrupt some of the very powerful people who visited Floss's Cross regularly. The pair, along with their lesser disavowed minions, have turned the place into a ticking time bomb.
They found a way to attract a caprathorn to the area. The potential devastation a caprathorn might cause gives the pair of disavowed leverage over the rich and powerful visitors to the town, and it would certainly bring powerful people to deal with it. Ruck and Yort also hold a powerful hostage in town, and they plan to either kill or corrupt the heroes who come to investigate.
\subsection{Set the Hook}
A powerful figure in Etharis (or your own world) came to Floss's Cross. Neither they, their bodyguards, nor other members of their retinue have been heard from since they first arrived. The disappearance of such a powerful figure has many people worried. Worse yet, rumors that a dreadful minion of the Great Beast has been hunting in the area stoke those worries. The characters are asked to find out what happened to this powerful figure.
\section{Lair Overview}
The provided map shows only the center of the hamlet of Floss's Cross. Other shops, businesses, taverns, and inns rest outside of the hamlet's center. However, all the danger and duplicity the characters find in Floss's Cross radiate out from this central area. Any building on the map that's not numbered is an empty residence—either fled by its normal occupants, or a building for rent that currently has no tenants. (Adding more interesting elements or enemies to those areas is recommended if you want to lengthen or increase the difficulty of the lair.)
The center of town has been attacked by the caprathorn once already, leaving some damage to the area. The population of the hamlet has locked themselves in their homes and businesses, waiting for the caprathorn to attack again. The people of the town are unaware that anything untoward is happening in and around the town square—at least in regards to the disavowed taking over and taking a hostage. They've been told to remain in their homes because of the threat of the caprathorn, but they think the problem is being handled.
The caprathorn and its minions presently haunt the forested area outside the hamlet. The caprathorn does not approach town until the magic of the Floss Memorial statue (see Area 1) is activated by Rusk or Yort. The two disavowed do that as soon as it becomes clear that the characters are too powerful for them and their lesser seraph allies to handle.
\subsection{1. Floss Memorial Statue}
\begin{DndReadAloud}{}
At the center of town, where the two well-traveled trade routes meet, stands a marble statue depicting a smiling woman, her arms spread wide as if casting a blessing over the entire area. A gold plaque on the base of the statue reads, "In Honor of Frannie Floss, Matron of Our Town, Honored by All Who Revere Home and Hearth." Unlit candles and small trinkets are placed on the statue and around its base.
\end{DndReadAloud}

\paragraph{Statue.}
The statue has been magically empowered by the disavowed seraphs Rusk and Yort to act as a beacon for the caprathorn. A successful DC 25 Intelligence (Arcana) check recognizes the summoning power of the magic infusing the statue. The magic of the statue cannot be dispelled with anything except a wish spell or powerful divine intervention.
A success DC 20 Intelligence (Investigation) reveals that none of the candles or trinkets have been lit or placed here in the last several days. Obviously, those who generally leave offerings in memory of Frannie Floss have failed to do so for a while.
\DndQuote%
{}
{''Rusk is attempting to charm Evi and gain a larger foothold in the hamlet.''}
{}
\subsection{2. Rusk's Residence}
The doors and windows to this building have been shuttered and locked. The locks are not strong, easily defeated with a successful DC 10 Dexterity (Thieves' Tools) check. Once inside, the characters can examine the interior.
\begin{DndReadAloud}{}
This small dwelling is nicely appointed with fine furnishings. Two women sit sipping wine and chatting in front of a fire.
\end{DndReadAloud}

Unless she is encountered elsewhere, Rusk (avarice seraph*) is here talking with Evi (noble), the proprietor of the Sideways Wayside, an establishment elsewhere in Floss's Cross. Rusk is attempting to charm Evi and gain a larger foothold in the hamlet.
Rusk realizes that the adventurers are quite powerful, so she does her best to put them at ease and answer their questions. She attempts to lie to them while still telling the truth. Her plan is to lead them to a place where she can get the assistance of the other seraphs in town, or even call the caprathorn to take out the characters.
\subsection{3. Private Tavern}
\begin{DndReadAloud}{}
This small tavern bustles with activity. An ogresh cook orders around several other workers and servers, preparing several plates of food. A party of six well-dressed diners sits at a table, sampling the delicious-smelling fare.
\end{DndReadAloud}

This private tavern caters to wealthy patrons looking for a unique meal cooked by a master chef residing in Floss's Cross. The chef, an ogresh named Hemina (noble), has been under suggestion spells from Yort, the gluttony seraph*, for the past several days. Under this magical compulsion, Hemina prepares and distributes food made with various herbs supplied by Yort, which have had magical effects on the people eating them.
\begin{figure}[h]
\includegraphics[width=0.5\textwidth]{images/adventure/GHLoE/128-18-004.gluttony-seraph.png}
\end{figure}
Hemina accepted the help supplied by Yort, in the form of seven Lesser Gluttony Seraph*, who are using magic to take the form of normal humanoids. A successful DC 20 Wisdom (Perception) check pierces their magical disguises and reveals something unusual about them.
\paragraph{Hemina.}
Between the magic and the innate charms of the disavowed, Hemina is muddled in her thinking. She strongly believes both Rusk and Yort are wonderful friends. Hemina remembers serving food to the powerful figure and their bodyguards, but she hasn't seen them in a few days. She assumes they, like many of the people in town, are staying hidden because of the threat of the caprathorn.
\paragraph{Diners.}
The dining party consists of six wealthy travelers (Noble) who couldn't turn down the chance to eat a meal at one of the premiere restaurants among the elite of Etharis. Unfortunately for them, the seraphs added poison to the meal, with the plan of killing them, stealing their coin and jewelry, and burning the bodies.
Not long after the characters enter the building, the diners start suffering the effects of the poison. If the characters attempt to help the diners, each needs a successful DC 10 Wisdom (Medicine) check to save them from dying. Magic that neutralizes poison also works. If the characters attempt to intercede, the seraphs attack them.
\paragraph{Pantry.}
If the characters search the pantry, they find many ingredients of a magical, rather than a culinary, nature. On a successful DC 20 Intelligence (Arcana or Nature) check, someone knows these ingredients, if ingested in the correct amounts and combinations, could cause anything from lethargy to intense pain to incapacitation to death. These ingredients are rare, and an alchemist specializing in rare ingredients and components would pay 750 gp for them.
\paragraph{Rewards.}
If the characters can save all six of the diners, the wealthy patrons give each character a reward of 500 gp, and they promise to owe the characters a favor (which you can determine what it means).
\subsection{4. Yort's Residence}
\begin{DndReadAloud}{}
A rotund man in a silk robe startles at your intrusion. "What's this all about," he blusters, as he falls out of his chair while reaching ineffectually for a dagger.
\end{DndReadAloud}

Yort, a gluttony seraph*, uses magic to disguise himself as human, pretending to be a doddering fool. A successful DC 20 Wisdom (Insight) check alerts a character that this person is acting like a fool but is very much in control of the situation and feigning surprise.
\paragraph{Ruse.}
Yort pretends to be a resident of the town, hiding out in his home. He knows better than to challenge the characters on his own, so he tries to get them to accompany him to an area where allies might help him. He also uses the statue to summon the caprathorn if the situation is dire.
\subsection{5. Hostage}
The doors and windows of this stable are boarded up, despite some structural damage. There are no locks to pick, but as this stable is not solidly built, a successful DC 20 Strength check can break open any door or window. If the characters gain access, they can see the situation.
\begin{DndReadAloud}{}
Four large angelic figures, their red eyes glowing and beautiful faces wearing a sneer, stand over a bound figure within this very clean stable.
\end{DndReadAloud}

Eight Lesser Avarice Seraph* occupy the stables. Only four are in the stables at a time, as they take turns watching over the powerful figure, whom they have kidnapped and plan to hold for ransom. The other four, however, are in the area and rush to help their associates if a battle breaks out. Their instructions are to make sure that their prisoner does not escape, and that no one snoops around and sees them. If a battle does begin, one of them rushes to retrieve Rusk.
\paragraph{Powerful Figure.}
This powerful figure can be an NPC from Etharis (if you are running a Grim Hollow campaign) or from your own world. They should have enough influence and power to make their survival something important to your campaign.
If the battle is going badly for the seraphs—which it most assuredly will at this level of play—one of the seraphs threatens the NPC if they characters do not back away. At this level, hostage situations are generally not an issue for characters, as they have magical capabilities that can bring people back to life with very little problem. If you want the hostage situation to present a true challenge, have the NPC hostage admit to the players that because of a family curse, no magic can bring them back to life once they're dead.
\subsection{The Caprathorn}
At whatever point you are ready to have the caprathorn make its appearance, it charges into the center of town.
\begin{DndReadAloud}{}
A huge goat-headed figure charges up the path toward the center of town. The four horns that adorn its head are covered with runes. Vines grow from the beast, and hanging from the ends of the vines are twisted-looking sheep. The lot are covered with blood.
\end{DndReadAloud}

In addition to the caprathorn*, nine Thornlamm* accompany their master. Driven to a frenzy by the call of the magic exuded by the statue, the caprathorn looks to kill anything within its line of sight.
\paragraph{Chaotic Battle.}
While the disavowed can summon the caprathorn, they cannot control it. This means that any battle involving all three sides (characters, disavowed, and caprathorn) can go in any direction. Many high-level groups could handle the battle even if the disavowed and caprathorn worked together, but regulating the action based on the success of the characters is highly recommended.
\section{Conclusion}
If the characters rescue the powerful figure being held in the stables before he is killed by his captors, that person rewards them with at least 1,000 gp worth of gems, coins, and other precious gifts—if the characters assist that person in getting home, of course.
If the characters rescue them but allow harm to come to them first, the reward is halved and the powerful figure acts much more stiffly around the characters, expecting that heroes of such power should have been able to mount a rescue more expertly than they did.
Regardless, the people of Floss's Cross are ecstatic that the characters not only eliminated the caprathorn, but also rooted out the corruption brought to the town by the disavowed. Being heroes of an area beloved by so many powerful people certainly could have its perks.
\section{Creatures}

\begin{itemize}
\begin{multicols}{3}
\item Caprathorn
\item Thornlamm
\item Lesser Avarice Seraph
\item Avarice Seraph
\item Lesser Gluttony Seraph
\item Gluttony Seraph
\end{multicols}

\end{itemize}

\chapter{19: Prison of Good Intent}
By Shawn Merwin
Prison of Good Intent is an angel of Empyreus lair for four or five 19th-level characters.
\section{Background}
The ancient elven vampire Aeilidrania terrorized communities around the world, gathering power, influence, information, and followers—and built a small undead fiefdom for herself. When the vampire was at the height of her power, an angel of Empyreus named Benitent swooped in and captured her. It dragged the powerful undead creature to a remote prison, far from any civilized areas that she might threaten.
Unable to kill Aeilidrania permanently and fearing she would simply escape if he destroyed her current form, Benitent secured the vampire in a magical dwelling on an island surrounded by an abjuration-powered moat. The vampire's many followers have scoured the world, looking for a way to rescue their absent master, but only a few have managed to find this area. They cannot, however, breach the moat and the bridge, so they hurl themselves ineffectively at the barrier, all the while trying to think of other ways to free their mistress.
Benitent is only able to keep this prison intact through the power of a mortal being. Fearing that bringing any living creature into the prison area might give the vampire someone to manipulate, the angel of Empyreus employed a blind and deaf caretaker named Habret to live on the island, taking care of pigs and living in a small house.
For her part, Aeilidrania has given up attempting to escape. She's formed a friendship with the old caretaker, sharing with him much of the knowledge and wisdom she's gathered across the many centuries of her unlife. She also has enacted plans to corrupt Benitent, tricking the single-focused angel of Empyreus into unconsciously doing her bidding even while she remains its prisoner.
It's into the precariously balanced situation that the characters thrust themselves.
\subsection{Set the Hook}
A world-threatening problem is rapidly hurtling toward the characters, and one piece of information they need to fight it rests firmly in the twisted but brilliant mind of the elven vampire Aeilidrania. After some research and investigation, they hear rumors that the vampire they seek lives in a secluded region far from any civilized areas. Whether or not they learn of the vampire's imprisonment by Benitent, the angel of Empyreus, depends on how much research they undertake and how many resources they use. With no other means of securing the information they need to deal with the incredible threat, the characters are forced to seek out Aeilidrania.
\section{Lair Overview}
Benitent has taken steps to protect the area from magical incursion. Attempts to teleport directly into the area fail, displacing the teleporting party into the wilderness a couple of miles from this site, forcing them to use more conventional means of transportation. Within the bounds of the magical moat, the following effects hold sway:

\begin{itemize}
\item Creatures (except Benitent and Habret) gain resistance to necrotic damage and vulnerability to radiant damage.
\item Magical flight is suppressed, although Benitent can still fly without interference. A creature who tries to fly using magic must succeed on a DC 25 Intelligence (Arcana) check. A failed check denies magical flight to that character until the start of its next turn.
\item Evil-aligned creatures cannot be brought back from death.
\end{itemize}

These effects end if either Benitent or Habret are slain.
\subsection{Lair Environs}
As the characters move closer and closer to the prison that Benitent created for Aeilidrania, they noticed the landscape turns from natural to unnatural in a noticeable way:
\begin{DndReadAloud}{}
The trek into the wilderness, far from any civilized area, brought you into contact with only small animals for the first part of your journey. The last mile, however, was rife with undead creatures that easily fell before you. Something must have drawn these vile creatures to this area.
\end{DndReadAloud}

A successful DC 20 Intelligence (Nature) check reveals that the change in the landscape has nothing to do with any large-scale magical effect in the area. It's simply a result of so many undead creatures coming to the area.
While not part of the lair specifically, you can run as many encounters as you want against undead foes as the characters travel toward the lair. Forcing the characters to use resources now means they have fewer resources later.
\subsection{1. Moat}
\begin{DndReadAloud}{}
A moat comprised of blue-grey water surrounds a plot of land that contains several large willow trees, a pigsty, and a small wooden dwelling. A hazy miasma rises from the moat, making it hard to discern further details about the contents of the island. A wooden bridge crosses the moat on the south side of the island.
\end{DndReadAloud}

When the characters approach, they notice a few lesser undead creatures hurling themselves into the miasma that rises from the moat's water. They instantly die when they touch it.
\paragraph{Magic Moat.}
The moat contains magic-infused water that acts to keep creatures from crossing it. The miasma acts as a curved wall of force which, if touched, does 66 (12d10) radiant damage (no save). The miasma seals at the top, 60 feet above the island, making the bridge the only means to enter safely. As mentioned earlier, teleportation, planar magic, or digging beneath the island fail to allow egress.
\subsection{2. Bridge}
\begin{DndReadAloud}{}
A wooden bridge crosses the moat. The highly polished wood of the bridge shines with a silver light, and small runes are etched into the planks.
\end{DndReadAloud}

A successful DC 25 Intelligence (Arcana) check reveals that the runes on the bridge empower it to deliver a disintegrate spell (Charisma saving throw instead of Dexterity) to anyone who gets halfway across the bridge. Good-aligned creatures are immune to the effect. Evil-aligned, undead, or fiendish creatures have disadvantage on the saving throw. Creatures don't have to be touching the bridge for this effect to occur. Creatures surviving the effect can access the island.
A successful DC 20 Intelligence (Investigation) check notices that dust has filled the cracks and knots of the wooden boards. Normal bird feathers rest among the dust. This is a clue that flying creatures are not immune to the disintegrating magic of the bridge.
\DndQuote%
{}
{''A hazy miasma rises from the moat, making it hard to discern further details about the contents of the island.''}
{}
\subsection{3. Sty}
\begin{DndReadAloud}{}
Several pigs wallow and root around in a fenced-in pigsty. Arcane runes decorate the fence.
\end{DndReadAloud}

\begin{figure}[h]
\includegraphics[width=0.5\textwidth]{images/adventure/GHLoE/135-19-001.empyrean-brazen-bull.png}
\end{figure}
Five of the pigs are Empyrean Brazen Bull* that the angel of Empyreus transformed to hide their presence.
If the characters attack Benitent or attempt to gain access to Aeilidrania, the brazen bulls take on their true form, break out of the pigsty, and attack immediately. A DC 20 Intelligence (Nature) check reveals that five of the pigs are not natural creatures.
\paragraph{Runed Fence.}
A successful DC 20 Intelligence (Arcana) check reveals that the runes on the fence power an ongoing magical effect. The runes siphon small amounts of the pig's life force and redirect it to fuel the power of the moat. Destroying the fence eliminates the moat-powered effects on the island within 24 hours.
\subsection{4. Tree}
\begin{DndReadAloud}{}
A large and voluminous willow tree stands out among its peers on the island. Silver runes adorn the trunk. A lead box sealed with silver chains rests at the base of the tree.
\end{DndReadAloud}

This tree has been blessed with the power of Empyreus, allowing it to hold the essence of Benitent. The Angel of Empyreus* can enter or exit the tree as a bonus action. Inside the lead box rests Popilo, a malikirian imp*.
\paragraph{Lead Box.}
A successful DC 15 Intelligence (Arcana) check reveals the purpose of the box: it's meant to hold a fiendish creature. The lock that holds the chains can be opened with a successful DC 15 Dexterity (Thieves' Tools) check, or the chain can be easily severed with a weapon attack.
\paragraph{Imp.}
Popilo, the Malikirian imp, serves Benitent. Its sole purpose is to tempt the characters to make a deal with it. Popilo offers to give the characters anything they request, particularly information, for a single piece of silver. Once the characters make any deal with Popilo, Benitent has the excuse needed to attack the characters for being evil and corrupt.
\DndQuote%
{}
{''Several pigs wallow and root around in a fenced-in pigsty. Arcane runes decorate the fence.''}
{}
\subsection{5. Dwelling}
\begin{DndReadAloud}{}
A simple wooden dwelling rests at the north side of the island. The building looks well built and sturdy, but plain. The windows are open, and the door stands ajar.
\end{DndReadAloud}

The dwelling houses Habret (a commoner), although Aeilidrania (an Elf Vampire (Ancient)*) spends her waking hours in the secret basement. The runes on the house can be examined successfully with a DC 15 Intelligence (Arcana) check. The runes gather life energy from around the island, bringing it to form a prison that keeps a specific undead creature within the dwelling.
\paragraph{Habret.}
Habret could be found anywhere on the island, as the old man spends every waking hour doing chores. Habret can sense the presence of the characters if they approach him, but he cannot hear or see them. He speaks to them, inviting them to share the meat, fruit, and vegetables that he's butchered, grown, or harvested.
If the characters use magic to communicate with Habret, he can convey the following information.

\begin{itemize}
\item A kind person who could talk in his mind asked Habret to come to this place and take care of it. He's been here for several months, and he takes care of the pigs, does chores, and otherwise just does what he was asked.
\item Not long after coming here, another person who could talk to his mind arrived. Habret talks to her every day, but he doesn't know where she lives. She just seems to appear and disappear at odd times. This person seems to be very learned and old, as she discusses things she's seen and done that indicate she's lived a long time and traveled the world.
\item Before taking on the task that the original person gave him, that person drew something on Habret's back. Habret isn't sure what it was, but it didn't hurt, and Habret has never felt better than he does now.
\end{itemize}

If the characters look at Habret's back, there's a large symbol of Empyreus drawn in silver ink there. A successful DC 15 Intelligence (Religion) check not only recognizes the symbol, but it also reveals that it is part of a complex network of energy that maintains all the magic of this island.
\begin{figure}[h]
\includegraphics[width=0.5\textwidth]{images/adventure/GHLoE/136-19-002.malkirian-imp.png}
\end{figure}
\paragraph{Secret Trapdoor.}
A successful DC 15 Intelligence (Investigation) or Wisdom (Perception) allows the characters to notice that the animal-skin rug on the floor covered a trapdoor. The door leads to a 10-foot-deep cellar, which houses Aeilidrania's prison.
\paragraph{Aeilidrania.}
The elven vampire spends most of her time in the basement, but when Habret is in the house, she comes up to talk with him. She can't leave the house because of the abjuration magic employed by Benitent. She also knows she cannot harm Habret, and she finds the old man curious and enjoys talking with him.
If the characters ask her to provide the information they seek, she lies and tells them she cannot provide it because of the magic of the island prison. If they could kill Habret for her, however, she would be able to leave—and she would be happy to provide that information. Just speaking with Aeilidrania is enough to bring Benitent to attack the characters.
\DndQuote%
{}
{''The elven vampire spends most of her time in the basement, but when Habret is in the house, she comes up to talk with him. She can't leave the house because of the abjuration magic employed by Benitent.''}
{}
\section{Conclusion}
It's likely that the characters must defeat Benitent to get the information they need from Aeilidrania; unfortunately, doing that releases all the abjurations that keep the vampire in and its minions out. A conscientious party would attempt to destroy the vampire instead of releasing her on the world.
If you're interested in giving the characters a way to get the information without having to kill everything on the island, here is an option: Since Aeilidrania does not see Habret as a threat, the vampire might have told the old man precisely the information that the characters are looking for. Habret would not have remembered the exact details, but by using magic, the characters might be able to retrieve that information.
\section{Creatures}

\begin{itemize}
\begin{multicols}{3}
\item Angel of Empyreus
\item Elf Vampire (Ancient)
\item Empyrean Brazen Bull
\item Malikirian Imp
\end{multicols}

\end{itemize}

\chapter{20: Top of the World}
By Shawn Merwin
Top of the World is Gegazol's lair for four or five 20th-level characters.
\section{Background}
The undead dragon known as Gegazol inhabits an infamous lair in the frozen north of Etharis, but it's also saddled with a corrupting link to a daemonic force that resides within its undead body. Gegazol is normally content to rest in its home for decades at a time before its bloodlust encourages it to gather an army and march upon the Valikan Clans.
The daemonic skull that infuses Gegazol with extra abilities has a sentience of its own, and it has plans. If Gegazol can rule the northern realms of Etharis, she can certainly rule the entirety of this world—and possibly other worlds as well. The daemonic entity has nudged Gegazol to leave her lair and take some steps to gain even more power, as well as look to spread that power into other realms.
An elemental pocket dimension floats through the universe, coming into contact with many worlds. The pocket dimension roils with elemental energy, and if someone can tap into these elemental wellsprings, they could learn to control the very elements themselves. And this is just what Gegazol's daemonic inner voice is pushing her to do.
Gegazol has spent months in this elemental pocket dimension, absorbing the elements and attuning to the magical obelisks.
\subsection{Set the Hook}
An enormous purple whirlpool of air and ether opens high above five of the largest cities of the world. While nothing has emerged from these whirlpools in the sky, the leaders of the world are naturally worried. The heroes of the world who've flown up and entered the whirlpool have not returned, except for the occasional corpses that plummet from the sky, scarred and maimed with elemental damage.
With the powers-that-be of the world running out of options, they turn to the most powerful and experienced adventurers in the world. They implore the characters to investigate any of these whirlpools. They can't promise the characters anything that the characters don't already have or can't obtain as level 20 characters, but use your imaginations to see what can tempt them.
If the characters fail to act in a timely matter, they wake up one morning to the horrible truth that every creature in the world is slowly being drained of power. (At sundown each evening, each creature loses one hit point from its maximum hit points, dying when their maximum reaches zero.) All signs point to the sources of this energy drain coming from the holes in the sky.
\section{Lair Overview}
This lair is not on this world, or really on any world. It's a realm between worlds, an elemental pocket dimension that Gegazol has taken the time to attune with. To defeat Gegazol and drive her away from this place and back to her own lair, they must first sever the connection between the undead dragon and the obelisks. Four smaller obelisks surround a larger central one.
\subsection{Lair Actions}
The lair actions for Gegazol in Grim Hollow: The Monster Grimoire represent her abilities when she is at her frozen lair in Etharis. She has attuned to the power of the elemental pockets dimension, giving her these new lair actions instead.
On initiative count 20 (losing initiative ties), Gegazol takes a lair action to cause one of the following effects. Gegazol can't use the same effect two rounds in a row.

\begin{itemize}
\item Geysers of elemental power burst from each of the obelisks. Each living creature not allied with Gegazol in the lair must attempt a DC 21 Constitution saving throw or take 27 (6d8) damage. The damage taken is of the type corresponding to the nearest obelisk. This damage ignores resistance or immunity.
\item Gegazol screeches, and the plateaus respond. Creatures in the lair must make a DC 21 Constitution saving throw. Those who fail and are in contact with the ground take 18 (4d8) force damage and are grappled (escape DC 21). On a success, the creature takes half the damage and isn't grappled. Those who are not in contact with the ground may succeed on a DC 21 Constitution saving throw or be pulled to the nearest spot on the ground.
\item Gegazol summons an elemental storm that affects the entire lair. The entire lair becomes difficult terrain for enemies of Gegazol. In addition, enemies of Gegazol must succeed on a DC 21 Charisma saving throw or gain a level of exhaustion until Gegazol uses a different lair action. While the storm is raging, Gegazol's fly speed doubles.
\end{itemize}

\subsection{Lich trolls, Plateaus, and Obelisks}
In her attempt to tame and control the elemental power of this pocket dimension, Gegazol has recruited four Lich Troll to her cause. Each of the lich trolls are spiritually connected to one of the four smaller elemental obelisks, while Gegazol is connected to the larger central obelisk.
\begin{figure}[h]
\includegraphics[width=0.5\textwidth]{images/adventure/GHLoE/143-20-002.lich-troll.png}
\end{figure}
The magic that binds all the creatures and elemental obelisks together is complex and intertwined. A lich troll cannot be killed until its associated elemental obelisk is powered down, and Gegazol cannot be killed until all four of the lich trolls are killed. This means none of the creatures can be reduced to less than 1 hit point.
The lich trolls cannot leave the plateau containing the obelisk they are connected to, but Gegazol is free to fly and move anywhere in the lair. While on the plateau with the obelisk that it's connected to, the spells it cast can be changed to do the corresponding damage instead of the spells normal damage. Additionally, when this happens, you should roll one extra damage die for the spell, and eliminate the low die. Each lich troll is also immune to the damage type of its associated obelisk. This immunity ends if the obelisk is deactivated.
If someone is pushed off any of the plateaus and have no way to fly or slow their descent, they fall thousands of feet to the ground of their home world.
\subsection{Bridges}
Rope bridges connect the plateaus to each other. These bridges aren't sturdy, but they hold up under heavy use. They are immune to elemental damage (fire, cold, thunder, lightning, force), but can be easily damaged with other types of damage. They have 20 hit points and an armor class of 15.
\section{Arrival}
As the characters move through the maelstrom in the sky, they are teleported to a random plateau. Roll 1d4 to see which plateau each character appears on.
\subsection{1. Thunder Obelisk}
\begin{DndReadAloud}{}
The obelisk on this plateau ripples with an internal energy that causes the monolith to quake and quiver. A troll-like creature standing near the base of the obelisk rumbles and trembles with barely contained fury.
\end{DndReadAloud}

The lich troll* connected to the thunder obelisk radiates tremors as it moves, and its voice rumbles with power.
\paragraph{Thunder Obelisk.}
A successful DC 20 Intelligence (Arcana or Nature) check reveals that this obelisk is associated with thunder. To deactivate the obelisk and sever it from connection to the lich troll and Gegazol, the characters must make the following steps. A successful DC 20 Intelligence (Arcana) check tells a character what these steps are, and what order they must be done in.

\begin{itemize}
\item Someone must touch the obelisk while casting a spell that does thunder damage, or attack the obelisk with thunder damage.
\item The obelisk is covered with runes, and one of the runes is associated with a divine entity representing thunder. A successful DC 20 Intelligence (Religion) notices this rune seems to be powering this obelisk.
\item If this rune is struck with a DC 20 Strength (Athletics) check, or 20 points of bludgeoning damage is done to that spot in a single blow, the obelisk cracks and loses its power. A character failing to break the obelisk on a single check takes 33 (6d10) thunder damage.
\end{itemize}

\subsection{2. Ice Obelisk}
\begin{DndReadAloud}{}
The obelisk on this plateau creaks and moans as an icy rime continually shifts on its surface. A troll-like creature covered with icy shards stalks the plateau.
\end{DndReadAloud}

The lich troll* connected to the ice obelisk is covered is sharp and frigid ice fragments, and frosty air puffs from its mouth when it speaks.
\paragraph{Ice Obelisk.}
A successful DC 20 Intelligence (Arcana or Nature) check reveals that this obelisk is associated with cold. To deactivate the obelisk and sever it from connection to the lich troll and Gegazol, the characters must make the following steps. A successful DC 20 Intelligence (Arcana) check tells a character what these steps are, and what order they must be done in.

\begin{itemize}
\item Someone must touch the obelisk while casting a spell that does cold damage, or attack the obelisk with cold damage.
\item A successful DC 20 Intelligence (Arcana) recognizes the pattern of icy veins running through the obelisk that carries the power through it.
\item A character who wants to depower the obelisk must succeed on a DC 20 Constitution saving throw to pull the icy veins apart. When that happen, the obelisk's icy surface melts and it loses its power. A creature failing the save takes 33 (6d10) cold damage and must let go before breaking the icy pillar.
\end{itemize}

\subsection{3. Flame Obelisk}
\begin{DndReadAloud}{}
The obelisk on this plateau is covered with small tongues of flickering flame. A troll-like creature covered with wearing cloaks of flame struts across the plateau.
\end{DndReadAloud}

The lich troll* connected to the flame obelisk is wreathed with flames, and noxious smoke escapes its lips when it speaks.
\paragraph{Flame Obelisk.}
A successful DC 20 Intelligence (Arcana or Nature) check reveals that this obelisk is associated with fire. To deactivate the obelisk and sever it from connection to the lich troll and Gegazol, the characters must make the following steps. A successful DC 20 Intelligence (Arcana) check tells a character what these steps are, and what order they must be done in.

\begin{itemize}
\item Someone must touch the obelisk while casting a spell that does fire damage, or attack the obelisk with fire damage.
\item A successful DC 20 Wisdom (Perception) checks reveals that there is a fire opal within the fiery obelisk that powers it.
\item If this gem can be plucked out from between the intense flames with a successful DC 20 Dexterity (Sleight of Hand) check, the obelisk's flames are snuffed and the magic of the obelisk ends. If this check is failed, the person attempting the checks takes 33 (6d10) fire damage.
\end{itemize}

\subsection{4. Lightning Obelisk}
\begin{DndReadAloud}{}
The obelisk on this plateau flickers as waves of bluish energy coruscate over its surface. A troll-like creature consumed in garments of static electricity howls in anticipation.
\end{DndReadAloud}

The lich troll* connected to the lightning obelisk crackles with blue static, and the wiry hair that covers its body stands on end.
\paragraph{Lightning Obelisk.}
A successful DC 20 Intelligence (Arcana or Nature) check reveals that this obelisk is associated with lightning. To deactivate the obelisk and sever it from connection to the lich troll and Gegazol, the characters must make the following steps. A successful DC 20 Intelligence (Arcana) check tells a character what these steps are, and what order they must be done in.

\begin{itemize}
\item Someone must touch the obelisk while casting a spell that does lightning damage, or attack the obelisk with lightning damage.
\item A successful DC 20 Intelligence (Investigation) finds a revolving metal rod that moves within the obelisk.
\item This rod can be grabbed and removed by a dexterous character who succeeds on a DC 20 Dexterity saving throw. Removing the rod breaks the magic of the obelisk. On a failed saving throw, the character takes 33 (6d10) lightning damage.
\end{itemize}

\subsection{5. Central Obelisk}
\begin{DndReadAloud}{}
An obelisk, larger and taller than the other four, dominates a plateau that rises 100 feet above the others. Dozens of corpses and cast off armor and weapons litter the ground all around this plateau. Perched atop the central pillar is an enormous draconic figure, but its body and wings show signs of rot and decay.
\end{DndReadAloud}

Gegazol* resides on her perch. As soon as the characters appear, she moves to attack whoever looks vulnerable.
\paragraph{Main Obelisk.}
This pillar is empowered by the smaller ones. As soon as the other pillars are deactivated, this one crumbles, making Gegazol able to be killed. Until that happens, no attack can reduce Gegazol to less than 1 hit point.
\begin{figure}[h]
\includegraphics[width=0.5\textwidth]{images/adventure/GHLoE/144-20-001.gegazol.png}
\end{figure}
\section{Conclusion}
Gegazol can't be killed permanently until the daemonic skull is destroyed, and since the skull is not here, Gegazol reforms back in her icy home. And she will not be pleased at what the characters did to her. Revenge becomes foremost on her mind.
\section{Creatures}

\begin{itemize}
\begin{multicols}{3}
\item Gegazol
\item Lich Troll
\item Zombie Troll
\end{multicols}

\end{itemize}

\appendix
\chapter{Magic and Salvage}
A monster hunter stalks through the overgrown forest, scanning the underbrush for signs of wild boar. By spreading a few silver coins at the local village's tavern to gain the trust and loosen the lips of reticent foresters, she learned of the potential location of a rare and dangerous gold mane boar.
The monster hunter can certainly use the meat from such a rare beast, but that's not her motive. The bristles of the gold mane boar can be suffused with magic and crafted into a light that pierces even the most powerful magical darkness. Adventurers traveling into the lairs of powerful fiends and spellcasters rely on such magic to gain an advantage.
\DndQuote%
{}
{''The most important tools a monster hunter can have are knowledge, skill in crafting, sufficient time, magical ingredients, and some valuable pieces of the monsters they've killed.''}
{}
\section{The Power of Salvage}
Each monster in Grim Hollow: Lairs of Etharis comes with a section that describes what can be salvaged from them—either as loot to be used or sold, or as components for powerful magic items to be created and used by adventurers. The salvage listed is only the first step, however. Creative DMs can use the examples provided to imagine and implement their own ideas of what can be salvaged from a monster, and more importantly, what can be created from that salvage.
In high-magic settings, one might imagine an economy built upon powerful crafters and arcanists creating magic items to be sold to adventurers for a tremendous profit. Magic in these settings would be easily available to all—and monsters, by extension, easier to defeat.
Grim Hollow is not one of those high-magic settings. The grim and dark fantasy supported by Grim Hollow eschews those notions of countless powerful creators making magic items willy-nilly, passing them out like candy at a midwinter festival. While magic is not necessarily rare in the setting, it's also not omnipresent and easy to obtain. And it's dangerous.
\section{Overcoming Incredible Odds}
A hallmark of dark fantasy is that the characters must rely on their wits, skill, luck, and choices to overcome challenges and achieve their goals. Survival is not a sure thing. Putting together the plans and tools for survival is an important part of the journey.
In such a world, nonplayer characters (NPCs) are not going to spend their valuable time and resources creating powerful magic items, on the off chance that some rich adventurers might come along with large amounts of gold to buy those goods. Gold can't buy safety in a world where the Great Beast may show up anytime, anywhere.
Part of the valor and resiliency of characters in a dark fantasy setting is their self-reliance—or of their dedication to finding, securing, and maintaining allies who can assist them in their struggles. That assistance is sometimes martial, but often it manifests as an ally with the ability to provide material support: a blacksmith, alchemist, armorer, sage, or some other artisan or expert.
\subsection{The Power of Survival in the Narrative}
Imagine a campaign where the characters are the only thing standing between a town and some terrible threat. The monsters threatening the area may seem much more powerful than the characters can handle, yet they have no choice but to try. Too many innocent lives are at stake.
When the characters first encounter the monsters, the heroes may barely survive, and fleeing is the only option to dying. Yet in that first encounter, the characters observe the monsters, maybe learning something about the monsters through successful ability checks or by locating clues in the area around the monsters' home.
Armed with that knowledge and the clues, the characters can rely on their own expertise—or the talents of the people they're protecting—to glean some information that may provide a spark of hope. For instance, maybe the monsters are vulnerable to a certain type of toxin.
The local alchemist knows of this very rare toxin, which can be distilled from the brain matter of a beast that lives in the nearby mountains. A local guide knows of that beast's lair.
It's easy to see a whole campaign, or at least a series of adventures, built naturally from this narrative construct. But more important than the plot set up by this is the dramas—big and small—that can emerge through play.
Maybe the alchemist is herself in danger from some other foe, so she requires the characters to take the time to solve her problem before she can make the necessary toxin. Maybe the local guide is a spy for the enemy, and instead of leading them to the lair, he leads the characters astray. Time factors may also force the characters to make hard decisions. All that narrative power comes when we, as game masters, don't simply hand the characters the answers to their problems.
\subsection{Keeping Friends Close}
In a setting where all the magic items and answers to problems must be worked for and crafted, NPCs become a much larger part of the campaign. A skilled weaponsmith with the necessary tools and training in arcane crafting is many times more valuable than any single magic weapon she could produce.
This heightens the characters (and their players) appreciation for, and attention to, the NPCs they encounter. NPCs are no longer cardboard cutouts that sit idly by, waiting for rescue from danger, while the heroes go about their business. The lives of these NPCs are much more important to the characters, and to the unfolding narrative of the campaign, when they have important information and assistance to offer.
NPCs, when they gain this increased importance in the success of the characters, enrich the campaign's setting. A calculating villain whose weak point can be exploited by the creations of an expert alchemist is much more likely to target the alchemist than the characters—and this brings a new dynamic to a campaign.
\section{Salvage and Crafting Rules}
The salvage rules for monsters includes not just what can be taken from the monsters; what can be crafted with those items is also stated.
Crafting these items takes into account four elements, in addition to the salvaged components:

\begin{itemize}
\item Crafting time
\item Necessary proficiencies or spells
\item Ability checks
\item Other components
\end{itemize}

Each of these elements is described in greater detail below. Remember: these elements are suggestions. Any or all of them can be changed to best fit your campaign.
\DndQuote%
{}
{''The monsters threatening the area may seem much more powerful than the characters can handle, yet they have no choice but to try.''}
{}
\subsection{Crafting Time}
The necessary time for item crafting is based on the type of item and its rarity. Items like weapons, armor, or clothing generally take longer than potions. However, powerful magic also adds complexity—and therefore time—to the crafting period. A simple potion might only take a few hours, while the most powerful arcane implement takes several days, weeks, or even months of constant work.
The length of time listed for crafting items from salvage passes whether the item is created successfully or not. The time assumes that the crafters work 16-hour days, and then take the mandatory 8 hours of a long rest.
Remember that during the time spent crafting, the crafters must still eat, rest, and take care of themselves. This might incur expenses renting space, buying food, etc. Those who aren't crafting can, of course, take steps to provide those resources.
\subsection{New Feats}
If crafting is going to play an important role in your campaign, providing characters (or even NPCs) with rules-based abilities to perform crafting might be just as valuable as combat-based abilities. For example, offering feats like the ones below might be valuable:
\subsubsection{Adroit Crafter}
You've trained in crafting items using a skill or tool that you're proficient with, and the steps in crafting have become second nature to you. You gain the following benefits:

\begin{itemize}
\item The time it takes to create an item using your proficient skills or tools is halved.
\item If you fail to successfully craft the item and would lose the components, you retain the components instead.
\end{itemize}

\subsubsection{Careful Crafter}
You've trained in crafting items using a skill or tool that you're proficient with, and your creations are known for being well made and reliable. You gain the following benefits:

\begin{itemize}
\item You make all ability checks for crafting with advantage.
\item Anyone assisting you in crafting an item gains advantage on their ability checks for crafting.
\item If you fail to successfully craft the item, the amount of time it took to complete the process is halved.
\end{itemize}

\DndQuote%
{}
{''The time assumes that the crafters work 16-hour days...''}
{}
\subsection{Necessary Proficiencies or Spells}
Crafting items sometimes requires a character or NPC with a certain proficiency to take part in the crafting. A person proficient with weaponsmith's tools must create the sword or mace if a magical weapon is being crafted. This proficiency is separate from the ability check that must be successfully made to craft the item.
For example, creating arcane oil from a xakalonus requires a successful DC 15 Intelligence (Arcana) check by a proficient alchemist. If a single person meets both conditions (proficient with alchemist's supplies AND proficient in the Arcana skill), they can craft the oil by themselves. However, two people can also work together if each is proficient in only one of the necessary skills or tools. In this case, they must each spend the entire time required for crafting.
\subsection{Ability Checks}
The ability check associated with a crafting should be made at the end of the crafting period. If the check is a success, the item is successful crafted. On a failure, the item is not crafted successful, and all the components that went into the crafting are consumed (unless otherwise stated).
Only people who take part in the entire process, and who are proficient in the skill called for, can attempt this check. If multiple people involved are proficient with the skill, they can roll individually, or they can assist someone else, giving them advantage on the check.
Magic that provides bonuses for a limited time—such as guidance—cannot be used to affect the ability check unless it is active for the entirety of the crafting process.
In some cases, rather than calling for a certain skill or ability check, the crafting of an item requires repeated casting of a particular spell. In these cases, someone must be available and willing to take part in the process for at least 1 hour of the day that the crafting takes place. This spell might be provided from a scroll, magic item, NPC, or some other source. Spells provided by a source other than a character might require payment, which could be in the form of coin, favors, or whatever else you can imagine that makes for a memorable story.
\subsubsection{Alternative Failure System}
Some players may bristle at a system where one failed check might have great repercussions in terms of time and components lost. If this is the case, a graduated failure system like the one below can take away some of the sting of failure.
1-3 lower than DC: No components are lost, and only half the time of crafting is spent.
4-6 lower than DC: Half the components are lost, and 75\% of the time are spent.
7-9 lower than DC: All components and time is lost.
10+ lower than DC: All components and time are lost, and all participants gain 2 levels of exhaustion (or alternative arcane mishap occurs).
\section{New Magic Items}
The following magic items can be crafted from salvage taken from monsters found in various Grim Hollow Lairs. The process for creating them is detailed in the salvage section of the monster entry.
You can, of course, simply present these items to your players as treasure taken from enemies—or even gifted or sold by NPCs or other elements of your game. However, presenting the items rather than encouraging the players (or their NPC associates) to craft them takes away from the objective of showing how valuable skills and proficiencies are in Etharis.
\subsection{Arcane Oil}
Potion, rare
This substance is created by distilling the jelly that encases the brain of a xakalonus. As an action, you can coat one weapon or 20 pieces of ammunition with the substance.
If the coated objects are not magical, for the next minute, your attacks with the coated weapon or ammunition gain a +2 bonus to attack and damage rolls.
If the coated objects are magical, the oil does not affect attack or damage rolls with them. However, your attacks against the hraptnon can fully damage that creature for one minute.
\subsection{Brazen Armor}
Armor (plate), very rare (requires attunement)
This armor, fashioned from the body of an Empyrean brazen bull, gleams brightly in direct sunlight. While wearing this armor, you gain a +1 bonus to AC. In addition, you have resistance to fire and radiant damage.
\subsection{Coat of Lies}
Wondrous item, rare (requires attunement by a spellcaster)
This coat made of the severed tongues of a Harvester of Lies's victims provides a +1 bonus to your AC, and you have advantage on Charisma (Deception) checks.
Additionally, as a bonus action three times per day, you can activate the magic of the cloak. If you do, the next spell you cast before the start of your next turn can be cast without any somatic or verbal components.
\subsection{Glass-Studded Armor}
Armor (studded leather), rare (requires attunement)
While wearing this armor, you have resistance to radiant damage. Each time a creature makes a melee attack against you, it takes 3 piercing damage. A creature can choose to make an attack with disadvantage to avoid this damage.
\subsection{Knifewing Cape}
Wondrous item, rare (requires attunement)
While wearing this cloak made from the skin of 20 Knifewing, you gain a +1 bonus to AC. In addition, when you fall while wearing this cloak, you descend 60 feet per round and take no damage from falling.
\subsection{Lindwyrm Venom}
Potion (poison), very rare
This terrible venom is extracted and distilled from the deadly lindwyrm. It can be delivered via ingestion or a wound. A creature exposed to the venom must succeed on a DC 18 Constitution saving throw or take 36 (8d8) poison damage and be poisoned for 1 minute. While poisoned, the creature has disadvantage on Strength and Dexterity skill checks and saving throws.
\subsection{Lycanthropy Antidote}
Potion, rare
This silvery potion, made with fzeg blood, is dotted with flecks of red. Drinking this liquid removes the curse of werewolf lycanthropy.
\subsection{Lycan Weapon}
Weapon (any), rare
Made from the claws, teeth, and bones of lycanthropes, these weapons are particularly effective against other lycanthropes or vampires. You gain a +2 bonus to attack and damage rolls made with this weapon.
Additionally, on a successful attack against a lycanthrope, the target is forced back into its human form unless it succeeds on a DC 10 Charisma saving throw. Also, damage done to vampires with lycan weapons cannot be regenerated until the creature finishes a short rest.
\subsection{Ready Gunk}
Wondrous item, rare
This sticky paste is fashioned from the powdered bone armor of a bone trader. When you treat a weapon, shield, or other item that can be wielded in one or two hands, you can use a bonus action to summon and equip it, as long as the item is within 30 feet of you. After using the magic of the paste three times, it loses its potency.
\subsection{Wand of Silence}
Wand, uncommon (requires attunement by a spellcaster)
This wand has 7 charges. While holding it, you can use an action to expend 1 charge to cast the silence spell (save DC 15). The wand regains 1d6 + 1 expended charges daily at dawn. If you expend the wand's last charge, roll a d20. On a 1, the wand crumbles into ashes and is destroyed.
\section{New Spells}

\begin{itemize}
\begin{multicols}{3}
\item Arboreal Curse
\item Consume Mind
\item Heartseeker
\item Hunter Sense
\item Magic Mirror
\item Suffocate
\item Wrack
\end{multicols}

\end{itemize}

\chapter{Etharis Creatures}

\begin{itemize}
\begin{multicols}{3}
\item Ancient Corpse Walker
\item Angel of Empyreus
\item Avarice Seraph
\item Awakened Chapped Brute
\item Beast Gnoll
\item Bloodbonded
\item Bone Trader
\item Candlelight Daemon
\item Caprathorn
\item Chapped Brute Abomination
\item Chapped Brute
\item Corpse Walker
\item Corpsejaw
\item Dawndrinker
\item Doomcaller
\item Downcast Apostate
\item Downcast Mercenary
\item Eldritch Herald
\item Eldritch Priest
\item Elf Vampire (Ancient)
\item Empyrean Brazen Bull
\item Eye Crawler
\item Eye Crow
\item Faevlin
\item Fzeglaich
\item Fzeg
\item Gegazol
\item Gluttony Seraph
\item Gnoll Brute
\item Harvester of Lies
\item Horror Flit Hunter
\item Hourglass Widow
\item Hraptnon
\item Infernal Tormentor
\item Ithjar
\item Knifewing
\item Lenchtahg
\item Lesser Avarice Seraph
\item Lesser Gluttony Seraph
\item Lich Troll
\item Lindwyrm
\item Lupilisk Elder
\item Lupilisk Whelp
\item Lupilisk
\item Malikirian Imp
\item Mjork Asher
\item Mjork Burner
\item Mjork Charger
\item Mjork Sootling Swarm
\item Mjork Sootling
\item Mjork
\item Mold Spider
\item Morbus Kobold
\item Oblivion Brute
\item Oblivion Juggernaut
\item Oblivion Leaper
\item Oblivion Whistler
\item Oozing Vulture
\item Ordeal Tree
\item Panjaian Ilharan
\item Psychic Fragment Swarm
\item Psychic Fragment
\item Scream Thief
\item Sea Drake
\item Shadowsteel Ghast
\item Shadowsteel Ghoul
\item Shatter Corpse
\item Shieldhead
\item Sitri Cat
\item Skeleton Cannoneer
\item Skeleton Commander
\item Skeleton Rifler
\item Skinweaver
\item Sky Drake
\item Sloth Galloper
\item Spythronar Sac
\item Spythronar Swarm
\item Spythronar Web
\item Stormborn Ithjar
\item Thornlamm
\item Torcheater
\item Venomous Gnoll
\item Vitebriate
\item Werebear Ascetic
\item Werewolf Ravager
\item Xakalonus
\item Young Blightscale Dragon
\item Zombie Troll
\end{multicols}

\end{itemize}

\chapter{Credits}

\begin{description}[nolistsep]
\item[Lead Designer.]
Shawn Merwin
\item[Designer Team.]
Thomas "Evhelm" Donovan, Greg Marks
\item[Lead Developer.]
Shawn Merwin
\item[Editor.]
Robert M. Everson
\item[Lairs Design.]
Thomas "Evhelm" Donovan, Greg Marks
\item[Production Assistant.]
Aurora Merwin
\item[Art Director.]
Suzanne Helmigh
\item[Graphic Design.]
Martin Hughes
\item[Interior Illustrators.]
Anastassia Grigorieva, Andreia Ugrai, Brent Hollowell, Damien Mammoliti, Dan Watson, Emily Ormerod, Helge C. Balzer, Josiah Cameron, Klaher Baklaher, Kurt Jakobi, Lucas Torquato, Marius Bota, Ona Kristensen, Quintin Gleim, Samuel Keiser, Suzanne Helmigh
\item[Lairs Cartographer.]
Joshua Orchard
\item[Marketing.]
Ian Gratton, Kathryn Griggs, Tyler Kempthorne
\item[Products.]
Matt Witbreuk, Simon Sherry, Nick Ingamells
\item[Playtesting.]
All the amazing people who took the time to playtest and provide feedback
\item[Special Thanks.]
To Wizards of the Coast and all the employees within. Thank you for the fantastic work
\end{description}

\subsection{Ghostfire Gaming Operations}

\begin{description}[nolistsep]
\item[Managing Director.]
Matt Witbreuk
\item[Financial Controller.]
James Atkins
\item[Strategy & Communications Manager.]
Hannah Peart
\item[General Manager of Operations.]
Nick Ingamells
\item[Media Content Manager.]
Ben Byrne
\item[Media Editor.]
Dante Szabo
\item[Ghostfire Gaming Discord Community Managers.]
Ian "Butters" Gratton, Nelson "Deathven" Di Carlo, Tom "Viking Walrus" Garland, Caleb Connendarf " Englehart, Cameron "C4Burgers" Brechin
\end{description}

\subsection{Ghostfire Production Studio}

\begin{description}[nolistsep]
\item[Head of Production.]
Simon Sherry
\item[Lead Producer.]
Joe Raso
\item[Producer.]
Kerstin Evans
\item[Principal Graphic Designer.]
Martin Hughes
\item[Graphic Designer.]
Josh Orchard
\item[Principal Art Director.]
Suzanne Helmigh
\item[Art Direction Team.]
Marius Bota, Ona Kristensen
\item[Lead Developer.]
Mark McIntyre
\item[Lead Game Designer.]
Shawn Merwin
\item[Senior Game Designer.]
James J. Haeck
\end{description}

\chapter{Creatures}

\begin{DndMonster}[float*=t]{Ancient Corpse Walker}
\DndMonsterType{Large undead (shapechanger), chaotic evil}
\DndMonsterBasics[
armor-class = {16 (natural armor)},
hit-points = {\DndDice{13d10 + 52}},
speed = {40 ft.}
]
\DndMonsterAbilityScores[
 str = 19,
 dex = 15,
 con = 19,
 int = 10,
 wis = 12,
 cha = 9
]
\DndMonsterDetails[
saving-throws = {Con +8, Wis +5},
skills = {Perception +5},
damage-resistances = {; bludgeoning, piercing, slashing from nonmagical attacks that aren't silvered},
damage-immunities = {necrotic, poison},
senses = {darkvision 120 ft., passive perception 15},
languages = {languages it knew in life},
challenge = 10,
]
\DndMonsterAction{Cursed Gaze}
A humanoid that starts its turn within 30 feet of and able to see the ancient corpse walker must succeed on a DC 16 Charisma saving throw, provided the corpse walker isn't incapacitated. On a failure, the humanoid is frightened for 1 minute. A frightened target can repeat the saving throw at the end of each of its turns, with disadvantage if not averting its eyes from the corpse walker, ending the effect on itself on a success. A creature failing two of these saving throws becomes cursed with disturbing thoughts. Such a creature can't benefit from a short or long rest until the curse is removed. On a successful save, the target is immune to any corpse walker's Cursed Gaze for the next 24 hours.
Unless surprised, a creature can avert its eyes to avoid the saving throw. If the creature does so, it can't see the corpse walker until the start of its next turn. A creature that looks at the corpse walker in the meantime must immediately make the save.
\DndMonsterAction{Spore Calling}
When the ancient corpse walker attacks, loyal undead within 120 feet of it come to its aid. Moderate or stronger wind can disperse these spores before they are effective.
\DndMonsterAction{Spore Cloud}
If the ancient corpse walker isn't in sunlight, it can use its action to polymorph into a Large cloud of spores or back into its true form. Anything it's wearing or carrying transforms with it. It reverts to its true form if it dies. In cloud form, the ancient corpse walker can't take any actions, speak, or manipulate objects. It's weightless, has a flying speed of 20 feet, can hover, and can enter a hostile creature's space and stop there. This cloud spreads around corners. If air can pass through a space, the cloud can do so without squeezing, but it can't pass through water. It has advantage on Strength, Dexterity, and Constitution saving throws, except those against wind. The cloud form is immune to nonmagical damage other than that from sunlight.
A creature that starts its turn in the cloud must make a DC 16 Constitution saving throw. If the save fails, the creature takes 13 3d8 poison and 10 (3d6) necrotic damage and becomes poisoned until the start of its next turn. On a successful save, a creature takes half the damage and isn't poisoned.
After 24 hours, corpses of beasts or humanoids that were in the spore area rise as zombies loyal to the ancient corpse walker. If such a corpse is that of an evil humanoid, it instead rises as a corpse walker loyal to the ancient corpse walker. These spores can animate a corpse only once.
\DndMonsterAction{Spore Fortitude}
If damage reduces the ancient corpse walker to 0 hit points, it must make a Constitution saving throw with a DC of 5 + the damage taken, unless the damage is radiant or from a critical hit. On a success, the corpse walker drops to 1 hit point instead.
When it drops to 0 hit points, the ancient corpse walker transforms into its Spore Cloud form. If it can't transform, it is destroyed. The spores of this cloud deal half their normal damage and can't animate corpses.
The ancient corpse walker can't transform back into its true form until it has at least 1 hit point. It must reach its preferred resting place within 2 hours or be destroyed. Once in its resting place, it reverts to its normal form.
It is then paralyzed and at 0 hit points for 1 hour, and then it regains 1 hit point and ceases being paralyzed.
If the ancient corpse walker takes damage while at 0 hit points, if able, it makes the same saving throw that would allow it to drop to 1 hit point. However, on a success, the ancient corpse walker remains at 0 hit points. If the ancient corpse walker fails or can't make this saving throw, it is destroyed.
\DndMonsterAction{Sunlight Hypersensitivity}
The ancient corpse walker takes 20 radiant damage when it starts its turn in sunlight. While in sunlight, it has disadvantage on attack rolls and ability checks.
\par
\DndMonsterSection{Actions}
\DndMonsterAction{Multiattack}
The ancient corpse walker makes two slam attacks.
\DndMonsterAction{Slam}
\textsl{Melee Weapon Attack:} +8 to hit, reach 5 ft., one target. \textsl{Hit:} 11 (2d6 + 4) bludgeoning damage and 10 (3d6) necrotic damage. If the target is a creature, it must succeed on a DC 15 Constitution saving throw, or its hit point maximum is reduced by an amount equal to the necrotic damage taken, and the ancient corpse walker regains hit points equal to that amount. This reduction lasts until the target finishes a long rest. The target dies if this effect reduces its hit point maximum to 0. A humanoid slain this way rises 24 hours later as a zombie, but it rises as a corpse walker if it was evil. These undead are loyal to the ancient corpse walker.
\end{DndMonster}



\begin{DndMonster}[float*=t]{Angel of Empyreus}
\DndMonsterType{Large celestial, lawful neutral}
\DndMonsterBasics[
armor-class = {21 (natural armor)},
hit-points = {\DndDice{33d10 + 264}},
speed = {40 ft., 90 ft. (key)}
]
\DndMonsterAbilityScores[
 str = 29,
 dex = 23,
 con = 26,
 int = 25,
 wis = 18,
 cha = 30
]
\DndMonsterDetails[
saving-throws = {Str +16, Int +14, Wis +11, Cha +17},
skills = {Arcana +14, Insight +11, Intimidation +17, Investigation +14, Perception +11, Religion +14},
damage-resistances = {fire, lightning, radiant; bludgeoning, piercing, slashing from nonmagical attacks},
damage-immunities = {cold, necrotic, poison},
senses = {truesight 120 ft., passive perception 21},
languages = {all, telepathy 120 ft.},
challenge = 22,
]
\DndMonsterAction{Celestial Arsenal}
The Angel of Empyreus' weapon attacks are magical. When it hits with any weapon, the weapon deals an extra 27 (6d8) radiant damage (included in the attack).
\DndMonsterAction{Authoritative Mien}
The Angel of Empyreus can cast the command spell at will (spell save DC 25).
\DndMonsterAction{Divine Awareness}
The Angel of Empyreus knows if it hears a lie.
\DndMonsterAction{Holy Aura}
The Angel of Empyreus can project a 60-ft. aura of holiness as a bonus action. The aura affects all creatures within its area unless the Angel of Empyreus chooses to exclude a creature from the effect. A creature remains under the aura's effect until it leaves the area, dies, or when the Angel of Empyrean chooses a new aura effect. Only one effect may be active at a time:

\begin{description}[nolistsep]
\item[Judgement.]
Creatures who start their turn in the aura must succeed on a DC 25 Charisma saving throw. On a failure, the creature gains a level of exhaustion as intensified feelings of guilt, doubt, remorse, and selfcriticism overwhelm the affected creature. Creatures who succeed on the saving throw are immune to the effect for one hour.
\item[Valor.]
Creatures who start their turn in the aura gain 15 temporary hit points, cannot be charmed or frightened, and have advantage on saving throws against spells and other magical effects
\end{description}

\par
\DndMonsterSection{Actions}
\DndMonsterAction{Multiattack}
The Angel of Empyreus can make three attacks with its holy avenger.
\DndMonsterAction{Holy Avenger}
\textsl{Melee Weapon Attack:} +24 to hit, reach 10 ft., one target. \textsl{Hit:} 23 (4d6 + 9) slashing damage plus 27 (6d8) radiant damage. If the target is a fiend or undead, it takes an additional 22 (4d10) radiant damage.
\DndMonsterAction{Blazing Radiance Recharge 5-6}
The Angel of Empyreus unleashes holy fire in a 60-foot cone. Each creature in that area must succeed on a DC 25 Dexterity saving throw or take 35 10d6 fire and 35 (10d6) radiant damage. Creatures take half damage on a success.
\par
\DndMonsterSection{Legendary Actions}
Angel of Empyreus can take 3 legendary actions, choosing from the options below. Only one legendary action can be used at a time and only at the end of another creature's turn. Angel of Empyreus regains spent legendary actions at the start of its turn.
\begin{DndMonsterLegendaryActions}
\DndMonsterLegendaryAction{Fly}{The Angel of Empyreus flies up to half its fly speed.}
\DndMonsterLegendaryAction{Sword}{The Angel of Empyreus makes a holy avenger attack.}
\DndMonsterLegendaryAction{Hurl Sword (2 actions)}{The Angel of Empyreus hurls its sword in a 30-foot line that is 5-feet wide. Each creature in that line must succeed on a DC 24 Dexterity saving throw or take 23 (4d6 + 9) slashing damage plus 27 (6d8) radiant damage. The sword immediately returns to its hand.}
\DndMonsterLegendaryAction{Recharge}{The Angel of Empyreus recharges its Blazing Radiance.}
\end{DndMonsterLegendaryActions}
\end{DndMonster}



\begin{DndMonster}[float*=t]{Archmage}
\DndMonsterType{Medium humanoid (any race), any alignment}
\DndMonsterBasics[
armor-class = {12 (15 with mage armor)},
hit-points = {\DndDice{18d8 + 18}},
speed = {30 ft.}
]
\DndMonsterAbilityScores[
 str = 10,
 dex = 14,
 con = 12,
 int = 20,
 wis = 15,
 cha = 16
]
\DndMonsterDetails[
saving-throws = {Int +9, Wis +6},
skills = {Arcana +13, History +13},
damage-resistances = {damage from spells; bludgeoning, piercing, slashing (from stoneskin)},
languages = {any six languages},
challenge = 12,
]
\DndMonsterAction{Magic Resistance}
The archmage has advantage on saving throws against spells and other magical effects.
\DndMonsterAction{Spellcasting}
The archmage is an 18th-level spellcaster. Its spellcasting ability is Intelligence (spell save DC 17, +9 to hit with spell attacks). The archmage can cast disguise self and invisibility at will and has the following wizard spells prepared:
\begin{DndMonsterSpells}
  \DndInnateSpellLevel{disguise self, invisibility}
  \DndMonsterSpellLevel{fire bolt, light, mage hand, prestidigitation, shocking grasp}
   \DndMonsterSpellLevel[1][4]{detect magic, identify, mage armor, magic missile}
   \DndMonsterSpellLevel[2][3]{detect thoughts, mirror image, misty step}
   \DndMonsterSpellLevel[3][3]{counterspell, fly, lightning bolt}
   \DndMonsterSpellLevel[4][3]{banishment, fire shield, stoneskin}
   \DndMonsterSpellLevel[5][3]{cone of cold, scrying, wall of force}
   \DndMonsterSpellLevel[6][1]{globe of invulnerability}
   \DndMonsterSpellLevel[7][1]{teleport}
   \DndMonsterSpellLevel[8][1]{mind blank}
\end{DndMonsterSpells}
\par
\DndMonsterSection{Actions}
\DndMonsterAction{Dagger}
\textsl{Melee or Ranged Weapon Attack:} +6 to hit, reach 5 ft. or range 20/60 ft., one target. \textsl{Hit:} 4 (1d4 + 2) piercing damage.
\end{DndMonster}



\begin{DndMonster}[float*=t]{Assassin}
\DndMonsterType{Medium humanoid (any race), any non-good alignment}
\DndMonsterBasics[
armor-class = {15 (studded leather armor)},
hit-points = {\DndDice{12d8 + 24}},
speed = {30 ft.}
]
\DndMonsterAbilityScores[
 str = 11,
 dex = 16,
 con = 14,
 int = 13,
 wis = 11,
 cha = 10
]
\DndMonsterDetails[
saving-throws = {Dex +6, Int +4},
skills = {Acrobatics +6, Deception +3, Perception +3, Stealth +9},
damage-resistances = {poison},
languages = {Thieves' cant plus any two languages},
challenge = 8,
]
\DndMonsterAction{Assassinate}
During its first turn, the assassin has advantage on attack rolls against any creature that hasn't taken a turn. Any hit the assassin scores against a Surprise creature is a critical hit.
\DndMonsterAction{Evasion}
If the assassin is subjected to an effect that allows it to make a Dexterity saving throw to take only half damage, the assassin instead takes no damage if it succeeds on the saving throw, and only half damage if it fails.
\DndMonsterAction{Sneak Attack (1/Turn)}
The assassin deals an extra 14 (4d6) damage when it hits a target with a weapon attack and has advantage on the attack roll, or when the target is within 5 feet of an ally of the assassin that isn't incapacitated and the assassin doesn't have disadvantage on the attack roll.
\par
\DndMonsterSection{Actions}
\DndMonsterAction{Multiattack}
The assassin makes two shortsword attacks.
\DndMonsterAction{Shortsword}
\textsl{Melee Weapon Attack:} +6 to hit, reach 5 ft., one target. \textsl{Hit:} 6 (1d6 + 3) piercing damage, and the target must make a DC 15 Constitution saving throw, taking 24 (7d6) poison damage on a failed save, or half as much damage on a successful one.
\DndMonsterAction{Light Crossbow}
\textsl{Ranged Weapon Attack:} +6 to hit, range 80/320 ft., one target. \textsl{Hit:} 7 (1d8 + 3) piercing damage, and the target must make a DC 15 Constitution saving throw, taking 24 (7d6) poison damage on a failed save, or half as much damage on a successful one.
\end{DndMonster}


\clearpage

\begin{DndMonster}[float*=t]{Avarice Seraph}
\DndMonsterType{Large celestial, neutral evil}
\DndMonsterBasics[
armor-class = {17 (natural armor)},
hit-points = {\DndDice{19d10 + 95}},
speed = {40 ft., 120 ft. (key)}
]
\DndMonsterAbilityScores[
 str = 20,
 dex = 20,
 con = 20,
 int = 14,
 wis = 17,
 cha = 20
]
\DndMonsterDetails[
saving-throws = {Con +10, Wis +8, Cha +10},
skills = {Deception +10, Insight +8, Perception +8},
damage-resistances = {radiant; bludgeoning, piercing, slashing from nonmagical attacks},
senses = {truesight 120 ft., passive perception 18},
languages = {all, telepathy 120 ft.},
challenge = 14,
]
\DndMonsterAction{Magic Resistance}
The seraph has advantage on saving throws against spells and other magical effects.
\DndMonsterAction{Seraphic Weapons}
The seraph's weapon attacks are magical. When the seraph hits with any weapon, the weapon deals an extra 4d8 radiant damage (included in the attack).
\DndMonsterAction{Innate Spellcasting}
The seraph's spellcasting ability is Charisma (spell save DC 18). The seraph can innately cast the following spells, using only verbal components:
\begin{DndMonsterSpells}
  \DndInnateSpellLevel{detect evil and good, detect magic, invisibility, light, minor illusion, thaumaturgy}
  \DndInnateSpellLevel[1]{zone of truth}
  \DndInnateSpellLevel[3e]{identify, suggestion}
\end{DndMonsterSpells}
\par
\DndMonsterSection{Actions}
\DndMonsterAction{Multiattack}
The seraph makes three attacks.
\DndMonsterAction{Scimitar}
\textsl{Melee Weapon Attack:} +10 to hit, reach 5 ft., one target. \textsl{Hit:} 12 (2d6 + 5) slashing damage, or 8 (1d6 + 5) slashing damage in Small or Medium form, and 18 (4d8) radiant damage.
\DndMonsterAction{Change Shape}
The seraph magically polymorphs into a Small or Medium humanoid, into a Large giant, or back into its true form. Other than its size, its statistics are the same in each form. It's clothing and scimitars change size to match its new form. If the seraph dies, it reverts to its true form, and its equipment to its normal size.
\DndMonsterAction{Voice of Avarice}
The seraph presents objects of desire, which can be illusions the seraph creates, to creatures of its choice. If those creatures are within 60 feet of the seraph, able to see it and the desirable objects, and able to hear the seraph, the creatures must make a DC 18 Wisdom saving throw or become charmed by the seraph until the seraph is incapacitated or fails to use a bonus action on its turn to continue to present and speak about the desirable objects.
While charmed in this way, a target can take no reactions and can use its action only to Dash. If the charmed target is more than 5 feet from the seraph, the target must move as far as it can toward the seraph, trying to come within 5 feet. The target avoids obvious danger, including opportunity attacks. If the target can't avoid such danger, the target can repeat the saving throw before entering the dangerous area. While an affected creature is within 5 feet of the seraph, the creature can't take reactions and spends its turns staring at the desirable objects.
The target also repeats the saving throw if the target is subjected to a harmful effect, if the target can't see or hear the seraph, and at the end of each of the target's turns. If a creature's saving throw is successful, the effect ends for it. A target that successfully saves is immune to the Voice of Avarice of this seraph or any lesser seraph for the next 24 hours.
\end{DndMonster}



\begin{DndMonster}[float*=t]{Awakened Chapped Brute}
\DndMonsterType{Large monstrosity, neutral evil}
\DndMonsterBasics[
armor-class = {14 (natural armor)},
hit-points = {\DndDice{12d10 + 60}},
speed = {40 ft.}
]
\DndMonsterAbilityScores[
 str = 20,
 dex = 10,
 con = 20,
 int = 5,
 wis = 10,
 cha = 5
]
\DndMonsterDetails[
skills = {Athletics +8},
damage-resistances = {necrotic},
senses = {blindsight 60 ft. (blind beyond this radius), passive perception 10},
languages = {understands one language but can't speak},
challenge = 7,
]
\DndMonsterAction{Keen Hearing and Smell}
The awakened chapped brute has advantage on Wisdom (Perception) checks that rely on hearing or smell. While deafened and unable to smell, the brute can't use its blindsight.
\DndMonsterAction{Sure-Footed}
The awakened chapped brute has advantage on Strength and Dexterity saving throws made against effects that would knock it prone.
\par
\DndMonsterSection{Actions}
\DndMonsterAction{Multiattack}
The awakened chapped brute makes two tentacle attacks or one tentacle slam.
\DndMonsterAction{Tentacle}
\textsl{Melee Weapon Attack:} +8 to hit, reach 15 ft., one target. \textsl{Hit:} 18 (3d8 + 5) bludgeoning damage. If the target is Medium or smaller, it is grappled (escape DC 16). Until this grapple ends, the target is restrained, and the brute can't use the tentacle against another target. The awakened chapped brute has two tentacles.
\DndMonsterAction{Tentacle Slam}
The awakened chapped brute slams one or two targets grappled by it into a solid object or each other. Each target must make a DC 16 Strength saving throw. If the saving throw fails, the target takes 18 (3d8 + 5) bludgeoning damage and is stunned until the end of the awakened chapped brute's next turn. On a successful saving throw, the target takes half the damage and isn't stunned.
\end{DndMonster}



\begin{DndMonster}[float*=t]{Beast Gnoll}
\DndMonsterType{Medium humanoid (gnoll), chaotic evil}
\DndMonsterBasics[
armor-class = {11},
hit-points = {\DndDice{2d8}},
speed = {30 ft.}
]
\DndMonsterAbilityScores[
 str = 13,
 dex = 12,
 con = 11,
 int = 6,
 wis = 10,
 cha = 7
]
\DndMonsterDetails[
senses = {darkvision 60 ft., passive perception 10},
languages = {Gnoll},
challenge = 1/8,
]
\DndMonsterAction{Rampage}
When the beast gnoll reduces a creature to 0 hit points with a melee attack on its turn, the gnoll can take a bonus action to move up to half its speed and make a bite attack.
\par
\DndMonsterSection{Actions}
\DndMonsterAction{Bite}
\textsl{Melee Weapon Attack:} +3 to hit, reach 5 ft., one target. \textsl{Hit:} 3 (1d4 + 1) piercing damage.
\DndMonsterAction{Claws}
\textsl{Melee Weapon Attack:} +3 to hit, reach 5 ft., one target. \textsl{Hit:} 4 (1d6 + 1) slashing damage.
\DndMonsterAction{Spine Spike}
\textsl{Ranged Weapon Attack:} +3 to hit, range 20/60 ft., one target. \textsl{Hit:} 3 (1d4 + 1) piercing damage.
\end{DndMonster}



\begin{DndMonster}[float*=t]{Bloodbonded}
\DndMonsterType{Medium monstrosity, chaotic evil}
\DndMonsterBasics[
armor-class = {16 (studded leather armor)},
hit-points = {\DndDice{12d8 + 12}},
speed = {{'number': 30, 'condition': '(or 10 ft. if missing a leg)'} ft.}
]
\DndMonsterAbilityScores[
 str = 12,
 dex = 15,
 con = 13,
 int = 15,
 wis = 7,
 cha = 18
]
\DndMonsterDetails[
saving-throws = {Con +3, Wis +0},
skills = {Deception +6, Insight +0, Persuasion +6},
damage-vulnerabilities = {psychic},
damage-resistances = {necrotic, radiant},
languages = {Celestial, Infernal, one national language},
challenge = 4,
]
\DndMonsterAction{Bloody Inspiration}
As a bonus action, the bloodbonded can inspire an ally within 30 feet to an enraged state of bloodlust for one minute. These followers gain 5 temporary hp, resistance to bludgeoning, piercing, and slashing damage, and vulnerability to psychic damage. All attacks made by the bloodbonded followers are reckless, gaining advantage to hit enemies but granting enemies advantage on attacks against them.
\DndMonsterAction{Resolute Followers}
Allies of the bloodbonded within 60 feet of it have advantage on Charisma and Wisdom saving throws.
\par
\DndMonsterSection{Actions}
\DndMonsterAction{Mace}
\textsl{Melee Weapon Attack:} +3 to hit, reach 5 ft., one target. \textsl{Hit:} 4 (1d6 + 1) bludgeoning damage.
\DndMonsterAction{Rage Bomb (2/Day)}
As a bonus action, the bloodbonded can infuse one of its bloodbonded followers who is under the effects of Bloody Inspiration with a fatal fury. This magic causes the follower to foam bloodily at the mouth. If the follower is not slain before the start of the bloodbonded's next turn, the follower dies and explodes in a shower of superheated blood. Creatures within 10 feet must make a DC 14 Dexterity saving throw, taking 12 (2d8 + 3) fire damage on a failure and taking half damage on a success.
\DndMonsterAction{Blood Boil}
The bloodbonded targets a creature it can see within 60 feet and utters words of madness. The target must succeed on a DC 14 Constitution saving throw or take 9 (2d8) fire damage as its blood superheats. The bloodbonded must be able to speak to use this ability.
\par
\DndMonsterSection{Reactions}
\DndMonsterAction{Loyal Followers}
When a bloodbonded would take damage from a melee or ranged attack it can see that targets only the bloodbonded, it can use its reaction to call on one of its followers within 5 feet to take the damage instead.
\end{DndMonster}


\clearpage

\begin{DndMonster}[float*=t]{Bone Trader}
\DndMonsterType{Medium fey, neutral evil}
\DndMonsterBasics[
armor-class = {12 (15 with bone armor)},
hit-points = {\DndDice{5d8}},
speed = {30 ft.}
]
\DndMonsterAbilityScores[
 str = 10,
 dex = 14,
 con = 11,
 int = 14,
 wis = 15,
 cha = 18
]
\DndMonsterDetails[
saving-throws = {Con +2},
skills = {Animal Handling +4, Perception +4},
senses = {darkvision 60 ft., passive perception 14},
languages = {Sylvan, two other languages},
challenge = 1,
]
\DndMonsterAction{Bone Armor (1/Day)}
A bone trader in its lair can use a bonus action to summon a temporary suit of bone armor that lasts for ten minutes or until dismissed by the bone trader. While wearing its bone armor, a bone trader has advantage on Wisdom saving throws and Charisma (Intimidation) checks.
\DndMonsterAction{Speak with Beasts}
The bone trader can communicate with beasts as if they shared a language.
\par
\DndMonsterSection{Actions}
\DndMonsterAction{Bone Club}
\textsl{Melee Weapon Attack:} +4 to hit, reach 5 ft., one creature. \textsl{Hit:} 6 (1d8 + 2) bludgeoning damage.
\DndMonsterAction{Skelesuck}
\textsl{Melee Weapon Attack:} +4 to hit, reach 5 ft., one creature. \textsl{Hit:} 4 (1d4 + 2) piercing damage. The target must succeed on a DC 10 Constitution saving throw or have one of its bones sucked out, resulting in the piercing damage becoming a reduction in maximum hit points until magical healing is applied.
\DndMonsterAction{Fey Charm}
The bone trader targets one humanoid, goblinoid, fey, or beast it can see within 30 feet of it. The target must succeed on a DC 14 Wisdom saving throw or be magically charmed. The charmed creature regards the bone trader as a trusted friend to be heeded and protected. Although the target isn't under the bone trader's control, it takes the bone trader's requests or actions in the most favorable way . Each time the bone trader or its allies do anything harmful to the target, it can repeat the saving throw, ending the effect on itself on a success. Otherwise, the effect lasts 24 hours or until the bone trader dies, is on a different plane of existence from the target, or ends the effect as a bonus action. If a target's saving throw is successful, the target is immune to the bone trader's fey charm for the next 24 hours.
The bone trader can have no more than one humanoid charmed at a time, along with up to three beasts.
\end{DndMonster}



\begin{DndMonster}[float*=t]{Candlelight Daemon}
\DndMonsterType{Medium fiend, neutral evil}
\DndMonsterBasics[
armor-class = {17 (natural armor)},
hit-points = {\DndDice{20d8 + 80}},
speed = {50 ft., 50 ft. (key)}
]
\DndMonsterAbilityScores[
 str = 16,
 dex = 17,
 con = 18,
 int = 12,
 wis = 14,
 cha = 17
]
\DndMonsterDetails[
skills = {Deception +7, Insight +6, Investigation +5, Perception +6, Stealth +7},
damage-resistances = {cold, fire, lightning; bludgeoning, piercing, slashing from nonmagical attacks},
damage-immunities = {poison},
senses = {darkvision 120 ft., passive perception 16},
languages = {Infernal},
challenge = 11,
]
\DndMonsterAction{Candle Dependence}
The daemon remains on the Material Plane only while the candle that summoned it burns.
\DndMonsterAction{Magic Resistance}
The daemon has advantage on saving throws against spells and other magical effects.
\DndMonsterAction{Nether Step}
As its movement for its turn, the candlelight daemon can teleport to an unoccupied space within 30 feet of it, provided the space it's teleporting to and from are both in dim light or darkness. The daemon doesn't need to see the destination.
\DndMonsterAction{Shadow Stealth}
While in dim light or darkness, the daemon can take the Hide action as a bonus action.
\DndMonsterAction{Slayer's Advantage}
The candlelight daemon has advantage on attack rolls against any creature that hasn't had a turn in combat.
\DndMonsterAction{Slayer's Attack (1/Turn)}
The candlelight daemon deals an extra 35 10d6 damage when it hits a target with an attack and has advantage on the attack roll. With the suffocate spell, the daemon deals this damage as extra bludgeoning damage.
\DndMonsterAction{Innate Spellcasting}
The daemon's spellcasting ability is Charisma (spell save DC 15, +7 to hit with spell attacks). The daemon can innately cast the following spells, requiring no components:
\begin{DndMonsterSpells}
  \DndInnateSpellLevel{detect magic, hunter sense, knock}
  \DndInnateSpellLevel[1e]{counterspell, darkness, dispel magic, shield, suffocate}
\end{DndMonsterSpells}
\par
\DndMonsterSection{Actions}
\DndMonsterAction{Multiattack}
The daemon makes two claw attacks and one gore attack. It can use hurl fire in place of any of these attacks.
\DndMonsterAction{Claw}
\textsl{Melee Weapon Attack:} +7 to hit, reach 5 ft., one target. \textsl{Hit:} 10 (2d6 + 3) slashing damage and 4 (1d8) fire damage.
\DndMonsterAction{Gore}
\textsl{Melee Weapon Attack:} +7 to hit, reach 5 ft., one target. \textsl{Hit:} 12 (2d8 + 3) piercing damage and 4 (1d8) fire damage.
\DndMonsterAction{Hurl Fire}
\textsl{Ranged Spell Attack:} +7 to hit, range 120 ft., one target. \textsl{Hit:} 13 (3d8) fire damage.
\end{DndMonster}



\begin{DndMonster}[float*=t]{Caprathorn}
\DndMonsterType{Huge fiend, chaotic evil}
\DndMonsterBasics[
armor-class = {19 (natural armor)},
hit-points = {\DndDice{20d12 + 120}},
speed = {40 ft., {'number': 40, 'condition': '(hover)'} ft. (key), True ft. (key)}
]
\DndMonsterAbilityScores[
 str = 25,
 dex = 14,
 con = 23,
 int = 14,
 wis = 17,
 cha = 10
]
\DndMonsterDetails[
saving-throws = {Int +9, Wis +10, Cha +7},
skills = {Perception +10},
damage-resistances = {cold, fire, lightning},
damage-immunities = {poison; bludgeoning, piercing, slashing from nonmagical attacks},
senses = {truesight 120 ft., passive perception 20},
languages = {all, telepathy 120 ft.},
challenge = 21,
]
\DndMonsterAction{Charge}
If the caprathorn moves at least 20 feet straight toward a target, and then hits it with a gore attack on the same turn, the target takes an extra 18 (4d8) bludgeoning damage. If the target is a creature, it must succeed on a DC 22 Strength saving throw or be pushed up to 10 feet away and knocked prone.
\DndMonsterAction{Implosion}
The caprathorn implodes when it dies, leaving behind a sphere of annihilation. Each creature within 60 feet of the caprathorn when this implosion occurs must succeed on a DC 21 Strength saving throw or be pulled 10 feet toward the sphere.
\DndMonsterAction{Legendary Resistance (3/Day)}
If the caprathorn fails a saving throw, it can choose to succeed instead.
\DndMonsterAction{Magic Resistance}
The caprathorn has advantage on saving throws against spells and other magical effects.
\DndMonsterAction{Magic Weapons}
The caprathorn's weapon attacks are magical.
\DndMonsterAction{Thornlamms}
A caprathorn typically has three inanimate thornlamms (see that stat block) attached to its horns with vines. Attached thornlamms don't take damage separately, but whenever the caprathorn takes 50 or more damage in a single turn, one of the attached thornlamms dies.
While it has at least one thornlamm attached to its vines, the caprathorn can use a bonus action to consume or detach a thornlamm. If the caprathorn consumes a thornlamm, the caprathorn regains 21 hit points. When the caprathorn detaches a thornlamm, that thornlamm lands in a space within 10 feet of the caprathorn and takes a turn immediately after the caprathorn. The thornlamm obeys the caprathorn's verbal or telepathic commands (no action for the caprathorn). If issued no commands, the thornlamm acts independently, often attacking a creature that isn't its ally. A thornlamm can survive indefinitely detached from a caprathorn.
\par
\DndMonsterSection{Actions}
\DndMonsterAction{Multiattack}
The caprathorn makes three attacks, one with its gore, one with its stomp, and one with its thorns.
\DndMonsterAction{Gore}
\textsl{Melee Weapon Attack:} +14 to hit, reach 10 ft., one target. \textsl{Hit:} 25 (4d8 + 7) bludgeoning damage.
\DndMonsterAction{Stomp}
\textsl{Melee Weapon Attack:} +14 to hit, reach 10 ft., one target. \textsl{Hit:} 25 (4d8 + 7) bludgeoning damage.
\DndMonsterAction{Thorns}
\textsl{Melee or Ranged Weapon Attack:} +14 to hit, reach 10 ft. or range 30/120 ft., one target. \textsl{Hit:} 21 (4d6 + 7) piercing damage. Being within 5 feet of a hostile creature doesn't cause the caprathorn to have disadvantage on the ranged attack roll.
\par
\DndMonsterSection{Legendary Actions}
Caprathorn can take 3 legendary actions, choosing from the options below. Only one legendary action can be used at a time and only at the end of another creature's turn. Caprathorn regains spent legendary actions at the start of its turn.
\begin{DndMonsterLegendaryActions}
\DndMonsterLegendaryAction{Attack}{The caprathorn makes one attack with its gore, stomp, or thorns.}
\DndMonsterLegendaryAction{Detach Thornlamm}{The caprathorn detaches a thornlamm as per the Thornlamms trait.}
\DndMonsterLegendaryAction{Grow Thornlamm}{Provided it has taken no damage from a silvered weapon this round, the caprathorn grows one thornlamm as per the Thornlamms trait. A caprathorn can have up to six attached thornlamms at a time.}
\DndMonsterLegendaryAction{Consume Thornlamm (Costs 2 Actions)}{The caprathorn consumes an attached thornlamm, regaining 21 hit points.}
\DndMonsterLegendaryAction{Entomb (Costs 2 Actions)}{The caprathorn chooses a creature that has 0 hit points that the caprathorn can see within 30 feet of it. The target must make a DC 18 Wisdom saving throw. On a success, the target is immune to this effect for 24 hours. If the save fails, the target is entombed deep within the earth in a sphere of magical force that is just large enough to contain the target. Nothing can pass through the sphere, nor can any creature teleport or use planar travel to get into or out of it. A caprathorn can entomb only one target at a time. This effect ends, returning the target to where it was before being entombed, if the target regains any hit points or the caprathorn dies.}
\DndMonsterLegendaryAction{Annihilation Field (Costs 3 Actions; Recharges if an Entombed Creature Dies or after a Long Rest)}{A sphere of energy ripples out in a 30-foot-radius from the caprathorn. Each creature in that area must make a DC 18 Constitution saving throw, taking 38 (7d10) force damage on a failed save, or half a s much damage on a successful one. For objects that aren't worn or carried, this emanation destroys Small and smaller ones and damages larger ones. The circle destroys a wall of force or similar structures of force in the area instantly. A creature or object is disintegrated if it drops to 0 hit points due to this magic. A disintegrated creature or object turns to fine ash, and nonmagical objects a disintegrated creature wears or carries disintegrate with it.}
\end{DndMonsterLegendaryActions}
\end{DndMonster}



\begin{DndMonster}[float*=t]{Chapped Brute}
\DndMonsterType{Medium humanoid, any non-good alignment}
\DndMonsterBasics[
armor-class = {13 (natural armor)},
hit-points = {\DndDice{10d8 + 50}},
speed = {30 ft.}
]
\DndMonsterAbilityScores[
 str = 18,
 dex = 12,
 con = 20,
 int = 7,
 wis = 10,
 cha = 7
]
\DndMonsterDetails[
skills = {Athletics +6},
damage-resistances = {necrotic},
senses = {blindsight 60 ft. (blind beyond this radius), passive perception 10},
languages = {understands one language but can't speak},
challenge = 4,
]
\DndMonsterAction{Keen Hearing}
The chapped brute has advantage on Wisdom (Perception) checks that rely on hearing. While deafened, the brute can't use its blindsight.
\par
\DndMonsterSection{Actions}
\DndMonsterAction{Multiattack}
The chapped brute makes two slam attacks.
\DndMonsterAction{Slam}
\textsl{Melee Weapon Attack:} +6 to hit, reach 5 ft., one target. \textsl{Hit:} 11 (2d6 + 4) bludgeoning damage.
\end{DndMonster}


\clearpage

\begin{DndMonster}[float*=t]{Chapped Brute Abomination}
\DndMonsterType{Huge monstrosity, chaotic evil}
\DndMonsterBasics[
armor-class = {16 (natural armor)},
hit-points = {\DndDice{14d12 + 70}},
speed = {40 ft.}
]
\DndMonsterAbilityScores[
 str = 23,
 dex = 10,
 con = 20,
 int = 3,
 wis = 10,
 cha = 5
]
\DndMonsterDetails[
skills = {Athletics +10},
damage-resistances = {necrotic},
senses = {blindsight 90 ft. (blind beyond this radius), passive perception 10},
challenge = 9,
]
\DndMonsterAction{Adaptive Hide}
After taking damage, the chapped brute abomination gains resistance to that type of damage until it takes another type of damage.
\DndMonsterAction{Keen Hearing and Smell}
The chapped brute abomination has advantage on Wisdom (Perception) checks that rely on hearing or smell. While deafened and unable to smell, the brute can't use its blindsight.
\DndMonsterAction{Unstable Life}
At the start of the chapped brute abomination's turn, roll 1d20. The abomination regains hit points equal to any even result and loses hit points equal to any odd result. The abomination can neither exceed its hit point maximum with a gain nor fall below 1 hit point with a loss.
\par
\DndMonsterSection{Actions}
\DndMonsterAction{Multiattack}
The chapped brute abomination makes two tentacle attacks.
\DndMonsterAction{Tentacle}
\textsl{Melee Weapon Attack:} +10 to hit, reach 20 ft., one target. \textsl{Hit:} 22 (3d10 + 6) bludgeoning damage. If the target is Medium or smaller, it is grappled (escape DC 18). Until this grapple ends, the target is restrained, and the brute can't use the tentacle against another target. The chapped brute abomination has four tentacles.
\DndMonsterAction{Engulf}
Each creature grappled by the chapped brute abomination must succeed on a DC 18 Strength saving throw or be pulled into the brute's body. An engulfed target is blinded, restrained, unable to breathe, and has total cover from effects that originate outside the brute. At the start of each of the brute's turns, each engulfed creature takes 14 (4d6) acid damage. The brute can hold only one Large creature or up to four Medium or smaller creatures inside it at a time.
An engulfed creature can try to escape by taking an action to make its choice of a DC 18 Strength (Athletics) or Dexterity (Acrobatics) check. On a success, the creature escapes and uses 10 feet of movement to enter a space of its choice within 5 feet of the brute.
If the brute dies, an engulfed creature is no longer restrained by it and can escape from the corpse by using 10 feet of movement.
If a creature dies inside the brute, the brute doesn't roll for Unstable Life at the start of its next turn and instead regains 20 hit points. The creature's body breaks up inside the brute, and the brute expels anything the creature wore or carried.
\end{DndMonster}



\begin{DndMonster}[float*=t]{Commoner}
\DndMonsterType{Medium humanoid (any race), any alignment}
\DndMonsterBasics[
armor-class = {10},
hit-points = {\DndDice{1d8}},
speed = {30 ft.}
]
\DndMonsterAbilityScores[
 str = 10,
 dex = 10,
 con = 10,
 int = 10,
 wis = 10,
 cha = 10
]
\DndMonsterDetails[
languages = {any one language (usually Common)},
challenge = 0,
]
\par
\DndMonsterSection{Actions}
\DndMonsterAction{Club}
\textsl{Melee Weapon Attack:} +2 to hit, reach 5 ft., one target. \textsl{Hit:} 2 (1d4) bludgeoning damage.
\end{DndMonster}



\begin{DndMonster}[float*=t]{Commoner}
\DndMonsterType{Medium humanoid (any race), any alignment}
\DndMonsterBasics[
armor-class = {10},
hit-points = {\DndDice{1d8}},
speed = {30 ft.}
]
\DndMonsterAbilityScores[
 str = 10,
 dex = 10,
 con = 10,
 int = 10,
 wis = 10,
 cha = 10
]
\DndMonsterDetails[
languages = {any one language (usually Common)},
challenge = 0,
]
\par
\DndMonsterSection{Actions}
\DndMonsterAction{Club}
\textsl{Melee Weapon Attack:} +2 to hit, reach 5 ft., one target. \textsl{Hit:} 2 (1d4) bludgeoning damage.
\end{DndMonster}



\begin{DndMonster}[float*=t]{Corpse Walker}
\DndMonsterType{Medium undead, neutral evil}
\DndMonsterBasics[
armor-class = {14 (natural armor)},
hit-points = {\DndDice{11d8 + 33}},
speed = {30 ft.}
]
\DndMonsterAbilityScores[
 str = 16,
 dex = 13,
 con = 17,
 int = 6,
 wis = 10,
 cha = 6
]
\DndMonsterDetails[
saving-throws = {Con +6, Wis +3},
damage-resistances = {; bludgeoning, piercing, slashing from nonmagical attacks that aren't silvered},
damage-immunities = {necrotic, poison},
senses = {darkvision 60 ft., passive perception 10},
languages = {understands languages it knew in life but can't speak},
challenge = 6,
]
\DndMonsterAction{Cursed Gaze}
A humanoid that starts its turn within 30 ft. of and able to see the corpse walker must succeed on a DC 14 Charisma saving throw, provided the corpse walker isn't incapacitated. On a failure, the humanoid is frightened for 1 minute. A frightened target can repeat the saving throw at the end of each of its turns, with disadvantage if not averting its eyes from the corpse walker, ending the effect on itself on a success. A creature failing two of these saving throws becomes cursed with disturbing thoughts. Such a creature can't benefit from a short or long rest until the curse is removed. On a successful save, the target is immune to any corpse walker's Cursed Gaze for the next 24 hours.
Unless surprised, a creature can avert its eyes to avoid the saving throw. If the creature does so, it can't see the corpse walker until the start of its next turn. A creature that looks at the corpse walker in the meantime must immediately make the save.
\DndMonsterAction{Sunlight Hypersensitivity}
The corpse walker takes 20 radiant damage when it starts its turn in sunlight. While in sunlight, it has disadvantage on attack rolls and ability checks.
\DndMonsterAction{Undead Fortitude}
If damage reduces the corpse walker to 0 hit points, it must make a Constitution saving throw with a DC of 5 + the damage taken, unless the damage is radiant or from a critical hit. On a success, the corpse walker drops to 1 hit point instead.
\par
\DndMonsterSection{Actions}
\DndMonsterAction{Multiattack}
The corpse walker makes two slam attacks.
\DndMonsterAction{Slam}
\textsl{Melee Weapon Attack:} +6 to hit, reach 5 ft., one target. \textsl{Hit:} 7 (1d6 + 4) bludgeoning damage and 5 (2d4) necrotic damage. If the target is a creature, it must succeed on a DC 14 Constitution saving throw, or its hit point maximum is reduced by an amount equal to the necrotic damage taken, and the corpse walker regains hit points equal to that amount. This reduction lasts until the target finishes a long rest. The target dies if this effect reduces its hit point maximum to 0. An evil humanoid slain this way rises 24 hours later as a corpse walker.
\end{DndMonster}


\clearpage

\begin{DndMonster}[float*=t]{Corpsejaw}
\DndMonsterType{Medium monstrosity, unaligned}
\DndMonsterBasics[
armor-class = {12},
hit-points = {\DndDice{4d8}},
speed = {40 ft.}
]
\DndMonsterAbilityScores[
 str = 15,
 dex = 14,
 con = 10,
 int = 3,
 wis = 13,
 cha = 7
]
\DndMonsterDetails[
senses = {darkvision 60 ft., passive perception 11},
challenge = 1/4,
]
\DndMonsterAction{Adamantine Jaws}
The corpsejaw's bite works as if made of adamantine. The bite scores a critical hit on a roll of 19 or 20.
\DndMonsterAction{Keen Smell}
The corpsejaw has advantage on Wisdom (Perception) checks that rely on smell.
\DndMonsterAction{Pounce}
If the corpsejaw moves at least 20 feet straight toward a creature and then hits it with a claw attack on the same turn, that target must succeed on a DC 12 Strength saving throw or fall prone. If the target is prone, the corpsejaw can make one bite attack against it as a bonus action.
\par
\DndMonsterSection{Actions}
\DndMonsterAction{Bite}
\textsl{Melee Weapon Attack:} +4 to hit, reach 5 ft., one target. \textsl{Hit:} 6 (1d8 + 2) piercing damage. If the corpsejaw scores a critical hit and the target wears nonmagical armor, the armor takes a permanent and cumulative -1 penalty to the AC it offers. Armor reduced to an AC of 10 is destroyed.
\DndMonsterAction{Claw}
\textsl{Melee Weapon Attack:} +4 to hit, reach 5 ft., one target. \textsl{Hit:} 4 (1d4 + 2) slashing damage.
\end{DndMonster}



\begin{DndMonster}[float*=t]{Cultist}
\DndMonsterType{Medium humanoid (any race), any non-good alignment}
\DndMonsterBasics[
armor-class = {12 (leather armor)},
hit-points = {\DndDice{2d8}},
speed = {30 ft.}
]
\DndMonsterAbilityScores[
 str = 11,
 dex = 12,
 con = 10,
 int = 10,
 wis = 11,
 cha = 10
]
\DndMonsterDetails[
skills = {Deception +2, Religion +2},
languages = {any one language (usually Common)},
challenge = 1/8,
]
\DndMonsterAction{Dark Devotion}
The cultist has advantage on saving throws against being charmed or frightened.
\par
\DndMonsterSection{Actions}
\DndMonsterAction{Scimitar}
\textsl{Melee Weapon Attack:} +3 to hit, reach 5 ft., one creature. \textsl{Hit:} 4 (1d6 + 1) slashing damage.
\end{DndMonster}



\begin{DndMonster}[float*=t]{Dawndrinker}
\DndMonsterType{Large monstrosity, unaligned}
\DndMonsterBasics[
armor-class = {15 (natural armor)},
hit-points = {\DndDice{12d10 + 48}},
speed = {40 ft.}
]
\DndMonsterAbilityScores[
 str = 18,
 dex = 14,
 con = 19,
 int = 3,
 wis = 11,
 cha = 7
]
\DndMonsterDetails[
skills = {Stealth +5},
damage-resistances = {fire},
damage-immunities = {radiant},
senses = {blindsight 120 ft., darkvision 240 ft., passive perception 10},
challenge = 5,
]
\DndMonsterAction{Drink Light}
The dawndrinker can take a bonus action to drink light within 30 feet of it. If it does so, nonmagical light sources in this area go out, extinguishing any flames that are the source of the light. Provided the light comes from an effect of 3rd level or lower, or a creature of CR 7 or lower, magical bright light becomes dim light, and magical dim light goes out. If a magical effect's light is extinguished, the effect is dispelled. A creature's or object's inherent light returns to normal after 1 minute. The dawndrinker can dim or extinguish stronger effects by succeeding on a Constitution check with a DC of 10 + the effect's level or 5 + the creator's CR, whichever is lower.
\DndMonsterAction{Light Sense}
The dawndrinker can sense any light source within 120 feet of it. This sense penetrates most barriers, but it is blocked by 3 feet of wood or dirt, 1 foot of stone, or 1 inch of metal.
\DndMonsterAction{Radiant Absorption}
Whenever the dawndrinker is subjected to radiant damage, it takes no damage and instead regains hit points equal to the radiant damage dealt.
\DndMonsterAction{Shadow Cunning}
In dim light or darkness, the dawndrinker can take a bonus action to use the Disengage or Hide actions.
\par
\DndMonsterSection{Actions}
\DndMonsterAction{Multiattack}
The dawndrinker makes two slam attacks.
\DndMonsterAction{Slam}
\textsl{Melee Weapon Attack:} +7 to hit, reach 5 ft., one target. \textsl{Hit:} 8 (1d8 + 4) bludgeoning damage and 7 (2d6) cold damage. If the target is a creature, it must succeed on a DC 13 Constitution saving throw or be blinded until the end of the dawndrinker's next turn.
\end{DndMonster}



\begin{DndMonster}[float*=t]{Doomcaller}
\DndMonsterType{Medium fiend, chaotic evil}
\DndMonsterBasics[
armor-class = {16 (natural armor)},
hit-points = {\DndDice{20d8 + 60}},
speed = {30 ft., {'number': 40, 'condition': '(hover)'} ft. (key), True ft. (key)}
]
\DndMonsterAbilityScores[
 str = 8,
 dex = 17,
 con = 17,
 int = 14,
 wis = 12,
 cha = 21
]
\DndMonsterDetails[
saving-throws = {Con +8, Int +7, Wis +6, Cha +10},
skills = {Arcana +7, Deception +10, Insight +6, Perception +6, Persuasion +10, Religion +7},
damage-resistances = {cold, fire, lightning; bludgeoning, piercing, slashing from nonmagical attacks},
damage-immunities = {acid, poison},
senses = {darkvision 120 ft., passive perception 16},
languages = {all},
challenge = 15,
]
\DndMonsterAction{Magic Resistance}
The doomcaller has advantage on saving throws against spells and other magical effects.
\DndMonsterAction{Spellcasting}
The doomcaller is an 18th-level spellcaster. Its spellcasting ability is Charisma (spell save DC 18, +10 to hit with spell attacks). It has the following spells prepared:
\begin{DndMonsterSpells}
  \DndMonsterSpellLevel{fire bolt, hunter sense, message, shocking grasp, true strike}
   \DndMonsterSpellLevel[1][4]{detect magic, magic missile, shield}
   \DndMonsterSpellLevel[2][3]{detect thoughts, misty step}
   \DndMonsterSpellLevel[3][3]{bestow curse, fireball}
   \DndMonsterSpellLevel[4][3]{banishment, consume mind}
   \DndMonsterSpellLevel[5][3]{magic mirror, scrying}
   \DndMonsterSpellLevel[6][1]{disintegrate}
   \DndMonsterSpellLevel[7][1]{arboreal curse}
   \DndMonsterSpellLevel[8][1]{mind blank}
\end{DndMonsterSpells}
\par
\DndMonsterSection{Actions}
\DndMonsterAction{Slam}
\textsl{Melee Weapon Attack:} +8 to hit, reach 5 ft., one target. \textsl{Hit:} 10 (2d6 + 3) bludgeoning damage, and the target has disadvantage on their next attack roll. Additionally, all attacks against the target by allies of the doomcaller are made with advantage until the start of the doomcaller's next turn.
\DndMonsterAction{Summon Fiend (1/Day)}
The doomcaller chooses what to summon and attempts a magical summoning. A doomcaller has a 60 percent chance of summoning 1d6 succubi or 1d4 barbed devils. A summoned fiend appears in an unoccupied space within 60 feet of its summoner, does as it pleases, and can't summon other fiends. The summoned fiend remains for 1 minute, until it or its summoner dies, or until its summoner takes a bonus action to dismiss it.
\par
\DndMonsterSection{Reactions}
\DndMonsterAction{Ill Fortune (3/Day)}
When a creature the doomcaller can see makes an attack roll, a saving throw, or an ability check, the doomcaller can roll a d20 and choose to use this roll in place of the attack roll, saving throw, or ability check.
\end{DndMonster}


\clearpage

\begin{DndMonster}[float*=t]{Downcast Apostate}
\DndMonsterType{Medium humanoid (downcast), any alignment}
\DndMonsterBasics[
armor-class = {12 (15 with mage armor)},
hit-points = {\DndDice{20d8 + 60}},
speed = {30 ft.}
]
\DndMonsterAbilityScores[
 str = 15,
 dex = 14,
 con = 17,
 int = 13,
 wis = 12,
 cha = 18
]
\DndMonsterDetails[
saving-throws = {Wis +5, Cha +8},
skills = {Deception +8, Intimidation +8, Perception +5, Persuasion +8, Religion +5, Stealth +6},
damage-resistances = {necrotic},
senses = {darkvision 120 ft., truesight 30 ft., passive perception 15},
languages = {Celestial and two other languages},
challenge = 11,
]
\DndMonsterAction{Apostate's Weapons}
The downcast apostate's weapon attacks are magical.
\DndMonsterAction{Innate Spellcasting}
The downcast apostate's spellcasting ability is Charisma. The apostate can innately cast the following spells, requiring no material components:
\begin{DndMonsterSpells}
  \DndInnateSpellLevel{alter self, detect magic, invisibility, jump, mage armor, thaumaturgy}
\end{DndMonsterSpells}
\DndMonsterAction{Spellcasting}
The downcast apostate is an 18th-level spellcaster. Its spellcasting ability is Charisma (spell save DC 16, +8 to hit with spell attacks). It regains its expended spell slots when it finishes a short or long rest. It knows the following warlock spells:
\begin{DndMonsterSpells}
  \DndInnateSpellLevel[1e]{glibness, mass suggestion, plane shift, power word kill}
  \DndMonsterSpellLevel{chill touch, light, mage hand, minor illusion, sacred flame}
   \DndMonsterSpellLevel[5][4]{counterspell, cure wounds, dimension door, dispel magic, flame strike, fly, greater restoration, guiding bolt, insect plague, lesser restoration, suggestion, tongues, vampiric touch, wall of fire}
\end{DndMonsterSpells}
\par
\DndMonsterSection{Actions}
\DndMonsterAction{Morningstar}
\textsl{Melee Weapon Attack:} +6 to hit, reach 5 ft., one target. \textsl{Hit:} 6 (1d8 + 2) bludgeoning damage and 18 (4d8) necrotic or radiant damage (downcast apostate's choice).
\end{DndMonster}



\begin{DndMonster}[float*=t]{Downcast Mercenary}
\DndMonsterType{Medium humanoid, any alignment}
\DndMonsterBasics[
armor-class = {17 (half plate armor)},
hit-points = {\DndDice{11d8 + 33}},
speed = {30 ft.}
]
\DndMonsterAbilityScores[
 str = 17,
 dex = 15,
 con = 16,
 int = 12,
 wis = 13,
 cha = 14
]
\DndMonsterDetails[
saving-throws = {Wis +4},
skills = {Athletics +6, Perception +4, Religion +4},
damage-resistances = {necrotic},
languages = {Celestial and two other languages},
challenge = 5,
]
\DndMonsterAction{Action Surge (Recharges after a Short or Long Rest)}
The downcast mercenary takes one extra action.
\DndMonsterAction{Indomitable (2/Day)}
If the downcast mercenary fails a saving throw, it can reroll the saving throw and use the new roll.
\DndMonsterAction{Thaumaturge}
The downcast mercenary can cast thaumaturgy at will.
\par
\DndMonsterSection{Actions}
\DndMonsterAction{Multiattack}
The downcast mercenary makes three attacks.
\DndMonsterAction{Greatsword}
\textsl{Melee Weapon Attack:} +6 to hit, reach 5 ft., one target. \textsl{Hit:} 10 (2d6 + 3) slashing damage.
\end{DndMonster}



\begin{DndMonster}[float*=t]{Dryad}
\DndMonsterType{Medium fey, neutral}
\DndMonsterBasics[
armor-class = {11 (16 with barkskin)},
hit-points = {\DndDice{5d8}},
speed = {30 ft.}
]
\DndMonsterAbilityScores[
 str = 10,
 dex = 12,
 con = 11,
 int = 14,
 wis = 15,
 cha = 18
]
\DndMonsterDetails[
skills = {Perception +4, Stealth +5},
senses = {darkvision 60 ft., passive perception 14},
languages = {Elvish, Sylvan},
challenge = 1,
]
\DndMonsterAction{Magic Resistance}
The dryad has advantage on saving throws against spells and other magical effects.
\DndMonsterAction{Speak with Beasts and Plants}
The dryad can communicate with beasts and plants as if they shared a language.
\DndMonsterAction{Tree Stride}
Once on her turn, the dryad can use 10 feet of her movement to step magically into one living tree within her reach and emerge from a second living tree within 60 feet of the first tree, appearing in an unoccupied space within 5 feet of the second tree. Both trees must be large or bigger.
\DndMonsterAction{Innate Spellcasting}
The dryad's innate spellcasting ability is Charisma (spell save DC 14). The dryad can innately cast the following spells, requiring no material components:
\begin{DndMonsterSpells}
  \DndInnateSpellLevel{druidcraft}
  \DndInnateSpellLevel[3e]{entangle, goodberry}
  \DndInnateSpellLevel[1e]{barkskin, pass without trace, shillelagh}
\end{DndMonsterSpells}
\par
\DndMonsterSection{Actions}
\DndMonsterAction{Club}
\textsl{Melee Weapon Attack:} +2 to hit (+6 to hit with shillelagh), reach 5 ft., one target. \textsl{Hit:} 2 (1d4) bludgeoning damage, or 8 (1d8 + 4) bludgeoning damage with shillelagh.
\DndMonsterAction{Fey Charm}
The dryad targets one humanoid or beast that she can see within 30 feet of her. If the target can see the dryad, it must succeed on a DC 14 Wisdom saving throw or be magically charmed. The charmed creature regards the dryad as a trusted friend to be heeded and protected. Although the target isn't under the dryad's control, it takes the dryad's requests or actions in the most favorable way it can.
Each time the dryad or its allies do anything harmful to the target, it can repeat the saving throw, ending the effect on itself on a success. Otherwise, the effect lasts 24 hours or until the dryad dies, is on a different plane of existence from the target, or ends the effect as a bonus action. If a target's saving throw is successful, the target is immune to the dryad's Fey Charm for the next 24 hours.
The dryad can have no more than one humanoid and up to three beasts charmed at a time.
\end{DndMonster}



\begin{DndMonster}[float*=t]{Eldritch Herald}
\DndMonsterType{Medium humanoid (any race), any non-good alignment}
\DndMonsterBasics[
armor-class = {16 (breastplate)},
hit-points = {\DndDice{12d8 + 24}},
speed = {30 ft.}
]
\DndMonsterAbilityScores[
 str = 14,
 dex = 11,
 con = 15,
 int = 10,
 wis = 17,
 cha = 12
]
\DndMonsterDetails[
skills = {Perception +6, Religion +3},
damage-resistances = {psychic},
languages = {any two languages},
challenge = 5,
]
\DndMonsterAction{Otherworldly Calm}
The eldritch herald has advantage on saving throws against being charmed or frightened. Also, attempts to read the herald's thoughts fail. The creature making the attempt must succeed on a DC 14 Wisdom saving throw or take 10 (3d6) psychic damage.
\DndMonsterAction{Spellcasting}
The herald is a 12th-level spellcaster. Its spellcasting ability is Wisdom (spell save DC 14, +6 to hit with spell attacks). It has the following cleric spells prepared:
\begin{DndMonsterSpells}
  \DndMonsterSpellLevel{light, sacred flame, thaumaturgy, vicious mockery}
   \DndMonsterSpellLevel[1][4]{bane, Tasha's hideous laughter, inflict wounds, sleep}
   \DndMonsterSpellLevel[2][3]{detect thoughts, hold person, see invisibility, wrack}
   \DndMonsterSpellLevel[3][3]{bestow curse, fear, spirit guardians, tongues}
   \DndMonsterSpellLevel[4][3]{confusion, death ward, phantasmal killer}
   \DndMonsterSpellLevel[5][2]{contact other plane, contagion, dream}
   \DndMonsterSpellLevel[6][1]{true seeing}
\end{DndMonsterSpells}
\par
\DndMonsterSection{Actions}
\DndMonsterAction{Handaxe}
\textsl{Melee Weapon Attack:} +5 to hit, reach 5 ft., one target. \textsl{Hit:} 5 (1d6 + 2) slashing damage.
\DndMonsterAction{Prophecy of Doom (Recharges after a Short or Long Rest)}
The eldritch herald chooses a point it can see within 120 feet of it and rolls on the Eldritch Effects table. Each creature within 15 feet of that point must succeed on a DC 14 Wisdom saving throw or be subjected to an eldritch effect for 1 minute. A creature can repeat the saving throw at the end of each of its turns, ending the effect on itself on a success.
\end{DndMonster}


\clearpage

\begin{DndMonster}[float*=t]{Eldritch Priest}
\DndMonsterType{Medium humanoid (any race), any non-good alignment}
\DndMonsterBasics[
armor-class = {13 (chain shirt)},
hit-points = {\DndDice{4d8 + 4}},
speed = {30 ft.}
]
\DndMonsterAbilityScores[
 str = 12,
 dex = 11,
 con = 13,
 int = 10,
 wis = 15,
 cha = 10
]
\DndMonsterDetails[
skills = {Perception +4, Religion +2},
languages = {any two languages},
challenge = 2,
]
\DndMonsterAction{Spellcasting}
The eldritch priest is a 4th-level spellcaster. Its spellcasting ability is Wisdom (spell save DC 12, +4 to hit with spell attacks). It has the following cleric spells prepared:
\begin{DndMonsterSpells}
  \DndMonsterSpellLevel{eldritch blast, light, sacred flame, thaumaturgy}
   \DndMonsterSpellLevel[1][4]{bane, Tasha's hideous laughter, inflict wounds, sleep}
   \DndMonsterSpellLevel[2][3]{detect thoughts, see invisibility, spiritual weapon}
\end{DndMonsterSpells}
\par
\DndMonsterSection{Actions}
\DndMonsterAction{Greatclub}
\textsl{Melee Weapon Attack:} +3 to hit, reach 5 ft., one target. \textsl{Hit:} 5 (1d8 + 1) bludgeoning damage.
\DndMonsterAction{Prophecy of Doom (Recharges after a Short or Long Rest)}
The eldritch priest chooses a point it can see within 120 feet of it and rolls on the Eldritch Effects table. Each creature within 15 feet of that point must succeed on a DC 12 Wisdom saving throw or be subjected to an eldritch effect for 1 minute. A creature can repeat the saving throw at the end of each of its turns, ending the effect on itself on a success.
\end{DndMonster}



\begin{DndMonster}[float*=t]{Elf Vampire (Ancient)}
\DndMonsterType{Medium undead, any evil alignment}
\DndMonsterBasics[
armor-class = {18 (natural armor)},
hit-points = {\DndDice{22d8 + 110}},
speed = {40 ft., 40 ft. (key)}
]
\DndMonsterAbilityScores[
 str = 20,
 dex = 20,
 con = 20,
 int = 15,
 wis = 16,
 cha = 16
]
\DndMonsterDetails[
saving-throws = {Dex +11, Wis +9, Cha +9},
skills = {Arcana +8, Athletics +11, History +8, Perception +9, Stealth +11},
damage-resistances = {necrotic, poison; bludgeoning, piercing, slashing from nonmagical attacks},
senses = {blindsight 30 ft., darkvision 120 ft., passive perception 19},
languages = {the languages known in life},
challenge = 19,
]
\DndMonsterAction{Legendary Resistance (3/Day)}
If the vampire fails a saving throw, it can choose to succeed instead.
\DndMonsterAction{Regeneration}
The vampire regains 25 hit points at the start of its turn if it has at least 1 hit point. If the vampire takes radiant damage, it regains 10 hit points from the next use of this trait instead.
\DndMonsterAction{Spider Climb}
The vampire can climb difficult surfaces, including upside down on ceilings, without needing to make an ability check.
\DndMonsterAction{Turn Resistance}
The vampire has advantage on saving throws against any effect that turns undead.
\DndMonsterAction{Vampire Weaknesses}
The vampire has the following flaws:
\DndMonsterAction{Forbiddance}
The vampire can't enter a residence without an invitation from one of the occupants.
\DndMonsterAction{Stake to the Heart}
The vampire dies if a piercing weapon made of wood is driven into its heart while it is incapacitated in its resting place.
\par
\DndMonsterSection{Actions}
\DndMonsterAction{Multiattack}
The vampire makes three attacks, only one of which can be a bite.
\DndMonsterAction{Claws}
\textsl{Melee Weapon Attack:} +11 to hit, reach 5 ft., one target. \textsl{Hit:} 15 (2d8 + 6) slashing damage. If the target is a creature, the vampire can grapple the target (escape DC 19).
\DndMonsterAction{Bite}
\textsl{Melee Weapon Attack:} +11 to hit, reach 5 ft., one creature. \textsl{Hit:} 15 (2d8 + 6) piercing damage and 17 (5d6) necrotic damage. The target's hit point maximum is reduced by an amount equal to the necrotic damage taken, and the vampire regains hit points equal to that amount. This reduction lasts until the target finishes a long rest. The target dies if this effect reduces its hit point maximum to 0. A humanoid slain in this way and then buried rises the following night as a vampire spawn under this vampire's control.
\DndMonsterAction{Fey Majesty}
The vampire targets one creature it can see within 30 feet of it. If the target can see or hear the vampire, the target must succeed on a DC 17 Charisma saving throw or be charmed. The charmed target regards the vampire as a friend to be heeded and protected. Although not under the vampire's control, it takes the vampire's requests in the most favorable way and is a willing target for the vampire's bite.
While the target is charmed, the vampire can communicate with the target telepathically within 300 feet. The vampire can use a bonus action to force the target to repeat the saving throw. If the save fails, the vampire can render the creature full of despair (disadvantage on ability checks and attack rolls) or confused (as the confusion spell) for 1 minute. The creature can repeat the saving throw at the end of each of its turns, ending the despair or confusion on itself on a success.
Each time the vampire or the vampire's allies do anything harmful to the target, it can repeat the saving throw, ending the effect on a success. Otherwise, the effect lasts 24 hours or until the vampire dies, is on a different plane of existence than the target, or takes a bonus action to end the effect.
\DndMonsterAction{Mist Form}
The vampire polymorphs into a Medium cloud of mist, or back into its true form. Anything it wears or carries transforms with it. While in mist form, the vampire can't take any actions, speak, or manipulate objects. It is weightless, has a flying speed of 30 feet, can hover, and can enter a hostile creature's space and stop there. In addition, if air can pass through a space, the mist can do so without squeezing, but it can't pass through water. It has advantage on Strength, Dexterity, and Constitution saving throws. A vampire in mist form is immune to nonmagical damage, except damage from sunlight, and has resistance to damage other than radiant. The vampire reverts to its true form if it dies.
\par
\DndMonsterSection{Reactions}
\DndMonsterAction{Misty Escape}
When it drops to 0 hit points outside its resting place, the vampire transforms into its Mist Form instead of falling unconscious. If the vampire can't transform, it dies.
While it has 0 hit points in mist form, the vampire dies if it takes damage three times or radiant damage once. It can't willingly revert to its vampire form, and it must reach its resting place within 2 hours or die. Once in its resting place, it reverts to its vampire form. It is then paralyzed until it regains at least 1 hit point. After spending 1 hour in its resting place with 0 hit points, it regains 1 hit point.
\par
\DndMonsterSection{Legendary Actions}
Elf Vampire (Ancient) can take 3 legendary actions, choosing from the options below. Only one legendary action can be used at a time and only at the end of another creature's turn. Elf Vampire (Ancient) regains spent legendary actions at the start of its turn.
\begin{DndMonsterLegendaryActions}
\DndMonsterLegendaryAction{Move and Claw}{The vampire moves up to its speed and attacks once with its claws.}
\DndMonsterLegendaryAction{Bite (2 Actions)}{The vampire makes one bite attack.}
\DndMonsterLegendaryAction{Purify Blood (2 Actions)}{Even if the vampire can't take actions, one condition ends on it.}
\end{DndMonsterLegendaryActions}
\end{DndMonster}



\begin{DndMonster}[float*=t]{Empyrean Brazen Bull}
\DndMonsterType{Huge celestial, lawful good}
\DndMonsterBasics[
armor-class = {16 (natural armor)},
hit-points = {\DndDice{12d12 + 84}},
speed = {50 ft., 30 ft. (key)}
]
\DndMonsterAbilityScores[
 str = 26,
 dex = 10,
 con = 25,
 int = 14,
 wis = 17,
 cha = 21
]
\DndMonsterDetails[
saving-throws = {Str +12, Con +11, Wis +7, Cha +9},
skills = {Athletics +12, Perception +7},
damage-resistances = {acid, cold, lightning; bludgeoning, piercing, slashing from nonmagical attacks},
damage-immunities = {fire, necrotic, poison, radiant},
senses = {darkvision 120 ft., passive perception 17},
languages = {Celestial},
challenge = 10,
]
\DndMonsterAction{Celestial Arsenal}
The Empyrean brazen bull's weapon attacks are magical. When it hits with any weapon, the weapon deals an extra 18 (4d8) radiant damage (included in the attack).
\DndMonsterAction{Death Throes}
When the Empyrean brazen bull dies, it explodes, and each creature within 20 feet of it must make a DC 17 Dexterity saving throw, taking 21 6d6 fire and 21 (6d6) radiant damage on a failed save, or half as much damage on a successful one.
\DndMonsterAction{Soulfire Nimbus}
At the start of each of the Empyrean brazen bull's turns, each creature within 5 feet of it takes 10 (3d6) radiant damage. A creature that touches the Empyrean brazen bull or hits it with a melee attack while within 5 feet of it takes 10 (3d6) radiant damage.
\DndMonsterAction{Siege Monster}
The Empyrean brazen bull deals double damage to objects and structures.
\DndMonsterAction{Unrelenting Charge}
If the Empyrean brazen bull moves at least 20 feet straight toward a creature and then hits it with a gore attack on the same turn, that target must succeed on a DC 20 Strength saving thrown or be knocked prone. If the target is prone, the Empyrean brazen bull can make one attack with its bite or hooves against it as a bonus action.
\par
\DndMonsterSection{Actions}
\DndMonsterAction{Bite}
\textsl{Melee Weapon Attack:} +12 to hit, reach 10 ft., one target., \textsl{Hit:} 24 (3d10 + 8) bludgeoning damage plus 18 (4d8) radiant damage. If the target is a creature, it is grappled (escape DC 20). Until this grapple ends, the target is restrained, and the Empyrean brazen bull can't bite another target.
\DndMonsterAction{Gore}
\textsl{Melee Weapon Attack:} +12 to hit, reach 10 ft., one target. \textsl{Hit:} 27 (3d12 + 8) piercing damage and 18 (4d8) radiant damage.
\DndMonsterAction{Hooves}
\textsl{Melee Weapon Attack:} +12 to hit, reach 5 feet., one target. \textsl{Hit:} 24 (3d10 + 8) bludgeoning damage and 18 (4d8) radiant damage.
\DndMonsterAction{Swallow}
The Empyrean brazen bull makes one bite attack against a Medium or smaller creature it is grappling. If the attack hits, that creature takes the bite's damage and is swallowed, and the grapple ends. While swallowed, the creature is blinded and restrained, it has total cover against attacks and other effects outside the Empyrean brazen bull, and it takes 21 (6d6) fire damage at the start of each of theEmpyrean brazen bull's turns. If the Empyrean brazen bull takes 30 damage or more on a single turn from a creature inside it, the Empyrean brazen bull must succeed on a DC 15 Constitution saving throw at the end of that turn or regurgitate all swallowed creatures, which fall prone in a space within 10 feet of the empyrean brazen bull. If the Empyrean brazen bull dies, a swallowed creature is no longer restrained by it and can escape from the corpse using 15 feet of movement, exiting prone.
\end{DndMonster}



\begin{DndMonster}[float*=t]{Eye Crawler}
\DndMonsterType{Tiny aberration, chaotic neutral}
\DndMonsterBasics[
armor-class = {11},
hit-points = {\DndDice{1d4}},
speed = {30 ft., 30 ft. (key)}
]
\DndMonsterAbilityScores[
 str = 3,
 dex = 13,
 con = 10,
 int = 6,
 wis = 11,
 cha = 5
]
\DndMonsterDetails[
skills = {Perception +2, Stealth +5},
senses = {darkvision 60 ft., passive perception 12},
challenge = 0,
]
\DndMonsterAction{Keen Sight}
The eye crawler has advantage on Wisdom (Perception) checks that rely on sight.
\end{DndMonster}


\clearpage

\begin{DndMonster}[float*=t]{Eye Crow}
\DndMonsterType{Tiny beast, unaligned}
\DndMonsterBasics[
armor-class = {12},
hit-points = {\DndDice{1d4}},
speed = {10 ft., 50 ft. (key)}
]
\DndMonsterAbilityScores[
 str = 2,
 dex = 14,
 con = 10,
 int = 3,
 wis = 12,
 cha = 7
]
\DndMonsterDetails[
skills = {Perception +3},
challenge = 0,
]
\DndMonsterAction{Mimicry}
The eye crow can mimic simple sounds it has heard, such as a person whispering, a baby crying, or an animal chittering. A creature that hears the sounds can tell they are imitations with a successful DC 10 Wisdom (Insight) check.
\DndMonsterAction{Pack Tactics}
The eye crow has advantage on an attack roll against a creature if at least one of the crow's allies is within 5 feet of the creature and the ally isn't incapacitated.
\par
\DndMonsterSection{Actions}
\DndMonsterAction{Beak}
\textsl{Melee Weapon Attack:} +4 to hit, reach 5 ft., one target. \textsl{Hit:} 1 piercing damage. If the eye crow scores a critical hit on a creature, the target is blinded until the end of the crow's next turn.
\end{DndMonster}



\begin{DndMonster}[float*=t]{Faevlin}
\DndMonsterType{Small fey, neutral evil}
\DndMonsterBasics[
armor-class = {15 (leather armor)},
hit-points = {\DndDice{3d6}},
speed = {30 ft.}
]
\DndMonsterAbilityScores[
 str = 8,
 dex = 14,
 con = 10,
 int = 10,
 wis = 12,
 cha = 11
]
\DndMonsterDetails[
skills = {Stealth +4},
senses = {darkvision 60 ft., passive perception 11},
languages = {Goblin, Sylvan},
challenge = 1/4,
]
\DndMonsterAction{Fey Touched}
Magic can't put the faevlin to sleep. Only a creature of the fey type can cause the faevlin to become charmed, and a faevlin always knows if a nonfey creature attempted to render it charmed.
\DndMonsterAction{Sneaky}
The faevlin can use a bonus action to take the Hide action.
\par
\DndMonsterSection{Actions}
\DndMonsterAction{Scimitar}
\textsl{Melee Weapon Attack:} +4 to hit, reach 5 ft., one target. \textsl{Hit:} 5 (1d6 + 2) slashing damage.
\DndMonsterAction{Shortbow}
\textsl{Ranged Weapon Attack:} +4 to hit, range 80/320 ft., one target. \textsl{Hit:} 5 (1d6 + 2) piercing damage.
\par
\DndMonsterSection{Reactions}
\DndMonsterAction{Fey Fade}
When an attack misses a faevlin, the faevlin can teleport up to 30 feet in a random direction, arriving in a safe, unoccupied space.
\end{DndMonster}



\begin{DndMonster}[float*=t]{Fzeg}
\DndMonsterType{Medium monstrosity, chaotic evil}
\DndMonsterBasics[
armor-class = {16 (natural armor)},
hit-points = {\DndDice{18d8 + 72}},
speed = {40 ft.}
]
\DndMonsterAbilityScores[
 str = 18,
 dex = 18,
 con = 18,
 int = 17,
 wis = 15,
 cha = 18
]
\DndMonsterDetails[
saving-throws = {Dex +9, Wis +7},
skills = {Perception +7, Stealth +9},
damage-resistances = {necrotic; bludgeoning, piercing, slashing from nonmagical attacks},
senses = {darkvision 120 ft., passive perception 17},
languages = {any four languages},
challenge = 16,
]
\DndMonsterAction{Keen Hearing and Smell}
The fzeg has advantage on Wisdom (Perception) checks that rely on hearing or smell.
\DndMonsterAction{Legendary Resistance (3/Day)}
If the fzeg fails a saving throw, it can choose to succeed instead.
\DndMonsterAction{Rampage}
When the fzeg reduces a creature to 0 hit points with a melee attack on its turn, the fzeg can take a bonus action to move up to half its speed and make a bite attack.
\DndMonsterAction{Regeneration}
The fzeg regains 20 hit points at the start of its turn if it has at least 1 hit point. If the fzeg takes radiant damage, this trait doesn't function at the start of the fzeg's next turn.
\DndMonsterAction{Shapechanger}
The fzeg can use its action to polymorph into a Small or Medium humanoid or back into its true form, that of a wolfhumanoid hybrid. It can transform its body parts as part of attacking with them and can choose whether the change remains after making the attack. Other than its size, its statistics are the same in each form. Any equipment it's wearing or carrying isn't transformed. It reverts to its true form if it dies.
\DndMonsterAction{Spider Climb}
The fzeg can climb difficult surfaces, including upside down on ceilings, without needing to make an ability check.
\par
\DndMonsterSection{Actions}
\DndMonsterAction{Multiattack}
The fzeg uses its howl, then attacks once with its bite and twice with its claws.
\DndMonsterAction{Bite}
\textsl{Melee Weapon Attack:} +9 to hit, reach 5 ft., one target. \textsl{Hit:} 11 (2d6 + 4) piercing damage plus 17 (5d6) necrotic damage and is stunned until the end of its next turn. A humanoid slain after being bitten rises the following night as a fzeglaich under the fzeg's control.
\DndMonsterAction{Claws}
\textsl{Melee Weapon Attack:} +9 to hit, reach 5 ft., one target. \textsl{Hit:} 13 (2d8 + 4) slashing damage.
\DndMonsterAction{Howl}
Each creature of the fzeg's choice that is within 300 feet of the fzeg and able to hear it must succeed on a DC 17 Wisdom saving throw or become frightened for 1 minute. A creature can repeat the saving throw at the end of each of its turns, ending the effect on itself on a success. If a creature's saving throw is successful or the effect ends for it, the creature is immune to the fzeg's Howl for the next 24 hours or until it is stunned by the fzeg.
\DndMonsterAction{Night Children (Recharges after a Short or Long Rest)}
The fzeg magically calls eight wolves. These wolves arrive in 2 rounds, acting as allies of the fzeg and obeying its spoken commands. The beasts remain for 1 hour, until the fzeg dies, or until the fzeg dismisses them as a bonus action.
\par
\DndMonsterSection{Legendary Actions}
Fzeg can take 3 legendary actions, choosing from the options below. Only one legendary action can be used at a time and only at the end of another creature's turn. Fzeg regains spent legendary actions at the start of its turn.
\begin{DndMonsterLegendaryActions}
\DndMonsterLegendaryAction{Move}{The fzeg moves up to its speed without provoking opportunity attacks.}
\DndMonsterLegendaryAction{Claws}{The fzeg attacks with its claws.}
\DndMonsterLegendaryAction{Incite}{The fzeg chooses one creature that is of a lower Challenge, has a bite attack, and is within 60 feet of the fzeg. Provided the target can see or hear the fzeg, that creature can use its reaction to move up to half its speed and make a bite attack.}
\DndMonsterLegendaryAction{Bite (Costs 2 Actions)}{The fzeg attacks with its bite.}
\end{DndMonsterLegendaryActions}
\end{DndMonster}



\begin{DndMonster}[float*=t]{Fzeglaich}
\DndMonsterType{Medium monstrosity, chaotic evil}
\DndMonsterBasics[
armor-class = {15 (natural armor)},
hit-points = {\DndDice{12d8 + 36}},
speed = {40 ft.}
]
\DndMonsterAbilityScores[
 str = 16,
 dex = 16,
 con = 16,
 int = 11,
 wis = 11,
 cha = 12
]
\DndMonsterDetails[
saving-throws = {Dex +6, Wis +3},
skills = {Perception +3, Stealth +6},
damage-resistances = {; bludgeoning, piercing, slashing from nonmagical attacks},
senses = {darkvision 60 ft., passive perception 13},
languages = {one language},
challenge = 7,
]
\DndMonsterAction{Keen Hearing and Smell}
The fzeglaich has advantage on Wisdom (Perception) checks that rely on hearing or smell.
\DndMonsterAction{Rampage}
When the fzeglaich reduces a creature to 0 hit points with a melee attack on its turn, the fzeglaich can take a bonus action to move up to half its speed and make a bite attack.
\DndMonsterAction{Regeneration}
The fzeglaich regains 10 hit points at the start of its turn if it has at least 1 hit point. If the fzeglaich takes radiant damage, this trait doesn't function at the start of the fzeglaich's next turn.
\par
\DndMonsterSection{Actions}
\DndMonsterAction{Multiattack}
The fzeglaich makes two attacks, only one of which can be a bite.
\DndMonsterAction{Bite}
\textsl{Melee Weapon Attack:} +6 to hit, reach 5 ft., one creature. \textsl{Hit:} 7 (1d8 + 3) piercing damage plus 7 (2d6) necrotic damage, and the target must succeed on a DC 14 Constitution saving throw or be incapacitated until the end of its next turn. The fzeglaich regains hit points equal to the amount of necrotic damage dealt.
\DndMonsterAction{Claws}
\textsl{Melee Weapon Attack:} +6 to hit, reach 5 ft., one creature. \textsl{Hit:} 8 (2d4 + 3) slashing damage.
\end{DndMonster}


\clearpage

\begin{DndMonster}[float*=t]{Gegazol}
\DndMonsterType{Gargantuan undead, chaotic evil}
\DndMonsterBasics[
armor-class = {22 (natural armor)},
hit-points = {\DndDice{25d20 + 200}},
speed = {40 ft., 40 ft. (key), 80 ft. (key), 40 ft. (key)}
]
\DndMonsterAbilityScores[
 str = 27,
 dex = 14,
 con = 26,
 int = 16,
 wis = 15,
 cha = 23
]
\DndMonsterDetails[
saving-throws = {Dex +9, Int +10, Wis +9},
skills = {Perception +9},
damage-immunities = {cold, necrotic, poison},
senses = {blindsight 60 ft., darkvision 120 ft., passive perception 19},
languages = {Draconic},
challenge = 24,
]
\DndMonsterAction{Ice Walk}
Gegazol can move across and climb icy surfaces without needing to make an ability check. Additionally, difficult terrain composed of ice or snow doesn't cost her extra movement.
\DndMonsterAction{Legendary Resistance (3/Day)}
If Gegazol fails a saving throw, she can choose to succeed instead.
\DndMonsterAction{Magic Resistance}
Gegazol has advantage on saving throws against spells and other magical effects.
\DndMonsterAction{Rejuvenation}
If destroyed while her soul vessel remains intact, Gegazol's unconscious body rejuvenates within 5 feet of the soul vessel, where she regains 78 (12d12) hit points each day. When she has regained half her hit points, she regains consciousness.
\DndMonsterAction{Stench}
Any creature not immune to the poisoned condition that starts its turn within 30 feet of Gegazol must succeed on a DC 23 Constitution saving throw or become poisoned until the start of that creature's next turn. On a successful saving throw, the creature is immune to Gegazol's stench for 1 hour.
\DndMonsterAction{Turning Defiance}
Gegazol and undead within 120 feet of her have advantage on saving throws against effects that turn undead.
\par
\DndMonsterSection{Actions}
\DndMonsterAction{Multiattack}
Gegazol can use her frightful presence. She then makes three attacks: one with her bite and two with her claws.
\DndMonsterAction{Bite}
\textsl{Melee Weapon Attack:} +15 to hit, reach 15 ft., one target. \textsl{Hit:} 21 (2d12 + 8) piercing damage, 9 (2d8) cold damage and 9 (2d8) necrotic damage.
\DndMonsterAction{Claw}
\textsl{Melee Weapon Attack:} +15 to hit, reach 10 ft., one target. \textsl{Hit:} 19 (2d10 + 8) piercing damage.
\DndMonsterAction{Tail}
\textsl{Melee Weapon Attack:} +15 to hit, reach 20 ft., one target. \textsl{Hit:} 19 (2d10 + 8) bludgeoning damage. If the target is a creature, it must succeed on a DC 23 Strength saving throw or fall prone.
\DndMonsterAction{Frightful Presence}
Each creature of Gegazol's choice within 120 feet of and aware of her must succeed on a DC 21 Wisdom saving throw or become frightened for 1 minute. A creature can repeat the saving throw at the end of each of its turns, ending the effect on itself on a success. If a creature's saving throw is successful or the effect ends for it, the creature is immune to Gegazol's Frightful Presence for the next 24 hours.
\DndMonsterAction{Breath of Long Night Recharge 5-6}
Gegazol exhales frozen material and rot in a 90-foot cone. Each creature in that area must make a DC 23 Constitution saving throw, taking 45 (10d8) cold damage and 45 (10d8) necrotic damage on a failed save, or half as much damage on a successful one. Those who fail the save also gain two levels of exhaustion until the end of Gegazol's next turn, and those who fail by 5 or more are also blinded for 1 minute. A blinded creature can repeat the saving throw at the end of each of its turns, ending the effect on itself on a success.
\par
\DndMonsterSection{Legendary Actions}
Gegazol can take 3 legendary actions, choosing from the options below. Only one legendary action can be used at a time and only at the end of another creature's turn. Gegazol regains spent legendary actions at the start of its turn.
\begin{DndMonsterLegendaryActions}
\DndMonsterLegendaryAction{Daemon Sight}{Gegazol gains truesight out to 120 feet until the end of her next turn and can make a Wisdom (Perception) check.}
\DndMonsterLegendaryAction{Tail}{Gegazol makes a tail attack.}
\DndMonsterLegendaryAction{Wings of Winter (Costs 2 Actions)}{Gegazol beats her wings. Each creature within 30 feet of her must succeed on a DC 23 Dexterity saving throw or take 17 (2d8 + 8) bludgeoning damage and 9 (2d8) cold damage, and then fall prone. Gegazol can then fly up to half her flying speed.}
\end{DndMonsterLegendaryActions}
\par \DndMonsterSection{Lair Actions}
On initiative count 20 (losing initiative ties), Gegazol takes a lair action to cause one of the following effects. Gegazol can't use the same effect two rounds in a row.

\begin{itemize}
\item Tendrils of necrotic power rise and reach for living creatures within the lair. Each living creature in the lair must attempt a DC 21 Constitution saving throw or take 13 (3d8) necrotic damage and gain one level of exhaustion until initiative count 20 on the next round.
\item Gegazol crashes her wings upon the ground, and rime spreads in a sphere with a 120-foot radius around her. This sphere spreads around corners and creates difficult terrain made of ice. Creatures in that area must make a DC 21 Constitution saving throw. Those who fail take 18 (4d8) cold damage and are grappled (escape DC 21). On a success, the creature takes half the damage and isn't grappled. The difficult terrain disappears on initiative count 20 on the next round.
\item Gegazol creates an opaque wall of black ice on a solid surface she can see within 120 feet of her. The wall is made up of twenty contiguous 10-foot square panels that are each 2 feet thick. When the wall appears, each creature in its area is pushed 5 feet out of the wall's space, appearing on whichever side of the wall the creature wants. Each 10-foot section of the wall has AC 5, 60 hit points, vulnerability to fire damage, and immunity to acid, cold, necrotic, poison, and psychic damage. The wall disappears when Gegazol uses this lair action again or when she dies.
\end{itemize}

\end{DndMonster}



\begin{DndMonster}[float*=t]{Giant Ape}
\DndMonsterType{Huge beast, unaligned}
\DndMonsterBasics[
armor-class = {12},
hit-points = {\DndDice{15d12 + 60}},
speed = {40 ft., 40 ft. (key)}
]
\DndMonsterAbilityScores[
 str = 23,
 dex = 14,
 con = 18,
 int = 7,
 wis = 12,
 cha = 7
]
\DndMonsterDetails[
skills = {Athletics +9, Perception +4},
challenge = 7,
]
\par
\DndMonsterSection{Actions}
\DndMonsterAction{Multiattack}
The ape makes two fist attacks.
\DndMonsterAction{Fist}
\textsl{Melee Weapon Attack:} +9 to hit, reach 10 ft., one target. \textsl{Hit:} 22 (3d10 + 6) bludgeoning damage.
\DndMonsterAction{Rock}
\textsl{Ranged Weapon Attack:} +9 to hit, range 50/100 ft., one target. \textsl{Hit:} 30 (7d6 + 6) bludgeoning damage.
\end{DndMonster}



\begin{DndMonster}[float*=t]{Giant Crocodile}
\DndMonsterType{Huge beast, unaligned}
\DndMonsterBasics[
armor-class = {14 (natural armor)},
hit-points = {\DndDice{9d12 + 27}},
speed = {30 ft., 50 ft. (key)}
]
\DndMonsterAbilityScores[
 str = 21,
 dex = 9,
 con = 17,
 int = 2,
 wis = 10,
 cha = 7
]
\DndMonsterDetails[
skills = {Stealth +5},
challenge = 5,
]
\DndMonsterAction{Hold Breath}
The crocodile can hold its breath for 30 minutes.
\par
\DndMonsterSection{Actions}
\DndMonsterAction{Multiattack}
The crocodile makes two attacks: one with its bite and one with its tail.
\DndMonsterAction{Bite}
\textsl{Melee Weapon Attack:} +8 to hit, reach 5 ft., one target. \textsl{Hit:} 21 (3d10 + 5) piercing damage, and the target is grappled (escape DC 16). Until this grapple ends, the target is restrained, and the crocodile can't bite another target.
\DndMonsterAction{Tail}
\textsl{Melee Weapon Attack:} +8 to hit, reach 10 ft., one target not grappled by the crocodile. \textsl{Hit:} 14 (2d8 + 5) bludgeoning damage. If the target is a creature, it must succeed on a DC 16 Strength saving throw or be knocked prone.
\end{DndMonster}



\begin{DndMonster}[float*=t]{Giant Crocodile}
\DndMonsterType{Huge beast, unaligned}
\DndMonsterBasics[
armor-class = {14 (natural armor)},
hit-points = {\DndDice{9d12 + 27}},
speed = {30 ft., 50 ft. (key)}
]
\DndMonsterAbilityScores[
 str = 21,
 dex = 9,
 con = 17,
 int = 2,
 wis = 10,
 cha = 7
]
\DndMonsterDetails[
skills = {Stealth +5},
challenge = 5,
]
\DndMonsterAction{Hold Breath}
The crocodile can hold its breath for 30 minutes.
\par
\DndMonsterSection{Actions}
\DndMonsterAction{Multiattack}
The crocodile makes two attacks: one with its bite and one with its tail.
\DndMonsterAction{Bite}
\textsl{Melee Weapon Attack:} +8 to hit, reach 5 ft., one target. \textsl{Hit:} 21 (3d10 + 5) piercing damage, and the target is grappled (escape DC 16). Until this grapple ends, the target is restrained, and the crocodile can't bite another target.
\DndMonsterAction{Tail}
\textsl{Melee Weapon Attack:} +8 to hit, reach 10 ft., one target not grappled by the crocodile. \textsl{Hit:} 14 (2d8 + 5) bludgeoning damage. If the target is a creature, it must succeed on a DC 16 Strength saving throw or be knocked prone.
\end{DndMonster}


\clearpage

\begin{DndMonster}[float*=t]{Gluttony Seraph}
\DndMonsterType{Huge celestial, chaotic evil}
\DndMonsterBasics[
armor-class = {15 (natural armor)},
hit-points = {\DndDice{19d12 + 95}},
speed = {30 ft., 30 ft. (key)}
]
\DndMonsterAbilityScores[
 str = 21,
 dex = 12,
 con = 20,
 int = 13,
 wis = 15,
 cha = 16
]
\DndMonsterDetails[
saving-throws = {Con +10, Wis +7, Cha +8},
skills = {Insight +7, Perception +7},
damage-resistances = {radiant; bludgeoning, piercing, slashing from nonmagical attacks},
damage-immunities = {poison},
senses = {truesight 120 ft., passive perception 17},
languages = {all, telepathy 120 ft.},
challenge = 14,
]
\DndMonsterAction{Magic Resistance}
The seraph has advantage on saving throws against spells and other magical effects.
\DndMonsterAction{Seraphic Weapons}
The seraph's weapon attacks are magical. When the seraph hits with any weapon, the weapon deals an extra (2d8) radiant damage (included in the attacks).
\DndMonsterAction{Innate Spellcasting}
The seraph's spellcasting ability is Charisma (spell save DC 16). The seraph can innately cast the following spells, using only verbal components:
\begin{DndMonsterSpells}
  \DndInnateSpellLevel{detect evil and good, detect magic, light, thaumaturgy}
  \DndInnateSpellLevel[1]{create food and water}
  \DndInnateSpellLevel[3]{suggestion}
\end{DndMonsterSpells}
\par
\DndMonsterSection{Actions}
\DndMonsterAction{Multiattack}
The seraph makes three attacks, only one of which can be a bite.
\DndMonsterAction{Bite}
\textsl{Melee Weapon Attack:} +10 to hit, reach 5 ft., one target. \textsl{Hit:} 24 (3d12 + 5) piercing damage and 9 (2d8) radiant damage. If the target is a Large or smaller creature, it is grappled (escape DC 18). Until this grapple ends, the target is restrained, and the seraph can't bite another target.
\DndMonsterAction{Claw}
\textsl{Melee Weapon Attack:} +10 to hit, reach 5 ft., one target. \textsl{Hit:} 12 (2d6 + 5) slashing damage and 9 (2d8) radiant damage.
\DndMonsterAction{Tongue}
\textsl{Melee Weapon Attack:} +10 to hit, reach 20 ft., one creature. \textsl{Hit:} 10 (2d4 + 5) bludgeoning damage and 9 (2d8) radiant damage, and the target is grappled (escape DC 18) and pulled to within 5 feet of the seraph. Until the grapple ends, the target is restrained, and the seraph can't use its tongue or bite on another target.
\DndMonsterAction{Swallow}
The seraph makes one bite attack against a Medium or smaller creature it holds grappled. If the attack hits, the target takes the bite's damage, the grapple ends, and the target is swallowed. While swallowed, a creature is blinded and restrained, has total cover against anything originating outside the seraph, and takes 18 (4d8) acid damage and 9 (2d8) radiant damage at the start of each of the seraph's turns. The seraph can have only one creature swallowed at a time.
If the seraph takes 35 damage or more on a single turn from a creature inside it, the seraph must succeed on a DC 20 Constitution saving throw at the end of that turn or regurgitate the swallowed creature, which falls prone in a space within 10 feet of the seraph. If the seraph dies, the swallowed creature is no longer restrained by it and can escape the corpse by using 10 feet of movement, exiting prone.
\end{DndMonster}



\begin{DndMonster}[float*=t]{Gnoll Brute}
\DndMonsterType{Medium humanoid (gnoll), chaotic evil}
\DndMonsterBasics[
armor-class = {13 (natural armor)},
hit-points = {\DndDice{8d8 + 16}},
speed = {40 ft.}
]
\DndMonsterAbilityScores[
 str = 16,
 dex = 12,
 con = 15,
 int = 6,
 wis = 13,
 cha = 7
]
\DndMonsterDetails[
skills = {Athletics +5},
senses = {darkvision 60 ft., passive perception 11},
languages = {Gnoll},
challenge = 2,
]
\DndMonsterAction{Brute}
A melee weapon deals one extra die of its damage when the gnoll brute hits with it (included in the attacks).
\DndMonsterAction{Charge}
If the gnoll brute moves at least 20 feet straight toward a target and then hits it with a melee weapon attack on the same turn, the target takes one extra die of the weapon's damage. If the target is a creature, it must succeed on a DC 13 Strength saving throw or fall prone.
\DndMonsterAction{Instinctive}
The gnoll brute has advantage on initiative rolls and saving throws against being charmed or frightened. If surprised, the gnoll can still act normally on its first turn.
\DndMonsterAction{Rampage}
When the gnoll brute reduces a creature to 0 hit points with a melee attack on its turn, the gnoll can take a bonus action to move up to half its speed and make a bite attack.
\par
\DndMonsterSection{Actions}
\DndMonsterAction{Multiattack}
The gnoll brute makes one bite attack and one claws attack.
\DndMonsterAction{Bite}
\textsl{Melee Weapon Attack:} +5 to hit, reach 5 ft., one target. \textsl{Hit:} 8 (2d4 + 3) piercing damage.
\DndMonsterAction{Claws}
\textsl{Melee Weapon Attack:} +5 to hit, reach 5 ft., one target. \textsl{Hit:} 10 (2d6 + 3) slashing damage.
\par
\DndMonsterSection{Reactions}
\DndMonsterAction{Shared Rampage}
If any creature allied with and within 100 feet of the gnoll brute uses its Rampage trait, the gnoll brute can use its Rampage trait, too, provided the brute can see that ally.
\end{DndMonster}



\begin{DndMonster}[float*=t]{Grimlock}
\DndMonsterType{Medium humanoid (grimlock), neutral evil}
\DndMonsterBasics[
armor-class = {11},
hit-points = {\DndDice{2d8 + 2}},
speed = {30 ft.}
]
\DndMonsterAbilityScores[
 str = 16,
 dex = 12,
 con = 12,
 int = 9,
 wis = 8,
 cha = 6
]
\DndMonsterDetails[
skills = {Athletics +5, Perception +3, Stealth +3},
senses = {blindsight 30 ft. or 10 ft. while deafened (blind beyond this radius), passive perception 13},
languages = {Undercommon},
challenge = 1/4,
]
\DndMonsterAction{Blind Senses}
The grimlock can't use its blindsight while deafened and unable to smell.
\DndMonsterAction{Keen Hearing and Smell}
The grimlock has advantage on Wisdom (Perception) checks that rely on hearing or smell.
\DndMonsterAction{Stone Camouflage}
The grimlock has advantage on Dexterity (Stealth) checks made to hide in rocky terrain.
\par
\DndMonsterSection{Actions}
\DndMonsterAction{Spiked Bone Club}
\textsl{Melee Weapon Attack:} +5 to hit, reach 5 ft., one target. \textsl{Hit:} 5 (1d4 + 3) bludgeoning damage plus 2 (1d4) piercing damage.
\end{DndMonster}



\begin{DndMonster}[float*=t]{Harvester of Lies}
\DndMonsterType{Medium undead, neutral evil}
\DndMonsterBasics[
armor-class = {17 (coat of lies)},
hit-points = {\DndDice{14d8 + 28}},
speed = {40 ft.}
]
\DndMonsterAbilityScores[
 str = 13,
 dex = 20,
 con = 14,
 int = 11,
 wis = 16,
 cha = 15
]
\DndMonsterDetails[
saving-throws = {Dex +9, Wis +7, Cha +6},
skills = {Acrobatics +9, Insight +7, Perception +7, Stealth +13},
damage-resistances = {necrotic; bludgeoning, piercing, slashing from nonmagical attacks},
damage-immunities = {poison},
senses = {darkvision 120 ft., passive perception 17},
languages = {the languages it knew in life},
challenge = 9,
]
\DndMonsterAction{Coat of Lies}
While wearing their specially prepared coat, the harvester's AC is increase by 2 and it can use its innate spellcasting.
\DndMonsterAction{Hear Lies}
The harvester hears all lies uttered within 1 mile of it and knows the liar's distance and direction while the harvester choses to concentrate.
\DndMonsterAction{Spider Climb}
The harvester can climb difficult surfaces, including upside down on ceilings, without needing to make an ability check.
\DndMonsterAction{Trackless}
The harvester leaves no tracks to indicate where it has been or where it's headed.
\DndMonsterAction{Innate Spellcasting}
The harvester's spellcasting ability is Charisma (spell save DC 14). It can innately cast the following spells, requiring only that the harvester be wearing its coat of lies:
\begin{DndMonsterSpells}
  \DndInnateSpellLevel{detect thoughts, zone of truth}
  \DndInnateSpellLevel[3e]{hold person, invisibility, misty step, silence}
  \DndInnateSpellLevel[1e]{levitate, tongues}
\end{DndMonsterSpells}
\par
\DndMonsterSection{Actions}
\DndMonsterAction{Multiattack}
The harvester makes three attacks, only one of which can be sever tongue.
\DndMonsterAction{Claw}
\textsl{Melee Weapon Attack:} +9 to hit, reach 5 ft., one target. \textsl{Hit:} 12 (2d6 + 5) slashing damage. If the target is Medium or smaller, it is grappled (escape DC 17).
\DndMonsterAction{Sever Tongue}
\textsl{Melee Weapon Attack:} +9 to hit, reach 5 ft., one humanoid grappled by the harvester of lies. \textsl{Hit:} 33 (8d6 + 5) slashing damage and the target must succeed on a DC 17 Constitution saving throw or have their tongue ripped out. A creature without a tongue cannot speak or cast spells with verbal components. A creature is immune to this effect if it is immune to slashing damage, doesn't have or need a tongue, or has legendary actions.
A tongue can be reattached within 1 hour of being severed with a successful DC 10 Wisdom (Medicine) check, or at can be regrown using magic that replaces lost limbs.
\end{DndMonster}


\clearpage

\begin{DndMonster}[float*=t]{Horror Flit Hunter}
\DndMonsterType{Small aberration, chaotic evil}
\DndMonsterBasics[
armor-class = {15 (natural armor)},
hit-points = {\DndDice{10d6 + 10}},
speed = {40 ft., 40 ft. (key), 40 ft. (key)}
]
\DndMonsterAbilityScores[
 str = 11,
 dex = 18,
 con = 12,
 int = 3,
 wis = 10,
 cha = 4
]
\DndMonsterDetails[
senses = {blindsight 60 ft. (blind beyond this radius), passive perception 10},
challenge = 2,
]
\DndMonsterAction{Amphibious}
The horror flit hunter can breathe air and water.
\DndMonsterAction{Breed}
When the horror flit hunter kills a Small or larger creature, a new horror flit swarm emerges from the corpse 1d4 hours later.
\DndMonsterAction{Spider Climb}
The horror flit hunter can climb difficult surfaces, including upside down on ceilings, without needing to make an ability check.
\par
\DndMonsterSection{Actions}
\DndMonsterAction{Multiattack}
The horror flit hunter makes two bite attacks.
\DndMonsterAction{Bite}
\textsl{Melee Weapon Attack:} +6 to hit, reach 5 ft., one target. \textsl{Hit:} 9 (2d4 + 4) piercing damage.
\end{DndMonster}



\begin{DndMonster}[float*=t]{Hourglass Widow}
\DndMonsterType{Medium undead, neutral evil}
\DndMonsterBasics[
armor-class = {14 (18 with mage armor)},
hit-points = {\DndDice{22d8 + 88}},
speed = {30 ft., {'number': 60, 'condition': '(hover)'} ft. (key), True ft. (key)}
]
\DndMonsterAbilityScores[
 str = 16,
 dex = 19,
 con = 18,
 int = 19,
 wis = 17,
 cha = 22
]
\DndMonsterDetails[
saving-throws = {Str +9, Con +10, Int +10, Wis +9, Cha +12},
skills = {Arcana +10, Athletics +9, Insight +9, Investigation +10, Perception +15},
damage-resistances = {cold, force, lightning, necrotic},
damage-immunities = {poison; bludgeoning, piercing, slashing from nonmagical attacks},
senses = {truesight 120 ft., passive perception 25},
languages = {the languages it knew in life},
challenge = 20,
]
\DndMonsterAction{Accelerate Through Time (2/Day)}
The widow takes one additional action.
\DndMonsterAction{Master of Time}
The widow can choose to be unaffected by effects that modify time or change its speed. The widow can affect others frozen in a time stop or other stasis-like effects.
\DndMonsterAction{Legendary Resistance (3/Day)}
If the widow fails a saving throw, it can choose to succeed instead.
\DndMonsterAction{Regeneration}
The widow regains 25 hit points at the start of its turn. If it takes fire or radiant damage, this trait doesn't function at the start of the widow's next turn. The widow's body is destroyed only if it starts its turn with 0 hit points and doesn't regenerate.
\DndMonsterAction{Turn Resistance}
The widow has advantage on saving throws against any effect that turns undead.
\DndMonsterAction{Innate Spellcasting}
The widow's spellcasting ability is Charisma (spell save DC 20, +12 to hit with spell attacks). It can innately cast the following spells, requiring no material components:
\begin{DndMonsterSpells}
  \DndInnateSpellLevel{eldritch blast, expeditious retreat, mage armor, mage hand}
  \DndInnateSpellLevel[3e]{blight, disintegrate, haste, misty step, shield, slow, see invisibility}
  \DndInnateSpellLevel[1e]{eyebite, foresight, freedom of movement, power word stun, teleport, time stop, wall of force}
\end{DndMonsterSpells}
\par
\DndMonsterSection{Actions}
\DndMonsterAction{Multiattack}
The widow uses disintegrating touch and withering gaze.
\DndMonsterAction{Disintegrating Touch}
\textsl{Melee Weapon Attack:} +10 to hit, reach 5 ft., one target. \textsl{Hit:} 16 (3d6 + 6) force damage and the target must succeed on a DC 20 Constitution saving throw or have their speed decreased by 15 feet for 1 minute. The target can repeat the saving throw at the end of each of its turns, ending the effect on itself on a success. The effects are cumulative.
\DndMonsterAction{Withering Gaze}
A creature within 30 feet that the widow can see suffers 21 (6d6) necrotic damage.
\par
\DndMonsterSection{Legendary Actions}
Hourglass Widow can take 3 legendary actions, choosing from the options below. Only one legendary action can be used at a time and only at the end of another creature's turn. Hourglass Widow regains spent legendary actions at the start of its turn.
\begin{DndMonsterLegendaryActions}
\DndMonsterLegendaryAction{Eldritch Blast}{The widow casts eldritch blast.}
\DndMonsterLegendaryAction{Sudden Rush (2 Actions)}{The widow teleports to a space it can see within 30 feet and uses disintegrating touch before or after it teleports.}
\DndMonsterLegendaryAction{Frozen in Time (3 Actions)}{All creatures within 20 ft. of the widow must succeed on a DC 20 Constitution saving throw or be paralyzed until the end of the creature's next turn.}
\end{DndMonsterLegendaryActions}
\end{DndMonster}



\begin{DndMonster}[float*=t]{Hraptnon}
\DndMonsterType{Gargantuan monstrosity, chaotic evil}
\DndMonsterBasics[
armor-class = {18 (natural armor)},
hit-points = {\DndDice{16d20 + 96}},
speed = {40 ft., 20 ft. (key), 30 ft. (key)}
]
\DndMonsterAbilityScores[
 str = 21,
 dex = 17,
 con = 22,
 int = 12,
 wis = 17,
 cha = 12
]
\DndMonsterDetails[
saving-throws = {Str +11, Dex +9, Con +12, Wis +9},
skills = {Athletics +11, Perception +9},
damage-immunities = {; bludgeoning, piercing, slashing from magical weapons; acid, bludgeoning, cold, fire, force, lightning, necrotic, piercing, poison, psychic, radiant, slashing, thunder from spells},
senses = {darkvision 120 ft., passive perception 19},
languages = {Draconic},
challenge = 18,
]
\DndMonsterAction{Legendary Resistance (3/Day)}
If the hraptnon fails a saving throw, it can choose to succeed instead.
\DndMonsterAction{Magic Resistance}
The hraptnon has advantage on saving throws against spells and other magical effects.
\DndMonsterAction{Undying}
Unless the hraptnon's Stillborn Heart is removed from its pool, the hraptnon's death is temporary. One of its countless xakalonus offspring grows into the hraptnon in 24 hours.
In addition, the hraptnon loses its immunities if the Stillborn Heart is removed from the pool.
\par
\DndMonsterSection{Actions}
\DndMonsterAction{Multiattack}
The hraptnon makes three melee attacks: one bite and two claws. Alternatively, the hraptnon can make two spine attacks.
\DndMonsterAction{Bite}
\textsl{Melee Weapon Attack:} +11 to hit, reach 5 ft., one target. \textsl{Hit:} 18 (3d8 + 5) piercing damage, and the target must succeed on a DC 20 Constitution saving throw or be incapacitated until the end of the hraptnon's next turn.
\DndMonsterAction{Claws}
\textsl{Melee Weapon Attack:} +11 to hit, reach 10 ft., one target. \textsl{Hit:} 27 (4d10 + 5) slashing damage.
\DndMonsterAction{Spine}
\textsl{Melee Weapon Attack:} +11 to hit, range 60 ft., one target. \textsl{Hit:} 23 (4d8 + 5) piercing damage, and the target must succeed on a DC 20 Dexterity saving throw. On a failed save, the target is knocked prone and restrained. The target must escape (DC 20) to remove the restrained condition and be able to stand.
\par
\DndMonsterSection{Reactions}
\DndMonsterAction{Spawn Xakalonus}
If the hraptnon takes 20 points of damage or more from a single attack, it can spawn a xakalonus, which appears in a space adjacent to the hraptnon and attacks on the same initiative count as the hraptnon.
\par
\DndMonsterSection{Legendary Actions}
Hraptnon can take 3 legendary actions, choosing from the options below. Only one legendary action can be used at a time and only at the end of another creature's turn. Hraptnon regains spent legendary actions at the start of its turn.
\begin{DndMonsterLegendaryActions}
\DndMonsterLegendaryAction{Claw}{The hraptnon makes a claw attack.}
\DndMonsterLegendaryAction{Baneful Howl (Costs 2 Actions)}{The hraptnon casts bane as a 3rd level spell (DC 20 to resist).}
\DndMonsterLegendaryAction{Life Drain (Costs 3 Actions)}{One xakalonus within 30 feet of the hraptnon dies. The hraptnon regains hit points equal to half the hit points of the xakalonus that died.}
\end{DndMonsterLegendaryActions}
\end{DndMonster}



\begin{DndMonster}[float*=t]{Incubus}
\DndMonsterType{Medium fiend (shapechanger), neutral evil}
\DndMonsterBasics[
armor-class = {15 (natural armor)},
hit-points = {\DndDice{12d8 + 12}},
speed = {30 ft., 60 ft. (key)}
]
\DndMonsterAbilityScores[
 str = 8,
 dex = 17,
 con = 13,
 int = 15,
 wis = 12,
 cha = 20
]
\DndMonsterDetails[
skills = {Deception +9, Insight +5, Perception +5, Persuasion +9, Stealth +7},
damage-resistances = {cold, fire, lightning, poison; bludgeoning, piercing, slashing from nonmagical attacks},
senses = {darkvision 60 ft., passive perception 15},
languages = {Abyssal, Common, Infernal, telepathy 60 ft.},
challenge = 4,
]
\DndMonsterAction{Telepathic Bond}
The fiend ignores the range restriction on its telepathy when communicating with a creature it has charmed. The two don't even need to be on the same plane of existence.
\DndMonsterAction{Shapechanger}
The fiend can use its action to polymorph into a Small or Medium humanoid, or back into its true form. Without wings, the fiend loses its flying speed. Other than its size and speed, its statistics are the same in each form. Any equipment it is wearing or carrying isn't transformed. It reverts to its true form if it dies.
\par
\DndMonsterSection{Actions}
\DndMonsterAction{Claw (Fiend Form Only)}
\textsl{Melee Weapon Attack:} +5 to hit, reach 5 ft., one target. \textsl{Hit:} 6 (1d6 + 3) slashing damage.
\DndMonsterAction{Charm}
One humanoid the fiend can see within 30 feet of it must succeed on a DC 15 Wisdom saving throw or be magically charmed for 1 day. The charmed target obeys the fiend's verbal or telepathic commands. If the target suffers any harm or receives a suicidal command, it can repeat the saving throw, ending the effect on a success. If the target successfully saves against the effect, or if the effect on it ends, the target is immune to this fiend's Charm for the next 24 hours.
The fiend can have only one target charmed at a time. If it charms another, the effect on the previous target ends.
\DndMonsterAction{Draining Kiss}
The fiend kisses a creature charmed by it or a willing creature. The target must make a DC 15 Constitution saving throw against this magic, taking 32 (5d10 + 5) psychic damage on a failed save, or half as much damage on a successful one. The target's hit point maximum is reduced by an amount equal to the damage taken. This reduction lasts until the target finishes a long rest. The target dies if this effect reduces its hit point maximum to 0.
\DndMonsterAction{Etherealness}
The fiend magically enters the Ethereal Plane from the Material Plane, or vice versa.
\end{DndMonster}


\clearpage

\begin{DndMonster}[float*=t]{Infernal Tormentor}
\DndMonsterType{Medium humanoid (any race, fiend), any evil alignment}
\DndMonsterBasics[
armor-class = {16 (half plate armor)},
hit-points = {\DndDice{13d8 + 39}},
speed = {30 ft.}
]
\DndMonsterAbilityScores[
 str = 18,
 dex = 12,
 con = 16,
 int = 12,
 wis = 13,
 cha = 15
]
\DndMonsterDetails[
skills = {Intimidation +5, Perception +4},
damage-resistances = {fire},
senses = {darkvision 60 ft., passive perception 14},
languages = {Infernal and one other language},
challenge = 6,
]
\par
\DndMonsterSection{Actions}
\DndMonsterAction{Multiattack}
The infernal tormentor makes two greatsword attacks. It can use infernal leash in place of one attack.
\DndMonsterAction{Greatsword}
\textsl{Melee Weapon Attack:} +7 to hit, reach 5 ft., one target. \textsl{Hit:} 11 (2d6 + 4) slashing damage and 7 (2d6) fire damage.
\DndMonsterAction{Infernal Leash Recharge 6}
A flaming whip appears in the infernal tormentor's hand or from its greatsword and lashes at a creature the tormentor can see within 30 feet of it. The target must succeed on a DC 13 Charisma saving throw. A creature that has signed a fiendish contract has disadvantage on this save. On a failure, the creature takes 14 (4d6) fire damage, is pulled 10 feet closer to the tormentor, and is cursed for 1 minute or until the tormentor uses this action again. This curse also ends if the target enters hallowed ground.
While cursed, the target can't move more than 30 feet from the tormentor. Also while the curse lasts, the target takes 7 (2d6) extra fire damage from the tormentor's greatsword attacks. The tormentor can use a bonus action to move the target 10 feet to an unoccupied space within 30 feet of the tormentor. If this space is dangerous, the target receives a new saving throw before moving, ending the effect on itself on a success.
\end{DndMonster}



\begin{DndMonster}[float*=t]{Ithjar}
\DndMonsterType{Large monstrosity, chaotic neutral}
\DndMonsterBasics[
armor-class = {16 (natural armor)},
hit-points = {\DndDice{12d10 + 36}},
speed = {30 ft., 90 ft. (key)}
]
\DndMonsterAbilityScores[
 str = 18,
 dex = 20,
 con = 17,
 int = 3,
 wis = 10,
 cha = 5
]
\DndMonsterDetails[
skills = {Acrobatics +8, Perception +3},
senses = {blindsight 10 ft., darkvision 60 ft., passive perception 13},
challenge = 5,
]
\DndMonsterAction{Keen Sight}
The ithjar has advantage on Wisdom (Perception) checks that rely on sight.
\DndMonsterAction{Slippery Serpent}
The ithjar has advantage on ability checks and saving throws made to escape a grapple or saving throws to avoid being paralyzed or restrained.
\par
\DndMonsterSection{Actions}
\DndMonsterAction{Multiattack}
The ithjar attacks once with its bite and once with its tail.
\DndMonsterAction{Bite}
\textsl{Melee Weapon Attack:} +7 to hit, reach 10 ft., one target. \textsl{Hit:} 15 (2d10 + 4) slashing damage.
\DndMonsterAction{Tail}
\textsl{Melee Weapon Attack:} +7 to hit, reach 10 ft., one target. \textsl{Hit:} 13 (2d8 + 4) bludgeoning damage. If the target is a Medium or smaller creature, it is grappled (escape DC 15). Until this grapple ends, the target is restrained, and the ithjar can't use its tail on another target.
\par
\DndMonsterSection{Reactions}
\DndMonsterAction{Writhing Escape}
When grappled, the ithjar can attempt to escape.
\end{DndMonster}



\begin{DndMonster}[float*=t]{Knifewing}
\DndMonsterType{Small beast, unaligned}
\DndMonsterBasics[
armor-class = {13 (natural armor)},
hit-points = {\DndDice{2d6 + 2}},
speed = {30 ft., 30 ft. (key), 30 ft. (key)}
]
\DndMonsterAbilityScores[
 str = 6,
 dex = 15,
 con = 13,
 int = 2,
 wis = 10,
 cha = 5
]
\DndMonsterDetails[
senses = {blindsight 10 ft., passive perception 10},
challenge = 1/8,
]
\DndMonsterAction{Flyby}
The knifewing doesn't provoke opportunity attacks when it flies out of an enemy's reach.
\DndMonsterAction{Glider}
The knifewing can't ascend while flying and its altitude decreases by 5 feet for every 30 feet it flies. If flying, the knifewing falls if it doesn't fly at least 30 feet during its turn.
\par
\DndMonsterSection{Actions}
\DndMonsterAction{Knifewing (Flying Only)}
\textsl{Melee Weapon Attack:} +4 to hit, reach 5 ft., one target. \textsl{Hit:} 5 (1d6 + 2) slashing damage.
\DndMonsterAction{Bite}
\textsl{Melee Weapon Attack:} +4 to hit, reach 5 ft., one target. \textsl{Hit:} 4 (1d4 + 2) piercing damage.
\end{DndMonster}



\begin{DndMonster}[float*=t]{Lenchtahg}
\DndMonsterType{Medium fiend, neutral evil}
\DndMonsterBasics[
armor-class = {15 (natural armor)},
hit-points = {\DndDice{13d8 + 52}},
speed = {40 ft.}
]
\DndMonsterAbilityScores[
 str = 18,
 dex = 11,
 con = 18,
 int = 10,
 wis = 11,
 cha = 13
]
\DndMonsterDetails[
skills = {Perception +3},
damage-resistances = {; bludgeoning, piercing, slashing from nonmagical attacks},
damage-immunities = {fire, poison, radiant},
senses = {darkvision 120 ft., passive perception 13},
languages = {Abyssal, Celestial, Infernal},
challenge = 5,
]
\DndMonsterAction{Berserk}
Whenever the lenchtahg starts its turn with 54 hit points or fewer, roll a d6. On a 1, the lenchtahg recalls its former life and current agony, and it goes berserk. On each of its turns while berserk, the lenchtahg attacks the nearest creature it can see. It treats all creatures as enemies and makes any opportunity attack it can. Also, it has advantage on attack rolls, but all attack rolls against it have advantage. If no creature is near enough to move to and attack, the lenchtahg moves as far as it can toward the closest creature.
Once the lenchtahg goes berserk, it remains so until destroyed or it has more than 54 hit points. Also, a fiend can calm the lenchtahg. The lenchtahg must be able to understand the fiend, who must take an action to make a DC 15 Charisma (Persuasion) check. If the check succeeds, the lenchtahg ceases being berserk.
\DndMonsterAction{Magic Resistance}
The lenchtahg has advantage on saving throws against spells and other magical effects.
\DndMonsterAction{Magic Weapons}
The lenchtahg's weapon attacks are magical.
\DndMonsterAction{Radiant Absorption}
Whenever the lenchtahg is subjected to radiant damage from a source other than a lenchtahg, it takes no damage and instead regains hit points equal to half the radiant damage dealt to it.
\DndMonsterAction{Rejuvenation}
If the lenchtahg's head survives its demise, the head must be doused in holy water and ritually consecrated in a 1-hour ceremony or be the target of a dispel evil and good spell. If it isn't, then the lenchtahg reappears in 6 days, emerging from a fiery explosion in a 20-foot-radius sphere centered on the head. Creatures in that area must make a DC 15 Dexterity saving throw, taking 28 8d6 damage on a failed saving throw, or half as much damage on a successful one. The head can't rejuvenate in a holy hallow spell's area.
\DndMonsterAction{Water Aversion}
If drenched in 5 gallons or more of water, or any amount of holy water, the lenchtahg has disadvantage on attack rolls and ability checks until the end of its next turn.
\par
\DndMonsterSection{Actions}
\DndMonsterAction{Multiattack}
The lenchtahg makes two claw attacks, makes one claw attack and uses radiant gaze once, or uses radiant gaze twice.
\DndMonsterAction{Claw}
\textsl{Melee Weapon Attack:} +7 to hit, reach 5 ft., one target. \textsl{Hit:} 9 (2d4 + 4) slashing damage and 9 (2d8) fire damage.
\DndMonsterAction{Radiant Gaze}
The lenchtahg targets one creature it can see within 60 feet of it with burning radiance. The target gains no benefit from cover and must succeed on a DC 15 Dexterity saving throw or take 13 (3d8) radiant damage.
\end{DndMonster}


\clearpage

\begin{DndMonster}[float*=t]{Lesser Avarice Seraph}
\DndMonsterType{Medium celestial, neutral evil}
\DndMonsterBasics[
armor-class = {15 (natural armor)},
hit-points = {\DndDice{14d8 + 42}},
speed = {30 ft., 90 ft. (key)}
]
\DndMonsterAbilityScores[
 str = 16,
 dex = 18,
 con = 17,
 int = 13,
 wis = 15,
 cha = 16
]
\DndMonsterDetails[
saving-throws = {Con +6, Wis +5, Cha +6},
skills = {Deception +6, Insight +5, Perception +5},
damage-resistances = {radiant; bludgeoning, piercing, slashing from nonmagical attacks},
senses = {darkvision 120 ft., passive perception 15},
languages = {all, telepathy 120 ft.},
challenge = 7,
]
\DndMonsterAction{Magic Resistance}
The seraph has advantage on saving throws against spells and other magical effects.
\DndMonsterAction{Seraphic Weapons}
The seraph's weapon attacks are magical. When the seraph hits with any weapon, the weapon deals an extra (2d8) radiant damage (included in the attack).
\DndMonsterAction{Innate Spellcasting}
The seraph's spellcasting ability is Charisma (spell save DC 14). The seraph can innately cast the following spells, using only verbal components:
\begin{DndMonsterSpells}
  \DndInnateSpellLevel{detect evil and good, detect magic, light, minor illusion, thaumaturgy}
  \DndInnateSpellLevel[1e]{identify, suggestion}
\end{DndMonsterSpells}
\par
\DndMonsterSection{Actions}
\DndMonsterAction{Multiattack}
The seraph makes three attacks.
\DndMonsterAction{Scimitar}
\textsl{Melee Weapon Attack:} +7 to hit, reach 5 ft., one target. \textsl{Hit:} 6 (1d6 + 3) slashing damage and 9 (2d8) radiant damage.
\DndMonsterAction{Tempting Call}
The seraph calls out a tempting promise to one creature that can hear it withing 30 feet. The creature must succeed on a DC 14 Charisma saving throw. On a failed save, the creature cannot take actions for 1 minute, as it dwells on the promise the seraph made. The target can repeat the saving throw at the end of each turn, ending the effect on a success. Each failed saving throw after the initial failure deals 9 (2d8) psychic damage to the target
\end{DndMonster}



\begin{DndMonster}[float*=t]{Lesser Gluttony Seraph}
\DndMonsterType{Large celestial, chaotic evil}
\DndMonsterBasics[
armor-class = {15 (natural armor)},
hit-points = {\DndDice{12d10 + 48}},
speed = {30 ft., 30 ft. (key)}
]
\DndMonsterAbilityScores[
 str = 18,
 dex = 14,
 con = 19,
 int = 12,
 wis = 14,
 cha = 15
]
\DndMonsterDetails[
saving-throws = {Con +7, Wis +5, Cha +5},
skills = {Insight +5, Perception +5},
damage-resistances = {radiant; bludgeoning, piercing, slashing from nonmagical attacks},
damage-immunities = {poison},
senses = {darkvision 120 ft., passive perception 15},
languages = {all, telepathy 120 ft.},
challenge = 7,
]
\DndMonsterAction{Magic Resistance}
The seraph has advantage on saving throws against spells and other magical effects.
\DndMonsterAction{Seraphic Weapons}
The seraph's weapon attacks are magical. When the seraph hits with any weapon, the weapon deals an extra (1d8) radiant damage (included in the attacks).
\DndMonsterAction{Innate Spellcasting}
The seraph's spellcasting ability is Charisma (spell save DC 13). The seraph can innately cast the following spells, using only verbal components:
\begin{DndMonsterSpells}
  \DndInnateSpellLevel{detect evil and good, detect magic, light, thaumaturgy}
  \DndInnateSpellLevel[1e]{create food and water, suggestion}
\end{DndMonsterSpells}
\par
\DndMonsterSection{Actions}
\DndMonsterAction{Multiattack}
The seraph makes three attacks, only one of which can be a bite.
\DndMonsterAction{Bite}
\textsl{Melee Weapon Attack:} +7 to hit, reach 5 ft., one target. \textsl{Hit:} 15 (2d10 + 4) piercing damage and 4 (1d8) radiant damage. If the target is a Medium or smaller creature, it is grappled (escape DC 15). Until this grapple ends, the target is restrained, and the seraph can't bite another target.
\DndMonsterAction{Claw}
\textsl{Melee Weapon Attack:} +7 to hit, reach 5 ft., one target. \textsl{Hit:} 9 (2d4 + 4) slashing damage and 4 (1d8) radiant damage.
\DndMonsterAction{Tongue}
\textsl{Melee Weapon Attack:} +7 to hit, reach 15 ft., one creature. \textsl{Hit:} 6 (1d4 + 4) bludgeoning damage and 4 (1d8) radiant damage, and the target is grappled (escape DC 15) and pulled to within 5 feet of the seraph. Until the grapple ends, the target is restrained, and the seraph can't use its tongue or bite on another target.
\end{DndMonster}



\begin{DndMonster}[float*=t]{Lich Troll}
\DndMonsterType{Large undead, chaotic evil}
\DndMonsterBasics[
armor-class = {16 (natural armor)},
hit-points = {\DndDice{14d10 + 56}},
speed = {30 ft.}
]
\DndMonsterAbilityScores[
 str = 18,
 dex = 14,
 con = 18,
 int = 15,
 wis = 13,
 cha = 12
]
\DndMonsterDetails[
saving-throws = {Con +9, Int +7, Wis +6},
damage-resistances = {cold, necrotic},
damage-immunities = {fire, poison},
senses = {truesight 60 ft., passive perception 16},
languages = {the languages it knew in life and Giant},
challenge = 14,
]
\DndMonsterAction{Legendary Resistance (2/Day)}
If the lich troll fails a saving throw, it can choose to succeed instead.
\DndMonsterAction{Regeneration}
The lich troll regains 15 hit points at the start of its turn. If the lich troll takes acid or radiant damage, this trait doesn't function until the start of the lich troll's next turn. The lich troll dies only if it starts its turn with 0 hit points and doesn't regenerate.
\DndMonsterAction{Spellcasting}
The lich troll is an 11th-level spellcaster. Its spellcasting ability is Intelligence (spell save DC 15, +7 to hit with spell attacks). The lich troll has the following wizard spells prepared:
\begin{DndMonsterSpells}
  \DndMonsterSpellLevel{chill touch, mage hand, prestidigitation, ray of frost}
   \DndMonsterSpellLevel[1][4]{detect magic, magic missile, shield, thunderwave}
   \DndMonsterSpellLevel[2][3]{detect thoughts, invisibility, mirror image, misty step}
   \DndMonsterSpellLevel[3][3]{animate dead, counterspell, dispel magic, haste}
   \DndMonsterSpellLevel[4][3]{blight, dimension door}
   \DndMonsterSpellLevel[5][2]{cloudkill, scrying}
   \DndMonsterSpellLevel[6][1]{create undead}
\end{DndMonsterSpells}
\par
\DndMonsterSection{Actions}
\DndMonsterAction{Multiattack}
The lich troll makes two claw attacks. Alternately, the lich troll can cast chill touch or ray of frost twice.
\DndMonsterAction{Claw}
\textsl{Melee Weapon Attack:} +9 to hit, reach 5 ft., one target. \textsl{Hit:} 11 (2d6 + 4) slashing damage plus 7 (2d6) cold damage.
\par
\DndMonsterSection{Legendary Actions}
Lich Troll can take 2 legendary actions, choosing from the options below. Only one legendary action can be used at a time and only at the end of another creature's turn. Lich Troll regains spent legendary actions at the start of its turn.
\begin{DndMonsterLegendaryActions}
\DndMonsterLegendaryAction{Cantrip}{The lich troll casts a cantrip.}
\DndMonsterLegendaryAction{Claw}{The lich troll makes a claw attack.}
\end{DndMonsterLegendaryActions}
\end{DndMonster}



\begin{DndMonster}[float*=t]{Lindwyrm}
\DndMonsterType{Gargantuan dragon, chaotic neutral}
\DndMonsterBasics[
armor-class = {16 (natural armor)},
hit-points = {\DndDice{12d20 + 60}},
speed = {40 ft., 30 ft. (key)}
]
\DndMonsterAbilityScores[
 str = 23,
 dex = 15,
 con = 21,
 int = 8,
 wis = 14,
 cha = 12
]
\DndMonsterDetails[
saving-throws = {Str +11, Dex +7, Con +10},
skills = {Perception +7, Stealth +7},
damage-resistances = {; bludgeoning, piercing, slashing from nonmagical attacks},
damage-immunities = {poison},
senses = {darkvision 120 ft., tremorsense 120 ft., passive perception 17},
languages = {Draconic},
challenge = 14,
]
\DndMonsterAction{Disease Immunity}
The lindwyrm is immune to all nonmagical diseases.
\DndMonsterAction{Legendary Resistance (3/Day)}
If the lindwyrm fails a saving throw, it can choose to succeed instead.
\par
\DndMonsterSection{Actions}
\DndMonsterAction{Multiattack}
The lindwyrm makes three attacks: one with its bite and two with its claws. The lindwyrm can substitute a swallow attack for one of its claw attacks.
\DndMonsterAction{Bite}
\textsl{Melee Weapon Attack:} +11 to hit, reach 15 ft., one target. \textsl{Hit:} 22 (3d10 + 6) piercing damage and the target must make a DC 18 Constitution saving throw, taking 18 (4d8) poison damage on a failed save, or half as much damage on a successful one. On a hit, the target is also grappled (escape DC 18). Until this grapple ends, the target is restrained, and the lindwyrm can't bite another target.
\DndMonsterAction{Claw}
\textsl{Melee Weapon Attack:} +11 to hit, reach 10 ft., one target. \textsl{Hit:} 15 (2d8 + 6) slashing damage.
\DndMonsterAction{Swallow}
\textsl{Melee Weapon Attack:} +11 to hit, reach 5 ft., one Huge or smaller creature the lindwyrm is grappling. \textsl{Hit:} 22 (3d10 + 6) piercing damage, the target is swallowed, and the grapple ends. The swallowed target is blinded and restrained, it has total cover against attacks and other effects outside the lindwyrm, and it takes 27 (6d8) acid damage at the start of each of the lindwyrm's turns. The lindwyrm can have only one target swallowed at a time.
If the lindwyrm takes 40 damage or more on a single turn from the swallowed creature, the lindwyrm must succeed on a DC 18 Constitution saving throw at the end of that turn or regurgitate the creature, which falls prone in a space within 10 feet of the lindwyrm. If the lindwyrm dies, a swallowed creature is no longer restrained by it and can escape from the corpse using 5 feet of movement, exiting prone.
\par
\DndMonsterSection{Legendary Actions}
Lindwyrm can take 2 legendary actions, choosing from the options below. Only one legendary action can be used at a time and only at the end of another creature's turn. Lindwyrm regains spent legendary actions at the start of its turn.
\begin{DndMonsterLegendaryActions}
\DndMonsterLegendaryAction{Claw Attack}{The lindwyrm makes a claw attack.}
\DndMonsterLegendaryAction{Burrowing Rush (Costs 2 Actions)}{The lindwyrm burrows its speed and comes up beneath a target touching the ground, making a bite attack. This movement does not provoke opportunity attacks.}
\end{DndMonsterLegendaryActions}
\end{DndMonster}


\clearpage

\begin{DndMonster}[float*=t]{Lupilisk}
\DndMonsterType{Large fiend, neutral evil}
\DndMonsterBasics[
armor-class = {13},
hit-points = {\DndDice{8d10 + 16}},
speed = {40 ft.}
]
\DndMonsterAbilityScores[
 str = 16,
 dex = 17,
 con = 14,
 int = 5,
 wis = 13,
 cha = 7
]
\DndMonsterDetails[
skills = {Perception +3, Stealth +5},
damage-resistances = {cold, fire},
damage-immunities = {poison},
senses = {darkvision 120 ft., passive perception 13},
languages = {can learn languages but can't speak},
challenge = 3,
]
\DndMonsterAction{Keen Hearing and Smell}
The lupilisk has advantage on Wisdom (Perception) checks that rely on hearing or smell.
\DndMonsterAction{Pack Tactics}
The lupilisk has advantage on an attack roll against a creature if at least one of the lupilisk's allies is within 5 feet of the creature and the ally isn't incapacitated.
\par
\DndMonsterSection{Actions}
\DndMonsterAction{Bite}
\textsl{Melee Weapon Attack:} +5 to hit, reach 5 ft., one target. \textsl{Hit:} 10 (2d6 + 3) piercing damage and 7 (2d6) poison damage. If the target is a creature, it must succeed on a DC 12 Constitution saving throw or become poisoned and restrained while poisoned in this way. The creature must repeat the saving throw at the end of its next turn. On a success, the effect ends. If the save fails again, the creature is petrified for 24 hours. At that time, it can repeat the saving throw, becoming permanently petrified on a failure or ending the effect on a success.
\end{DndMonster}



\begin{DndMonster}[float*=t]{Lupilisk Elder}
\DndMonsterType{Large fiend, neutral evil}
\DndMonsterBasics[
armor-class = {15 (natural armor)},
hit-points = {\DndDice{13d10 + 26}},
speed = {50 ft.}
]
\DndMonsterAbilityScores[
 str = 18,
 dex = 16,
 con = 15,
 int = 6,
 wis = 14,
 cha = 8
]
\DndMonsterDetails[
skills = {Perception +5, Stealth +6},
damage-resistances = {cold, fire},
damage-immunities = {poison},
senses = {darkvision 120 ft., passive perception 15},
languages = {can learn languages but can't speak},
challenge = 6,
]
\DndMonsterAction{Keen Hearing and Smell}
The lupilisk elder has advantage on Wisdom (Perception) checks that rely on hearing or smell.
\DndMonsterAction{Pack Tactics}
The lupilisk elder has advantage on an attack roll against a creature if at least one of the lupilisk's allies is within 5 feet of the creature and the ally isn't incapacitated.
\par
\DndMonsterSection{Actions}
\DndMonsterAction{Bite}
\textsl{Melee Weapon Attack:} +7 to hit, reach 5 ft., one target. \textsl{Hit:} 13 (2d8 + 4) piercing damage and 9 (2d8) poison damage. If the target is a creature, it must succeed on a DC 13 Constitution saving throw or become poisoned and restrained while poisoned in this way. The creature must repeat the saving throw at the end of its next turn. On a success, the effect ends. If the save fails again, the creature is permanently petrified.
\DndMonsterAction{Pack Attack Recharge 5-6}
One lupilisk or lupilisk whelp the lupilisk elder can see within 30 feet of it can use a reaction to move up to half its speed and make a bite attack. The target must be able to hear or see the elder. The lupilisk elder then makes one bite attack.
\end{DndMonster}



\begin{DndMonster}[float*=t]{Lupilisk Whelp}
\DndMonsterType{Medium fiend, neutral evil}
\DndMonsterBasics[
armor-class = {12},
hit-points = {\DndDice{6d8 + 6}},
speed = {40 ft.}
]
\DndMonsterAbilityScores[
 str = 13,
 dex = 14,
 con = 12,
 int = 3,
 wis = 11,
 cha = 6
]
\DndMonsterDetails[
skills = {Perception +2, Stealth +4},
damage-resistances = {cold, fire},
damage-immunities = {poison},
senses = {darkvision 120 ft., passive perception 12},
languages = {can learn languages but can't speak},
challenge = 1,
]
\DndMonsterAction{Keen Hearing and Smell}
The lupilisk whelp has advantage on Wisdom (Perception) checks that rely on hearing or smell.
\DndMonsterAction{Pack Tactics}
The lupilisk whelp has advantage on an attack roll against a creature if at least one of the lupilisk's allies is within 5 feet of the creature and the ally isn't incapacitated.
\par
\DndMonsterSection{Actions}
\DndMonsterAction{Bite}
\textsl{Melee Weapon Attack:} +4 to hit, reach 5 ft., one target. \textsl{Hit:} 7 (2d4 + 2) piercing damage and 5 (2d4) poison damage. If the target is a creature, it must succeed on a DC 11 Constitution saving throw or become poisoned and restrained while poisoned in this way. The creature must repeat the saving throw at the end of its next turn. On a success, the effect ends. If the save fails again, the creature is petrified for 24 hours.
\end{DndMonster}



\begin{DndMonster}[float*=t]{Mage}
\DndMonsterType{Medium humanoid (any race), any alignment}
\DndMonsterBasics[
armor-class = {12 (15 with mage armor)},
hit-points = {\DndDice{9d8}},
speed = {30 ft.}
]
\DndMonsterAbilityScores[
 str = 9,
 dex = 14,
 con = 11,
 int = 17,
 wis = 12,
 cha = 11
]
\DndMonsterDetails[
saving-throws = {Int +6, Wis +4},
skills = {Arcana +6, History +6},
languages = {any four languages},
challenge = 6,
]
\DndMonsterAction{Spellcasting}
The mage is a 9th-level spellcaster. Its spellcasting ability is Intelligence (spell save DC 14, +6 to hit with spell attacks). The mage has the following wizard spells prepared:
\begin{DndMonsterSpells}
  \DndMonsterSpellLevel{fire bolt, light, mage hand, prestidigitation}
   \DndMonsterSpellLevel[1][4]{detect magic, mage armor, magic missile, shield}
   \DndMonsterSpellLevel[2][3]{misty step, suggestion}
   \DndMonsterSpellLevel[3][3]{counterspell, fireball, fly}
   \DndMonsterSpellLevel[4][3]{greater invisibility, ice storm}
   \DndMonsterSpellLevel[5][1]{cone of cold}
\end{DndMonsterSpells}
\par
\DndMonsterSection{Actions}
\DndMonsterAction{Dagger}
\textsl{Melee or Ranged Weapon Attack:} +5 to hit, reach 5 ft. or range 20/60 ft., one target. \textsl{Hit:} 4 (1d4 + 2) piercing damage.
\end{DndMonster}


\clearpage

\begin{DndMonster}[float*=t]{Malikirian Imp}
\DndMonsterType{Tiny fiend, lawful evil}
\DndMonsterBasics[
armor-class = {11},
hit-points = {\DndDice{6d4}},
speed = {20 ft., 40 ft. (key)}
]
\DndMonsterAbilityScores[
 str = 7,
 dex = 13,
 con = 10,
 int = 11,
 wis = 12,
 cha = 14
]
\DndMonsterDetails[
saving-throws = {Wis +3, Cha +4},
skills = {Deception +4, Insight +3, Stealth +3},
damage-resistances = {cold; bludgeoning, piercing, slashing from nonmagical attacks},
damage-immunities = {fire, poison},
senses = {darkvision 120 ft., passive perception 11},
languages = {Infernal and 2 other languages},
challenge = 1,
]
\DndMonsterAction{Creeping Corruption}
Creatures within 30 ft. that the Malikirian imp can see must make a DC 12 Charisma saving throw at the start of their turn. On a failure, the creature acts in accordance with its most selfish, venal, and ruthless impulses to the best of its ability. After 24 hours, the creature must make another Charisma saving throw, ending the compulsion on a success. Failure indicates the compulsion lasts indefinitely until the affected is subject to remove curse or similar magic. Creatures that cannot be charmed are immune to this effect.
\DndMonsterAction{Magic Resistance}
The Malikirian imp has advantage on saving throws against spells and other magical effects.
\DndMonsterAction{Innate Spellcasting}
The Malikirian imp's spellcasting ability is Charisma (spell save DC 12). It can innately cast the following spells without material components:
\begin{DndMonsterSpells}
  \DndInnateSpellLevel{minor illusion}
  \DndInnateSpellLevel[1e]{detect thoughts, sleep, suggestion}
\end{DndMonsterSpells}
\par
\DndMonsterSection{Actions}
\DndMonsterAction{Bite}
\textsl{Melee Weapon Attack:} -1 to hit, reach 0 ft., one target. \textsl{Hit:} 1 (1d4 - 2) piercing damage and the target must succeed on a DC 10 Constitution saving throw or take 7 (2d6) poison damage.
\DndMonsterAction{Invisibility}
The Malikirian imp becomes invisible until its concentration ends or it makes an attack.
\end{DndMonster}



\begin{DndMonster}[float*=t]{Mjork}
\DndMonsterType{Medium elemental, chaotic neutral}
\DndMonsterBasics[
armor-class = {14 (natural armor)},
hit-points = {\DndDice{8d8 + 24}},
speed = {30 ft.}
]
\DndMonsterAbilityScores[
 str = 15,
 dex = 12,
 con = 16,
 int = 6,
 wis = 10,
 cha = 6
]
\DndMonsterDetails[
damage-vulnerabilities = {thunder},
damage-resistances = {; bludgeoning, piercing, slashing from nonmagical attacks},
damage-immunities = {fire, poison},
senses = {darkvision 60 ft., passive perception 10},
languages = {Ignan, Terran},
challenge = 3,
]
\DndMonsterAction{Fiery Body}
A creature that touches the mjork or hits it with a melee attack while within 5 feet of it takes 3 (1d6) fire damage.
\DndMonsterAction{Illumination}
The mjork sheds dim light in a 10-foot radius.
\DndMonsterAction{Smoke Sense}
A mjork in contact with smoke has blindsight in that smoke's area in a radius of up to 120 feet.
\DndMonsterAction{Water Susceptibility}
For every 5 feet the mjork moves in water, or for every gallon of water splashed on it, it takes 1 cold damage.
\par
\DndMonsterSection{Actions}
\DndMonsterAction{Multiattack}
The mjork makes two slam attacks.
\DndMonsterAction{Slam}
\textsl{Melee Weapon Attack:} +4 to hit, reach 5 ft., one target. \textsl{Hit:} 6 (1d8 + 2) bludgeoning damage and 3 (1d6) fire damage.
\end{DndMonster}



\begin{DndMonster}[float*=t]{Mjork Asher}
\DndMonsterType{Large elemental, chaotic neutral}
\DndMonsterBasics[
armor-class = {14 (natural armor)},
hit-points = {\DndDice{10d10 + 30}},
speed = {30 ft.}
]
\DndMonsterAbilityScores[
 str = 15,
 dex = 15,
 con = 16,
 int = 7,
 wis = 13,
 cha = 7
]
\DndMonsterDetails[
damage-vulnerabilities = {thunder},
damage-resistances = {; bludgeoning, piercing, slashing from nonmagical attacks},
damage-immunities = {fire, poison},
senses = {darkvision 60 ft., passive perception 11},
languages = {Ignan, Terran},
challenge = 6,
]
\DndMonsterAction{Blinding Soot}
A creature that takes fire damage from a mjork asher's attacks must make a DC 14 Constitution saving throw or be blinded until the end of the mjork's next turn.
\DndMonsterAction{Fiery Body}
A creature that touches the mjork asher or hits it with a melee attack while within 5 feet of it takes 7 (2d6) fire damage.
\DndMonsterAction{Smoke Sense}
A mjork asher in contact with smoke has blindsight in that smoke's area in a radius of up to 120 feet.
\DndMonsterAction{Smoke Shroud}
While the mjork asher has 41 or more hit points, smoke renders a sphere within 20 feet of it heavily obscured. This smoke goes around corners. Even while the mjork has 40 or fewer hit points, tendrils of this smoke extend to 60 feet from the mjork, providing smoke for the Smoke Sense of all mjorks. However, the whole shroud dissipates during any round the mjork is exposed to moderate (at least 10 miles per hour) or stronger wind.
\DndMonsterAction{Water Susceptibility}
For every 5 feet the mjork asher moves in water, or for every gallon of water splashed on it, it takes 1 cold damage.
\par
\DndMonsterSection{Actions}
\DndMonsterAction{Multiattack}
The mjork asher makes two attacks.
\DndMonsterAction{Slam}
\textsl{Melee Weapon Attack:} +5 to hit, reach 5 ft., one target. \textsl{Hit:} 9 (2d6 + 2) bludgeoning damage and 7 (2d6) fire damage and Blinding Soot.
\DndMonsterAction{Hurl Soot}
\textsl{Ranged Spell Attack:} +5 to hit, ranged 30/60 ft., one target. \textsl{Hit:} 12 (3d6 + 2) fire damage and Blinding Soot.
\DndMonsterAction{Create Sootlings (1/Day)}
The mjork asher chooses a point within 60 feet of it that it can sense. Smoke fills a 20-foot radius sphere centered on that point. The sphere spreads around corners, and its area is heavily obscured. This smoke lasts until the start of the mjork's next turn or until a wind of moderate or greater speed (at least 10 miles per hour) disperses it. Unless dispersed before the start of the mjork's next turn, a sootling swarm forms in a space in the area, favoring a space a creature hostile to the mjork occupies. The swarm takes its turn just after the mjork and obeys the mjork's verbal commands (no action for the mjork). If issued no commands, the swarm acts independently, often attacking a creature that isn't its ally. The swarm can survive indefinitely after being created.
\end{DndMonster}



\begin{DndMonster}[float*=t]{Mjork Burner}
\DndMonsterType{Large elemental, chaotic neutral}
\DndMonsterBasics[
armor-class = {14 (natural armor)},
hit-points = {\DndDice{10d10 + 40}},
speed = {40 ft.}
]
\DndMonsterAbilityScores[
 str = 15,
 dex = 15,
 con = 18,
 int = 7,
 wis = 12,
 cha = 7
]
\DndMonsterDetails[
damage-vulnerabilities = {thunder},
damage-resistances = {; bludgeoning, piercing, slashing from nonmagical attacks},
damage-immunities = {fire, poison},
senses = {darkvision 60 ft., passive perception 11},
languages = {Ignan, Terran},
challenge = 6,
]
\DndMonsterAction{Fiery Body}
A creature that touches the mjork burner or hits it with a melee attack while within 5 feet of it takes 7 (2d6) fire damage.
\DndMonsterAction{Fire Path}
The mjork burner leaves fire behind in any space it moves through. Until the start of the mjork's next turn, any creature that enters this area for the first time on a turn or ends its turn there takes 5 (2d4) fire damage. The fire can ignite flammable objects (not creatures), causing the mjork's Fire Path to last longer.
\DndMonsterAction{Ignite}
A creature or flammable object ignites if it takes fire damage from the mjork burner's attacks. Until a creature takes an action to douse this fire, the target takes 5 (2d4) fire damage at the start of each of its turns.
\DndMonsterAction{Illumination}
The mjork burner sheds bright light in a 10-foot radius and dim light for another 10 feet.
\DndMonsterAction{Smoke Sense}
A mjork burner in contact with smoke has blindsight in that smoke's area in a radius of up to 120 feet.
\DndMonsterAction{Water Susceptibility}
For every 5 feet the mjork burner moves in water, or for every gallon of water splashed on it, it takes 1 cold damage.
\par
\DndMonsterSection{Actions}
\DndMonsterAction{Multiattack}
The mjork burner makes two attacks.
\DndMonsterAction{Slam}
\textsl{Melee Weapon Attack:} +5 to hit, reach 5 ft., one target. \textsl{Hit:} 9 (2d6 + 2) bludgeoning damage, 7 (2d6) fire damage, and Ignite.
\DndMonsterAction{Hurl Fire}
\textsl{Ranged Spell Attack:} +5 to hit, ranged 30/60 ft., one target. \textsl{Hit:} 13 (3d6 + 3) fire damage and Ignite.
\end{DndMonster}


\clearpage

\begin{DndMonster}[float*=t]{Mjork Charger}
\DndMonsterType{Large elemental, chaotic neutral}
\DndMonsterBasics[
armor-class = {13 (natural armor)},
hit-points = {\DndDice{11d10 + 44}},
speed = {40 ft.}
]
\DndMonsterAbilityScores[
 str = 18,
 dex = 11,
 con = 18,
 int = 6,
 wis = 10,
 cha = 6
]
\DndMonsterDetails[
damage-vulnerabilities = {thunder},
damage-resistances = {; bludgeoning, piercing, slashing from nonmagical attacks},
damage-immunities = {fire, poison},
senses = {darkvision 60 ft., passive perception 10},
languages = {Ignan, Terran},
challenge = 7,
]
\DndMonsterAction{Charge}
If the mjork charger moves at least 20 feet straight toward a target and then hits it with a slam attack on the same turn, the target takes an extra 4 (1d8) bludgeoning damage and 3 (1d6) fire damage. If the target is a creature, it must succeed on a DC 15 Strength saving throw or fall prone.
\DndMonsterAction{Fiery Body}
A creature that touches the mjork charger or hits it with a melee attack while within 5 feet of it takes 7 (2d6) fire damage.
\DndMonsterAction{Illumination}
The mjork charger sheds dim light in a 10-foot radius.
\DndMonsterAction{Momentum}
If the mjork charger moves on its turn before attacking, creatures have disadvantage on attack rolls for opportunity attacks against it.
\DndMonsterAction{Smoke Sense}
A mjork charger in contact with smoke has blindsight in that smoke's area in a radius of up to 120 feet.
\DndMonsterAction{Water Susceptibility}
For every 5 feet the mjork charger moves in water, or for every gallon of water splashed on it, it takes 1 cold damage.
\par
\DndMonsterSection{Actions}
\DndMonsterAction{Multiattack}
The mjork makes two slam attacks.
\DndMonsterAction{Slam}
\textsl{Melee Weapon Attack:} +7 to hit, reach 5 ft., one target. \textsl{Hit:} 13 (2d8 + 4) bludgeoning damage and 7 (2d6) fire damage.
\end{DndMonster}



\begin{DndMonster}[float*=t]{Mjork Sootling}
\DndMonsterType{Tiny elemental, chaotic neutral}
\DndMonsterBasics[
armor-class = {12},
hit-points = {\DndDice{3d4 + 3}},
speed = {25 ft.}
]
\DndMonsterAbilityScores[
 str = 4,
 dex = 15,
 con = 12,
 int = 4,
 wis = 10,
 cha = 3
]
\DndMonsterDetails[
damage-vulnerabilities = {thunder},
damage-immunities = {fire, poison},
senses = {darkvision 60 ft., passive perception 10},
languages = {understands Ignan and Terran but can't speak},
challenge = 1/4,
]
\DndMonsterAction{Heated Body}
A creature that ends its turn grappling the sootling takes 1 fire damage.
\DndMonsterAction{Smoke Sense}
A sootling in contact with smoke has blindsight in that smoke's area in a radius of up to 120 feet.
\DndMonsterAction{Water Susceptibility}
For every 5 feet the sootling moves in water, or for every gallon of water splashed on it, it takes 1 cold damage.
\par
\DndMonsterSection{Actions}
\DndMonsterAction{Burn}
\textsl{Melee Weapon Attack:} +4 to hit, reach 5 ft., one target. \textsl{Hit:} 4 (1d4 + 2) fire damage.
\end{DndMonster}



\begin{DndMonster}[float*=t]{Mjork Sootling Swarm}
\DndMonsterType{Medium elemental, chaotic neutral}
\DndMonsterBasics[
armor-class = {12},
hit-points = {\DndDice{6d8 + 6}},
speed = {25 ft.}
]
\DndMonsterAbilityScores[
 str = 9,
 dex = 15,
 con = 12,
 int = 4,
 wis = 10,
 cha = 3
]
\DndMonsterDetails[
damage-vulnerabilities = {thunder},
damage-resistances = {bludgeoning, piercing, slashing},
damage-immunities = {fire, poison},
senses = {darkvision 60 ft., passive perception 10},
languages = {understands Ignan and Terran but can't speak},
challenge = 2,
]
\DndMonsterAction{Hot Soot}
A creature that ends its turn in the sootling swarm's space takes 5 (2d4) fire damage, or 2 (1d4) fire damage if the swarm has half of its hit points or fewer. The creature must also succeed on a DC 11 Constitution saving throw or be blinded and unable to breathe until the end of its next turn.
\DndMonsterAction{Smoke Sense}
A sootling swarm in contact with smoke has blindsight in that smoke's area in a radius of up to 120 feet.
\DndMonsterAction{Swarm}
The sootling swarm can occupy another creature's space and vice versa, and the swarm can move through any opening large enough for a Tiny creature. The swarm can't regain hit points or gain temporary hit points
\DndMonsterAction{Water Susceptibility}
For every 5 feet the sootling swarm moves in water, or for every gallon of water splashed on it, it takes 2 cold damage.
\par
\DndMonsterSection{Actions}
\DndMonsterAction{Burns}
\textsl{Melee Weapon Attack:} +4 to hit, reach 0 ft., one target in the swarm's space. \textsl{Hit:} 10 (4d4) fire damage, or 5 (2d4) fire damage if the swarm has half of its hit points or fewer.
\end{DndMonster}



\begin{DndMonster}[float*=t]{Mold Spider}
\DndMonsterType{Tiny plant, unaligned}
\DndMonsterBasics[
armor-class = {11},
hit-points = {\DndDice{1d4 + 1}},
speed = {20 ft., 20 ft. (key)}
]
\DndMonsterAbilityScores[
 str = 2,
 dex = 13,
 con = 13,
 int = 1,
 wis = 8,
 cha = 2
]
\DndMonsterDetails[
damage-vulnerabilities = {fire},
senses = {blindsight 60 ft. (blind beyond this radius), passive perception 9},
challenge = 0,
]
\DndMonsterAction{Regrowth}
When a mold spider dies, unless fire destroyed it, it releases spores that regrow a mold spider in 3 days.
\DndMonsterAction{Spider Climb}
The mold spider can climb difficult surfaces, including upside down on ceilings, without needing to make an ability check.
\par
\DndMonsterSection{Actions}
\DndMonsterAction{Claw}
\textsl{Melee Weapon Attack:} +3 to hit, reach 5 ft., one target. \textsl{Hit:} 1 slashing damage.
\end{DndMonster}


\clearpage

\begin{DndMonster}[float*=t]{Morbus Kobold}
\DndMonsterType{Small humanoid (kobold), lawful evil}
\DndMonsterBasics[
armor-class = {12},
hit-points = {\DndDice{2d6 + 2}},
speed = {30 ft., 30 ft. (key)}
]
\DndMonsterAbilityScores[
 str = 6,
 dex = 14,
 con = 13,
 int = 10,
 wis = 9,
 cha = 8
]
\DndMonsterDetails[
saving-throws = {Con +3},
skills = {Medicine +1, Stealth +4},
damage-resistances = {poison},
senses = {darkvision 60 ft., passive perception 9},
languages = {Draconic and one other language},
challenge = 1/8,
]
\DndMonsterAction{Carrier}
The morbus kobold can contract a disease but is immune to effects other than cosmetic symptoms.
\DndMonsterAction{Infection}
A creature that touches the morbus kobold or anything it has touched in the past 24 hours is exposed to a disease, usually sewer plague.
\DndMonsterAction{Pack Tactics}
The morbus kobold has advantage on an attack roll against a creature if at least one of the kobold's allies is within 5 feet of the creature and the ally isn't incapacitated.
\DndMonsterAction{Sunlight Sensitivity}
While in sunlight, the morbus kobold has disadvantage on attack rolls, as well as on Wisdom (Perception) checks that rely on sight.
\par
\DndMonsterSection{Actions}
\DndMonsterAction{Dagger}
\textsl{Melee or Ranged Weapon Attack:} +4 to hit, reach 5 ft. or range 20/60 ft., one target. \textsl{Hit:} 4 (1d4 + 2) piercing damage and Infection.
\DndMonsterAction{Sling}
\textsl{Ranged Weapon Attack:} +4 to hit, range 30/120 ft., one target. \textsl{Hit:} 4 (1d4 + 2) bludgeoning damage and Infection.
\DndMonsterAction{Gore Bag}
The morbus kobold throws a sack filled with diseased gore at a point it can see within 20 feet of it. Each creature within 5 feet of the chosen point is subjected to Infection. A kobold usually carries only one such bag at a time.
\end{DndMonster}



\begin{DndMonster}[float*=t]{Noble}
\DndMonsterType{Medium humanoid (any race), any alignment}
\DndMonsterBasics[
armor-class = {15 (breastplate)},
hit-points = {\DndDice{2d8}},
speed = {30 ft.}
]
\DndMonsterAbilityScores[
 str = 11,
 dex = 12,
 con = 11,
 int = 12,
 wis = 14,
 cha = 16
]
\DndMonsterDetails[
skills = {Deception +5, Insight +4, Persuasion +5},
languages = {any two languages},
challenge = 1/8,
]
\par
\DndMonsterSection{Actions}
\DndMonsterAction{Rapier}
\textsl{Melee Weapon Attack:} +3 to hit, reach 5 ft., one target. \textsl{Hit:} 5 (1d8 + 1) piercing damage.
\par
\DndMonsterSection{Reactions}
\DndMonsterAction{Parry}
The noble adds 2 to its AC against one melee attack that would hit it. To do so, the noble must see the attacker and be wielding a melee weapon.
\end{DndMonster}



\begin{DndMonster}[float*=t]{Noble}
\DndMonsterType{Medium humanoid (any race), any alignment}
\DndMonsterBasics[
armor-class = {15 (breastplate)},
hit-points = {\DndDice{2d8}},
speed = {30 ft.}
]
\DndMonsterAbilityScores[
 str = 11,
 dex = 12,
 con = 11,
 int = 12,
 wis = 14,
 cha = 16
]
\DndMonsterDetails[
skills = {Deception +5, Insight +4, Persuasion +5},
languages = {any two languages},
challenge = 1/8,
]
\par
\DndMonsterSection{Actions}
\DndMonsterAction{Rapier}
\textsl{Melee Weapon Attack:} +3 to hit, reach 5 ft., one target. \textsl{Hit:} 5 (1d8 + 1) piercing damage.
\par
\DndMonsterSection{Reactions}
\DndMonsterAction{Parry}
The noble adds 2 to its AC against one melee attack that would hit it. To do so, the noble must see the attacker and be wielding a melee weapon.
\end{DndMonster}



\begin{DndMonster}[float*=t]{Oblivion Brute}
\DndMonsterType{Large undead, neutral evil}
\DndMonsterBasics[
armor-class = {12},
hit-points = {\DndDice{8d10 + 16}},
speed = {40 ft.}
]
\DndMonsterAbilityScores[
 str = 16,
 dex = 14,
 con = 15,
 int = 10,
 wis = 10,
 cha = 8
]
\DndMonsterDetails[
skills = {Stealth +5},
damage-vulnerabilities = {radiant},
damage-resistances = {cold; bludgeoning, piercing, slashing from nonmagical attacks},
damage-immunities = {necrotic, poison},
senses = {darkvision 60 ft., passive perception 10},
challenge = 5,
]
\DndMonsterAction{Shadow Stealth}
While in dim light or darkness, the oblivion brute can take the Hide action as a bonus action.
\DndMonsterAction{Sunlight Weakness}
While in sunlight, the oblivion brute has disadvantage on attack rolls, ability checks, and saving throws.
\par
\DndMonsterSection{Actions}
\DndMonsterAction{Soul Drain}
\textsl{Melee Weapon Attack:} +6 to hit, reach 5 ft., one creature. \textsl{Hit:} 16 (3d8 + 3) necrotic damage, and the target must make a DC 13 Constitution saving throw. On a failure, the target's available Hit Dice are reduced by 1. If this effect reduces the target's Hit Dice to 0, the target falls unconscious. The target can repeat the saving throw at the end of each of its turns, regaining consciousness on a success. Hit Dice lost this way return at the end of a short rest, in time to spend them to regain hit points.
If an oblivion brute kills a living creature with this attack when the creature is at 0 Hit Dice, the target's soul is lost to the realm of shadow. This lost soul can be raised from the dead only with mighty magic, such as wish or divine intervention.
\DndMonsterAction{Unbearable Howl (Recharges after a Short or Long Rest)}
The oblivion brute emits a howl in a 30-foot cone. Each creature that isn't deafened or undead in the area must succeed on a DC 13 Wisdom saving throw or become frightened for 1 minute. While frightened this way, a creature's speed is halved. A frightened target can repeat the saving throw at the end of each of its turns, ending the effect on itself on a success. If a target's saving throw is successful or the effect ends for it, the target is immune to any oblivion brute's Unbearable Howl for the next 24 hours.
\end{DndMonster}


\clearpage

\begin{DndMonster}[float*=t]{Oblivion Juggernaut}
\DndMonsterType{Huge undead, neutral evil}
\DndMonsterBasics[
armor-class = {12},
hit-points = {\DndDice{10d12 + 20}},
speed = {40 ft.}
]
\DndMonsterAbilityScores[
 str = 18,
 dex = 14,
 con = 15,
 int = 10,
 wis = 10,
 cha = 8
]
\DndMonsterDetails[
skills = {Stealth +5},
damage-vulnerabilities = {radiant},
damage-resistances = {cold; bludgeoning, piercing, slashing from nonmagical attacks},
damage-immunities = {necrotic, poison},
senses = {darkvision 60 ft., passive perception 10},
challenge = 8,
]
\DndMonsterAction{Shadow Stealth}
While in dim light or darkness, the oblivion juggernaut can take the Hide action as a bonus action.
\DndMonsterAction{Sunlight Weakness}
While in sunlight, the oblivion juggernaut has disadvantage on attack rolls, ability checks, and saving throws.
\par
\DndMonsterSection{Actions}
\DndMonsterAction{Multiattack}
The oblivion juggernaut uses unbearable roar and then uses soul drain.
\DndMonsterAction{Soul Drain}
\textsl{Melee Weapon Attack:} +7 to hit, reach 5 ft., one creature. \textsl{Hit:} 26 (4d10 + 4) necrotic damage, and the target must make a DC 13 Constitution saving throw. On a failure, the target's available Hit Dice are reduced by 2. If this effect reduces the target's Hit Dice to 0, the target falls unconscious. The target can repeat the saving throw at the end of each of its turns, regaining consciousness on a success. Hit Dice lost this way return at the end of a short rest, in time to spend them to regain hit points.
If an oblivion juggernaut kills a living creature with this attack when the creature is at 0 Hit Dice, the target's soul is lost to the realm of shadow. This lost soul can be raised from the dead only with mighty magic, such as wish or divine intervention.
\DndMonsterAction{Unbearable Roar}
Each creature of the oblivion juggernaut's choice that isn't deafened and is within 120 feet of the juggernaut must succeed on a DC 18 Wisdom saving throw or become frightened for 1 minute. A creature can repeat the saving throw at the end of each of its turns, ending the effect on itself on a success. If a creature's saving throw is successful or the effect ends for it, the creature is immune to any oblivion juggernaut's Unbearable Roar for 24 hours.
\end{DndMonster}



\begin{DndMonster}[float*=t]{Oblivion Leaper}
\DndMonsterType{Small undead, neutral evil}
\DndMonsterBasics[
armor-class = {13},
hit-points = {\DndDice{4d6}},
speed = {30 ft.}
]
\DndMonsterAbilityScores[
 str = 10,
 dex = 16,
 con = 11,
 int = 10,
 wis = 10,
 cha = 8
]
\DndMonsterDetails[
skills = {Stealth +5},
damage-vulnerabilities = {radiant},
damage-resistances = {cold; bludgeoning, piercing, slashing from nonmagical attacks},
damage-immunities = {necrotic, poison},
senses = {darkvision 60 ft., passive perception 10},
challenge = 1,
]
\DndMonsterAction{Shadow Leap}
Starting in an area of dim light or darkness, an oblivion leaper can teleport to an unoccupied space also in dim light or darkness. Every foot of teleportation costs the oblivion leaper 1 foot of movement.
\DndMonsterAction{Shadow Stealth}
While in dim light or darkness, the oblivion leaper can take the Hide action as a bonus action.
\DndMonsterAction{Sunlight Weakness}
While in sunlight, the oblivion leaper has disadvantage on attack rolls, ability checks, and saving throws.
\par
\DndMonsterSection{Actions}
\DndMonsterAction{Soul Drain}
\textsl{Melee Weapon Attack:} +5 to hit, reach 5 ft., one creature. \textsl{Hit:} 7 (2d4 + 3) necrotic damage, and the target must make a DC 11 Constitution saving throw. On a failure, the target's available Hit Dice are reduced by 1. If this effect reduces the target's Hit Dice to 0, the target falls unconscious. The target can repeat the saving throw at the end of each of its turns, regaining consciousness on a success. Hit Dice lost this way return at the end of a short rest, in time to spend them to regain hit points.
If an oblivion leaper kills a living creature with this attack when the creature is at 0 Hit Dice, the target's soul is lost to the realm of shadow. This lost soul can be raised from the dead only with mighty magic, such as wish or divine intervention.
\end{DndMonster}



\begin{DndMonster}[float*=t]{Oblivion Whistler}
\DndMonsterType{Medium undead, neutral evil}
\DndMonsterBasics[
armor-class = {13},
hit-points = {\DndDice{6d8 + 6}},
speed = {40 ft.}
]
\DndMonsterAbilityScores[
 str = 12,
 dex = 16,
 con = 13,
 int = 10,
 wis = 10,
 cha = 8
]
\DndMonsterDetails[
skills = {Stealth +5},
damage-vulnerabilities = {radiant},
damage-resistances = {cold; bludgeoning, piercing, slashing from nonmagical attacks},
damage-immunities = {necrotic, poison},
senses = {darkvision 60 ft., passive perception 10},
challenge = 2,
]
\DndMonsterAction{Shadow Stealth}
While in dim light or darkness, the oblivion whistler can take the Hide action as a bonus action.
\DndMonsterAction{Sunlight Weakness}
While in sunlight, the oblivion whistler has disadvantage on attack rolls, ability checks, and saving throws.
\DndMonsterAction{Unbearable Whistling}
Creatures of the oblivion whistler's choice that aren't deafened and start their turn within 30 feet of the oblivion whistler must make a DC 11 Wisdom saving throw. A target that fails is frightened until the start of its next turn and unable to hide the distress of hearing the whistling, which is the source of the fear and considered to be within line of sight while the target can hear and is within 30 feet of it. On a successful saving throw, the creature is immune to any oblivion whistler's Unbearable Whistling for 1 hour.
\par
\DndMonsterSection{Actions}
\DndMonsterAction{Soul Drain}
\textsl{Melee Weapon Attack:} +5 to hit, reach 5 ft., one creature. \textsl{Hit:} 10 (2d6 + 3) necrotic damage, and the target must make a DC 11 Constitution saving throw. On a failure, the target's available Hit Dice are reduced by 1. If this effect reduces the target's Hit Dice to 0, the target falls unconscious. The target can repeat the saving throw at the end of each of its turns, regaining consciousness on a success. Hit Dice lost this way return at the end of a short rest, in time to spend them to regain hit points.
If an oblivion whistler kills a living creature with this attack when the creature is at 0 Hit Dice, the target's soul is lost to the realm of shadow. This lost soul can be raised from the dead only with mighty magic, such as wish or divine intervention.
\end{DndMonster}



\begin{DndMonster}[float*=t]{Oozing Vulture}
\DndMonsterType{Medium monstrosity, unaligned}
\DndMonsterBasics[
armor-class = {13 (natural armor)},
hit-points = {\DndDice{10d8 + 10}},
speed = {20 ft., 50 ft. (key)}
]
\DndMonsterAbilityScores[
 str = 16,
 dex = 13,
 con = 12,
 int = 5,
 wis = 14,
 cha = 8
]
\DndMonsterDetails[
skills = {Perception +4, Stealth +3},
senses = {darkvision 60 ft., passive perception 14},
challenge = 2,
]
\DndMonsterAction{Blood Frenzy}
The oozing vulture has advantage on melee attack rolls against any creature that doesn't have all its hit points.
\DndMonsterAction{Keen Sight and Smell}
The oozing vulture has advantage on Wisdom (Perception) checks that rely on sight or smell.
\par
\DndMonsterSection{Actions}
\DndMonsterAction{Bite}
\textsl{Melee Weapon Attack:} +5 to hit, reach 5 ft., one target. \textsl{Hit:} 10 (2d6 + 3) slashing damage. If the target is a creature, it must succeed on a DC 13 Strength saving throw or be knocked prone.
\end{DndMonster}


\clearpage

\begin{DndMonster}[float*=t]{Ordeal Tree}
\DndMonsterType{Gargantuan plant, neutral evil}
\DndMonsterBasics[
armor-class = {17 (natural armor)},
hit-points = {\DndDice{12d20 + 60}},
speed = {20 ft.}
]
\DndMonsterAbilityScores[
 str = 24,
 dex = 7,
 con = 21,
 int = 12,
 wis = 16,
 cha = 12
]
\DndMonsterDetails[
damage-vulnerabilities = {fire},
damage-resistances = {bludgeoning, piercing},
languages = {Sylvan},
challenge = 12,
]
\DndMonsterAction{Corrupt Desire}
A humanoid may bind themselves to the ordeal tree. Each day the petitioner is bound to the tree, they must succeed on a Constitution saving throw with a DC equal to 15 + number of days bound or suffer a level of exhaustion. If they survive for 7 days, the ordeal tree may choose to grant a blessing. The tree only does this if the humanoid is evil, once per humanoid. If the humanoid is neutral or good, the tree attacks them.
The blessing is tailored to the petitioner's desires (GM discretion) but might include +1 to an ability score, a feat, expertise in a skill, or the ability to innately cast a 3rd-level (or lower) spell once per day using Charisma as the petitioner's spellcasting ability.
The petitioner also contracts the curse of foul blight (Grim Hallow Campaign Guide). If the petitioner commits a profoundly evil act every day, the curse never advances beyond stage 1. If the petitioner fails to commit the act, the curse automatically escalates.
\DndMonsterAction{False Appearance}
While the ordeal tree remains motionless, it is indistinguishable from a normal tree.
\DndMonsterAction{Heart Sight}
The ordeal tree knows the alignment and desires of any creature that binds itself to the tree.
\par
\DndMonsterSection{Actions}
\DndMonsterAction{Multiattack}
The ordeal tree makes two slam attacks and one attack with a vine or enervate.
\DndMonsterAction{Slam}
\textsl{Melee Weapon Attack:} +11 to hit, reach 10 ft., one target. \textsl{Hit:} 29 (4d10 + 7) bludgeoning damage.
\DndMonsterAction{Vine}
\textsl{Melee Weapon Attack:} +11 to hit, reach 20 ft., one target. \textsl{Hit:} 20 (3d8 + 7) slashing damage, and the creature is grappled (escape DC 17).
\DndMonsterAction{Enervate (2/Day)}
A creature bound to the ordeal tree or grappled by it must succeed on a DC 17 Constitution saving throw or suffer 44 (8d10) necrotic damage. The ordeal tree gains temporary hit points equal to half the damage done.
\end{DndMonster}



\begin{DndMonster}[float*=t]{Panjaian Ilharan}
\DndMonsterType{Medium humanoid (any race, elemental), any alignment}
\DndMonsterBasics[
armor-class = {15 (natural armor)},
hit-points = {\DndDice{10d8 + 10}},
speed = {30 ft., 30 ft. (key)}
]
\DndMonsterAbilityScores[
 str = 10,
 dex = 16,
 con = 13,
 int = 12,
 wis = 15,
 cha = 13
]
\DndMonsterDetails[
saving-throws = {Dex +5},
skills = {Acrobatics +5, Deception +3, Perception +4, Sleight Of Hand +5, Stealth +5},
damage-resistances = {lightning, thunder},
senses = {darkvision 60 ft., passive perception 14},
languages = {Primordial and two other languages},
challenge = 3,
]
\DndMonsterAction{Returning Wind}
The panjaian ilharan can take a bonus action to choose up to four objects it can see within 60 feet of it. These objects can't be carried or worn by another creature, and they can weigh up to 5 pounds combined. Those objects fly up to 60 feet toward the ilharan, delivered into its hands or open containers, such as sheathes, once within 5 feet.
\DndMonsterAction{Innate Spellcasting}
The panjaian ilharan's innate spellcasting ability is Wisdom (spell save DC 12). It can innately cast the following spells, requiring no material components:
\begin{DndMonsterSpells}
  \DndInnateSpellLevel{mage hand}
  \DndInnateSpellLevel[3e]{gust of wind, invisibility}
\end{DndMonsterSpells}
\par
\DndMonsterSection{Actions}
\DndMonsterAction{Multiattack}
The panjaian ilharan makes four dagger attacks. It can draw the daggers as part of this action.
\DndMonsterAction{Dagger}
\textsl{Melee or Ranged Weapon Attack:} +5 to hit, reach 5 ft. or range 20/60 ft., one target. \textsl{Hit:} 5 (1d4 + 3) piercing damage.
\DndMonsterAction{Storm Wave Recharge 5-6}
The panjaian ilharan unleashes lightning, thunder, and strong wind in a 15-foot cone. Each creature in that area must make a DC 13 Constitution saving throw, taking 6 (2d4 + 1) lightning damage and 6 (2d4 + 1) thunder damage on a failed save, or half as much damage on a successful one. Those who fail the saving throw are also pushed up to 15 feet to the edge of the area.
\DndMonsterAction{Gift of Flight (1/Day)}
The panjaian ilharan touches up to four willing creatures. Each target gains a flying speed of 30 feet for 1 hour. When the effect ends, the target descends at 60 feet per round for 1 minute, taking no falling damage and landing on its feet if on the ground after that minute.
\end{DndMonster}



\begin{DndMonster}[float*=t]{Psychic Fragment}
\DndMonsterType{Tiny undead, chaotic evil}
\DndMonsterBasics[
armor-class = {12},
hit-points = {\DndDice{1d4 + 1}},
speed = {0 ft., {'number': 30, 'condition': '(hover)'} ft. (key), True ft. (key)}
]
\DndMonsterAbilityScores[
 str = 3,
 dex = 15,
 con = 13,
 int = 2,
 wis = 12,
 cha = 10
]
\DndMonsterDetails[
damage-resistances = {acid, cold, fire; bludgeoning, piercing, slashing from nonmagical attacks},
damage-immunities = {necrotic, poison},
senses = {darkvision 60 ft., passive perception 11},
languages = {can repeat one phrase it knew in life},
challenge = 0,
]
\DndMonsterAction{Incorporeal Movement}
The psychic fragment can move through other creatures and objects as if they were difficult terrain. It takes 5 (1d10) force damage if it ends its turn inside an object.
\DndMonsterAction{Mimicry}
The psychic fragment can mimic a single word or phrase it spoke in life. A creature that hears the sounds can tell they are imitations with a successful DC 11 Wisdom (Insight) check.
\par
\DndMonsterSection{Actions}
\DndMonsterAction{Withering Touch}
\textsl{Melee Weapon Attack:} +4 to hit, reach 5 ft., one target. \textsl{Hit:} 2 (1d4) necrotic damage.
\end{DndMonster}



\begin{DndMonster}[float*=t]{Psychic Fragment Swarm}
\DndMonsterType{Medium undead, chaotic evil}
\DndMonsterBasics[
armor-class = {13},
hit-points = {\DndDice{10d8 + 10}},
speed = {0 ft., 30 ft. (key)}
]
\DndMonsterAbilityScores[
 str = 3,
 dex = 16,
 con = 13,
 int = 7,
 wis = 14,
 cha = 12
]
\DndMonsterDetails[
damage-resistances = {acid, cold, fire; bludgeoning, piercing, slashing from nonmagical attacks},
damage-immunities = {necrotic, poison},
senses = {darkvision 60 ft., passive perception 14},
languages = {can repeat one phrase it knew in life},
challenge = 3,
]
\DndMonsterAction{Incorporeal Movement}
The swarm can move through other creatures and objects as if they were difficult terrain. It takes 5 (1d10) force damage if it ends its turn inside an object.
\DndMonsterAction{Maddening Voices}
The swarm constantly emits a babble of words and phrases spoken in life by its many undead spirits. Each creature that starts its turn within 20 ft. of the swarm must succeed on a DC 13 Wisdom saving throw. On a failure, the creature is stunned until the start of its next turn. A creature that succeeds on a saving throw becomes immune to maddening voices for 24 hours.
\DndMonsterAction{Swarm}
The swarm can occupy another creature's space and vice versa. The swarm can't regain hit points or gain temporary hit points.
\par
\DndMonsterSection{Actions}
\DndMonsterAction{Withering Touch}
\textsl{Melee Weapon Attack:} +5 to hit, reach 5 ft., one target. \textsl{Hit:} 21 (6d6) necrotic damage, or 10 (3d6) necrotic damage if the swarm has half of its hit points or fewer.
\end{DndMonster}


\clearpage

\begin{DndMonster}[float*=t]{Riding Horse}
\DndMonsterType{Large beast, unaligned}
\DndMonsterBasics[
armor-class = {10},
hit-points = {\DndDice{2d10 + 2}},
speed = {60 ft.}
]
\DndMonsterAbilityScores[
 str = 16,
 dex = 10,
 con = 12,
 int = 2,
 wis = 11,
 cha = 7
]
\DndMonsterDetails[
challenge = 1/4,
]
\par
\DndMonsterSection{Actions}
\DndMonsterAction{Hooves}
\textsl{Melee Weapon Attack:} +5 to hit, reach 5 ft., one target. \textsl{Hit:} 8 (2d4 + 3) bludgeoning damage.
\end{DndMonster}



\begin{DndMonster}[float*=t]{Scout}
\DndMonsterType{Medium humanoid (any race), any alignment}
\DndMonsterBasics[
armor-class = {13 (leather armor)},
hit-points = {\DndDice{3d8 + 3}},
speed = {30 ft.}
]
\DndMonsterAbilityScores[
 str = 11,
 dex = 14,
 con = 12,
 int = 11,
 wis = 13,
 cha = 11
]
\DndMonsterDetails[
skills = {Nature +4, Perception +5, Stealth +6, Survival +5},
languages = {any one language (usually Common)},
challenge = 1/2,
]
\DndMonsterAction{Keen Hearing and Sight}
The scout has advantage on Wisdom (Perception) checks that rely on hearing or sight.
\par
\DndMonsterSection{Actions}
\DndMonsterAction{Multiattack}
The scout makes two melee attacks or two ranged attacks.
\DndMonsterAction{Shortsword}
\textsl{Melee Weapon Attack:} +4 to hit, reach 5 ft., one target. \textsl{Hit:} 5 (1d6 + 2) piercing damage.
\DndMonsterAction{Longbow}
\textsl{Ranged Weapon Attack:} +4 to hit, ranged 150/600 ft., one target. \textsl{Hit:} 6 (1d8 + 2) piercing damage.
\end{DndMonster}



\begin{DndMonster}[float*=t]{Scream Thief}
\DndMonsterType{Medium undead, chaotic evil}
\DndMonsterBasics[
armor-class = {12},
hit-points = {\DndDice{8d8 + 16}},
speed = {0 ft., {'number': 60, 'condition': '(hover)'} ft. (key), True ft. (key)}
]
\DndMonsterAbilityScores[
 str = 6,
 dex = 15,
 con = 14,
 int = 10,
 wis = 12,
 cha = 12
]
\DndMonsterDetails[
damage-resistances = {acid, cold, fire; bludgeoning, piercing, slashing from nonmagical attacks},
damage-immunities = {necrotic, poison},
senses = {darkvision 60 ft., passive perception 11},
languages = {languages it knew in life},
challenge = 3,
]
\DndMonsterAction{Incorporeal Movement}
The scream thief can move through other creatures and objects as if they were difficult terrain. It takes 5 (1d10) force damage if it ends its turn inside an object.
\par
\DndMonsterSection{Actions}
\DndMonsterAction{Life Drain}
\textsl{Melee Weapon Attack:} +4 to hit, reach 5 ft., one target. \textsl{Hit:} 18 (4d8) necrotic damage. The target must succeed on a DC 12 Constitution saving throw or its hit point maximum is reduced by an amount equal to the damage taken. This reduction lasts until the target finishes a long rest. The target dies if this effect reduces its hit point maximum to 0.
\DndMonsterAction{Silence the Living Recharge 5-6}
The scream thief forces a humanoid within 30 feet of it to make a DC13 Wisdom saving throw. On a failure, the target takes 19 (3d12) psychic damage, is frightened, and is unable to vocalize a sound of any kind. The target may repeat the saving throw at the end of its turns to end these effects. On a success the creature takes half as much damage and is not frightened or silenced.
\par
\DndMonsterSection{Reactions}
\DndMonsterAction{Fragment the Spirit}
The scream thief targets a creature that died from one of its attacks. The target's spirit rises as 3 (1d6) psychic fragments in the space of its corpse or in the nearest unoccupied space. The psychic fragments are under the scream thief's control.
\end{DndMonster}



\begin{DndMonster}[float*=t]{Sea Drake}
\DndMonsterType{Medium dragon, neutral evil}
\DndMonsterBasics[
armor-class = {14 (natural armor)},
hit-points = {\DndDice{9d8 + 18}},
speed = {20 ft., 50 ft. (key)}
]
\DndMonsterAbilityScores[
 str = 15,
 dex = 16,
 con = 14,
 int = 7,
 wis = 14,
 cha = 8
]
\DndMonsterDetails[
skills = {Perception +4, Stealth +5},
damage-resistances = {cold},
senses = {darkvision 120 ft., passive perception 14},
languages = {Draconic},
challenge = 2,
]
\DndMonsterAction{Amphibious}
The sea drake can breathe air and water.
\par
\DndMonsterSection{Actions}
\DndMonsterAction{Multiattack}
The sea drake makes one bite attack and one tail attack.
\DndMonsterAction{Bite}
\textsl{Melee Weapon Attack:} +5 to hit, reach 5 ft., one target. \textsl{Hit:} 10 (2d6 + 3) piercing damage.
\DndMonsterAction{Tail}
\textsl{Melee Weapon Attack:} +5 to hit, reach 10 ft., one target. \textsl{Hit:} 6 (1d6 + 3) bludgeoning damage. If the target is a Medium or smaller creature, it is pulled within 5 feet of the sea drake and grappled (escape DC 13). Until this grapple ends, the drake can't use its tail on another target.
\end{DndMonster}


\clearpage

\begin{DndMonster}[float*=t]{Shadow}
\DndMonsterType{Medium undead, chaotic evil}
\DndMonsterBasics[
armor-class = {12},
hit-points = {\DndDice{3d8 + 3}},
speed = {40 ft.}
]
\DndMonsterAbilityScores[
 str = 6,
 dex = 14,
 con = 13,
 int = 6,
 wis = 10,
 cha = 8
]
\DndMonsterDetails[
skills = {Stealth +4},
damage-vulnerabilities = {radiant},
damage-resistances = {acid, cold, fire, lightning, thunder; bludgeoning, piercing, slashing from nonmagical attacks},
damage-immunities = {necrotic, poison},
senses = {darkvision 60 ft., passive perception 10},
challenge = 1/2,
]
\DndMonsterAction{Amorphous}
The shadow can move through a space as narrow as 1 inch wide without squeezing.
\DndMonsterAction{Shadow Stealth}
While in dim light or darkness, the shadow can take the Hide action as a bonus action. Its stealth bonus is also improved to +6.
\DndMonsterAction{Sunlight Weakness}
While in sunlight, the shadow has disadvantage on attack rolls, ability checks, and saving throws.
\par
\DndMonsterSection{Actions}
\DndMonsterAction{Strength Drain}
\textsl{Melee Weapon Attack:} +4 to hit, reach 5 ft., one creature. \textsl{Hit:} 9 (2d6 + 2) necrotic damage, and the target's Strength score is reduced by 1d4. The target dies if this reduces its Strength to 0. Otherwise, the reduction lasts until the target finishes a short or long rest.
If a non-evil humanoid dies from this attack, a new shadow rises from the corpse 1d4 hours later.
\end{DndMonster}



\begin{DndMonster}[float*=t]{Shadowsteel Ghast}
\DndMonsterType{Medium undead, chaotic evil}
\DndMonsterBasics[
armor-class = {16 (natural armor)},
hit-points = {\DndDice{17d8 + 42}},
speed = {40 ft., 40 ft. (key)}
]
\DndMonsterAbilityScores[
 str = 16,
 dex = 14,
 con = 17,
 int = 14,
 wis = 13,
 cha = 8
]
\DndMonsterDetails[
saving-throws = {Int +5, Wis +4},
skills = {Arcana +5, Perception +4},
damage-resistances = {necrotic},
damage-immunities = {poison},
senses = {darkvision 120 ft., passive perception 14},
languages = {two languages known in life},
challenge = 7,
]
\DndMonsterAction{Magic Resistance}
The Shadowsteel ghast has advantage on saving throws against spells and other magical effects.
\DndMonsterAction{Shadowsteel Coat}
The Shadowsteel ghast is coated in metal armor made of Shadowsteel.
\DndMonsterAction{Shadowsteel Miasma}
Creatures of the Shadowsteel ghast's choice that start their turn within 30 feet of the ghast must roll a d4 and subtract the number rolled from any attack roll or saving throw they make until the start of their next turn. Only one miasma can affect a given creature at a time. A creature that rolls a natural 20 on any attack roll or saving throw while in a Shadowsteel miasma is immune to the miasma of any Shadowsteel ghast for 1 hour.
\DndMonsterAction{Turning Defiance}
The Shadowsteel ghast and any Shadowsteel ghouls within 30 feet of it have advantage on saving throws against effects that turn undead.
\DndMonsterAction{Innate Spellcasting}
The Shadowsteel ghast's innate spellcasting ability is Intelligence (spell save DC 13, +5 to hit with spell attacks). It can innately cast the following spells, requiring no components:
\begin{DndMonsterSpells}
  \DndInnateSpellLevel{chill touch, shocking grasp}
  \DndInnateSpellLevel[2]{bestow curse}
  \DndInnateSpellLevel[3e]{counterspell, misty step, shield}
\end{DndMonsterSpells}
\par
\DndMonsterSection{Actions}
\DndMonsterAction{Multiattack}
The Shadowsteel ghast makes two claw attacks or casts one at-will spell and makes one claw attack.
\DndMonsterAction{Claw}
\textsl{Melee Weapon Attack:} +6 to hit, reach 5 ft., one target. \textsl{Hit:} 12 (2d8 + 3) slashing damage and 7 (2d6) necrotic damage. A cursed target must make an escalation check. If a target drops to 0 hit points due to this damage and isn't already cursed, they are randomly afflicted with one from among any of the curses in the Grim Hollow Campaign Guide (GM's choice).
\end{DndMonster}



\begin{DndMonster}[float*=t]{Shadowsteel Ghoul}
\DndMonsterType{Medium undead, chaotic evil}
\DndMonsterBasics[
armor-class = {15 (natural armor)},
hit-points = {\DndDice{10d8 + 20}},
speed = {30 ft.}
]
\DndMonsterAbilityScores[
 str = 16,
 dex = 12,
 con = 14,
 int = 9,
 wis = 10,
 cha = 6
]
\DndMonsterDetails[
saving-throws = {Wis +2},
skills = {Perception +2},
damage-immunities = {poison},
senses = {darkvision 120 ft., passive perception 12},
languages = {one language known in life},
challenge = 4,
]
\DndMonsterAction{Magic Resistance}
The Shadowsteel ghoul has advantage on saving throws against spells and other magical effects.
\DndMonsterAction{Shadowsteel Coat}
The Shadowsteel ghoul is coated in metal armor made of Shadowsteel.
\DndMonsterAction{Shadowsteel Hunger}
The Shadowsteel ghoul has advantage on attack rolls against creatures carrying or wearing Shadowsteel.
\par
\DndMonsterSection{Actions}
\DndMonsterAction{Multiattack}
The Shadowsteel ghoul makes two claw attacks.
\DndMonsterAction{Claw}
\textsl{Melee Weapon Attack:} +5 to hit, reach 5 ft., one target. \textsl{Hit:} 10 (2d6 + 3) slashing damage and 5 (2d4) necrotic damage. A cursed target must make an escalation check. If a target drops to 0 hit points due to this damage and isn't already cursed, roll 1d4, and the target is afflicted with a curse based on the result: 1-2 Curse of Damned Aging, 3-4 Curse of Foul Blight. (See Grim Hollow Campaign Guide for information on curses and escalation.)
\end{DndMonster}



\begin{DndMonster}[float*=t]{Shambling Mound}
\DndMonsterType{Large plant, unaligned}
\DndMonsterBasics[
armor-class = {15 (natural armor)},
hit-points = {\DndDice{16d10 + 48}},
speed = {20 ft., 20 ft. (key)}
]
\DndMonsterAbilityScores[
 str = 18,
 dex = 8,
 con = 16,
 int = 5,
 wis = 10,
 cha = 5
]
\DndMonsterDetails[
skills = {Stealth +2},
damage-resistances = {cold, fire},
damage-immunities = {lightning},
senses = {blindsight 60 ft. (blind beyond this radius), passive perception 10},
challenge = 5,
]
\DndMonsterAction{Lightning Absorption}
Whenever the shambling mound is subjected to lightning damage, it takes no damage and regains a number of hit points equal to the lightning damage dealt.
\par
\DndMonsterSection{Actions}
\DndMonsterAction{Multiattack}
The shambling mound makes two slam attacks. If both attacks hit a Medium or smaller target, the target is grappled (escape DC 14), and the shambling mound uses its Engulf on it.
\DndMonsterAction{Slam}
\textsl{Melee Weapon Attack:} +7 to hit, reach 5 ft., one target. \textsl{Hit:} 13 (2d8 + 4) bludgeoning damage.
\DndMonsterAction{Engulf}
The shambling mound engulfs a Medium or smaller creature grappled by it. The engulfed target is blinded, restrained, and unable to breathe, and it must succeed on a DC 14 Constitution saving throw at the start of each of the mound's turns or take 13 (2d8 + 4) bludgeoning damage. If the mound moves, the engulfed target moves with it. The mound can have only one creature engulfed at a time.
\end{DndMonster}


\clearpage

\begin{DndMonster}[float*=t]{Shatter Corpse}
\DndMonsterType{Medium undead, chaotic neutral}
\DndMonsterBasics[
armor-class = {10 (natural armor)},
hit-points = {\DndDice{6d8 + 18}},
speed = {20 ft.}
]
\DndMonsterAbilityScores[
 str = 13,
 dex = 6,
 con = 16,
 int = 3,
 wis = 6,
 cha = 5
]
\DndMonsterDetails[
saving-throws = {Wis +0},
damage-immunities = {piercing, poison},
senses = {darkvision 60 ft., passive perception 8},
languages = {understands the languages it knew in life but can't speak},
challenge = 1,
]
\DndMonsterAction{Glass Spikes}
Each time a creature makes a melee attack against a shatter corpse, it takes 3 piercing damage. A creature can choose to make an attack with disadvantage to avoid this damage.
\DndMonsterAction{Loud}
Creatures have advantage on Wisdom (Perception) checks made to hear the approach of a shatter corpse.
\DndMonsterAction{Turn Resistance}
The shatter corpse has advantage on saving throws against any effect that turns undead.
\DndMonsterAction{Undead Fortitude}
If damage reduces the shatter corpse to 0 hit points, it must make a Constitution saving throw with a DC of 5 + the damage taken, unless the damage is radiant or from a critical hit. On a success, the shatter corpse drops to 1 hit point instead.
\par
\DndMonsterSection{Actions}
\DndMonsterAction{Slam}
\textsl{Melee Weapon Attack:} +3 to hit, reach 5 ft., one target. \textsl{Hit:} 4 (1d6 + 1) bludgeoning damage and 2 (1d4) piercing damage.
\DndMonsterAction{Volley of Glass Recharge 6}
The shatter corpse flicks a volley of glass shards in a 15-foot cone. Each creature in that area must make a DC 13 Dexterity saving throw, taking 10 (3d6) slashing damage on a failed save, or half as much damage on a successful one.
\end{DndMonster}



\begin{DndMonster}[float*=t]{Shieldhead}
\DndMonsterType{Large beast, unaligned}
\DndMonsterBasics[
armor-class = {13 (natural armor)},
hit-points = {\DndDice{4d10 + 8}},
speed = {50 ft.}
]
\DndMonsterAbilityScores[
 str = 16,
 dex = 11,
 con = 14,
 int = 2,
 wis = 12,
 cha = 7
]
\DndMonsterDetails[
challenge = 1/2,
]
\DndMonsterAction{Charge}
If the shieldhead moves at least 20 feet straight toward a target and then hits it with a ram attack on the same turn, the target takes an extra 7 (2d6) bludgeoning damage. If the target is a creature, it must succeed on a DC 13 Strength saving throw or fall prone.
\par
\DndMonsterSection{Actions}
\DndMonsterAction{Ram}
\textsl{Melee Weapon Attack:} +5 to hit, reach 5 ft., one target. \textsl{Hit:} 10 (2d6 + 3) bludgeoning damage.
\end{DndMonster}



\begin{DndMonster}[float*=t]{Sitri Cat}
\DndMonsterType{Tiny monstrosity, unaligned}
\DndMonsterBasics[
armor-class = {13},
hit-points = {\DndDice{2d4}},
speed = {40 ft., 30 ft. (key)}
]
\DndMonsterAbilityScores[
 str = 3,
 dex = 16,
 con = 10,
 int = 3,
 wis = 12,
 cha = 7
]
\DndMonsterDetails[
skills = {Perception +3, Stealth +5},
challenge = 0,
]
\DndMonsterAction{Keen Smell}
The Sitri cat has advantage on Wisdom (Perception) checks that rely on smell.
\DndMonsterAction{Seven Shades of Sitri}
A humanoid who touches a Sitri cat must succeed on a DC 13 Constitution saving throw or become infected with the Seven Shades of Sitri. If a humanoid succeeds on the saving throw, they're immune to the disease for 24 hours. Carnal urges overwhelm an infected humanoid. When the infected interacts with another non-hostile humanoid, the infected becomes charmed by that humanoid until the subject of the infatuation or their friends harm the infected. The infected can repeat the saving throw once every 24 hours. Current infatuations end when one of these saves succeeds. The disease ends for the infected when they have succeeded on five saving throws.
\par
\DndMonsterSection{Actions}
\DndMonsterAction{Claws}
\textsl{Melee Weapon Attack:} +5 to hit, reach 5 ft., one target. \textsl{Hit:} 1 slashing damage and Seven Shades of Sitri.
\end{DndMonster}



\begin{DndMonster}[float*=t]{Skeleton Cannoneer}
\DndMonsterType{Medium undead, lawful evil}
\DndMonsterBasics[
armor-class = {14 (rotting buff coat)},
hit-points = {\DndDice{8d8 + 16}},
speed = {30 ft.}
]
\DndMonsterAbilityScores[
 str = 12,
 dex = 16,
 con = 15,
 int = 8,
 wis = 13,
 cha = 5
]
\DndMonsterDetails[
skills = {Perception +3},
damage-vulnerabilities = {bludgeoning},
damage-immunities = {poison},
senses = {darkvision 120 ft., passive perception 13},
languages = {understands languages it knew in life but can't speak},
challenge = 2,
]
\DndMonsterAction{Dead Firearms}
The cannon and pistol a skeleton cannoneer uses are nonfunctional as ranged weapons for a creature other than a skeleton cannoneer. For the skeleton cannoneer, the pistol never needs to be reloaded.
\par
\DndMonsterSection{Actions}
\DndMonsterAction{Load}
The skeleton cannoneer loads the dead cannon.
\DndMonsterAction{Dead Cannon}
The skeleton cannoneer fires its loaded cannon at a point it can see within 500 feet of it, but no closer than 50 feet from it. Each creature within 15 feet of that point must make a DC 13 Dexterity saving throw, taking 7 (2d6) bludgeoning damage and 7 (2d6) necrotic damage on a failed save, or half as much damage on a successful one.
Alternatively, the cannoneer can choose the cannon as the point of fire. Doing so means the shot goes off inside the cannon, destroying it and increasing each type of damage by 3 (1d6).
\DndMonsterAction{Pistol Club}
\textsl{Melee Weapon Attack:} +5 to hit, reach 5 ft., one target. \textsl{Hit:} 5 (1d6 + 2) bludgeoning damage and 3 (1d6) necrotic damage.
\DndMonsterAction{Dead Pistol}
\textsl{Ranged Weapon Attack:} +5 to hit, range 30/90 ft., one target. \textsl{Hit:} 8 (1d10 + 3) piercing damage plus 3 (1d6) necrotic damage.
\end{DndMonster}


\clearpage

\begin{DndMonster}[float*=t]{Skeleton Commander}
\DndMonsterType{Medium undead, lawful evil}
\DndMonsterBasics[
armor-class = {17 (battered splint mail)},
hit-points = {\DndDice{10d8 + 20}},
speed = {30 ft.}
]
\DndMonsterAbilityScores[
 str = 16,
 dex = 10,
 con = 15,
 int = 11,
 wis = 12,
 cha = 14
]
\DndMonsterDetails[
saving-throws = {Wis +3},
skills = {Insight +3, Perception +3},
damage-vulnerabilities = {bludgeoning},
damage-immunities = {poison},
senses = {darkvision 60 ft., passive perception 13},
languages = {understands languages it knew in life but can't speak},
challenge = 3,
]
\par
\DndMonsterSection{Actions}
\DndMonsterAction{Multiattack}
The skeleton commander makes two pike attacks.
\DndMonsterAction{Pike}
\textsl{Melee Weapon Attack:} +5 to hit, reach 10 ft., one target. \textsl{Hit:} 8 (1d10 + 3) piercing damage and 5 (2d4) necrotic damage.
\DndMonsterAction{Warlord Strike Recharge 5-6}
The skeleton commander makes one pike attack. The commander can use a bonus action to direct another undead who can see or hear the commander to attack the target the commander attacked this turn. That undead can use a reaction to make the attack, adding the commander's proficiency bonus to the damage roll.
\par
\DndMonsterSection{Reactions}
\DndMonsterAction{Brandish Standard}
If an undead the skeleton commander can see within 30 feet of the commander makes an attack roll or saving throw, the commander can brandish the tattered banner on its pike. If the commander does so, the target can add a d4 to its roll, provided it can see the commander and isn't receiving this benefit from another commander.
\end{DndMonster}



\begin{DndMonster}[float*=t]{Skeleton Rifler}
\DndMonsterType{Medium undead, lawful evil}
\DndMonsterBasics[
armor-class = {12 (rotting buff coat)},
hit-points = {\DndDice{4d8 + 8}},
speed = {30 ft.}
]
\DndMonsterAbilityScores[
 str = 10,
 dex = 15,
 con = 15,
 int = 7,
 wis = 10,
 cha = 5
]
\DndMonsterDetails[
skills = {Perception +2},
damage-vulnerabilities = {bludgeoning},
damage-immunities = {poison},
senses = {darkvision 60 ft., passive perception 12},
languages = {understands languages it knew in life but can't speak},
challenge = 1/2,
]
\DndMonsterAction{Dead Rifle}
The rifle a skeleton rifler carries is nonfunctional as a ranged weapon for a creature other than a skeleton rifler. For the skeleton rifler, the rifle never needs to be reloaded.
\par
\DndMonsterSection{Actions}
\DndMonsterAction{Bayonet}
\textsl{Melee Weapon Attack:} +4 to hit, reach 5 ft., one target. \textsl{Hit:} 5 (1d6 + 2) piercing damage and 3 (1d6) necrotic damage.
\DndMonsterAction{Dead Rifle}
\textsl{Ranged Weapon Attack:} +4 to hit, range 40/120 ft., one target. \textsl{Hit:} 8 (1d12 + 2) piercing damage and 3 (1d6) necrotic damage.
\end{DndMonster}



\begin{DndMonster}[float*=t]{Skinweaver}
\DndMonsterType{Large aberration, chaotic evil}
\DndMonsterBasics[
armor-class = {17 (natural & tailored leather)},
hit-points = {\DndDice{8d10 + 24}},
speed = {30 ft., 15 ft. (key)}
]
\DndMonsterAbilityScores[
 str = 12,
 dex = 14,
 con = 16,
 int = 14,
 wis = 13,
 cha = 18
]
\DndMonsterDetails[
skills = {Stealth +4},
damage-resistances = {fire},
senses = {darkvision 60 ft., passive perception 11},
languages = {any three languages},
challenge = 3,
]
\DndMonsterAction{Supernaturally Sharp}
Skinweaver's tanning knives overcome slashing damage resistance (though not immunity).
\DndMonsterAction{Shocking Visage}
Any humanoid that starts its turn within 30 feet of the skinweaver and can see it must make a DC 11 Wisdom saving throw. On a failed save, the creature is frightened until the end of the creature's next turn). Whether a creature succeeds or fails, it cannot be affected by this ability again until it loses sight of the skinweaver (as when the skinweaver moves behind total cover or successfully hides).
\par
\DndMonsterSection{Actions}
\DndMonsterAction{Multiattack}
The skinweaver makes two tanning knife attacks or one tanning knife and one web whip attack.
\DndMonsterAction{Tanning Knife}
\textsl{Melee Weapon Attack:} +4 to hit, reach 5 ft., one target. \textsl{Hit:} 11 (2d8 + 2) slashing damage.
\DndMonsterAction{Web Whip}
\textsl{Ranged Weapon Attack:} Attack: +4 to hit, reach 25 ft., one target. \textsl{Hit:} 7 (2d4 + 2) piercing damage, and the target is pulled up to 20 feet toward the skinweaver.
\DndMonsterAction{Web Net Recharge 5-6}
\textsl{Ranged Weapon Attack:} +4 to hit, range 30/60 ft., one creature. \textsl{Hit:} The target is restrained by leathery webbing. As an action, the restrained target can attempt a DC 13 Strength check, bursting the webbing on a success. The webbing can also be attacked and destroyed (AC 10; hp 10; vulnerability to slashing damage; immunity to bludgeoning, poison, and psychic damage).
\end{DndMonster}



\begin{DndMonster}[float*=t]{Sky Drake}
\DndMonsterType{Large dragon, neutral evil}
\DndMonsterBasics[
armor-class = {14 (natural armor)},
hit-points = {\DndDice{10d10 + 30}},
speed = {20 ft., 60 ft. (key)}
]
\DndMonsterAbilityScores[
 str = 18,
 dex = 11,
 con = 17,
 int = 8,
 wis = 14,
 cha = 11
]
\DndMonsterDetails[
skills = {Athletics +6, Perception +4, Survival +4},
senses = {darkvision 120 ft., passive perception 14},
languages = {Draconic and one other language},
challenge = 4,
]
\DndMonsterAction{Flyby}
The sky drake doesn't provoke opportunity attacks when it flies out of an enemy's reach.
\DndMonsterAction{Keen Sight}
The sky drake has advantage on Wisdom (Perception) checks that rely on sight.
\par
\DndMonsterSection{Actions}
\DndMonsterAction{Multiattack}
The sky drake makes one bite attack and one attack with its claws.
\DndMonsterAction{Bite}
\textsl{Melee Weapon Attack:} +6 to hit, reach 10 ft., one target. \textsl{Hit:} 11 (2d6 + 4) piercing damage and 5 (2d4) lightning damage.
\DndMonsterAction{Claws}
\textsl{Melee Weapon Attack:} +6 to hit, reach 5 ft., one target. \textsl{Hit:} 13 (2d8 + 4) slashing damage.
\end{DndMonster}


\clearpage

\begin{DndMonster}[float*=t]{Sloth Galloper}
\DndMonsterType{Large beast, unaligned}
\DndMonsterBasics[
armor-class = {11},
hit-points = {\DndDice{4d10 + 4}},
speed = {25 ft., 25 ft. (key)}
]
\DndMonsterAbilityScores[
 str = 17,
 dex = 13,
 con = 12,
 int = 4,
 wis = 12,
 cha = 7
]
\DndMonsterDetails[
senses = {darkvision 60 ft., passive perception 11},
challenge = 1,
]
\DndMonsterAction{Speed Burst}
Provided it isn't moving over difficult terrain, the sloth galloper can use a bonus action to take the Dash action.
\par
\DndMonsterSection{Actions}
\DndMonsterAction{Multiattack}
The sloth galloper makes two slam attacks.
\DndMonsterAction{Slam}
\textsl{Melee Weapon Attack:} +5 to hit, reach 5 ft., one target. \textsl{Hit:} 8 (2d4 + 3) bludgeoning damage.
\end{DndMonster}



\begin{DndMonster}[float*=t]{Snapjaw}
\DndMonsterType{Gargantuan beast, unaligned}
\DndMonsterBasics[
armor-class = {17 (natural armor)},
hit-points = {\DndDice{9d12 + 27}},
speed = {30 ft., 50 ft. (key)}
]
\DndMonsterAbilityScores[
 str = 21,
 dex = 9,
 con = 17,
 int = 10,
 wis = 10,
 cha = 7
]
\DndMonsterDetails[
skills = {Stealth +5},
challenge = 5,
]
\DndMonsterAction{Hold Breath}
Snapjaw can hold its breath for 30 minutes.
\par
\DndMonsterSection{Actions}
\DndMonsterAction{Multiattack}
Snapjaw makes three attacks: two with its bite and one with its tail.
\DndMonsterAction{Bite}
\textsl{Melee Weapon Attack:} +11 to hit, reach 5 ft., one target. \textsl{Hit:} 31 (3d10 + 15) piercing damage, and the target is grappled (escape DC 16). Until this grapple ends, the target is restrained, and Snapjaw can't bite another target.
\DndMonsterAction{Tail}
\textsl{Melee Weapon Attack:} +11 to hit, reach 10 ft., one target not grappled by Snapjaw. \textsl{Hit:} 24 (2d8 + 15) bludgeoning damage. If the target is a creature, it must succeed on a DC 16 Strength saving throw or be knocked prone.
\end{DndMonster}



\begin{DndMonster}[float*=t]{Spythronar Sac}
\DndMonsterType{Tiny aberration, unaligned}
\DndMonsterBasics[
armor-class = {5},
hit-points = {\DndDice{1d4 - 1}},
speed = {0 ft.}
]
\DndMonsterAbilityScores[
 str = 1,
 dex = 1,
 con = 8,
 int = 1,
 wis = 3,
 cha = 1
]
\DndMonsterDetails[
senses = {tremorsense (blind beyond this radius), passive perception 6},
challenge = 0,
]
\DndMonsterAction{False Appearance}
The spythronar sac appears to be a tangled ball of string, twigs, and dirt. Someone who can see the sac can identify it with a successful DC 15 Intelligence (Arcana or Nature) check.
\DndMonsterAction{Fragile}
A creature who enters the spythronar sac's space must succeed on a DC 10 Dexterity saving throw, or the sac is destroyed.
\DndMonsterAction{Lightning Release}
When the spythronar sac is destroyed, it releases lightning in a 10-foot radius. A creature who destroyed the sac by entering its space receives no saving throw. Other creatures in that area must succeed on a DC 10 Dexterity saving throw or take 4 (1d8) lightning damage. Each spythronar swarm and web in this area instead gains advantage on its next attack roll.
\DndMonsterAction{Shocking Birth}
When a spythronar sac takes lightning damage from a source other than another spythronar, it hatches, transforming into a spythronar swarm with half the normal hit points. This swarm rolls initiative and enters the combat.
\end{DndMonster}



\begin{DndMonster}[float*=t]{Spythronar Swarm}
\DndMonsterType{Medium aberration, neutral evil}
\DndMonsterBasics[
armor-class = {13},
hit-points = {\DndDice{10d8}},
speed = {30 ft., 30 ft. (key)}
]
\DndMonsterAbilityScores[
 str = 5,
 dex = 17,
 con = 10,
 int = 3,
 wis = 11,
 cha = 3
]
\DndMonsterDetails[
skills = {Stealth +5},
damage-resistances = {bludgeoning, piercing, slashing},
damage-immunities = {lightning},
senses = {darkvision 60 ft., tremorsense 30 ft., passive perception 10},
challenge = 2,
]
\DndMonsterAction{Furtive}
The spythronar swarm can use a bonus action to take the Hide action.
\DndMonsterAction{Spider Climb}
The spythronar swarm can climb difficult surfaces, including upside down on ceilings, without needing to make an ability check.
\DndMonsterAction{Swarm}
The spythronar swarm can occupy another creature's space and vice versa, and the swarm can move through any opening large enough for a Tiny arachnid. The swarm can't regain hit points or gain temporary hit points.
\DndMonsterAction{Web Sense}
While in contact with a web, the spythronar swarm knows the exact location of any other creature in contact with the same web.
\DndMonsterAction{Web Walker}
The spythronar swarm ignores movement restrictions caused by webbing.
\par
\DndMonsterSection{Actions}
\DndMonsterAction{Bites}
\textsl{Melee Weapon Attack:} +5 to hit, reach 0 ft., one creature in the swarm's space. \textsl{Hit:} 10 (4d4) piercing damage and 9 (2d8) lightning damage, or 5 (2d4) piercing damage and 4 (1d8) lightning damage if the swarm has half of its hit points or fewer.
\DndMonsterAction{Weave (4/Day)}
The spythronar swarm weaves webbing, covering a 5-foot square. The swarm can then give up 1 hit point to create a spythronar sac in that space. If the swarm's webbing covers a 10-foot square, it can give up 23 hit points as part of the same action, creating a spythronar web with the same number of hit points. The swarm can't drop to 0 hit points to create a sac or web.
If the swarm weaves an incomplete web, that web is a normal web rather than a sentient, monstrous one. A creature that enters this normal web's space must succeed on a DC 10 Strength saving throw or become grappled (escape DC 10) and restrained while grappled in this way. Each 5-foot square of this web has AC 10; the damage vulnerabilities, immunities, and resistances of a spythronar web; and 5 hit points.
\end{DndMonster}


\clearpage

\begin{DndMonster}[float*=t]{Spythronar Web}
\DndMonsterType{Large aberration, neutral evil}
\DndMonsterBasics[
armor-class = {13},
hit-points = {\DndDice{10d10 + 10}},
speed = {30 ft., 30 ft. (key)}
]
\DndMonsterAbilityScores[
 str = 17,
 dex = 16,
 con = 13,
 int = 3,
 wis = 8,
 cha = 1
]
\DndMonsterDetails[
damage-vulnerabilities = {fire},
damage-resistances = {bludgeoning, piercing},
damage-immunities = {lightning},
senses = {tremorsense 60 ft. (blind beyond this radius), passive perception 9},
challenge = 2,
]
\DndMonsterAction{Adhesive}
The spythronar web adheres to anything that touches it. To avoid adhering when touching the web, a creature must succeed on a DC 13 Strength saving throw, but a creature hit by the web's crush attack doesn't receive a saving throw. A Huge or smaller creature adhered to the web is also grappled by it (escape DC 13). If adhered to a target larger than it can grapple, the web can release the adhesive or remain adhered and be dragged along as the target moves. Neither choice requires an action from the web. Removing an object adhered to the web or the web from anything it's adhered to requires a successful DC 13 Strength (Athletics) check as an action.
\DndMonsterAction{Amorphous}
The spythronar web can move through a space as narrow as 1 inch wide without squeezing.
\DndMonsterAction{False Appearance}
While the spythronar web remains motionless and takes no action, it is indistinguishable from a large, thick spider web.
\DndMonsterAction{Spider Climb}
A spythronar web can climb difficult surfaces, including upside down on ceilings, without needing to make an ability check.
\DndMonsterAction{Web Sense}
While in contact with a web, the spythronar web knows the exact location of any other creature in contact with the same web.
\DndMonsterAction{Web Walker}
The spythronar web ignores movement restrictions caused by webbing.
\par
\DndMonsterSection{Actions}
\DndMonsterAction{Crush}
\textsl{Melee Weapon Attack:} +5 to hit, reach 5 ft., one creature. \textsl{Hit:} 10 (2d6 + 3) bludgeoning damage and Adhesive. The target is also restrained while grappled in this way, and the spythronar web can't crush another target.
\DndMonsterAction{Lightning Surge}
\textsl{Melee Weapon Attack:} +5 to hit, reach 30 ft., one target. \textsl{Hit:} 9 (2d8) lightning damage. If the spythronar web has a creature grappled and attacks another target, the grappled creature also takes this damage.
\end{DndMonster}



\begin{DndMonster}[float*=t]{Stormborn Ithjar}
\DndMonsterType{Huge monstrosity, chaotic neutral}
\DndMonsterBasics[
armor-class = {17 (natural armor)},
hit-points = {\DndDice{18d12 + 72}},
speed = {40 ft., 120 ft. (key)}
]
\DndMonsterAbilityScores[
 str = 21,
 dex = 20,
 con = 19,
 int = 3,
 wis = 12,
 cha = 5
]
\DndMonsterDetails[
skills = {Acrobatics +9, Perception +5},
damage-immunities = {lightning, thunder},
senses = {blindsight 10 ft., darkvision 60 ft., passive perception 15},
challenge = 9,
]
\DndMonsterAction{Keen Sight}
The stormborn ithjar has advantage on Wisdom (Perception) checks that rely on sight.
\DndMonsterAction{Slippery Serpent}
The stormborn ithjar has advantage on ability checks or saving throws made to escape a grapple and saving throws to avoid being paralyzed or restrained.
\par
\DndMonsterSection{Actions}
\DndMonsterAction{Multiattack}
The stormborn ithjar attacks once with its bite and once with its tail.
\DndMonsterAction{Bite}
\textsl{Melee Weapon Attack:} +9 to hit, reach 15 ft., one target. \textsl{Hit:} 21 (3d10 + 5) slashing damage and 9 (2d8) lightning damage.
\DndMonsterAction{Tail}
\textsl{Melee Weapon Attack:} +9 to hit, reach 15 ft., one target. \textsl{Hit:} 18 (3d8 + 5) bludgeoning damage and 7 (2d6) thunder damage. If the target is a Large or smaller creature, it is grappled (escape DC 17). Until this grapple ends, the target is restrained, and the stormborn ithjar can't use its tail on another target.
\DndMonsterAction{Storm Breath Recharge 5-6}
The stormborn ithjar exhales a 30-foot cone of lightning and thunder. Each creature in that area must make a DC 16 Dexterity saving throw. On a failure, a creature takes 31 (7d8) lightning damage and 31 (7d8) thunder damage and falls prone. If the save succeeds, the creature takes half the lightning and thunder damage and doesn't fall prone. The discharge of this breath weapon is audible for 1 mile outside and 300 feet inside.
\end{DndMonster}



\begin{DndMonster}[float*=t]{Succubus}
\DndMonsterType{Medium fiend (shapechanger), neutral evil}
\DndMonsterBasics[
armor-class = {15 (natural armor)},
hit-points = {\DndDice{12d8 + 12}},
speed = {30 ft., 60 ft. (key)}
]
\DndMonsterAbilityScores[
 str = 8,
 dex = 17,
 con = 13,
 int = 15,
 wis = 12,
 cha = 20
]
\DndMonsterDetails[
skills = {Deception +9, Insight +5, Perception +5, Persuasion +9, Stealth +7},
damage-resistances = {cold, fire, lightning, poison; bludgeoning, piercing, slashing from nonmagical attacks},
senses = {darkvision 60 ft., passive perception 15},
languages = {Abyssal, Common, Infernal, telepathy 60 ft.},
challenge = 4,
]
\DndMonsterAction{Telepathic Bond}
The fiend ignores the range restriction on its telepathy when communicating with a creature it has charmed. The two don't even need to be on the same plane of existence.
\DndMonsterAction{Shapechanger}
The fiend can use its action to polymorph into a Small or Medium humanoid, or back into its true form. Without wings, the fiend loses its flying speed. Other than its size and speed, its statistics are the same in each form. Any equipment it is wearing or carrying isn't transformed. It reverts to its true form if it dies.
\par
\DndMonsterSection{Actions}
\DndMonsterAction{Claw (Fiend Form Only)}
\textsl{Melee Weapon Attack:} +5 to hit, reach 5 ft., one target. \textsl{Hit:} 6 (1d6 + 3) slashing damage.
\DndMonsterAction{Charm}
One humanoid the fiend can see within 30 feet of it must succeed on a DC 15 Wisdom saving throw or be magically charmed for 1 day. The charmed target obeys the fiend's verbal or telepathic commands. If the target suffers any harm or receives a suicidal command, it can repeat the saving throw, ending the effect on a success. If the target successfully saves against the effect, or if the effect on it ends, the target is immune to this fiend's Charm for the next 24 hours.
The fiend can have only one target charmed at a time. If it charms another, the effect on the previous target ends.
\DndMonsterAction{Draining Kiss}
The fiend kisses a creature charmed by it or a willing creature. The target must make a DC 15 Constitution saving throw against this magic, taking 32 (5d10 + 5) psychic damage on a failed save, or half as much damage on a successful one. The target's hit point maximum is reduced by an amount equal to the damage taken. This reduction lasts until the target finishes a long rest. The target dies if this effect reduces its hit point maximum to 0.
\DndMonsterAction{Etherealness}
The fiend magically enters the Ethereal Plane from the Material Plane, or vice versa.
\end{DndMonster}



\begin{DndMonster}[float*=t]{Swarm of Quippers}
\DndMonsterType{Medium beast, unaligned}
\DndMonsterBasics[
armor-class = {13},
hit-points = {\DndDice{8d8 - 8}},
speed = {0 ft., 40 ft. (key)}
]
\DndMonsterAbilityScores[
 str = 13,
 dex = 16,
 con = 9,
 int = 1,
 wis = 7,
 cha = 2
]
\DndMonsterDetails[
damage-resistances = {bludgeoning, piercing, slashing},
senses = {darkvision 60 ft., passive perception 8},
challenge = 1,
]
\DndMonsterAction{Blood Frenzy}
The swarm has advantage on melee attack rolls against any creature that doesn't have all its hit points.
\DndMonsterAction{Swarm}
The swarm can occupy another creature's space and vice versa, and the swarm can move through any opening large enough for a Tiny quipper. The swarm can't regain hit points or gain temporary hit points.
\DndMonsterAction{Water Breathing}
The swarm can breathe only underwater.
\par
\DndMonsterSection{Actions}
\DndMonsterAction{Bites}
\textsl{Melee Weapon Attack:} +5 to hit, reach 0 ft., one creature in the swarm's space. \textsl{Hit:} 14 (4d6) piercing damage, or 7 (2d6) piercing damage if the swarm has half of its hit points or fewer.
\end{DndMonster}


\clearpage

\begin{DndMonster}[float*=t]{Thornlamm}
\DndMonsterType{Small plant, neutral evil}
\DndMonsterBasics[
armor-class = {14 (natural armor)},
hit-points = {\DndDice{6d6}},
speed = {30 ft.}
]
\DndMonsterAbilityScores[
 str = 11,
 dex = 14,
 con = 11,
 int = 3,
 wis = 8,
 cha = 3
]
\DndMonsterDetails[
damage-resistances = {cold, fire, lightning, poison},
senses = {blindsight 60 ft. (blind beyond this radius), passive perception 9},
challenge = 1/4,
]
\par
\DndMonsterSection{Actions}
\DndMonsterAction{Hooves}
\textsl{Melee Weapon Attack:} +4 to hit, reach 5 ft., one target. \textsl{Hit:} 7 (2d4 + 2) bludgeoning damage.
\DndMonsterAction{Needles}
\textsl{Ranged Weapon Attack:} +4 to hit, range 30/60 ft., one target. \textsl{Hit:} 7 (2d4 + 2) piercing damage.
\end{DndMonster}



\begin{DndMonster}[float*=t]{Thug}
\DndMonsterType{Medium humanoid (any race), any non-good alignment}
\DndMonsterBasics[
armor-class = {11 (leather armor)},
hit-points = {\DndDice{5d8 + 10}},
speed = {30 ft.}
]
\DndMonsterAbilityScores[
 str = 15,
 dex = 11,
 con = 14,
 int = 10,
 wis = 10,
 cha = 11
]
\DndMonsterDetails[
skills = {Intimidation +2},
languages = {any one language (usually Common)},
challenge = 1/2,
]
\DndMonsterAction{Pack Tactics}
The thug has advantage on an attack roll against a creature if at least one of the thug's allies is within 5 feet of the creature and the ally isn't incapacitated.
\par
\DndMonsterSection{Actions}
\DndMonsterAction{Multiattack}
The thug makes two melee attacks.
\DndMonsterAction{Mace}
\textsl{Melee Weapon Attack:} +4 to hit, reach 5 ft., one creature. \textsl{Hit:} 5 (1d6 + 2) bludgeoning damage.
\DndMonsterAction{Heavy Crossbow}
\textsl{Ranged Weapon Attack:} +2 to hit, range 100/400 ft., one target. \textsl{Hit:} 5 (1d10) piercing damage.
\end{DndMonster}



\begin{DndMonster}[float*=t]{Torcheater}
\DndMonsterType{Tiny monstrosity, unaligned}
\DndMonsterBasics[
armor-class = {12},
hit-points = {\DndDice{1d4 + 1}},
speed = {5 ft., 30 ft. (key)}
]
\DndMonsterAbilityScores[
 str = 2,
 dex = 14,
 con = 12,
 int = 2,
 wis = 11,
 cha = 6
]
\DndMonsterDetails[
damage-immunities = {fire},
senses = {darkvision 60 ft., passive perception 10},
challenge = 0,
]
\DndMonsterAction{Eat Fire}
If the torcheater flies through fire, it reduces the fire's area by a 1-foot cube. A fire goes out if its area is reduced to nothing.
\par
\DndMonsterSection{Actions}
\DndMonsterAction{Bite}
\textsl{Melee Weapon Attack:} +4 to hit, reach 5 ft.one target. \textsl{Hit:} 1 fire damage.
\end{DndMonster}



\begin{DndMonster}[float*=t]{Venomous Gnoll}
\DndMonsterType{Medium humanoid (gnoll), chaotic evil}
\DndMonsterBasics[
armor-class = {13 (natural armor)},
hit-points = {\DndDice{6d8 + 6}},
speed = {30 ft.}
]
\DndMonsterAbilityScores[
 str = 12,
 dex = 14,
 con = 13,
 int = 6,
 wis = 13,
 cha = 7
]
\DndMonsterDetails[
saving-throws = {Con +3},
skills = {Stealth +4},
damage-resistances = {poison},
senses = {darkvision 60 ft., passive perception 11},
languages = {Gnoll},
challenge = 1,
]
\DndMonsterAction{Rampage}
When the venomous gnoll reduces a creature to 0 hit points with a melee attack on its turn, the gnoll can take a bonus action to move up to half its speed and make a bite attack.
\par
\DndMonsterSection{Actions}
\DndMonsterAction{Multiattack}
The venomous gnoll makes two attacks, only one of which can be a bite.
\DndMonsterAction{Bite}
\textsl{Melee Weapon Attack:} +4 to hit, reach 5 ft., one target. \textsl{Hit:} 4 (1d4 + 2) piercing damage and 2 (1d4) poison damage.
\DndMonsterAction{Claws}
\textsl{Melee Weapon Attack:} +4 to hit, reach 5 ft., one target. \textsl{Hit:} 5 (1d6 + 2) slashing damage and 2 (1d4) poison damage.
\DndMonsterAction{Spine Spike}
\textsl{Ranged Weapon Attack:} +4 to hit, range 20/60 ft., one target. \textsl{Hit:} 4 (1d4 + 2) piercing damage and 2 (1d4) poison damage.
\end{DndMonster}


\clearpage

\begin{DndMonster}[float*=t]{Vitebriate}
\DndMonsterType{Medium monstrosity, neutral evil}
\DndMonsterBasics[
armor-class = {14 (natural armor)},
hit-points = {\DndDice{7d8 + 21}},
speed = {30 ft.}
]
\DndMonsterAbilityScores[
 str = 11,
 dex = 14,
 con = 16,
 int = 12,
 wis = 12,
 cha = 16
]
\DndMonsterDetails[
damage-immunities = {poison},
senses = {darkvision 60 ft., passive perception 20},
languages = {languages it knew previously},
challenge = 2,
]
\DndMonsterAction{Unnatural Life}
Vitebriates are immune to the effects of old age and death by natural causes.
\DndMonsterAction{Vitebriance}
While on the same plane of existence as a person charmed by the vitebriate's Dominate Person action, the vitebriate may leech 1 hp per day permanently from the thrall. When a thrall is reduced to 0 hp through this process, it is killed. Drained hit points are regained after 7 days of rest.
\par
\DndMonsterSection{Actions}
\DndMonsterAction{Multiattack}
The vitebriate makes two melee attacks.
\DndMonsterAction{Dagger}
\textsl{Melee or Ranged Weapon Attack:} +4 to hit, reach 5 ft. or range 20/60 ft., one creature. \textsl{Hit:} 4 (1d4 + 2) piercing damage.
\DndMonsterAction{Dominate Person}
A vitebriate can target a single humanoid with dominate person as the spell, with the following differences: once the vitebriate has chosen the target, it can use this ability only on the chosen target (until the next day when it can choose a different target and drop the focus on the current target). The spell does not require concentration but is broken if the vitebriate is knocked unconscious involuntarily or killed (though not if it trances or sleeps). If the person has been charmed for at least an hour, damage does not force a new saving throw.
\par
\DndMonsterSection{Reactions}
\DndMonsterAction{Life Shield (1/Day)}
The vitebriate uses its psychic link with its thrall to transfer all damage it would take to its thrall. Damage that goes beyond what would kill the thrall is dealt to the vitebriate.
\end{DndMonster}



\begin{DndMonster}[float*=t]{Werebear Ascetic}
\DndMonsterType{Medium humanoid (any, shapechanger), any alignment}
\DndMonsterBasics[
armor-class = {12 in humanoid form (15 (natural armor) in hybrid form)},
hit-points = {\DndDice{20d8 + 60}},
speed = {{'number': 30, 'condition': '(50 ft. and 30 ft. climb in hybrid and kindred form)'} ft.}
]
\DndMonsterAbilityScores[
 str = 16,
 dex = 14,
 con = 17,
 int = 10,
 wis = 16,
 cha = 11
]
\DndMonsterDetails[
saving-throws = {Wis +7},
skills = {Animal Handling +7, Nature +4, Perception +7},
senses = {darkvision 60 ft., passive perception 17},
languages = {any three languages},
challenge = 12,
]
\DndMonsterAction{Shapechanger}
The werebear ascetic can use its action to polymorph into a bear-humanoid hybrid or into its kindred bear form, or back into its true form, which is humanoid. Its statistics, unless noted, are the same in each form. Any armor it is wearing merges into its alternate forms. It reverts to its true form if it dies.
\DndMonsterAction{Kindred Form}
The werebear's kindred form takes the shape of a Large bear. While in this form, it can speak to and understand bears, cannot take any actions requiring hands, adds its Constitution modifier to the result of any saving throw, and gains 25 temporary hit points the first time it takes this form each day.
\DndMonsterAction{Spellcasting}
The werebear is a 13th-level spellcaster. Its spellcasting ability is Wisdom (spell save DC 15, +6 to hit with spell attacks). It can cast spells in any form, regardless of verbal, somatic, or material components provided they have no gold cost. It has the following druid spells prepared:
\begin{DndMonsterSpells}
  \DndMonsterSpellLevel{druidcraft, guidance, mending, resistance}
   \DndMonsterSpellLevel[1][4]{animal friendship, faerie fire, healing word}
   \DndMonsterSpellLevel[2][3]{hold person, lesser restoration}
   \DndMonsterSpellLevel[3][3]{conjure animals, dispel magic, plant growth}
   \DndMonsterSpellLevel[4][3]{freedom of movement, ice storm}
   \DndMonsterSpellLevel[5][2]{commune with nature, tree stride}
   \DndMonsterSpellLevel[6][1]{wall of thorns}
   \DndMonsterSpellLevel[7][1]{arboreal curse}
\end{DndMonsterSpells}
\par
\DndMonsterSection{Actions}
\DndMonsterAction{Multiattack}
The werebear makes three melee attacks, only one of which can be a bite.
\DndMonsterAction{Bite (Hybrid or Kindred Form Only)}
\textsl{Melee Weapon Attack:} +7 to hit, reach 5 ft., one creature. \textsl{Hit:} 16 (3d8 + 3) piercing damage. If the target is a humanoid, it must succeed on a DC 15 Constitution saving throw or be cursed with werebear lycanthropy.
\DndMonsterAction{Claw (Hybrid or Kindred Form Only)}
\textsl{Melee Weapon Attack:} +7 to hit, reach 5 ft., one creature. \textsl{Hit:} 10 (2d6 + 3) slashing damage.
\end{DndMonster}



\begin{DndMonster}[float*=t]{Werewolf Ravager}
\DndMonsterType{Medium humanoid (any, shapechanger), any alignment}
\DndMonsterBasics[
armor-class = {13 in humanoid form (16 (natural armor) in hybrid form)},
hit-points = {\DndDice{14d8 + 28}},
speed = {{'number': 30, 'condition': '(50 ft. in hybrid form)'} ft.}
]
\DndMonsterAbilityScores[
 str = 17,
 dex = 16,
 con = 14,
 int = 9,
 wis = 15,
 cha = 10
]
\DndMonsterDetails[
saving-throws = {Str +6, Dex +6},
skills = {Perception +5},
damage-vulnerabilities = {; bludgeoning, piercing, slashing from nonmagical attacks that are silvered},
damage-immunities = {; bludgeoning, piercing, slashing from nonmagical attacks that aren't silvered},
senses = {darkvision 60 ft., passive perception 15},
languages = {any two languages},
challenge = 7,
]
\DndMonsterAction{Grappler}
While a creature is grappled by the werewolf ravager, the werewolf has advantage on attack rolls against grappled targets, and the target takes an additional 3 1d6 damage whenever it is hit with the werewolf's bite or claw attack.
\DndMonsterAction{Shapechanger}
The werewolf can use its action to polymorph into a wolf-humanoid hybrid or back into its true form, which is humanoid. Any armor it is wearing merges into its hybrid form. It reverts to its true form if it dies.
\par
\DndMonsterSection{Actions}
\DndMonsterAction{Multiattack}
The werewolf ravager makes two melee attacks. It can only make one bite attack per turn.
\DndMonsterAction{Bite (Hybrid Form Only)}
\textsl{Melee Weapon Attack:} +6 to hit, reach 5 ft., one creature. \textsl{Hit:} 12 (2d8 + 3) piercing damage. The target must succeed on a DC 14 Strength saving throw or be grappled.
Additionally, if the target is a humanoid, it must succeed on a DC 13 Constitution saving throw or be cursed with werewolf lycanthropy.
\DndMonsterAction{Claw (Hybrid Form Only)}
\textsl{Melee Weapon Attack:} +6 to hit, reach 5 ft., one creature. \textsl{Hit:} 10 (2d6 + 3) slashing damage.
\DndMonsterAction{Short Sword}
\textsl{Melee Weapon Attack:} +6 to hit, reach 5 ft., one creature. \textsl{Hit:} 6 (1d6 + 3) piercing damage.
\end{DndMonster}



\begin{DndMonster}[float*=t]{Winter Wolf}
\DndMonsterType{Large monstrosity, neutral evil}
\DndMonsterBasics[
armor-class = {13 (natural armor)},
hit-points = {\DndDice{10d10 + 20}},
speed = {50 ft.}
]
\DndMonsterAbilityScores[
 str = 18,
 dex = 13,
 con = 14,
 int = 7,
 wis = 12,
 cha = 8
]
\DndMonsterDetails[
skills = {Perception +5, Stealth +3},
damage-immunities = {cold},
languages = {Common, Giant, Winter Wolf},
challenge = 3,
]
\DndMonsterAction{Keen Hearing and Smell}
The wolf has advantage on Wisdom (Perception) checks that rely on hearing or smell.
\DndMonsterAction{Pack Tactics}
The wolf has advantage on an attack roll against a creature if at least one of the wolf's allies is within 5 feet of the creature and the ally isn't incapacitated.
\DndMonsterAction{Snow Camouflage}
The wolf has advantage on Dexterity (Stealth) checks made to hide in snowy terrain.
\par
\DndMonsterSection{Actions}
\DndMonsterAction{Bite}
\textsl{Melee Weapon Attack:} +6 to hit, reach 5 ft., one target. \textsl{Hit:} 11 (2d6 + 4) piercing damage. If the target is a creature, it must succeed on a DC 14 Strength saving throw or be knocked prone.
\DndMonsterAction{Cold Breath Recharge 5-6}
The wolf exhales a blast of freezing wind in a 15-foot cone. Each creature in that area must make a DC 12 Dexterity saving throw, taking 18 (4d8) cold damage on a failed save, or half as much damage on a successful one.
\end{DndMonster}


\clearpage

\begin{DndMonster}[float*=t]{Wraith}
\DndMonsterType{Medium undead, neutral evil}
\DndMonsterBasics[
armor-class = {13},
hit-points = {\DndDice{9d8 + 27}},
speed = {0 ft., {'number': 60, 'condition': '(hover)'} ft. (key), True ft. (key)}
]
\DndMonsterAbilityScores[
 str = 6,
 dex = 16,
 con = 16,
 int = 12,
 wis = 14,
 cha = 15
]
\DndMonsterDetails[
damage-resistances = {acid, cold, fire, lightning, thunder; bludgeoning, piercing, slashing from nonmagical attacks that aren't silvered},
damage-immunities = {necrotic, poison},
senses = {darkvision 60 ft., passive perception 12},
languages = {the languages it knew in life},
challenge = 5,
]
\DndMonsterAction{Incorporeal Movement}
The wraith can move through other creatures and objects as if they were difficult terrain. It takes 5 (1d10) force damage if it ends its turn inside an object.
\DndMonsterAction{Sunlight Sensitivity}
While in sunlight, the wraith has disadvantage on attack rolls, as well as on Wisdom (Perception) checks that rely on sight.
\par
\DndMonsterSection{Actions}
\DndMonsterAction{Life Drain}
\textsl{Melee Weapon Attack:} +6 to hit, reach 5 ft., one creature. \textsl{Hit:} 21 (4d8 + 3) necrotic damage. The target must succeed on a DC 14 Constitution saving throw or its hit point maximum is reduced by an amount equal to the damage taken. This reduction lasts until the target finishes a long rest. The target dies if this effect reduces its hit point maximum to 0.
\DndMonsterAction{Create Specter}
The wraith targets a humanoid within 10 feet of it that has been dead for no longer than 1 minute and died violently. The target's spirit rises as a specter in the space of its corpse or in the nearest unoccupied space. The specter is under the wraith's control. The wraith can have no more than seven specters under its control at one time.
\end{DndMonster}



\begin{DndMonster}[float*=t]{Xakalonus}
\DndMonsterType{Medium monstrosity, neutral evil}
\DndMonsterBasics[
armor-class = {16 (natural armor)},
hit-points = {\DndDice{10d8 + 30}},
speed = {30 ft.}
]
\DndMonsterAbilityScores[
 str = 18,
 dex = 15,
 con = 17,
 int = 9,
 wis = 12,
 cha = 9
]
\DndMonsterDetails[
skills = {Perception +4},
senses = {darkvision 120 ft., passive perception 14},
languages = {understands Draconic but can't speak},
challenge = 8,
]
\DndMonsterAction{Magic Absorption}
If a xakalonus succeeds on a saving throw against a magical effect that would damage it, it takes no damage. Instead, the xakalonus regains 11 (2d10) hit points. Alternatively, if the xakalonus is at full hit points, the next successful melee attack by the xakalonus deals an additional 11 (2d10) force damage.
\DndMonsterAction{Magic Resistance}
The xakalonus has advantage on saving throws against spells and other magical effects.
\par
\DndMonsterSection{Actions}
\DndMonsterAction{Multiattack}
The xakalonus makes one bite and one claws attack.
\DndMonsterAction{Bite}
\textsl{Melee Weapon Attack:} +7 to hit, reach 5 ft., one target. \textsl{Hit:} 13 (2d8 + 4) piercing damage.
\DndMonsterAction{Claws}
\textsl{Melee Weapon Attack:} +7 to hit, reach 5 ft., one target. \textsl{Hit:} 15 (2d10 + 4) slashing damage.
\DndMonsterAction{Sonic Bark Recharge 5-6}
All creatures within 30 ft. who can hear the xakalonus must succeed on a DC 14 Constitution saving throw or take 14 (4d6) thunder damage and be stunned until the end of their next turn. Other xakalonus and the hraptnon are immune to this.
\end{DndMonster}



\begin{DndMonster}[float*=t]{Young Blightscale Dragon}
\DndMonsterType{Large dragon (cursed), neutral evil}
\DndMonsterBasics[
armor-class = {18 (natural armor)},
hit-points = {\DndDice{16d10 + 64}},
speed = {40 ft., 80 ft. (key)}
]
\DndMonsterAbilityScores[
 str = 19,
 dex = 12,
 con = 19,
 int = 12,
 wis = 11,
 cha = 8
]
\DndMonsterDetails[
saving-throws = {Dex +4, Wis +3, Cha +2},
skills = {Insight +3, Perception +6, Stealth +4},
damage-immunities = {necrotic, poison},
senses = {blindsight 30 ft., darkvision 120 ft., passive perception 16},
languages = {Draconic and one national language},
challenge = 8,
]
\DndMonsterAction{Accursed Death}
When the dragon dies, it exudes a cloud of sickening, accursed gas. Each creature within 20 ft. of the dragon must succeed on a DC 15 Constitution saving throw or become poisoned for 1 minute. A creature can repeat the saving throw at the end of each of its turns, ending the effect on itself on a success. However, a creature that fails two of these saving throws becomes cursed. A creature cursed in this way takes 10 (3d6) necrotic damage and can't regain hit points, and the creature's hit point maximum is reduced by an amount equal to the necrotic damage taken. Every 24 hours, the cursed creature takes 10 (3d6) necrotic damage with similar reduction to the creature's hit point maximum. If the curse reduces the target's hit point maximum to 0, the target dies.
\par
\DndMonsterSection{Actions}
\DndMonsterAction{Multiattack}
The dragon makes three attacks: one with its bite and two with its claws.
\DndMonsterAction{Bite}
\textsl{Melee Weapon Attack:} +7 to hit, reach 10 ft., one target. \textsl{Hit:} 13 (2d8 + 4) piercing damage and 5 (2d4) necrotic damage.
\DndMonsterAction{Claw}
\textsl{Melee Weapon Attack:} +7 to hit, reach 5 ft., one target. \textsl{Hit:} 9 (2d4 + 4) slashing damage and 2 (1d4) poison damage.
\DndMonsterAction{Swarm Breath Recharge 5-6}
The dragon exhales a 30-foot cone of poisonous insects and vile gas. Each creature in that area must make a DC 15 Constitution saving throw, taking 14 (4d6) piercing damage and 28 (8d6) poison damage on a failed save, or half as much damage on a successful one.
\end{DndMonster}



\begin{DndMonster}[float*=t]{Zombie}
\DndMonsterType{Medium undead, neutral evil}
\DndMonsterBasics[
armor-class = {8},
hit-points = {\DndDice{3d8 + 9}},
speed = {20 ft.}
]
\DndMonsterAbilityScores[
 str = 13,
 dex = 6,
 con = 16,
 int = 3,
 wis = 6,
 cha = 5
]
\DndMonsterDetails[
saving-throws = {Wis +0},
damage-immunities = {poison},
senses = {darkvision 60 ft., passive perception 8},
languages = {understands all languages it spoke in life but can't speak},
challenge = 1/4,
]
\DndMonsterAction{Undead Fortitude}
If damage reduces the zombie to 0 hit points, it must make a Constitution saving throw with a DC of 5 + the damage taken, unless the damage is radiant or from a critical hit. On a success, the zombie drops to 1 hit point instead.
\par
\DndMonsterSection{Actions}
\DndMonsterAction{Slam}
\textsl{Melee Weapon Attack:} +3 to hit, reach 5 ft., one target. \textsl{Hit:} 4 (1d6 + 1) bludgeoning damage.
\end{DndMonster}


\clearpage

\begin{DndMonster}[float*=t]{Zombie Troll}
\DndMonsterType{Large undead, neutral evil}
\DndMonsterBasics[
armor-class = {14 (natural armor)},
hit-points = {\DndDice{9d10 + 36}},
speed = {30 ft.}
]
\DndMonsterAbilityScores[
 str = 19,
 dex = 6,
 con = 18,
 int = 3,
 wis = 6,
 cha = 5
]
\DndMonsterDetails[
damage-resistances = {cold, necrotic},
damage-immunities = {poison},
senses = {darkvision 60 ft., passive perception 8},
languages = {understands Giant but can't speak},
challenge = 4,
]
\DndMonsterAction{Regeneration}
The zombie troll regains 10 hit points at the start of its turn. If the zombie troll takes acid or radiant damage, this trait doesn't function at the start of the zombie troll's next turn. The zombie troll dies only if it starts its turn with 0 hit points and doesn't regenerate.
\par
\DndMonsterSection{Actions}
\DndMonsterAction{Multiattack}
The zombie troll makes two attacks.
\DndMonsterAction{Morningstar}
\textsl{Melee Weapon Attack:} +6 to hit, reach 5 ft., one target. \textsl{Hit:} 13 (2d8 + 4) bludgeoning damage.
\DndMonsterAction{Unrelenting Brute}
If a creature is hit by both morningstar attacks, the creature must take a DC 14 Strength saving throw or be knocked prone.
\end{DndMonster}


\end{document}
